% ============================================================
%  R. Nielsen, Natur og Aand.
%  Bidrag til en med Physiken stemmende Naturphilosophie.
%  Kjøbenhavn: Forlagt af J. H. Schubothes Boghandel, 1873.
%  (xiv + 549 pp., ANTIQUA)
%  TRANSCRIPTION (Danish)
%
%  ------------------------------------------------------------------
%  TWO WITNESSES. Both are in bibliotek/Nielsen, Rasmus/.
%
%  WORKING SCAN — Natur_og_aand.pdf, Google Books, 689 PDF pages,
%      600 ppi 1-bit JBIG2, one image per leaf, with a Google OCR text
%      layer. Every script reads this one.
%      It is a SAMMELBAND: PDF 572--689 are a second work, Grundideernes
%      Logik i kort Begreb (1870), transcribed separately at
%      texts/nielsen/grundideernes-logik-kort-begreb/. Nothing above
%      PDF 571 belongs to this edition. The "689 pp." that stood in
%      catalog.yaml was this PDF's page count, not the book's extent.
%  KB COLOUR SCAN — natur-og-aand-kb-color.pdf, Royal Danish Library
%      digitise-on-demand (www.kb.dk/e-mat/dod/130001514629-color.pdf),
%      575 PDF pages, 300 ppi RGB JPEG, no text layer, A DIFFERENT
%      PHYSICAL COPY.
%
%  WHY BOTH. The working scan is JBIG2, a symbol-substitution codec: it
%      can replace a glyph image with a visually similar one from its
%      dictionary, silently, and zooming in only renders the substitution
%      larger. So a doubtful SINGLE GLYPH is never settled on the working
%      scan — it is settled on the KB copy, which is continuous-tone and
%      cannot substitute a symbol. The KB copy also answers the second
%      question: where both copies agree, a defect is the compositor's
%      and not this copy's. crop.sh defaults to the KB scan for this
%      reason.
%
%  PAGE MAPS — see pagemap.py. Both are UNIFORM; no leaf is duplicated
%      or dropped anywhere in either scan.
%          printed 1--549  =  working PDF + 21    (PDF  22--570)
%          printed 1--549  =  KB PDF      + 19    (PDF  20--568)
%      Working scan verified by reading the folio out of the text layer
%      of all 689 PDF pages: 539 of the 549 body pages yield a folio and
%      every one gives +21, a contiguous run from PDF 22 to PDF 570 with
%      no sustained exception; the ten silent pages are bridged by their
%      neighbours. KB scan verified by CONTENT rather than folio digits —
%      each sampled leaf OCR'd and its words matched against the working
%      scan's text for the printed page the offset predicts — at 23
%      samples from printed p. 2 to p. 549, all matching at +19 with a
%      word overlap of 0.95--1.00, both endpoints included.
%
%  THREE MISPRINTED FOLIOS, in the type and not in the digitisation:
%          printed 357  prints as  "257"
%          printed 394  prints as  "294"
%          printed 395  prints as  "295"
%      The neighbours are correct on both sides (355, 356, [257], 358,
%      359 and 393, [294], [295], 396, 397) and the KB copy also prints
%      "295" on p. 395, so both copies carry it. The pages are marked
%      here by their true number and the misprint is logged at the site.
%
%  THE BOOK'S OWN ERRATA. A Trykfeil leaf follows p. 549 (working PDF
%      571), unfoliated. It is transcribed in full at the end of this
%      file. Per the house rule the text is transcribed AS PRINTED and
%      Nielsen's corrections are NOT silently applied; where a batch
%      reaches a page the Trykfeil touches, the correction is noted in a
%      % comment at the site so a reader can see both.
%
%  INDHOLD/HEAD COLLATION. The contents pages and the printed chapter
%      heads disagree in one place, and the disagreement is recorded
%      rather than normalised, per the house rule:
%          Anden Deel, first chapter — the Indhold reads "Tredie
%          Kapitel", the head on p. 197 reads "Første Kapitel".
%          Checked on the image. The head is followed here.
%      All nine chapter heads were verified against their printed pages
%      (pp. 1, 62, 125, 197, 250, 318, 375, 445, 509), not taken from
%      the Indhold. Two of the nine do NOT open a page: Tredie Kapitel
%      of Første Deel stands part-way down p. 125, and Tredie Kapitel of
%      Tredie Deel part-way down p. 509. The custom heading macros below
%      therefore never force a page break — using \chapter would have
%      invented one at both.
%      One further variant is NOT yet adjudicated: the Indhold reads
%      "Det elektro-chemiske System" where the run-in head on p. 300
%      appears to read "Det elektrochemiske System". Check on the image
%      when that batch is done and record both readings if they differ.
%
%  MATHEMATICS. Unlike every other book in this repo, this one contains
%      real mathematics — Kepler's laws derived in a footnote on p. 160
%      with polar coordinates and integrals, and lighter algebra through
%      the mechanics chapters. Roughly 83 pages carry it. The OCR is NOT
%      a witness on those pages and the batch markers below say so.
%      Set formulas in LaTeX math mode preserving the printed notation,
%      including Nielsen's own spellings (cos, sin, tang, Radius vector)
%      and his variables; do not modernise, do not tidy his notation.
%
%  OCR. This is the first book in this repository whose OCR genuinely
%      carries the words: it is antiqua, the scan is clean, and
%      tesseract -l dan on the native 600 ppi leaf returns correct
%      Danish with the line breaks and the printed word-break hyphens
%      intact. ocr.sh therefore has an EMPTY sed table, deliberately —
%      seeding it from a Fraktur book's harness would rewrite text that
%      is already right. One mechanical fix IS safe and is applied:
%      Å/å do not exist in this orthography (the letter entered Danish
%      only with the 1948 reform; this book writes aa throughout), so
%      any Å the OCR produces — it likes "Ånden Deel" for "Anden Deel"
%      and "Å." for "A." — is certainly an error.
%
%  CONVENTIONS. 1873 orthography exactly as printed: aa not å,
%      capitalised nouns, Christendommen, Videnskaben, høiere. Danish
%      quotation marks „…“ low-high. Em-dash ---. Letterspaced emphasis
%      -> \emph{}; run spacing.py, which is calibrated for this scan.
%      Latin and some proper names are set in italic against the
%      surrounding roman -> \textit{}, and each occurrence follows its
%      own page. The a)/b)/c) division heads are RUN-IN — printed on
%      the same line as the paragraph they open -> \runin{}.
%      Printer's defects are transcribed as printed and logged in a %
%      comment at the site; the quote balance is expected to drift off
%      zero as a result and the running total belongs in RESUME-NOTES.
%      Word-breaks: the compositor's line-break hyphen must be resolved
%      when the text reflows — write  Cau\-%  (joints.py enforces this).
%      The printed word-division mark is a single raised hyphen throughout
%      — settled in batch 1; see WORD-DIVISION MARK below.
%
%  QUOTATION HABIT OF THIS BOOK — found in batch 1, and expected to recur.
%      Where a quotation is broken by an interpolation („hedder det“,
%      „hedder det i Thesis“), the compositor RE-OPENS with „ afterwards and
%      still closes only once, at the end. Seen on pp. 10, 11 and 16, in both
%      copies. It is the house habit, not a slip, and it is transcribed as
%      printed — so each instance pushes the running „/“ balance up by one.
%      DO NOT "fix" these. The balance is expected to climb through the book;
%      every unbalanced page must carry a % comment at the site, and a
%      quotation that merely runs across a page boundary (p. 9 into p. 10) is
%      not a defect and gets none.
%
%  WORD-DIVISION MARK — settled. A single, slightly raised hyphen, uniform
%      through the pages examined so far, verified on the KB copy. tesseract
%      renders it as "=" fairly often; that is an OCR artefact and not the
%      type. (The "Andet=" noticed on p. 549 while surveying is the same
%      artefact, so the doubt recorded above is resolved: the book does not
%      vary its mark.)
%
%  FOOTNOTES. *) is confirmed as the book's mark. p. 12 carries TWO notes,
%      printed *) and **), handled there with a local \renewcommand around
%      the second. If further two-note pages turn up, follow that pattern.
%
%  WHERE THE RULES ACTUALLY FALL, now that enough openings have been seen:
%      * between a Deel head and the Kapitel head below it — always (see below);
%      * before an A./B./C. division head that opens PART-WAY DOWN a leaf;
%      * at the close of a chapter, and at the close of Første Deel, where it
%        is printed DOUBLE (pp. 61, 125, 196, 249);
%      * NOT above a chapter head that opens its own page — checked at p. 250,
%        where the head sits directly under the folio. The p. 197 case is not a
%        counter-example: what precedes the rule there is the Deel head.
%
%  A RULE STANDS BETWEEN EVERY DEEL HEAD AND THE KAPITEL HEAD BELOW IT.
%      The skeleton originally set \deel and \kapitel adjacent, which dropped
%      it. Verified on the KB copy at p. 1 (Første Deel) and p. 197 (Anden
%      Deel) and restored at all three Deel openings, p. 375 included. Noticed
%      because the p. 197 agent saw the rule, judged it to fall above its own
%      marker, and said so rather than silently ignoring it.
%
%  PRINTED RULES. A short centred rule stands between some divisions, and its
%      length varies between occurrences. That variation is not reproduced:
%      \streg serves throughout (see the macro's comment). Note each occurrence
%      in a % comment at the site.
%
%  GREEK. The book has both monotonic and polytonic Greek (ἅπτεσθαι, p. 26).
%      textalpha handles both — twenty-odd books in this repository already
%      carry polytonic Greek under textalpha with their PDFs built, ἅ among
%      them — so no extra package is needed. The sandbox compile test cannot
%      show this, because stripping textalpha is what that test does; map
%      Greek to a placeholder to test the rest of the file, and never delete
%      Greek from the source.
%
%  MATHEMATICAL NOTATION — precedent set in batch 4 (pp. 39--47), to be
%      followed by every later mathematical batch:
%        * Nielsen writes "tg", not "tang" or "tan" -> \operatorname{tg}.
%        * The period between "tg α" and "x" is HIS multiplication mark and
%          is kept.
%        * He prints ÷ where a modern text prints a minus sign (p. 46, second
%          equilibrium equation, against + in the first). ÷ was the ordinary
%          Danish minus of the period, so this is his notation, not a
%          compositor's fault; the Trykfeil leaf does not touch it. Set \div
%          and note it at the site.
%        * Printed variables are upright roman. This is NOT reproduced:
%          forcing \mathrm throughout would clutter past usefulness, so the
%          default math italic stands. Isolated variables in running prose
%          (x, y, s, t, g) stay as plain text, as printed.
%        * spacing.py flags algebra as letterspacing. Ignore it on formula
%          lines entirely; it is only a witness on running prose.
%        * ENUMERATIONS KEEP ONE FACE. Where a member of an enumeration carries
%          a prime or an index, the whole enumeration is set in math, so that
%          m, m', m'' does not change face part-way along the list. Set in
%          batch 5 (the Lagrange footnote, pp. 59--61) and logged at the site.
%        * Minus is a true minus in most formulas; ÷ is Nielsen's minus only
%          where he actually prints it (p. 46 so far). Do not normalise either
%          way -- read the sign off the page.
%
%  THE WORKING COPY HAS BEEN CORRECTED IN PENCIL BY A PREVIOUS OWNER, AND THE
%  CORRECTIONS ARE NIELSEN'S OWN ERRATA. This is the most consequential fact
%      about the working scan, and it was established at p. 160 by putting the
%      two copies side by side.
%      The Trykfeil leaf directs that p. 160's numerator "c/μ" be read "c²/μ".
%      The KB copy prints a plain c. The working copy carries a hand-drawn
%      superscript 2 beside it — lighter than the type, flat-topped, sitting
%      high and to the right. Someone worked through this copy applying the
%      author's errata by hand.
%      The same hand accounts for every earlier anomaly at an errata site:
%        p. 113 — a full stop "already there" where the leaf says to add one;
%        p. 163 — the two added commas and the "o. s. v." after "= 0;";
%        p.  79 — what an agent reported as an "ink speck mimicking a prime"
%                 on Q_k is, on this reading, the pencil prime the leaf asks
%                 for.
%      CONSEQUENCE, and it is a large one: at ANY of the leaf's 32 errata
%      sites the working scan is liable to show the CORRECTED reading, which
%      would silently defeat this edition's whole policy — transcribe as
%      printed, note the correction. **Every errata site must be read on the
%      KB copy**, without exception. The pages are listed in pagemap.py.
%      The pencil is not confined to errata sites either: it underlines
%      passages (pp. 35, 37, 63, 64, 66, 142, 152), brackets notes (pp. 144,
%      145, 153) and marks margins. None of it is text.
%      AND IT IS NOT ALWAYS A SMALL MARK. At p. 224 the hand strikes the
%      printed word "mechaniske" through and writes "chemiske" IN THE MARGIN —
%      and BOTH OCR witnesses transcribe the marginal word as if it were text.
%      At p. 233 it strikes the division letter B through and writes C beside
%      it. So the hazard is not only a forged accent or stop: whole words in
%      the margin can enter the OCR. Read the KB copy.
%      TALLY so far: of the errata sites reached, the pencil had got to
%      pp. 79, 113, 160, 163 (three entries), 176, 178 (two), 180, 183, 184,
%      206, 224, 233, 256 and 257; it had NOT got to pp. 164 or 169, where the
%      fault is genuinely in the type. So it is thorough but not exhaustive,
%      and each site must be checked rather than assumed either way.
%
%  EMPHASIS IS NEVER SETTLED ON THE WORKING SCAN. This is the single most
%      important methodological point in this file, and it was established by
%      putting the two copies of p. 64 side by side.
%      The working scan is 1-bit JBIG2. Thresholding FATTENS every stroke, so
%      ordinary roman on that scan reads as bold — the whole page looks
%      emphatic. On the KB colour copy the same passage is plain light roman.
%      Two consequences:
%        1. A word that "looks emphasised" on the working scan usually is not.
%           Confirm EVERY emphasis call on the KB copy before marking it —
%           INCLUDING A RUN. An earlier version of this note said a run of
%           adjacent flagged words was "nearly always real emphasis". That is
%           NOT safe here: the previous reader's pencil underline runs beneath
%           whole clauses, merges with the descenders when the scan is
%           thresholded, and inflates the measured widths of every word above
%           it. On p. 142 it produced three consecutive spacing.py RUNs across
%           a four-line passage that the KB copy shows as plain roman
%           throughout. A run raises the question; only KB answers it.
%           Agents have independently rejected calls this way on pp. 24, 71,
%           76, 79, 82, 85 — all justification or thresholding artefacts.
%        2. spacing.py measures advance widths on that same fattened scan, so
%           its false-positive rate is structural, not incidental. A RUN it
%           reports is still good evidence; a single flagged word is not, and
%           loose justification mimics letterspacing. It also MISSES real
%           emphasis — inside quotation marks, and two-letter words.
%      Separately, the working copy carries a previous reader's PENCIL
%      underlining and marginal strokes (pp. 35, 37, 63, 64, 66 so far), which
%      tesseract transcribes and which can read as emphasis. The KB copy, being
%      a different physical copy, does not have them. Another reason the
%      emphasis question goes to KB.
%
%  FOOTNOTES THAT SPAN LEAVES. Several notes run across a page break, and one
%      (the Lagrange interpolation note) begins on p. 59, fills p. 60 beneath
%      two lines of body text, and ends on p. 61. The book does NOT repeat the
%      mark on the continuation. A reflowed transcription cannot preserve that
%      layout, so the whole note is placed at its mark on the opening page,
%      with a % comment recording which leaves it covered in the print.
%
%  HOW THE HEADS ARE SET — a pattern, now that eleven batches have seen them.
%      The a) b) c) run-in heads are printed HALVFED (semi-bold); the Greek
%      α) β) γ) sub-heads are LETTERSPACED. Confirmed at pp. 79, 101 and 107
%      for the first, and throughout for the second. This edition does not
%      reproduce the distinction — \runin and \subrunin both render as \emph —
%      but it is logged at the sites, so a later editor can restore it from the
%      macros alone rather than re-reading every page.
%
%  ± OCCURS, AND SO DOES ÷ ON THE SAME PAGE. p. 92 prints ± (a plus over a
%      bar, no dots) in the very formula that p. 91 sets with ÷, and p. 106
%      prints both. Both copies agree in every case. Set as printed --- do not
%      harmonise Nielsen's signs, in either direction.
%
%  FOOTNOTE APPARATUS. Two or three notes on one leaf are common (pp. 12, 118,
%      134, 136, 137, 145, 146), set *) **) ***) and handled with a local
%      \renewcommand{\thefootnote} around each non-default mark, restored
%      after. p. 145 carries a genuine fault in the apparatus itself: the mark
%      in the body reads **) while its note at the foot reads ***). That is
%      reproduced with \footnotemark/\footnotetext rather than silently
%      regularised.
%
%  ITALIC AND LETTERSPACED TOGETHER. „Thema Coeli“ (p. 142) is set in italic
%      AND letterspaced. Use \textit and log the spacing in a % comment: \emph
%      inside an italic context flips the text back upright, which would lose
%      the italic the page actually prints.
%
%  A TYPE-SIZE CHANGE, p. 94. The first six lines of that leaf are set in
%      SMALLER type; the change falls exactly at the leaf break, mid-sentence,
%      and the small type runs to the end of that paragraph. Both copies agree.
%      Set with {\small ...} and logged. Watch for more.
%
%  DIVISION BOUNDARIES CAN FALL MID-PAGE, AND ONE DOES SO ACROSS A BATCH.
%      Division A of Andet Kapitel does not end with p. 86: it runs two more
%      paragraphs onto p. 87, where a centred rule closes it and the
%      B. Mechanisk Realitet head follows. So the batch beginning at p. 87
%      opens with the TAIL OF A, not with B, and that agent writes both the
%      rule and the \afdeling{B.} head. Checked on the page. Expect the same
%      shape at other division boundaries -- the Indhold's spans overlap by a
%      page precisely because divisions open part-way down a leaf.
%
%  STATE. Batches 1--23 (pp. 1--258) and the Trykfeil leaf transcribed and
%      verified. FØRSTE DEEL COMPLETE (pp. 1--196). Anden Deel: Første Kapitel
%      complete (pp. 197--249), Andet Kapitel begun. Nearly half the book.
%      25 batches outstanding, plus the Indhold and the Indledning. See RESUME-NOTES.md for the production log and
%      BATCH-AGENT.md for the batch brief.
% ============================================================

\documentclass[12pt,a4paper]{book}
\usepackage[utf8]{inputenc}
\usepackage[T1]{fontenc}
\usepackage{libertinus}
\usepackage{libertinust1math}
\usepackage{amsmath}     % this book has real mathematics; see the header
\usepackage{textalpha}   % Greek typed directly (θ, ω, π occur in the physics)
\usepackage{graphicx}    % for \reflectbox (the old Danish »ɔ:« id-est mark)
\DeclareUnicodeCharacter{0254}{\reflectbox{c}}  % ɔ (U+0254) = reversed c
\usepackage[danish]{babel}
\usepackage[protrusion=true,expansion=true]{microtype}
\usepackage{setspace}
\usepackage[margin=1.2in]{geometry}
\usepackage{enumitem}
\usepackage{fancyhdr}
\setlength{\headheight}{28pt}
\addtolength{\topmargin}{-13.5pt}

% Footnotes are marked *) in this book (p. 160 and the other mathematical
% footnotes). \fnsymbol is NOT used: it is a math-mode command and gives
% U+2217. footmisc[perpage] is NOT used either: it resets on the TYPESET page,
% not the printed one, so a second note on a reflowed page silently becomes a
% dagger. If a page ever turns out to carry two notes, deal with it there.
\renewcommand{\thefootnote}{*)}

% The book has NO running heads: the folio stands alone at the top of each
% page. Reproduced.
\pagestyle{fancy}
\fancyhf{}
\fancyhead[C]{\thepage}
\renewcommand{\headrulewidth}{0pt}

% ---------------------------------------------------------------------------
% Heading macros. None of them forces a page break, because two of the nine
% chapter heads do not open a page (pp. 125 and 509) — see the header.
% ---------------------------------------------------------------------------
\newcommand{\deel}[2]{%
  \par\vspace{2.2\baselineskip}%
  {\centering\large #1\par\vspace{0.3\baselineskip}\Large #2\par}%
  \addcontentsline{toc}{part}{#1 #2}%
  \vspace{1.4\baselineskip}\par\nointerlineskip}
\newcommand{\kapitel}[2]{%
  \par\vspace{1.8\baselineskip}%
  {\centering #1\par\vspace{0.25\baselineskip}\large #2\par}%
  \addcontentsline{toc}{chapter}{#1 #2}%
  \vspace{1.1\baselineskip}\par\nointerlineskip}
\newcommand{\afdeling}[2]{%
  \par\vspace{1.3\baselineskip}%
  {\centering #1\par\vspace{0.2\baselineskip}\normalsize #2\par}%
  \addcontentsline{toc}{section}{#1 #2}%
  \vspace{0.8\baselineskip}\par\nointerlineskip}
% A short centred rule is printed between some divisions. Its printed LENGTH
% varies from one occurrence to the next, and that variation is NOT reproduced
% typographically -- one rule serves throughout, following the decision taken
% in brochner/det-religioese. Two batch agents independently invented
% 0.12\textwidth and 0.20\textwidth before this macro existed; both were
% normalised. Where a rule occurs, note it in a % comment at the site.
\newcommand{\streg}{\begin{center}\rule{2cm}{0.4pt}\end{center}}

% Run-in head: a) b) c) are printed on the same line as the paragraph they
% open, letterspaced. \runin{a) Det leddeelte Hele.} at the head of the
% paragraph, NOT on a line of its own.
\newcommand{\runin}[1]{\emph{#1}\hspace{0.4em}}
% Greek-letter sub-heads α) β) γ) are run-in likewise.
\newcommand{\subrunin}[1]{\emph{#1}\hspace{0.4em}}

\setstretch{1.15}

\title{Natur og Aand\\[0.4em]
  \large Bidrag til en med Physiken stemmende Naturphilosophie}
\author{R.\ Nielsen}
\date{Kjøbenhavn: Forlagt af J.\ H.\ Schubothes Boghandel, 1873}

\begin{document}
\maketitle
\tableofcontents


% ============================================================
%  FRONT MATTER
% ============================================================
%
% The Indhold (printed V--X, working PDF 12--18) is transcribed LAST, from
% the image, once the body is done and every chapter and division head has
% been read off its own page. It is not reconstructed from the OCR: the
% contents-page OCR of this scan already carries known damage — it renders
% the Greek sub-heads α β γ as "a", "B"/"ß"/"P", "y"/"7", and it has dropped
% or garbled several leaders — and a chapter argument taken from it would
% import those errors into the edition. The one collation already made from
% it is recorded in the header.

% [text to be added: Indhold, printed pp. V--X]

% The Indledning is printed XI--XIV. Its first leaf, printed XI, carries no
% folio in either scan; XII--XIV do. The Indhold lists it as "XI--XIV".

% [text to be added: Indledning, printed pp. XI--XIV]


% ======================================================================
%  Første Deel — Første Kapitel: Grundlagsbegreber   (printed pp. 1–61)
% ======================================================================

\deel{Første Deel.}{Den mechaniske Natur.}
% Mellem Deel-Overskriften og Kapitel-Overskriften staaer i Satsen en kort,
% centreret Streg. Efterprøvet paa KB-Exemplaret for S. 1 og S. 197.
\streg

\kapitel{Første Kapitel.}{Grundlagsbegreber.}

%  --- A. Rum og Tid  (pp. 1–18) ---

% --- p. 1 ---
\afdeling{A.}{Rum og Tid.}

Da alle udvortes Ting ere til i Rummet og umulig kunde være til,
dersom Rummet ikke var; da Tingene opstaae og forgaae i Tiden og
umulig kunde enten opstaae eller forgaae, hvis der ingen Tid var, saa
følger jo vistnok heraf, at Rum og Tid maae høre med til
Naturtingenes Væsen, at de, med andre Ord, maae betragtes som de
første almene, altomfattende Naturformer. Men uagtet Rum og Tid ere
almeengyldige Naturformer, skal det dog i det Følgende vise sig, at
disse Formers naturlige Bestaaen har en aandelig Grund; at de kun ere
Gjenstandsformer ved at være Selvhedsformer, kun objective ved at
objectiveres af den evige Aand.

\runin{a) Rummet.} De vigtigste Bidrag til dybere Forstaaelse af
Rummets Natur vil man i den nyere Philosophie
% --- p. 2 ---
fornemmelig søge hos Kant og Hegel; men de to Synsmaader maae
behandles med Kritik, thi, overfladisk seet, er den ene i Modsigelse
med den anden.

Kant stiller sig, som bekjendt, paa et ubetinget subjectivt Standpunkt
og betragter Rummet som en apriorisk Sandseform, en blot menneskelig
Bevidsthedsform. Rummet er intet fra Erfaringen laant Begreb; for at
kunne betragte Tingene i Rummet, maae vi selv medbringe
Rumsanskuelsen. Man kan umulig forestille sig, at der intet Rum er;
men man kan meget vel forestille sig, at der ingen Gjenstande ere i
Rummet. Naar der tales om mange Rum, er det ikke at forstaae, som var
der Tale om Dele, hvoraf et Hele var sammensat: de mange Dele flyde
sammen i et eneste med vor nødvendige Sandseanskuelse givet Rum. Heraf
uddrages saa den Slutning, at det ikke er Tingen i sig men kun
Phænomenet af Tingen, d. v. s. Tingen som Synsting, der opfattes i
Rummet; hvad ogsaa kan udtrykkes saaledes, at \emph{Rummet kun er en
Selvhedsform, men ingen i sig gyldig Andethedsform.}

I skarp afgjørende Modsætning til Kant staaer Hegel paa et absolut
objectivt Standpunkt. For ham er Rummet i sin „abstracte
Almindelighed“ den første Naturform, den umiddelbart nødvendige Form
for \emph{udvortes Væren}. For at tydeliggjøre, hvad der ligger i
Begrebet af Rummets „ideelle Vedsidenafhinanden“,
% Trykkefeil: det andet Anførselstegn aabnes med » (Gaaseøie) og
% lukkes med “; begge Exemplarer (Google og KB) sætte det saaledes.
»Værenudenforhinanden“, osv., vælger han at paavise, hvorledes Rummet
seer ud i vor Forestilling. For at bevise, at Rummet svarer til vor
Tanke, maae vi, siger han, sammenligne vore Rumsforestillinger med de
Bestemmelser, der ligge i vort Begreb. De mange \emph{Her} og
\emph{Der} i Rummet ere ved Siden af hverandre; det ene \emph{Her} er
fuldkomment det samme som det andet, og denne abstracte Fleerhed ---
uden sand Afbrydelse og Grændse --- er netop Udvorteshed. Rummet med
de utallige Punkter er lutter Discretion; men denne Discretion er igjen
flydende
% --- p. 3 ---
Eenhed, Eenhed i „Udenforhverandre“, fuldkommen Continuitet. Punktet
har kun Betydning, forsaavidt det indtager et Rum, altsaa er et
Udvortes mod sig selv og Andet; i Punktet selv er der da igjen et
Foroven, et Forneden, et Høire, et Venstre, thi et sidste \emph{Her}
gives ikke. Udenfor sin Grændse er Rummet bestandig „hos sig selv“:
det Continuerende er discret; det Discrete continuerende. Eenheden af
begge Momenter, Discretion og Continuitet, er „Rummets objectivt
bestemte Begreb“. Det er da ifølge denne Synsmaade saa langt fra, at
det objective Rum skulde beroe paa nogensomhelst subjectiv Anskuelse,
at det tvertimod som „en usandselig Sandselighed“, hvilende paa sit
eget Begreb, afgiver Naturtingenes Grundlag: \emph{Rummet er
Andethedsform uden at være Selvhedsform.}

Ved at holde de to Opfattelser skarpt mod hinanden vil man let
overbevise sig om, at de begge ere eensidige.

At Kant med Rette gjør Rumsformen til en menneskelig
Bevidsthedsform, kan udtrykkelig sees af Hegels Beviisførelse; thi
denne Beviisførelse støtter sig Punkt for Punkt til menneskelig
Anskuen, menneskelig Tænken og Viden. Vor Indsigt i det objective Rums
Gyldighed afhænger absolut af dets Gyldighed for os. Rummet bestemmes
som den „\emph{abstracte}“ Form for udvortes Væren; men det er os ---
ikke Naturen --- der abstrahere Formen. De mange forskjellige
„\emph{Her}“ og „\emph{Der}“ maae sættes og udslettes; men det er os
--- ikke Naturen --- der sætte dem og slette dem ud. Eenheden af
„Discretion“ og „Continuitet“ er at tænke i Forskjel, og Forskjellen
igjen i Eenhed; men det er os --- ikke Naturen --- der maae tænke denne
Forskjelseenhed: for at Rummet kan vise sig som objectivt Begreb, er
det os, der maae begribe.

Men at Hegel desuagtet har Ret, idet han gjør Rummet til et objectivt,
af os uafhængigt Begreb, bliver omvendt kjendeligt paa den Kantiske
Beviisførelse. Kant fastholder,
% --- p. 4 ---
at Phænomenerne, de sandselige Ting, umulig kunne opfattes uden
Rummet; hermed har han da tillige indrømmet, at Rummet er en for
Synstingene nødvendig, objectiv Form. Tingene i Naturen blive jo ikke
sandselige derved, at vi sandse dem; men vi sandse dem, fordi de i sig
ere sandselige. Saa lidt som det er vort Syn, der gjør Tingene synlige,
eller vor Tænkning, der gjør dem tænkelige, ligesaa lidt er det vor
Rumsanskuelse, der giver dem Tilværelse i Rummet. Dersom den
Skillevæg, Kant har reist mellem Tingenes Væsen i sig og deres
Phænomener, var mere end en kritisk Indbildning, da var det aldeles
utænkeligt, at dette ubekjendte, mod Rummet ufordragelige Væsen skulde
være istand til at afgive noget Erfaringsindhold og udenfra fylde den
tomme Anskuelse med virkelige Gjenstande.

Altsaa: hver Gang Rumsformen udelukkende gjøres til Andethedsform,
forvandler den sig og bliver en Selvhedsform; og naar den udelukkende
fastholdes som Selvhedsform, forvandler den sig igjen og bliver en Form
for Andet. Dette viser da tydelig nok, at Rummet umulig kunde være en
virkelig sand Naturform, dersom denne Naturform ikke fra Evighed af fik
objectiv Væren i guddommelig Anskuen. Det Kantiske Feilsyn bestaaer
ikke deri, at han gjør Rummet til en Anskuelsesform, men deri, at han
vilkaarlig indskrænker Rumsformen til at være en blot menneskelig
Selvhedsform, og netop derved kommer til at benegte dens objective
Gyldighed. Hegels Feil er ikke den, at han hævder Rummets
Objectivitet; men Grundfeilen er, at han i sin Naturphilosophie ligesom
i sin Logik bestandig overseer, at intet Objectivt er til, uden det,
som objectiveres, og at ingen Objectivering er mulig uden et
objectiverende Selv. At den menneskelige Bevidsthed tilkjender Rummet
i Naturen udvortes Væren, beroer i Grunden paa den stiltiende
Forudsætning, at det Verdensrum, vi i vor Naturbetragtning antage for
objectivt, fra Begyndelsen af maa være --- objecti\-%
% --- p. 5 ---
veret af den skabende Aand. Stjernehimlen gjør et symbolsk Indtryk;
det er den menneskelige Anskuen, der under sin Selvudviden i Rummets
Uendelighed mødes med den guddommelige som Øie med Øie.

Hvorledes det, hvad Rumsobjectiveringen angaaer, nærmere forholder sig
med Vexelbestemmelsen af Selvhed og Andethed, kan sees ved en
Betragtning af Rummets Dimensioner. Længde, Brede og Høide ere i
Rummet forskjellige og dog ikke forskjellige. Der maa nødvendig gjøres
Forskjel mellem de ovenanførte Hovedretninger; thi uden en saadan
Forskjel vilde der ingen Bestemmelser være i Anskuelsen, og da heller
intet objectivt bestemt Rum. Men hvorpaa beroer Forskjellen? Længde er
Høide, og Høide er igjen Brede: den ene Dimension er, hvad den anden
er, og kan ombyttes med den anden. De sættes samtidig og udslettes
samtidig; de forudsætte hinanden i det Uendelige. Heraf sees, at de tre
Dimensioner i deres udvortes Gjensidighed ere dobbeltsidige
Bestemmelser, Selvhedsbestemmelser og Bestemmelser for Andet. Som
Selvhedsbestemmelser ere de vilkaarlige, som Andethedsbestemmelser
tilfældige, som objectiverede, altsaa objective
Gjensidighedsbestemmelser indbyrdes nødvendige. Vælge vi f. Ex. et
retvinklet Coordinatsystem med Axerne $X$, $Y$ og $Z$, er
Objectiveringen let at belyse. De tre Axer maae have et
Skjæringspunkt. Hvor er nu dette Skjæringspunkt i Rummet selv? Overalt
og ingensteds. Punktet er kun objectivt ved at objectiveres. For at
objectiveres, maa det sættes; for at sættes, maa det sættes vilkaarlig,
altsaa ved en Selvbestemmelse. Vilkaarligheden er nødvendig; thi
hvilket Sted, vi end vælge, saa kunde vi dog lige saa godt have valgt
et andet Sted. Ved at være sat er Punktet blevet objectivt; det har
faaet en forestillet Tilværelse og afgiver nu en Bestemmelse for Andet,
nemlig for Axerne. I denne objective Bestemthed er Punktet at tage som
givet; nu er det netop her, ikke der.
% --- p. 6 ---
Dets Tilværelse er tilfældig, thi Axerne kunde lige saa godt have
skaaret hverandre i et hvilketsomhelst andet Punkt. Men Tilfældigheden
er nødvendig, thi hvilket Punkt der end var givet, kunde der jo med
lige Grund være givet et andet. Hvad der gjælder om Skjæringspunktet,
gjælder deelvis om Axerne. Det er vilkaarligt at gjøre $Z$ til
Høide-, $X$ til Længde- og $Y$ til Brede-Axe; men har Vilkaarligheden
en Gang givet dem denne Betydning, saa have Høide, Længde og Brede som
Bestemmelser for hverandre indbyrdes, altsaa som
Andethedsbestemmelser, ogsaa derved erholdt en objectiv, tilfældig
Tilværen. Tilfældigheden i denne Tilværen er lige saa uundgaaelig
d. v. s. lige saa nødvendig som Vilkaarligheden. Det Vilkaarlige og
det Tilfældige gjøre da heller ikke Nødvendigheden Afbræk; tvertimod,
disse Grundlagsbestemmelser kunne med Rette siges at afgive
Forudsætninger, uundværlige Betingelser for en exact Objectivering af
den betingede Nødvendighed. Antages et Punkt $M$, hvis Coordinater i
det givne System ere $x$, $y$ og $z$, og hvis Afstand fra
Skjæringspunktet kan betegnes ved $r$, da vil en Ligning som:
\[
  r^2 = x^2 + y^2 + z^2
\]
uvilkaarlig minde om den Mangfoldighed af Rumsbestemmelser, som i
deres exacte Nødvendighed alle ere betingede, subjectivt af
vilkaarlige, objectivt af tilfældige Forudsætninger.

De tre samværende Axer ere i deres udvortes, umiddelbare Tilfældighed
uden indbyrdes Relationer; først med Viden komme Relationerne, og med
Relationerne Nødvendigheden. Her er, naar vi tage al Viden bort, ikke
engang tre Axer, thi Trehed fordrer, at en Viden af de to forener sig
med en Viden af den tredie. Forestillingen om en abstract objectiv
Trehed hidrører fra en Samviden, som den objectiverende Bevidsthed,
uden at tænke derpaa, selv underskyder. Det er uimodsigelig vist, at
den anførte, til Axerne
% --- p. 7 ---
svarende Ligning har objectiv Gyldighed, naar den vel at mærke,
objectiveres i Viden; men naar der ingen Viden er, saa er Ligningen
meningsløs. \emph{Al objectiv Relation er subjectivt begrundet i
Reflexion}; det er en Objectiveringens Lov, saa streng som nogen
Naturlov.

I den Schellingske Naturphilosophie tales der mystisk dybsindig om en
Naturen iboende Anskuen, en objectiv Naturens Phantasie, en ubevidst
Naturviden o. desl. Hvor grundløse slige Talemaader ere, viser sig
hver Gang Blikket fæstes paa noget Bestemt. Skulde Rummets tre
Dimensioner være satte af en ubevidst Viden, da maatte der i denne
Viden jo være en Samviden af dem alle tre, en Viden, hvori hver enkelt
af dem blev bestemt med Hensyn paa de to andre. Den ubevidste Viden
skulde da være \emph{en Viden af Hensyn}; her er Modsigelsen. En Viden
af Hensyn er en de enkelte Bestemmelser oversvævende Viden, d. e. en
Selvviden. For at $X$ kan bestemmes med Hensyn paa $Y$ og $Y$ med
Hensyn paa $X$, maa den bestemmende Viden nødvendig oversvæve baade
$X$ og $Y$. En saadan Oversvæven er kun mulig, naar det samme Selv,
der er vidende om $X$, ogsaa er vidende om $Y$. Men naar Viden af
Andet, ifølge Sagens Natur, er uadskillelig fra Selvviden, saa ere jo
saadanne Dybsindigheder som: „en ubevidst Naturviden“, „et sig
anskuende Rum“ i Lighed med „en sig tænkende Luft“, „en sig begribende
Materie“, at henføre, ikke til „den strenge Videnskab“, men til de
utopiske Speculationer.

At bestemme Rummet er at sætte Grændser i Rummet. Her viser sig da en
qvalitativ Forskjel mellem Grændsen, det Negative, og det begrændsede
Rum, det Positive. Det Negative er for Tanken, det Positive for
Anskuelsen. Det første absolut Negative er Punktet. Ved Bestemmelsen
af Punktet opkommer en eiendommelig Strid mellem Anskuen og Tænken.
Tanken fordrer, at Punktet skal være uden Udstrækning; Anskuelsen
fordrer, at det skal have en, om
% --- p. 8 ---
end nok saa lille, Udstrækning. Hvis Tanken ikke skeer Fyldest, bliver
Punktet en Flade, intet Punkt; hvis Anskuelsen ikke tilfredsstilles,
kan Punktet ikke indtage noget Sted: et reent Tankepunkt er ikke noget
Punkt i Rummet. Modsigelsens Løsning er en Reflexion, hvori Anskuelse
og Tanke vexelviis ophæve og bekræfte hinanden. Forsaavidt Tanken
ophæver Anskuelsen, bliver Rumselementet, det Positive, fortæret af det
Negative og sat som Punkt. Forsaavidt Anskuelsen ophæver Tanken,
bliver det Negative omsat i det Positive, og Punktet sat som Element.
For Tanken er Punktet Liniens Grændse, men for Anskuelsen er det
Liniens Element; for Tanken er Linien Fladens Grændse, men for
Anskuelsen Fladens Element; saaledes er da ogsaa Fladen igjen baade
Legemets Grændse og dets Element. Liniens Element kan integreres til
Linie, Fladens til Flade, Legemets til Legeme; men Grændsen selv, det
absolut Negative, kan ikke integreres.

Som det Negative, det Ikke-Værende, kan Grændsen altsaa ikke sætte
Rummets Væren nogen Skranke; der er Rum paa begge Sider af Grændsen;
Grændsen er selv den negative Eenhed, hvori de Begrændsede enes og
skilles. Overalt hvor Tanken sætter en Grændse i Rummet, gaaer
Anskuelsen ud over Grændsen og opdager et nyt Rum. Rummet er fuldt af
Grændser, thi Grændsen er allevegne, og dog grændseløst, thi Grændsen
overskrides allevegne. At ombytte en Grændse, som overskrides, med en
ny Grændse, som igjen overskrides, er Hemmeligheden: det uendelige Rum
er kun objectivt for en uendelig Objectivering.

\runin{b) Tiden.} Rummets uendelige Udstrækning er begrundet i
Grændsens uendelige Vorden. Denne Vorden, der sætter Forskjel mellem
den svundne, den værende og den kommende Grændse, viser hen til en
anden fra Rummet forskjellig Naturform: den i Udstrækningen skjulte
Succession
% --- p. 9 ---
tilhører Tiden. Ligegyldigt mod sin Grændse, er Rummet ligegyldigt mod
sin Vorden; det Rum, som kommer, naar Grændsen overskrides, var der
allerede forud, og det, som har været, er der endnu; Forskjellen er kun
en Skygge af Tiden. I Tiden selv er Forskjellen mellem det, som har
været, det som er, og det, som vil komme, derimod en væsenlig, om end
endnu udvortes Forskjel.

At Tidsformen oprindelig er en Selvhedsform, kan sees af den Reflexion,
hvori Anskuelsen opløses. Tiden har kun een Dimension,
\emph{Længden}; men denne er i Tredelingen af Fortid, Nutid og Fremtid
saa brudt, at kun en gjennemgaaende Tænkning er istand til at fastholde
Eenheden. Fortiden er til i Erindring, men Erindring er en
Selvhedsact; Fremtiden er til i Forventning, men Forventning er en
Selvhedsact. Den Tidsdeel, der skulde kunne siges at have en
umiddelbar, af erindrende og forventende Subjectivitet uafhængig Væren,
skulde da være Nutiden. Men Nutiden er, naar vi see nærmere til, lige
saa reflecteret som Fortid og Fremtid. Opfattes Nutiden som det
absolute Øieblik, saa opfattes den ikke som Tid, men som Tidsgrændse,
den Grændse, hvori Fortid og Fremtid berøre hinanden. Denne Grændse,
det rene Nu, er et Ikke-Værende, et mathematisk Nulpunkt. Opfattes
Nutiden som relativt Øieblik, som Tidselement, da er den Vedvaren,
Anskuelsen underskyder, i uophørlig Flugt fra Forbigangenhed til
Tilkommelse d. e. fra Ikke-Væren til Ikke-Væren. Vel kan et
Tidselement opfattes constant, $dt = dt$, naar det anvendes som Maal
for et Rumselement; men Sætningen: Nu er Nu, bliver en modsigende Dom,
thi Subjectets Nu er forbi, naar vi komme til Prædicatets. Fastholdes
Nutiden som endelig, værende Tid, da er det alene Subjectivitetens
Erindring og Fremsyn, som i reflecteret Eenhed formaae at lade
Anskuelsen strække sig hen ad en vilkaarlig nærværende Tidslængde.

Hegels Dialektik af hvad han kalder „Tidens Dimen\-%
% --- p. 10 ---
sioner“, er fra den objective Side fiin og skarp. Det er Vorden, hvis
abstracte Eenhedsmomenter, hvert for sig, ere satte som det Hele, men
under modsatte Bestemmelser. De to Bestemmelser, Fortid og Fremtid,
ere saaledes hver især en Eenhed af Væren og Intet; men de ere ogsaa
forskjellige. Denne Forskjellighed kan kun være den mellem Opstaaen og
Forgaaen. „Den ene Gang, i Forgangenheden (Hades), er Væren
Grundlaget, hvorfra der begyndes; Forbigangenheden har virkelig været
som Verdenshistorie, Naturbegivenheder; men sat under Bestemmelsen af
Ikke-Væren, der træder til. Den anden Gang er det omvendt; i den
tilkommende Tid er Ikke-Væren den første Bestemmelse, Væren den senere,
om end ikke efter Tiden. Midten er begges indifferente Eenhed“ o. s.
fremdl. Det er fortræffeligt. Men i skjærende Modsigelse med disse
objectiverede Begrebsbestemmelser, Bestemmelser, der aabenbart skyldes
en selvbevidst, objectiverende Subjectivitet, staaer en Erklæring som
den, at Forskjellen mellem Objectiviteten og en mod samme subjectiv
Bevidsthed aldeles ikke vedkommer Tidens Begreb. „Tiden, hedder det,
% Anførselstegnet aabnes paany efter Indskuddet „hedder det“; kun eet
% lukkende “ staaer, ved Citatets Slutning. Saaledes trykt.
„er det samme Princip som den rene Selvbevidstheds Jeg $=$ Jeg: men
det samme eller det eenfolde Begreb endnu i sin hele Udvorteshed og
Abstraction som den anskuede blotte \emph{Vorden} --- den rene
I-sig-Væren, som en slet og ret Udenfor-sig-Kommen.“ --- Under slige
Vendinger bliver den klare Objectivering af Tiden, en Objectivering,
som kun er mulig for den vidende Aand, skudt tilside, og Fixideen om en
sig anskuende, tænkende Naturform, en sig selv begribende Tid omtaager
Bevidstheden.

Som Vordens Form har Tiden et Vordende til Forudsætning.
Vexelbestemmelsen af Tidens Vorden og det i Tiden Vordende kaster Lys
paa Forholdet mellem det Objective og det Subjective. Fra den
objective Side gaaer Tiden i Eet med det, som foregaaer i Tiden: det
vordende Indhold udlægger sit Væsen i Vordens Form og meddeler Formen
% --- p. 11 ---
sine Prædicater. Tiden vorder, Tiden iler, løber; Tiderne omskiftes
(tempora mutantur), der er gode Tider, onde Tider o. s. v. Fra den
subjective Side bliver Formen derimod fastholdt i abstract
Almindelighed for sig til Forskjel fra Indholdet. Alt Endeligt
begynder i Tiden og ender i Tiden, men Tiden selv har hverken
Begyndelse eller Ende. Sættes et første, absolut Begyndelsespunkt, da
opdager Tanken let den Modsigelse, at der altsaa maa have været en Tid,
da der endnu ingen Tid var. Sættes et sidste, absolut Endepunkt, maa
Følgen blive, at der tilsidst vil komme en Tid, da der ingen Tid er.
Gjør Hegel Uret i at fastholde den objective Tid uden Hensyn til en
objectiverende Selvviden, da gjør Kant ogsaa Uret i at lade det
Subjective og det Objective falde fra hinanden i to Stykker. Hvorledes
en saadan Udstykning maa blive en Snare, hvori den endelige Forstand
fanger sig selv, er paa det Klareste viist i „Antinomien“ af Verdens
Oprindelse.

% Anførselstegnet aabnes paany efter Indskuddet „hedder det i Thesis“;
% kun eet lukkende “ staaer, ved Citatets Slutning. Saaledes trykt i
% begge Exemplarer (Google og KB).
„\emph{Verden}, hedder det i Thesis, „\emph{har en Begyndelse i
Tiden}“; thi hvis man antog det Modsatte, da maatte der til ethvert
givet Tidspunkt være forløbet en heel Evighed, og en uendelig Række af
paa hinanden følgende Tilstande være afsluttet. \emph{Verden}, siger
Antithesis, \emph{har ingen Begyndelse i Tiden}; thi i modsat Fald
maatte der forud for Verdens Begyndelse være en tom Tid, men i en tom,
for alle Betingelser blottet Tid kan ingen Ting begynde.

Ifølge Thesis bestaaer en Rækkes Uendelighed deri, at den aldrig lader
sig „fuldende ved successiv Synthese“. Hvad betyder dette Aldrig? At
Fuldendelsen udsættes i en uendelig Tid. Sætningen kan altsaa
udtrykkes positivt saaledes: for at afslutte en uendelig Række ved
successiv Synthese fordres en uendelig Tid; eller: ved at fortsætte den
successive Synthese i en uendelig Tid danner Forestillingen en Række af
uendelig mange Led, en uendelig Række.
% --- p. 12 ---
% Denne Side bærer TO Fodnoter, mærkede *) og **).
Rækkedannelsen bliver den samme, enten vi begynde i Endeligheden og
nærme os Uendeligheden, eller vi ud fra Uendeligheden nærme os
Endeligheden.\footnote{% trykt Mærke: *)
I Rækken $\dfrac{1}{1-x} = 1 + x + x^2 + \ldots x^{n-1} + x^n +
\dfrac{x^n + 1}{1 - x}$ er $x^n$ det almindelige Led; sættes $n = o$,
% „n = o“ er trykt med Bogstavet o, ikke med Nullet; saaledes gjengivet.
faae vi $x^0 = 1$; sættes $n = \infty$, faae vi $x^\infty$; men Veien
fra 1 til $x^\infty$ og fra $x^\infty$ til 1 er lige lang: den er
uendelig.} Imellem et i Timeligheden givet Tidspunkt og en Begyndelse
i Evigheden ligger der da en uendelig Tid. Der er altsaa slet ingen
Modsigelse i, at en uendelig Række af Tilstande kan „til ethvert givet
Tidspunkt“ være afløbet. Skulde Beviset for Thesis have Noget at
betyde, da maatte det være beviist, enten: at en uendelig Række umulig
kan dannes i en uendelig Tid, eller: at der mellem en Begyndelse i
Evigheden og et timeligt Tidspunkt umulig kan ligge en uendelig Tid.
Men at føre Beviis for slige Umuligheder er umuligt.

Hvad Antithesis angaaer, maa Eftertrykket falde paa Begyndelsens
Relativitet. Enhver endelig Begyndelse forudsætter en foregaaende og
afgiver Forudsætning for en efterfølgende. Dette skulde nu ifølge
Kants Tankegang bevise Umuligheden af en \emph{første} Begyndelse. Men
nei. Lad de endelige Begyndelser kun svare til Leddene i en uendelig
Række; der er slet Intet i Veien for, at Rækken kan begynde med et
første endeligt Led, thi dette Led har jo netop samme almindelige
Forudsætning som alle de følgende: den uendelige Række er en successiv
Udfoldning af den Function, der indeslutter Uendelighedens Grund.%
\renewcommand{\thefootnote}{**)}%
\footnote{% trykt Mærke: **)
Den uendelige Række: $1 + x^2 + x^3 + \ldots$ har sin Grund og
% Trykfeil i Rækken: Leddet $x$ mangler mellem 1 og $x^2$; saaledes trykt.
Forudsætning i Functionen $F(x) = \dfrac{1}{1-x}$; den uendelige Række
$1 + \dfrac{x}{1} + \dfrac{x^2}{1.2} + \dfrac{x^3}{1.2.3} + \ldots$
har sin Grund og Forudsætning i Functionen $f(x) = e^x$; der er
Exempler nok i Mathematiken.}%
\renewcommand{\thefootnote}{*)}
Skulde Be\-%
% --- p. 13 ---
viset for Antithesis have nogen sand Gyldighed, da maatte det være
beviist, at en Begyndelse i Tiden umulig kan bestaae med en uendelig,
være sig fortidig eller fremtidig, Succession af endelige
Begyndelser. Men dette er, af gode Grunde, ikke blevet beviist.

At Striden mellem det Endelige og Uendelige, hvorpaa Antinomien er
bygget, væsenlig grunder sig paa et kritisk Brud mellem Bevidstheden
selv og den objective Væren, er let at indsee. Hvortil svarer vel
Endeligheden? Til den Objectivitet, som Sandsningen kan fatte og
Forestillingen fastholde. I denne Henseende er det da ganske rigtigt,
at en uendelig Tidsrække hverken lader sig sandse eller forestille.
Hvorledes komme vi da til Uendeligheden? Ad Tankens Vei. Vi gaae, som
det hedder i Mathematikens Sprog, til Grændsen; det vil sige: Tanken
løser Endeligheden op i dens yderste Grændse, sætter Negationen af det
Endelige. Det skeer ved et Spring. Skal der til den saaledes satte
Uendelighed nu ogsaa svare en egen Objectivitet, da er det vistnok
klart, at Objectiviteten i det Uendelige maa komme i Strid med
Objectiviteten i det Endelige. Men istedenfor at skyde Skylden paa
Tidsanskuelsen, skulle vi tvertimod erkjende, at det netop er Tiden,
der i naturlig Forstand stræber at forlige det Endelige med det
Uendelige. Tiden er Succession, og Successionen en udvortes Overgang,
der i lige Grad tilfredsstiller Forestilling og Tanke. Ved
Successionen følger Endeligt paa Endeligt; hvert endeligt Led bliver
successivt til for Forestillingen. Forestillingen kan da ikke forlange
Mere end, at Successionen fremdeles skal vedblive. Men Successionens
Vedbliven er netop Uendelighedens Tilbliven; den tilblivende
Uendelighed maa da ligeledes tilfredsstille Tanken. Objectivt seet, er
Successionen en Vorden; men naar Vorden i Tid er den udvortes Maade,
hvorpaa det Endelige uophørlig gaaer over i det Uendelige,
% --- p. 14 ---
og omvendt: saa have jo begge, det Endelige og det Uendelige i Tiden
som Tid, lige megen Objectivitet.

For at skærpe Antinomiens Braad har Kant slaaet Tid og Evighed sammen
i Eet og dog stillet dem op i den skarpeste Modsætning. Da Verden, for
den uendelige Tidsrækkes Skyld, ikke kan være fra Evighed af, maa den
have en Begyndelse i Tiden; det er Thesis. Da der ellers forud for en
Begyndelse i Tiden maatte ligge en uendelig, tom Tid (en tom Evighed),
maa Verden være fra Evighed af; det er Antithesis. Her alterneres
mellem Tid og Evighed, og saa er Tiden dog Evighed, og Evigheden Tid
--- i det Uendelige. Sagen er, at der til Grund for al Vorden maa
ligge en Væren, som ikke vorder. Denne i sig reflecterede Væren,
Væsenet, reflecteres nemlig i sit Andet, i Phænomenet. I Kraft af denne
dobbeltsidige Reflexion har Væsenet sin indre Uendelighed i Selvhedens
Form, der er Andethedsformernes ydre Uendelighed modsat. Den ydre
Uendelighed svarer til Tiden, den indre til Evigheden. Der er en
qvalitativ Forskjel mellem Tid og Evighed.

Paa denne Forskjel grunder sig igjen Forskjellen mellem
\emph{Begyndelse} og \emph{Oprindelse}. Begyndelsen skeer i Tiden, men
Oprindelsen udgaaer fra Evigheden. Det absolute Væsen er evigt,
uoprundet, ophøiet over Tiden; det relative Væsen har derimod sin
Oprindelse fra det Absolute og er ifølge sin Relativitet Tiden
underlagt. Oprindelse viser hen paa Selvvæsenets Objectivering af sit
væsenlige Andet, paa Aandens Objectivering af Naturen, Guds
Objectivering af Verden; Oprindelse er Overgang fra Evighed til Tid.
Men Begyndelse er en Tidskategorie; Begyndelsen forudsætter
Oprindelsen, thi den tilhører Tiden, den uendelige lige saa vel som den
endelige: den uendelige Tid er en Succession af lutter endelige
Begyndelser. Aanden er evig i Tiden over Tiden; men den timelige
Verden er ikke evig, om end dens Vorden er fra evig Tid.

% --- p. 15 ---
% I Overskriften staaer der i Satsen et løst Mellemrum inde i Ordet:
% „Mellemvære nde.“ Begge Exemplarer (Google og KB) vise det; her sat
% som eet Ord.
\runin{c) Rummets og Tidens Mellemværende.} Dersom Naturbegreberne:
Rum og Tid, i ubevidst Viden vare istand til at begribe sig selv, hvert
især, og hinanden indbyrdes; dersom de ved fælles begribende
Anstrengelse virkelig evnede at frembringe noget saa Haandgribeligt som
det, vi kalde Materie, da vilde den moderne Naturpantheisme unegtelig
kunne rose sig af et sikkert Fodfæste i „den strenge Videnskab“; men
Fixideen bliver kjendeligere med hvert Skridt, Begrebsudviklingen gjør.

Den Hegelske Deduction af Sted, Bevægelse og Materie er i denne
Henseende lærerig. Ordene ere dunkle, men Tanken er klar. Rummet med
sin ligegyldige Discretion og sin lige saa ligegyldige Continuitet er
en Modsigelse, der løser sig op i Tiden. Men Tiden for sig er ogsaa en
Modsigelse; thi det Nærværende er en ustandselig Omsætning af det
Forbigangne i det Tilkommende: de adskilte Momenter opløse sig i
umiddelbar Indifferens. Det concrete Punkt, hvori de fra begge Sider
sammenfaldende Modsigelser falde sammen, er „\emph{Stedet}“. Stedet er
Eenhed af „\emph{Her}“ og „\emph{Nu}“ og derved udtrykkelig sat som
Eenhed af Rum og Tid. Det ved Her og Nu bestemte Sted viser imidlertid
ud over sig selv hen til et andet ved Her og Nu bestemt Sted; denne
Stedets Vorden er „\emph{Bevægelsen}“. Men Her og Nu er allevegne; det
bestemte Sted er i al Ubestemthed kun det almindelige Sted:
Bevægelsens modsigende Momenter falde sammen i umiddelbar Eenhed;
denne umiddelbart tilværende Eenhed er „\emph{Materien}“. „Materien er
\emph{uigjennemtrængelig} og gjør \emph{Modstand}, er et Føleligt,
Synligt o. s. v. Disse Prædicater ere ikke Andet, end at Materien
deels er for den bestemte Iagttagelse, overhovedet \emph{for et Andet},
men deels ogsaa \emph{for sig}. Begge ere Bestemmelser, hvilke den
just har som Identiteten af Rum og Tid, af det umiddelbare
\emph{Udenforhverandre} og \emph{Negativiteten} eller af den for sig
værende Enkelthed.“

% --- p. 16 ---
Denne Deduction og dens Dialektik er characteristisk for
% Anførselstegnet aabnes paany efter Indskuddet „hedder det“; kun eet
% lukkende “ staaer, ved Citatets Slutning (efter „Existens.“).
% Saaledes trykt i begge Exemplarer (Google og KB).
Immanensphilosophien. „Rummet, hedder det, „er ikke adæqvat med sit
Begreb; det er derfor Rummets Begreb selv, der i Materien skaffer sig
Existens.“ Tydeligere kan Læren om det objective Begrebs
mirakelgjørende Evne ikke udtales. Men Trylleriet forsvinder, og
Miraklet kaldes igjen tilbage ved en Forklaring som følgende: „Man
har,“ tilføies der, „ofte begyndt med Materien, og da anseet Rum og Tid
for dens Former. Det Rigtige heri er, at Materien er det Reale ved Rum
og Tid. Men disse maae, paa Grund af deres Abstraction, her forekomme
som det Første, og saa maa det vise sig, at Materien er deres
Sandhed.“ Den sidste af disse Udtalelser tiltræde vi ubetinget; men
den staaer i skrigende Modsigelse med den første. Er det virkelig
Rummets Begreb, som i Forening med Tidens har formaaet at give Materien
Tilværelse, da kan Materien, det Reale, umulig være det Første; den maa
saavel i subjectiv som i objectiv Forstand være det Andet, det
Afledede. Skal Materien derimod i sig selv være det Første, fordi den
er det Concrete, der i sig optager Rummet og Tiden som abstracte
Former, da er det jo dog en vanvittig Paastand, at det Abstracte skulde
have frembragt det, hvoraf det selv er abstraheret, at Rummet,
tilskyndet af sit eget urolige, fra Materien abstraherede Begreb,
skulde have skabt Materien.

Abstraction, Modsigelse, Negativitet, uendelig Negativitet o. s. v.
ere Selvhedskategorier, der tilvisse have Gyldighed i den uendelige
Objectivering; men de misforstaaes, misbruges og tages forfængelig,
naar de uden videre lægges ind i det Objective, saa at der tillægges
dem udvortes Væren. Rum og Tid ere Former, hvis abstracte Væren maa
integreres ved en fuldstændigere Tilværen: de maae have et
\emph{Mellemværende}, der optages i dem, og hvori de optages. Men
Abstraction og Integration ere Tankehandlinger, Functioner af
objectiverende Selvhed. Hvilken Realitet, det
% --- p. 17 ---
fordrede Mellemværende skal have, vil da beroe paa Objectiveringen.
For den abstracte Objectivering, forsaavidt det endnu kun gjælder om en
almindelig Bestemmen af Sted og mechanisk Bevægelse, er et forestillet
Mellemværende tilstrækkeligt.

Ingen Bevægelse er mulig uden et Noget, som kan bevæge sig. Bevægelsen
foregaaer i Rummet, men det sig Bevægende er forskjelligt fra Rummet,
thi Rummet bevæger sig ikke. Stedforandringen er en Forandring i
Rummet, men uden at Rummet forandres. Forandringen skeer med det
Bevægede, der gaaer over fra Sted til Sted; og dog er det ikke det
Bevægede selv, men kun Stedet, der forandres. Stedet tilhører ikke det
Bevægede, thi det tilhører Rummet; og dog tilhører det ikke Rummet for
sig, thi det bliver først i Bevægelsen til Sted ved at indtages af det
Bevægede. Stedet er hverken ligefrem værende eller ikkeværende; det er
en objectiv Relationsbestemmelse, der er objectivt reflecteret.
Bevægelsen foregaaer i Tiden, og det i Tiden Bevægede er forskjelligt
fra Tiden; men Tiden indgaaer i Bevægelsen paa en anden Maade end
Rummet. Rummets Steder ere samtidig værende, nemlig efter Mulighed;
ethvert tilbagelagt Sted kan desaarsag igjen indtages. Men Tidens
Øieblikke svinde, og ethvert forsvundet Øieblik er uigjenkaldelig tabt.
Med det Bevægelige til Mellemværende faae da Rum og Tid hver sin
reflecterede Andeel i Bevægelsen. Den bestemte Bevægelse er Hastighed,
og Hastigheden en Qvotient, hvori Rumselementet som Tæller og
Tidselementet som Nævner reflectere hinanden objectivt,
$\dfrac{ds}{dt} = v$; men paa Objectiveringens Vilkaar: den
objectiverende Reflexion er subjectiv.

Hvilken Begrebsforvirring der kan opstaae, naar man i Bevægelsens
Phænomen absolut vil tænke det Objective uden at betænke dets Forhold
til den objectiverende Subjec\-%
% --- p. 18 ---
tivitet, er bekjendt. Eleaternes Piil, der flyver i det uendelige Rum,
er nu her, nu her, altid nu her! og kommer ikke af Stedet. Det er et
objectivt Extrem, der ophæver sig selv. For at Bevægelsen kan blive
til, maa Betragteren erindre, sammenfatte og anskue, nemlig de Steder,
det Bevægede har tilbagelagt og fremdeles tilbagelægger. Intet Under,
at det objective Extrem hos Eleaterne slog over i et subjectivt, saa at
hine ubøielige Tænkere reentud erklærede al Bevægelse for Skin og
Bedrag. Men at ogsaa det Modsatte kan skee, at nemlig det Subjective
ved et logisk Selvbedrag kan slaae over i det Objective, er af den
Hegelske Dialektik indlysende. Medens Zeno, for at faae Pilen til at
staae stille, viser hen til det uendelige Rum, hvori al Relativitet er
forsvundet, holder Hegel Stedbestemmelsen fast i Tanken. „Her er tre
Steder: det, som er, det næstfølgende og det forladte. Men det er
tillige kun eet Sted, hine Steders almene Sted, et i alle Forandringer
Uforandret“. Man lærer heraf, hvorledes Abstractionen, Modsigelsen og
Negationen drive hinanden frem. At de tre Steder ere eet Sted, beroer
paa en Abstraction, idet der abstraheres fra den Forskjel i udvortes
Væren, der netop er betegnende for Rummet. Det er tillige en
Modsigelse; men denne af Abstractionen betingede Modsigelse bliver
opløst derved, at den i Treheden negerede Eenhed igjen negeres i den
almene Eenhed, det Blivende, saa at „Vedvaren (die Dauer) efter sit
Begreb er Bevægelsen“. At Naturen ikke paa egen Haand anstiller slige
Reflexioner, men overlader det til Subjectiviteten, maa vel
indrømmes. Ikke desto mindre bliver det subjectivt reflecterede
Bevægelsesbegreb ligefrem erklæret for et ikke subjectivt, men blot
objectivt Naturbegreb og ovenikjøbet gjort til Eet med Bevægelsen selv.

Som abstracte, mellem Væren og Ikke-Væren svævende Former ere Rum og
Tid uden særlig Selvstændighed; deres Tilværen afhænger af det
Mellemværende, hvori de forenes,
%  --- B. Materie og Kraft  (pp. 19–38) ---
% --- p. 19 ---
altsaa af det Indhold, hvormed de fyldes: i denne Afhængighed sees
deres Idealitet. Fyldes Rummet og Tiden med Drømmebilleder, have de
kun en drømt Tilværen; fyldes de med en Verden af poetiske Skikkelser,
opnaae de en æsthetisk Tilværen; først med et Indhold af virkelige
Ting blive de virkelige. Gjælder det, som vi have seet om Formerne,
at deres objective Væren beroer paa Objectivering, da maa det Samme
nødvendig gjælde om det Indhold, hvormed de fyldes.

% Kort centreret Streg trykt her, mellem Afdeling A og Afdeling B.
\streg

\afdeling{B.}{Materie og Kraft.}

Rummets Virkelighed falder sammen med dets virkelige Udfyldning. Det
Rumudfyldende, Materien, og det i Rummet Virksomme, Kraften, tilhøre
Naturen, men sættes af Aanden; Opgaven er at vise, hvorledes de to
Grundlagsbestemmelser: Oprindelse fra Aand og naturlig Realitet
forudsætte hinanden.

\runin{a) Materien.} Naturens første umiddelbare Reale er ikke et
Begreb, men et sandseligt Tilværende, hvori et Begreb existerer.
Existensbegrebet er et Realitetsbegreb; det gjælder da først og
fornemmelig om at indsee Forholdet mellem Begreb og Realitet.

At forudsætte Realiteten og udelukke Begrebet er betegnende for
Materialismen. Det Eensidige i denne Synsmaade bestaaer ikke deri, at
den tillægger Materien Realitet, thi hvis Materien kun var et Skin,
saa var Naturgrunden udhulet; men Eensidigheden er, at Materien gjøres
til noget
% --- p. 20 ---
i sig Absolut, til et evigt, ubetinget Værende, af hvis Grunddele Alt
sammensættes og dannes. Disse Grunddele, Atomerne, opfattes paa een
Gang baade som Massedele og som absolute Realpunkter. Det er en
Modsigelse, thi dersom Atomerne ere Massedele, saa ere de relative; ere
de derimod absolute, saa er al virkelig Relation imellem dem umulig.
At tillægge Atomerne Egenskaber er at gjøre dem relative, thi det ene
Atoms Egenskab faaer i Tiltrækning og Frastødning kun Betydning ved det
andets Egenskab. At forudsætte Atomernes absolute Eenhed med deres
Egenskaber er at indrømme deres absolute Relativitet og nedsætte dem
til Momenter for hverandre indbyrdes. Med denne Conseqvens har
Materialismen opløst sig selv.

At sætte Begrebet som Realitet og Realiteten som Begreb er betegnende
for Idealismens Eensidighed. Idealismen er magtesløs; den miskjender
Modsætningen af det Ideelle og det Reelle, af Aand og Materie, og
indseer ikke, at der til et mægtigt Indre maa svare et mægtigt Ydre.
Idealismen vil forvandle Materiens Begreb til virkelig Materie; det er
en ufrugtbar Speculation. Modsætningen af Begreb og Materie er en
Modsætning i Magt, en Modsætning saa haard og stærk, at
Eenhedsbegrebet falder sammen med Begrebet af Almagt, absolut
overgribende Selvmagt. Selvmagtens indre oprindelige Modsætning er
Gjenstandsmagt, Magt i Andet; det første Ydre for Magt i Andet er
Materien. Derfra Atomernes Vrimmel: lutter Andethedspunkter,
Realitetspunkter. Den qvalitative Modsætning af Aand og Materie er
betegnet ved Magtens Udslag. Udslaget er ikke blindt, det er en
Grundhandling i Kraft af almægtig Begriben: Begrebet er Realitetens
Grund, Realiteten Begrebets Existens.

En umiddelbar Gjenstandsmagt uden tilsvarende Begreb er lige saa
umulig, som et existerende Magtbegreb uden Gjenstand. Man kan ikke
dele det Existerende i to absolut forskjellige Arter: et Existerende
\emph{med} Grund og et Existe\-%
% --- p. 21 ---
rende \emph{uden} Grund; thi Existensens Væsen er at have en Grund.
Den reale Grund er en i Magtbegrebet opløst Existens; det Existerende,
der ingen Grund har, existerer ikke, dets Tilværen er et huult
Phænomen, et Skin.

Materialismens Grundvildfarelse viser sig nu deri, at den gjør den
„fra Evighed af existerende“ Materie til et Existerende uden Grund.
Atomerne, de absolute Existenser, paa hvis Adskillelse og Forbindelse
alle endelige, materielle Dannelser beroe, forholde sig da til de
virkelige Ting i deres Endelighed, som det Grundløse til det
Begrundede! En unaturligere Adskillelse af Grund og Existens kan ikke
gives. Materialismen har et skarpt Blik for Realiteten som
Gjenstandsmagt. Den begynder med et Magtsprog, peger paa
Kjendsgjerningen, tager det Existerende som et ufravigelig Givet; det
er sundt og rigtigt. Men at oversee Magtbegrebet og fornegte
Existensens Væsen for at bygge paa et Existerende uden Grund, er
usundt og falsk.

Den objective Idealisme holder paa Begrundelsen, idet den holder sig
til Begrebet; men Grundfeilen er, at den overseer Gjenstandsmagten og
udsletter Forskjellen mellem Reflexion i Viden og Reflexion i Magt.
Hvad dette Misgreb har at betyde, kan i Materien sees paa Atomer og
Masse. Den existerende Masse har Atomerne til Grundlag; men Atomernes
Existens har igjen den forudsatte Masse til Grundlag. For at danne
Grundlag, maa et Existerende gaae til Grunde. Det Existerende
forestilles; men Grunden og Grundlaget tænkes. At Atomerne existere
betyder, at de sættes i Forestillingen, medens den underforstaaede
Masse sættes i Tanken; er Massen derimod i Forestillingen, idet
Atomerne underforstaaes, saa er det Massen, der existerer. Atomernes
Undergang i Massen og Massens Opløsning i Atomerne ere imaginære
Bevægelser, der foregaae indenfor Tanken og Forestillingen; det er en
Reflexion, som, hvormange Begrebsbestemmelser den end optager, aldrig
kommer ud over en
% --- p. 22 ---
forestillet Existens med en tænkt Grund, aldrig naaer Virkeligheden.

I Magtreflexionen begyndes der udtrykkelig ud fra Virkeligheden,
forsaavidt der begyndes i det Existerende med Gjenstandsmagten;
Grunden er et Magtbegreb, og Grundlagsbestemmelserne Realmomenter af
opløst Existens. Overgangen fra Existens til Grund, og fra Grund til
Existens, er en væsenlig Forandring, en objectiv Magthandling, en
Naturbegivenhed. Massen existerer, og Atomerne existere; og dog er her
ikke paa een Gang to Slags Existerende eller to Lag Existens. Falder
Existensen i Massen, saa have Atomerne ophørt at existere; de ere som
Grundlagsmomenter optagne i Massens Existens. Men da Atomerne gik
under, foregik der virkelig Noget: de nærmede sig hverandre, opgave i
Sammenhæng (Cohæsion) deres Selvstændighed og bleve Momenter i Massen.
Falder Existensen i Atomerne, saa er Massen opløst: den er gaaet
under. Men da Massen gik under, foregik der virkelig Noget: den
opløste sig og tabte i Opløsningen sin Selvstændighed. At Massen tabte
sin Selvstændighed er Grunden til, at Atomerne fik Selvstændighed og
kom til at existere. Atomernes Existens har lige saa vist sin reale
Grund i Massens Undergang, som Massens Existens har sin reale Grund i
Atomernes Undergang. I denne Magtreflexion er det Ene ikke et
ubetinget Første, og det Andet et Afledet, men begge Existenser,
Atomernes og Massens, ere lige oprindelige og lige afledede. Undergang
er da heller ikke Tilintetgjørelse; det Existerende, der gik under i et
andet Existerende, kommer igjen til Existens, naar den paa dets
Undergang beroende Existens gaaer under.

Den Reflexion af Magtbegreb og Gjenstandsmagt, der afgiver den
mechaniske Naturs almindelige Grundlag, er overalt kjendelig paa en
brat Omsætten af existerende Led i Momenter, af Momenter i existerende
Led. Magtbegrebets
% --- p. 23 ---
reflecterede Eenhed med Gjenstandsmagten er udtrykt i den haardeste
Modsætning. Det Existerende er intet Begreb; det er Gjenstand,
virkelig Kjendsgjerning. Begrebet existerer ikke; det reflecteres i
opløst Existens som Grund og Grundlag. Men Eenheden er igjen lige saa
mægtig som Modsætningen: det Existerende har Grunden i sig:
Gjenstanden er sit virkeliggjorte Begrebs materielle Udtryk. Fra
Reflexionens Side er Begrebet det Overgribende; fra Gjenstandssiden er
det Kjendsgjerning og Kjendsgjerning først og sidst. Den uendelige
Vexelbestemmelse af Reflexion og Kjendsgjerning, Grund og Existens, er
Actualitetens Form for \emph{Magt i Andet}. Men Magt i Andet
forudsætter Selvmagt; Kjendsgjerningernes iboende Begreber maae
begribes af overgribende Enemagt: \emph{den mechaniske Naturmagt bæres
af Aandens Almagt.}

I hvilken Betydning den existerende Materie maa siges at være den
første mechaniske Fremstilling af Magt i Andet, kan sees af de
Egenskaber, Physiken særlig har tillagt den som Materie. Materien er
\emph{sandselig}, ikke blot for os, men i sig selv. At tænke Materien
er at forvandle den til et Tankevæsen og fornegte dens Existens. Er
Materien for Aanden kun et Tankevæsen, saa er hele Naturen et Skin.
Men Sandheden er, at Materiens Væsen er et Magtvæsen; dens Begreb er et
i Magten omsat Vidensbegreb, dens Existens en Gjenstandsmagt, en
Modstand for Tanken, en umiddelbar Gjenstand, ogsaa for Aanden. At
være Aandens umiddelbare Gjenstand, Anskuelsens reale Andet, en
Grændse for Tanken, er at være sandselig. Materien er
\emph{rumopfyldende}. Den begrændsede Materie indtager sit Sted og
gjør Modstand mod anden Materie, der vil indtage det samme Sted; denne
Modstandsmagt er Andethedsmagt, Magt til at udelukke Andet fra Andet.
Materien er \emph{deelbar}. Deelbarhed beroer paa Andethed. Massen er
Et, Atomerne et Andet. Massens Eenhed er ingen Selveenhed, men
% --- p. 24 ---
Eenhed af Andet med Andet; Massen bestaaer ved sit Andet, thi den
bestaaer af Atomer. Skulde Atomerne absolut bestaae ved sig selv, da
kunde det dog umulig være som mathematiske Punkter; det maatte da være
som smaa Masser. Men hvad der gjælder om de større Masser, gjælder om
de mindre, mindste, allermindste: Massen bestaaer i Andet ved Andet.
Hvert enkelt Atom har samme Bestaaen som alle andre. Realiteten af eet
Atom er Realiteten af alle; men Realiteten af alle er Massen: Atomerne
bestaae i Massen, altsaa i Andet. Materien er \emph{ligegyldig mod
Hvile og Bevægelse}; denne Ligegyldighed er ret et charakteristisk
Tegn paa Andethed. At Materien er underlagt Inertiens Lov, bliver
ogsaa betegnet ved, at den er villieløs, hvilket med andre Ord vil
sige, at den er uden Selvbestemmelse. Den kan ikke, naar den er i
Hvile, bestemme sig til Bevægelse, men maa bestemmes af Andet; den kan
ikke, naar den er i Bevægelse, bestemme sig til Hvile, men maa
bestemmes af Andet. Denne Lære om Materiens selvløse Selvstændighed
(Inertie) har afgjørende Betydning for Erkjendelsen saavel af Materiens
Væsen som af Bevægelsens Love.

Ikke desto mindre vil Immanenslæren som speculativ Naturphilosophie
endnu til Trods for empirisk Physik og rationel Mechanik hævde den
mystiske Forestilling om en Materien iboende Selvhed, en stille Længsel
efter Befrielse, en skjult Evne til at sætte sig selv i Bevægelse.
Materien, siger Hegel, er Bortfjernelse i Rummet. Den gjør Modstand,
støder sig selv bort fra sig selv: det er \emph{Repulsionen}, hvorved
den sætter sin Realitet og fylder Rummet. Men de Enkeltled, der
repelleres fra hverandre, ere alle Eet. At det er det samme Ene, som
ved at støde sig selv fra sig selv ophæver Enkeltleddenes Fjernelse,
deri bestaaer \emph{Attractionen}. Begge tilsammen udgjøre som Tyngde
Materiens Begreb. Tyngdens Eenhed er kun „en Skullen, en Længsel, den
% Citatet aabnes her og lukkes først paa næste Side („naaes ikke.“);
% en Anførsel, der løber over et Sideskifte, ikke en Trykfeil.
ulyksaligste Stræben, hvortil Materien evig er
% --- p. 25 ---
fordømt; thi Eenheden kommer ikke til sig selv, den opnaaes ikke.“
Materien kan imidlertid trøste sig med, „at den ikke er saa dum“ ---
som Physikerne --- „at den skulde holde Eet og Flere ud fra
hinanden.“ Det er kun i den endelige Sphære, i Kredsen af „de selvløse
Legemer“, at Materien lider af Træghed; i sit „frie Begreb“, i den
„himmelske Sphære“ har Materien Selvhed og bevæger sig selv.

Idet Immanenslæren paa denne Maade træder polemisk op mod Physiken,
kommer den med det Samme i en mislig Stilling til Aandslæren. For at
bøde paa Savnet af Aandens begribende Selvmagt og lade Naturen begribe
% Trykfeil: der staaer „sin egne Begreber“ i begge Exemplarer.
% Bogens egen Trykfeil-Fortegnelse retter: „Side 25 Linie 11 f. o.
% sin --- læs sine“. Her gjengivet som trykt.
sin egne Begreber, anbringes Personificationer: til Erstatning for den
udelukkede Aand indføres Aandrigheder. „Ligesom Tiden er den eenfolde
formelle Natursjæl, og Rummet ifølge Newton Guds Sensorium, saaledes er
Bevægelsen Begrebet af den sande Verdenssjæl. Vi ere vante til at
ansee den for Prædicat, Tilstand, men den er i Virkeligheden Selvet,
Subjectet som Subject, det Blivende i al Forsvinden.“ Ligesom
Bevægelsen er det væsenlige Selv, saaledes er Materien i det Uendeliges
Mechanik naturligviis det virkelige Selv. Den Kjendsgjerning, at
Materie bevæger Materie i det Uendelige, forklares ved en logisk
Omskrivning derhen, at Materien altsaa \emph{bevæger sig selv}. Ved
slige Forklaringer bliver Naturudviklingens Begreb fra Grunden af
mystificeret.

\runin{b) Kraften.} Materien bevæger og bevæges; den ene Masse bevæger
den anden: Materien har Bevægelsens Grund og Aarsag i sig selv, men
under Andethedens Form. Først naar det tydelig indsees, at det
nødvendig hører med til dens rumopfyldende Realitet at kunne bevæge og
lade sig bevæge, er Materiens Begreb udtømt. Det er ikke Bevægelsens
abstracte Formbestemmelse, men den virkelige Bevægelse, hvorom her
handles. Vi maae see os om efter et
% --- p. 26 ---
Grundphænomen, hvoraf det klarest kan vise sig, hvad det Bevægende
egenlig er.

Man har fra et endelig mechanisk Synspunkt henviist til
\emph{Stødet}. Under to sammenstødende Legemers Berøring foranlediger
den gjensidige Modstand en Kamp om Stedet, der synes at afgive
tilstrækkelig Grund til Stedforandring, altsaa til Bevægelse. Grunden
er saa meget naturligere, som den øiensynlig ligger i den Realitet,
hvormed enhver Masse indtager sit Sted. Umiddelbar Berøring
% Græsk, sat med kursiv græsk Skrift: ἅπτεσθαι (spiritus asper og
% akut over alfa, læst paa KB-Exemplaret).
(ἅπτεσθαι) er da ogsaa af Aristoteles sat som Betingelse for, at den
ene Materie kan indvirke bevægende paa den anden. Men at Stødet ikke
kan være det Bevægende selv, er indlysende. Stødet forudsætter nemlig,
at de to sammenstødende Legemer enten begge maae have bevæget sig i
modsat Retning, eller, hvis Retningen er den samme, maa enten det ene
have bevæget sig med større Hastighed end det andet eller det ene have
bevæget sig og det andet været i Hvile. Det sees da fra alle Sider, at
Stødet ikke kan være det egenlig Bevægende; det er jo selv et Resultat
af Bevægelse.

For at undgaae den forudsatte Bevægelse, kunde man forsøge at sætte det
Bevægende i et Tryk og gjøre \emph{Trykket} til et Grundphænomen. Tryk
og Modtryk udtrykke nu ligefrem den med Materiens Realitet identiske
Modstandsmagt; Bevægelsen indtræder som en nødvendig Følge, hvis
Forudsætning er, at Ligevægten mellem Tryk og Modtryk ophæves.
Bevægelsen begynder med en Stræben til Bevægelse; en saadan Stræben
ligger virkelig i Trykket, og Forklaringen har forsaavidt en Grund; men
det er en utilstrækkelig Grund: her forudsættes et Noget, som endnu
ikke er begrundet. Trykkets Phænomen kan tydes deels saaledes, at
begge Legemer, hvert fra sin Side, stræbe at trænge ind i hinanden,
medens de dog gjensidig modstaae hinanden; deels saaledes, at kun det
ene er paatrængende, det andet blot modstaaende. I begge Tilfælde
falder Bestræbelsen sammen med Paatræn\-%
% --- p. 27 ---
genheden; og Spørgsmaalet bliver, hvorfra Paatrængenheden da hidrører.
Paatrængenhed forudsætter Nærmelse: to adskilte \emph{Massers
gjensidige Nærmelse} maa altsaa være det søgte Grundphænomen.

To i Rummet adskilte Masser, \textit{A} og \textit{B}, ville, naar
uvedkommende Hindringer tænkes borte, nærme sig hinanden, altsaa komme
i Bevægelse; men hvad er Grunden? Hegel beraaber sig paa Tyngden og
sætter Tyngden i en fælles Søgen efter et fælles, udenfor Masserne
liggende Midtpunkt. „Midtpunktet, siger han, er ikke at tage som
% Indskuddet „siger han“ bryder Citatet, men her aabnes der IKKE paany
% med „; Anførselstegnene balancere altsaa paa denne Side.
materielt; thi det Materielle er netop det, at have sit Midtpunkt
\emph{udenfor sig}. Ikke Midtpunktet, men en Stræben efter det, er
Materien immanent. Tyngden er saa at sige den Bekjendelse, Materien
aflægger om Intetheden af sin udvortes Væren for sig, om sin
Uselvstændighed, sin Modsigelse.“ Men at tillægge Masserne et saadant
indre Savn, en Higen mod et udenfor værende Ideelt er vel poetisk. men
% Trykfeil: Punktum, hvor Meningen kræver Komma --- „poetisk. men ikke
% videnskabeligt“. Saaledes trykt i begge Exemplarer (Google og KB);
% ikke rettet i Bogens egen Trykfeil-Fortegnelse.
ikke videnskabeligt; ligesom det da ogsaa udtrykkelig strider mod den
haarde Andethed, den selvløse Selvstændighed, at forklare Massernes
gjensidige Nærmelse af dunkel Sympathie og digte en --- naturligviis
ubevidst --- Længsel ind i \textit{A} efter \textit{B}, besvaret af en
--- naturligviis ubevidst --- Længsel i \textit{B} efter \textit{A}.
Den strenge Mechanik har en anden Udtydning. \textit{A}, hedder det,
tiltrækker \textit{B}, fordi \textit{B} tiltrækker \textit{A}: det i
begge Masser, ja i hvert Atom af begge Masser, Tiltrækkende er
\emph{Kraften}.

Kraften er, mechanisk forstaaet, Bevægelsens Aarsag; men --- tilføier
Physiken --- en ubekjendt Aarsag. Vi skulle nu nærmere undersøge,
hvorledes det egenlig forholder sig med „den ubekjendte Aarsag, vi
kalde Kraft“. Naar det hedder, at vi ikke vide, hvad Kraft er, saa
bliver der stiltiende underforstaaet, at vi endda nok vide, hvad
Materie er. Men dersom Materien virkelig er en bekjendt Størrelse, da
er det umuligt, at Kraften i sin eiendommelige Sammen\-%
% --- p. 28 ---
hæng med Materien kan være os aldeles ubekjendt. Kraften, det
Usandselige, svarer til Materien, det Sandselige, som Magtreflexionen
til den umiddelbare Gjenstandsmagt: Materien er det Ydre, Kraften det
Indre. Men et Indre kan aldrig vise sig ligefrem; at det viser sig,
vil sige, at det yttrer sig; det reflecterer sig i sit Ydre og maa
kjendes paa sine Yttringer. En Sætning som den: vi vide ikke, hvad
Kraft er i sig selv, vi kjende kun dens Yttringer, dens Virkninger,
udsiger da slet intet Andet, end hvad der ligger i Reflexionens Natur,
at nemlig det Indre nødvendig maa kjendes paa sit Ydre. Det abstract
logiske Væsen reflecterer sig i Skin af Skin, men Magtens Reflecteren
er en Virken, dens Phænomener ere Kjendsgjerninger: som
Magtreflexionens kategoriske Udtryk er Kraften ikke en i menneskelig
Indbildning „hypostaseret Tanke“, men en i sine Virkninger reflecteret
Kjendsgjerning. Kraften som Indre er Intet i sig selv absolut; det
Indre er en Mulighed, som først ved at yttre sig og virke bliver til
Virkelighed. Den ovenanførte Sætning kan da ikke i objectiv Forstand
betragtes som nogen sikker Betegnelse for den menneskelige Videns
absolute Grændse, men vel som en subjectiv Antydning af Uklarhed over
Tænkningens, Reflexionens, særlig Magtreflexionens Væsen.

Hvor uklar og svævende den sædvanlige Forestilling om Kraft er,
fremlyser af den Maade, hvorpaa dens Forhold til Materien nærmere
bestemmes. Paa den ene Side taler man, som om Materien var Et og
Kraften noget ganske Andet. Det er da ikke Massen \textit{A}, der ved
Tiltrækning sætter \textit{B} i Bevægelse; nei, det er Kraften, den i
\textit{A} værende Tiltrækningskraft, hvorved Bevægelsen af \textit{B}
foraarsages, og omvendt: det er ikke \textit{B}, men Kraften i
\textit{B}, der bevæger \textit{A}. Masserne bevæges, men Kræfterne
bevæge: Materien er \emph{passiv}, Kraften \emph{activ}. Paa den anden
Side taler man igjen, som om Kraften kun var en Egenskab ved Materien.
% --- p. 29 ---
Saa bliver det da virkelig Massen \textit{A}, der i Egenskab af Masse
tiltrækker \textit{B}, og omvendt. Den sidste Synsmaade staaer i
øiensynlig Strid med den første. Ved at gjøre Kraften til Materiens
Egenskab og, hvad den almindelige Tiltrækning angaaer, ovenikjøbet til
Grundegenskab lader man Kraftens Begreb uden videre gaae op i
Materiens. Ifølge den første Synsmaade skulde her være to Væsener, men
ifølge den sidste er her kun eet Væsen. Som Egenskab ved Materien
% „høist“: begge OCR-Vidner læse „heist“, men ø'et staaer tydeligt paa
% KB-Exemplaret. Læsningen „heist“ forkastet.
kommer Kraften i en høist betænkelig Stilling til Materiens øvrige
Egenskaber, navnlig til Inertien. Ifølge Inertien kan Materien kun
bevæges, men ikke bevæge; ifølge Kraften kan den kun bevæge, ikke
bevæges. Det maa da synes, som var Materien en sig selv modsigende
Realitet. Som det Rumopfyldende har Materien desuden en Deel
Egenskaber tilfælles med Rummet, for ikke at tale om en Mængde
umiddelbare Sandselighedsegenskaber, der ikke vel lade sig forlige med
Kraften som Kraft. Thi man vil dog ligesaa lidt kunne tale om en bred
Kraft, en rund Kraft, en fiirkantet Kraft, som om en tynd eller tyk
eller graa Kraft.

De anførte Vanskeligheder bortfalde, naar vi holde os til den ovenfor
angivne Vexelbestemmelse af Magtreflexion og umiddelbar
Gjenstandsmagt. Materien er paa sin Maade reflecteret, nemlig i
gjensidig Begrundelse af Atomer og Masser; Kraften er paa sin Maade
% „iøinefaldende“: begge OCR-Vidner læse „ieinefaldende“; ø'et staaer
% tydeligt paa KB-Exemplaret.
umiddelbar, nemlig i sine Yttringer. Og dog er Forskjellen
iøinefaldende; thi for Materien er Umiddelbarheden, men for Kraften er
Reflexionen det Afgjørende. Reflexionen af Magt er noget Mere end en
logisk Speilen af det Ene i det Andet, den er en virkelig Bestemmen, en
Virken, hvori det viser sig, hvad det Ene formaaer over det Andet. Det
er netop en saadan Formaaen, der mechanisk udtaler sig i Bevægelsen.
Bevægelsen af \textit{B} reflecterer Kraften i \textit{A}, idet
Bevægelsen af \textit{A} reflecterer Kraften i \textit{B}. Den
gjensidige Bevægelse er et Phænomen af Magtreflexionens Væsen: den i
Væsen overgribende Reflexion
% --- p. 30 ---
beviser, at det Almene i Leddenes Relation behersker de særskilte Led.
Physiken er forsaavidt i sin gode Ret, naar den opfatter Kraften, det
Active, som Væsenhed for sig til Forskjel fra den træge Materie, det
Passive. Men herved overseer Physiken dog ikke, at de to Væsener i
Grunden ere Sidebestemmelser af samme Væsen. Magtreflexionens Eenhed
med Gjenstandsmagten fremhæves udtrykkelig derved, at Kraften opfattes
som Materiens egen reflecterede Bestemthed, altsaa som dens Egenskab.
Som Materiens gjennemgaaende Egenskab er Kraften noget Andet end en
umiddelbar Qvalitet eller Beskaffenhed, den er et kategorisk Udtryk,
ikke for Materiens sandselige Realitet, men for dens reflecterede
Væsen, dens iboende Energie, dens Indre. Hermed forsvinder
Modsigelsen mellem Kraften som reflecteret Egenskab og Materiens
øvrige, umiddelbare Egenskaber. De umiddelbare Egenskaber tilhøre
\emph{det Extensive}, men Bestemmelserne ved Kraften tilhøre \emph{det
Intensive}, thi de tilhøre Reflexionen, vel at mærke: Reflexionen i
Magt. „Man kan“, hedder det om Tyngdepunktet i den mechaniske
Naturlære, „forestille sig Tingene, som om hele Legemets Vægt var
forenet i Tyngdepunktet, og alle de øvrige Dele ingen Vægt havde.“
Denne Forestilling er mere end en Forestilling; det er en virkelig
Sandhed. Den paa Delene fordeelte Tyngde, Tyngden i det Extensive,
forvandles i Magtreflexionen til Tyngde-Eenhed, en Tyngde i det
Intensive: Tyngdepunktet \emph{er} det intensive
Eenhedspunkt.\footnote{% trykt Mærke: *)
Jevnfr. Mathematik og Dialektik. S. 129--131.}

Gyldigheden af Ovenstaaende bliver endnu nærmere at prøve, forsaavidt
det gjælder en endelig Tydning af Tiltrækningens ellers saa gaadefulde
Phænomen. Masserne \textit{A} og \textit{B}, der paa Afstand udfylde
hver sit Rum, synes slet ikke at have Noget med hinanden at gjøre; og
dog have de med hinanden at gjøre, thi de nærme sig hinanden. Det er
en
% --- p. 31 ---
forunderlig Modsætning af Væsen og Existens, Magtreflexion og
Gjenstandsmagt, som i Phænomenet lader sig tilsyne; Existensen siger
Et, Væsenet et Andet. Gaadens Løsning ligger i den Indsigt, at den
materielle Virkelighed gjennem Væsen og Existens har udtalt to
modsatte, men lige uafviselige Fordringer, og at
Tiltrækningsphænomenet netop er begge Fordringers Opfyldelse. Den
første Fordring er, at de materielle Led maae have umiddelbar
Selvstændighed; denne Selvstændighed er anskueliggjort i Massernes
sandselige Existens. Existensen svarer til Stedet. Hver Masse
indtager sit Sted, og det kan let forekomme Beskueren, som var den
rumopfyldende Realitet det fuldstændige Udtryk for Materiens
Virkelighed. Det er imidlertid et Sandsebedrag; det, der saaledes sees
i umiddelbar Existens, er kun en halv Virkelighed. Existensen skjuler
Noget i de existerende Led, der kun kan komme tilsyne i Reflexionens
Form som Væsen. Den anden Fordring er, at Væsenet maa reflectere sig i
Leddenes Relation, i Massernes indbyrdes Forhold. Forholdsled uden
Forhold er en lige saa ufuldbaaren Tanke, som Forhold uden Forholdsled.
Ere Leddene virkelige, maae Forholdene ogsaa være det; deres
Virkelighed er i deres Virksomhed, og denne Virksomhed er Forholdets
Realitet. Naar Magten reflecteres i Forhold af virkelige Led, viser
Væsenet sig som Kraft. Massernes Afstandsforhold bliver ved Kraften et
virksomt, altsaa virkeligt Forhold. I Masserne er Materiens Existens
begrændset, indskrænket til Stedet; Masserne ere som virkelige Led
Momenter i Virkeligheden. Som Væsen derimod gaaer Kraften fra Masse
til Masse ud over Grændsen, altsaa ud over Stedet, og
Virkelighedsmomentet sees i den gjensidige Stedforandring. At Kraften
gaaer gjennem Rummet, at Masserne virke paa Afstand, er et mechanisk
betegnende Udtryk for den Vexelbestemmelse af Magtreflexion og
Gjenstandsmagt, der gjør Virkeligheden fuldstændig.
% --- p. 32 ---
Den tilsyneladende Modsigelse, at en Masse har Kraft til at bevæge en
anden Masse, men ikke Kraft til at bevæge sig selv, har med det Samme
fundet sin Løsning. Materiens Virkelighed er en Selvstændighed i
Andethedens Form, en selvløs Selvstændighed. Selvløshed er
Uselvstændighed; den materielle Virkelighed er en uselvstændig
Selvstændighed. Uselvstændigheden er Materiens Inertie,
Selvstændigheden dens Kraft. Inertien er, at Materien bevæges af
Andet, Kraften er, at den bevæger Andet. I de hinanden tiltrækkende
Masser ere Inertie og Kraft paa begge Sider forenede: Kraft er
Inertiens Complement, Inertie er Receptivitet for Kraften. Saaledes
ere da ogsaa Selvstændighed og Uselvstændighed paa begge Sider
forenede. Massen \textit{A} har Selvstændighed, den sætter ud fra sit
Sted \textit{B} i Afhængighed; men \textit{B} har ogsaa Selvstændighed
og sætter ud fra sit Sted \textit{A} i Afhængighed. Relativiteten er,
at begge Masser i deres Selvstændighed ere uselvstændige. Her behøves
ingen udvortes Betragter for at sammenholde de to Steder i
Forestillingen og bestemme deres Afstand; thi saa gjennemgaaende er
Relativiteten, at det i Virkeligheden er Masserne, der gjensidig
bestemme hinandens Steder: Materiens Relativitet er i Bevægelsen
virkelig formedelst Kraften.

\runin{c) Materiens og Kraftens reflecterede Eenhed: Aand og Materie.}
Den hertil førte Udvikling har Magtens Idealitet, den reflecterede
Magt, til Grundlag; dens Betydning staaer og falder med Besvarelsen af
det Spørgsmaal: \emph{hvorledes er Magtreflexionen mulig?} At
graviterende Masser optræde som Magthavere, der i Forhold til Størrelse
udøve et vist Herredømme over hverandre, er kun en figurlig Talemaade,
saalænge Kjendsgjerningens Mulighed ikke er indseet. Det er, som
allerede viist, Vexelbestemmelsen af Led og Momenter, hvorpaa det
kommer an. Efter det Ydre, i Udstykningens og Deelthedens Form, ere
Mo\-%
% --- p. 33 ---
menterne virkelige Masser; deres Reflexion er Relation, og Relationen
Vexelvirkning. Men i Begrebet, efter det Indre, ere de virkelige Masser
Momenter af samme Eenhed, Reflexer af een og samme Begriben. At begribe
er at magte. Det er da klart, at den almægtige Begriben, der forener
alle særskilte Magtbegreber i det almene Begrebs Eenhed, maa foregaae i
en Alt begribende Selvmagt. Kun en uendelig Selvmagt, et vidende Selv,
der objectiverer i Magt, formaaer at give de objective Begrebers System
Virkelighed i Massernes Existens og hæve sig over Stykværket.
Reflexionen af Magt i Andet har sin Grund i reflecterende Selvmagt;
hermed er Materiens absolute Afhængighed af Aanden erkjendt i Begrebet.

Videnskabens Historie lærer, at alle naturphilosophiske Forsøg paa at
opklare Forholdet mellem Materie og Bevægelse, Materie og Kraft, hidtil
ere strandede paa en total Miskjendelse af Forholdet mellem Aand og
Materie. Materiens Eenhed med Kraften er nemlig reflecteret saaledes,
at den, foruden at være en logisk Eenhed i Væsen, tillige er en virkelig
Eenhed i Kjendsgjerning. Reflexionen, hvori den dobbeltsidige Eenhed
begrundes, er selv dobbeltsidig; det er en Reflexion af Materiens og
Kraftens Væsen; men det er med det Samme en Reflexion af den existerende
Materie og Kraftens Yttring. Sættes Eenheden i Væsen, som naar man
siger: Materie og Kraft ere i Grunden det Samme, saa brydes den i
Existens; thi den existerende Materie er ikke Kraft, og den virkende
Kraft er ikke Materie. Sættes Eenheden i Existens, som naar man siger:
Kraften er materiel (den atomistiske Opfattelse) eller: Materien er et
Product af Kræfter (den dynamiske Opfattelse), saa kommer
Væsensforskjellen tilsyne i Reflexionen; thi Materien er et Ydre,
Kraften et Indre. Skal Eenheden erkjendes i Dobbeltreflexion af Væsen i
Væsen, af Kjendsgjerning i Kjendsgjerning, da er det kun muligt derved,
at Selvmagt og Magt i Andet, Væ\-%
% --- p. 34 ---
sensmagt i Idealitet og Gjenstandsmagt i Realitet, afgive saa stærke,
saa haarde Extremer, at Aandens Almagt kan vise sig som absolut
overgribende Enemagt.

Uden Selvhed ingen Reflecteren, og omvendt; uden et reflecterende Selv
ingen Reflexion; uden Reflexion ingen Relation; uden Reflexion i Magt
ingen virksom Relation, ingen Aarsag og Virkning, ingen Vexelvirkning.
I Begribningens Form er Magtreflexionen subjectiv; i Begrebets, det
Begrebnes, Form er den objectiv; i Reflexion af subjectiv og objectiv
Reflexion er Objectiveringen fuldendt.

\subrunin{α) Subjectiv Magtreflexion.} Selvmagt er ikke Magtens tomme,
% Bogens første græske Underoverskrift. Bogstavet er læst paa Siden selv
% (KB-Exemplaret): et utvetydigt α, ikke det „a“, som Indholdets OCR giver.
% Overskriften er indløbende, sat i samme Linie som det Afsnit, den aabner,
% og spærret; Punktum efter „Magtreflexion“, ikke Komma.
blot formelle Lighed med sig selv uden indre Forskjel og Modsætning, som
naar den uendelige Magt uden videre antages for værende og bestemmes ved
en tautologisk Sætning: Magt $=$ Magt, $A = A$. Nei, Selvmagten maa
udtrykkelig sætte sin egen Lighed og derved gjøre den til Selvlighed.
Dette kan i Virkeligheden kun skee derved, at Magteenheden viser sig
forskjellig fra sig selv, fordobler sig gjennem Andet og bekræfter sig
ved Modsætning til Andet. Kun den Bekræftelse, hvori det samme Væsen er
i Resultatet sin egen Forudsætning og i Forudsætningen sit eget
Resultat, kan i Sandhed kaldes \emph{Selv}bekræftelse, en Bekræftelse af
\emph{Selvet}.
% Kun Forleddet „Selv“ er spærret i „Selvbekræftelse“; „bekræftelse“ er
% sat med almindelig Skrift. Ordet er ét Ord, ikke to.

Jo større Selvstændighed det Andet har, desto mægtigere er det sig
derved bekræftende Selv, thi Selvet reflecterer sig subjectivt ved at
magte sit Andet: jo stærkere Modstanden er, desto større er den Magt,
hvormed den ophæves. Heraf sees, at Selvmagten ikke kan opfattes som et
umiddelbart Værende, en absolut, mod alt Andet ligegyldig Substans.
Selvmagtens Væren er Virken, og denne Virken en uendelig Reflecteren af
sig selv i Andet, af Andet i sig selv. Den indre Bevægelse, bvori den
% TRYKFEIL: „bvori“ trykt for „hvori“ — en urigtig Type, b for h. Bugten
% er lukket og svarer nøiagtig til b'et i „bestaaer“ i samme Linie; begge
% Exemplarer (Arbeidsscanningen og KB's Farvescanning) har samme Form, saa
% Feilen er Sætterens. Trykfeilsbladet retter den ikke. Gjengivet som trykt.
evig bestaaer, er en Kredsbevægelse. Den absolute Selvmagt forlader ikke
sig selv, fordi den gaaer ud fra sig selv; den naaer ikke hen
% --- p. 35 ---
til sit Andet som til et udvortes Givet, thi den har sit Andet i sig
selv; den forandres ikke, thi den magter sit Andet; Bevægelsen er ingen
Uro, thi Selvet er i sin Bevægelse hos sig selv, dets Bevægelse er
Hvile; det hviler i sig selv paa sig selv: \emph{den absolute Selvmagt
er Aand.}

For Aanden som Aand er hele den materielle Verden et Intet og dog et
virkeligt Noget. Begge Bestemmelser forliges i den Grundbestemmelse, at
den materielle Verden er Aandens Andet. Som Andet for Aanden er det
Materielle Moment for den almægtige Selvreflecteren, men Intet, slet
Intet i sig selv: kun Aanden, det Ene, er, alt Andet er Intet. Som
Andet for Aanden er det Materielle virkeligt, det er jo
Gjenstandsmagt, Selvbekræftelsens indre Modstandsmagt, Maalestok for
Aandens Almagt. Var Materien et Skin, og Verden en Drøm, saa var Aanden
selv en Drøm, saa var der ingen, enten indvortes eller udvortes,
Realitet, ingen, enten ideel eller reel, Virkelighed.

\subrunin{β) Objectiv Magtreflexion.} Den umiddelbare og dog
% Bogstavet læst paa Siden (KB-Exemplaret): β, ikke „B“, „ß“ eller „P“,
% som Indholdets OCR giver. Indløbende Overskrift som α) ovenfor.
reflecterede Eenhed af Materie og Kraft er en Magteenhed. Men den
objective Eenhed af Magt er tillige en brudt Magt, Magt i Andet,
Mangfoldighedsmagt; herpaa beroer dens eiendommelige Reflexion.
Reflexionen er kjendelig paa Tyngden. Physiken opfatter Tyngden fra tre
Synspunkter: som \emph{Egenskab} ved Materien, som særegen \emph{Kraft}
og som en Universet omspændende \emph{Magt}. Som Egenskab er Tyngden
reflecteret i Massernes Selvstændighed, som Kraft i deres indbyrdes
Vexelvirkning, som Verdensmagt i de store Love. Det er fra alle tre
Sider samme Væsen, samme Magt, samme Substantialitet; hver Side af
Magten reflecterer de to andre saaledes, at den selv maa siges at være
den hele Magt. Have vi for os de virkelige Masser, saa have vi i deres
fælles Egenskab Kraften, i Kraften Tiltrækningen og i Tiltrækningen den
som Magt almindelige Tyngde. Ere de tiltrækkende Kræfter givne, saa
have vi med det Samme
% --- p. 36 ---
Masserne, thi Masserne ere jo i denne Belysning intet Andet end
Kræfternes Existens, og Kraft og Masse igjen Momenter af Tyngden som
Magt. Have vi endelig Tyngden som Magt, da er Massernes Existens med
den tilsvarende Kræfternes Fordeling netop dens egen Realitet, Udslag af
det Særskilte, hvori den bestaaer som det Almene. Fra alle Sider er
altsaa den objective Magt reflecteret.

Men den objective Magt kan ikke reflectere sig selv ved sig selv; det
kan kun Selvmagten. Lad os igjen see paa de to Masser. Det bevægende
$A$ er Aarsag, Bevægelsen af $B$ Virkning. Aarsagen, der er Realitet i
$A$, er Idealitet i $B$, thi Aarsagens Idealitet er at være ---
reflecteret --- i sin Virkning. Da nu $A$ er fjernet fra $B$, saa har
det jo som Aarsag sin Virkning udenfor sig, er som Realitet fjernet fra
sin Idealitet. Og dog er Idealiteten i Aarsagen, og Aarsagen er i $A$.
$A$ kan da umulig være fjernet fra, hvad det selv er; det maa, saa sandt
det har Aarsagen i sig, nødvendig ogsaa have Idealiteten i sig. Denne
Modsigelse, at nemlig hver Masse har sin Idealitet i sig ved at have den
udenfor sig, er en Modsigelse i den objective Reflexion, som klarlig
beviser, at den ikke kan bære sig selv, men maa bæres af Aanden. Fra
Aanden, i og ved Beaandelsen, faaer Massens Realitet en Tilgift, hvori
det Ideelle objectivt og væsenlig ener sig med det Reelle, denne
ideel-reelle Tilgift er Kraften. Formeldelst Kraften er Materien i $A$
% TRYKFEIL: „Formeldelst“ trykt for „formedelst“ — et Bogstav l for meget.
% Overstaven mellem „Forme“ og „delst“ staaer tydelig i begge Exemplarer.
% Trykfeilsbladet retter den ikke. Gjengivet som trykt.
efter Mulighed mere end Materie, den er Aarsag. At Virksomheden
indtræder, at Aarsagen bliver virkelig og faaer sin Virkning, er
objectiv Reflexion i Magt. Dersom den objective Magt ikke blev
reflecteret i Aanden, saa var Tyngden ikke den almene Magt, saa var der
ingen almene Magter i Naturen, saa vare Lovene umulige. Den objective,
almene Magt er ikke en Abstraction af Masser og Kræfter, men deres egen
substantielle Væsenhed, deres iboende Grundlag. Den objective Magt
reflecterer sig --- ved at blive reflec\-%
% --- p. 37 ---
teret; den er i Virkeligheden objectiv, men --- ved at være
objectiveret.

\subrunin{γ) Objectivering i Magt.} Naturbegrebene have dobbeltsidig
% Bogstavet læst paa Siden (KB-Exemplaret): γ, ikke „y“ eller „7“, som
% Indholdets OCR giver. Indløbende Overskrift som α) og β) ovenfor.
% „Naturbegrebene“ staaer saaledes trykt her; længere nede i samme Afsnit
% skriver Bogen „Naturbegreberne“.
Bestaaen, Bestaaen i virkelige Ting og Bestaaen i Selvets Begriben. Det
Sidste betinger det Første; Begreberne kunde ikke give sig Afpræg i
tilværende Gjenstande, dersom de ikke selv vare prægede af uendelig
Begriben. Det er jo ikke de sandselige Naturting, som i objectiv
Forstand bestemme Begreberne, men det er omvendt Naturbegreberne, som
usandseligt bestemme de sandselige Ting. For at være Tingens
bestemmende Væsen, maa Begrebet i sig selv nødvendig have Bestemthed,
maa være forud bestemt til at afgive Bestemmelse. Massens bestemte
Begreb bliver ikke først til ved Hjælp af existerende Masser, ei heller
Kraftens Begreb ved allerede tilværende Kræfter. Hvis Masser og Kræfter
havde en af Begreber uafhængig Tilværelse, saa vilde Tingene med det
Samme kunne undvære alle Begreber; der vilde slet ingen Fornuft være i
Naturen. Hvis Begrebernes Tilblivelse faldt umiddelbart sammen med
Tingenes, vilde det være aldeles umuligt at sige, hvorfra det
Bestemmende kom. Det kunde da ikke være Begrebet, der bestemte Tingen,
thi Begrebet skulde jo netop blive bestemt ved Tingen; men at det heller
ikke kunde være den i sig bestemte Ting, der bestemte Begrebet, er en
Selvfølge, eftersom det for en Ting Bestemmende just er dens Begreb.
Det staaer altsaa fast, at Naturbegreberne maae have en af Tingenes
Existens uafhængig, for Tingene selv afgjørende Bestemthed i sig.
Hvorfra faae nu Begreberne denne indre, for al ydre Tilværelse
uundværlige Bestemthed? Fra en Begriben. At bestemme Begreber er at
begribe; at fremstille de i sig bestemte og desaarsag bestemmende
Begreber i Existeren og Existens er ikke at \emph{begribe}, det er at
\emph{virke}. Naturen fremstiller Begreber, udformer Begreber, men den
begriber ikke, den virker. Aandens ideelle Virken
% --- p. 38 ---
er Begriben, Naturens Gjengivelse af det Begrebne er ikke Begriben, det
% Komma, ikke Punktum, efter det første „Begriben“; efterprøvet paa
% KB-Exemplaret. Sætningsrækken er komma-føiet, saaledes som trykt.
er reel Virken, materiel Virksomhed, Begrebsbestemmelsernes Udfoldning i
sandselig-virkelige Kjendsgjerninger.

Modsætningen af Aandens Begriben og Naturens Virken er begrundet i
Reflexionen af subjectiv og objectiv Magtreflexion, et alsidigt Udtryk
for Objectivering i Magt. I den subjective Magtreflexion, den ideelle
Objectivering, have Naturbegreberne Tilværen som aandelige
Magtbestemmelser, afgjørende for al materiel Existeren. Men under
Idealitetens Form, i Selvmagtens Begriben, have Begreberne ikke selv
Existens, deres Tilværen, en Væren i Viden, er Magtbestemmelsernes
uendelige Vexelskin af hverandre indbyrdes; det Enkelte og Særskilte
reflectere sig her i det klare, frie Almene: Magt i Andet er Moment for
Selvmagten. Fra Existenssiden, altsaa fra Natursiden, falder
Eftertrykket derimod paa Magten i Andet, paa Gjenstandsmagten, paa
Realiteten, og den reelle Objectivering afsætter Begreber i tilsvarende
Objecter. Objecterne ere Gjenstande, en Vrimmel af Enkeltled, Atomer og
Masser. Det som Begreb bestemmende Almene gaaer her umiddelbart op i
det Enkelte, bindes i Magt af Enkeltled, af udvortes Objecter. Saa
hedder det i Physikens Sprog: kun Enkeltleddene, Grunddelene med de
deraf sammensatte Ting, have i Naturen Existens; det Almene existerer
ikke. Og dog er det just det Almene, der kommer til Existens; det
bliver jo netop gjennem Udslag af Enkeltled objectiveret i Magt.
Enkeltleddene virke, antage ved deres Vexelvirkning Præg af det
Særskilte og hjemfalde til det Almene; saa viser det sig, at det Enkelte
og det Særskilte ere det Almene underlagte. Den objectivt reflecterede
Magt, Naturens, Objectivitetens almene Magt, kommer til Existens i
Tingenes Kræfter og de Love, hvorunder Kræfterne virke.

% Efter sidste Linie paa s. 38 staaer en kort, centreret Streg, som
% afslutter Afdeling B. Findes i begge Exemplarer. Gjengivet.
\streg

%  --- C. Kræfter og Love  (pp. 39–61) ---

%  MATHEMATICAL — work from the image; the OCR is not a witness here.
% --- p. 39 ---
\afdeling{C.}{Kræfter og Love.}

I objectiv Reflexion er den almene Magt baade Existensen selv og
Existensens Væsen; men Forskjellen mellem Væsen og Existens er lige saa
mægtig som Eenheden. Magten i Andet, der er Eet med sig selv i Væsen,
er forskjellig fra sig selv i Existensen; den bestaaer med denne
Forskjellighed som Gjenstandsmagt i Leddenes Masser og som virkende
Kræfter i Leddenes Vexelvirkning; men den har ogsaa Bestaaen i
Væsenhedens almene Form som hvilende, Alt bestemmende Magt, nemlig i
Lovene. Kræfternes Yttringer ere Kjendsgjerninger, men Lovene see ud som
Tanker, mægtige Tanker, Naturtanker. At kjende de virkende Kræfter og
Betingelserne for deres Virksomhed er at kjende Lovene; og dog ere
Kræfterne ikke Love, og Lovene ikke Kræfter. For at komme til Indsigt i
den almene Forskjelseenhed af Kræfter og Love, skulle vi nu fra et
mechanisk Synspunkt undersøge Lovenes
% Spærret Sats, efterprøvet paa KB-Exemplaret: hele Leddet fra
% „Oprindelse“ til „Sammenhæng“ er spærret. Punktum, ikke Komma, efter
% „Sammenhæng“; tesseract læser her Komma.
\emph{Oprindelse, Opfyldelse og indbyrdes Sammenhæng}.

\runin{a) Forestillede Kræfter: Lovenes Oprindelse.}%
„Forstanden“, siger Kant, „laaner ikke sine (aprioriske) Love fra
Naturen, men den foreskriver den disse.“ Denne Sætning viser hen til
Lovenes subjective Oprindelse og retfærdiggjøres ved talløse Exempler
fra den analytiske Mechanik. Dersom to lige store Kræfter virke paa et
materielt Punkt i modsat Retning, maae de ophæve hinanden, og Punktet
forblive ubevæget; det er en Lov. At denne Lov har objectiv Gyldighed,
er lige saa indlysende, som at dens Gyldighed er uafhængig af al
Erfaring. Lovens almene Gyldighed ligger i den Nødvendighed, hvormed den
opstillede Forudsætning afgiver sin Følgesætning: Loven er i logisk Form en
% --- p. 40 ---
hypothetisk Dom, i mathematisk Form en Ligning. Der ligger ingen Vægt
paa, hvorvidt det antagne Tilfælde forekommer i den ydre Natur eller
ikke, thi Loven er Sagens Natur; den har sin Natur i sig selv. Det er
det tænkende Selv, der opstiller Forudsætningen; det er det tænkende
Selv, der drager Conseqvensen; det er det tænkende Selv, der giver sig
selv Loven og forvisser sig om dens Objectivitet. Der er hverken i
Himlen eller paa Jorden nogen udvortes Magt, der formaaer at omstøde den
i Tankens Medfør givne Lov eller gjøre den tvivlsom; den eneste
Magthaver, der er i Stand til at ophæve Loven, er den lovgivende
Subjectivitet. Vel kan Subjectiviteten ikke modsige Loven uden at
modsige sig selv; men den kan ophæve den ved at ophæve Forudsætningen,
den kan forandre den ved at forandre Forudsætningen.

At en Lov maa forandre sig, naar dens Forudsætning forandres, er selv en
Lov. Dersom de to i modsat Retning virkende Kræfter ere ulige store, maa
Punktet komme i Bevægelse, og Kræfternes Differens vil da være den
bevægende Kraft; men dersom de to Kræfter danne en hvilkensomhelst
Vinkel med hinanden, vil det ved første Øiekast see ud, som havde vi
faaet en ny Lov. Den Nødvendighed, hvormed Loven gjælder, er nemlig
altid bunden til en Forudsætning af een eller flere Betingelser, og
følgelig ikke en absolut, men hypothetisk Nødvendighed. Ved at indføre
snart færre, snart flere Betingelser i Forudsætningen og lade
Betingelserne vexle kan den lovgivende Subjectivitet lade den ene Lov
fremgaae af den anden, kan paa denne Maade anlægge og udvikle Systemer
af Love. Vi see det paa Kræfternes Parallelogram. Man skulde troe, at
her blev opstillet en egen Lov for et eget specielt Tilfælde; men den
Tænkende opdager snart, at det tilsyneladende specielle Tilfælde netop
er det almindelige Tilfælde, thi to paa eet
% --- p. 41 ---
Punkt virkende Kræfter maae jo altid danne en Vinkel med hinanden.%
% Fodnotemærket er *), som Bogens øvrige. Een Fodnote paa denne Side.
% Formlerne ere satte efter Billedet, ikke efter OCR'en. Rodtegnets
% Streg dækker hele Udtrykket, ogsaa „cos (P. Q)“; i KB-Exemplaret er
% Stregen sat af to Stykker, hvilket ikke er gjengivet.
\footnote{De to Kræfter være $P$ og $Q$ og deres Resultant $R$. For
Kræfternes Parallelogram gjælder da Ligningen:
$R = \sqrt{P^2 + Q^2 + 2\,P\,Q \cos (P.\,Q)}$. Er nu Vinkelen
$(P.\,Q) = 0^{\circ}$, saa er $\cos = 1$, og $R = P + Q$; er Vinkelen
$(P.\,Q) = 180^{\circ}$, saa er $\cos = -1$, og $R = P-Q$.
Parallelogrammets Lov er en Grundlov for Sammensætning og Opløsning af
mechaniske Kræfter, hvoraf et uendeligt System af særskilte Love lader
sig udvikle.}

Bevægelsen efter Diagonalen i Parallelogrammet er, væsenlig seet, en
sammensat Bevægelse, thi den Kraft, hvorved den bestemmes, er en
Samvirken af to Sidekræfter; men ligesom Forskjellen mellem enkelt og
sammensat Kraft kun er relativ, saaledes er ogsaa Forskjellen mellem
enkelt og sammensat Bevægelse relativ. Den ved Diagonalen betegnede
Resultant kan i sin Fortsættelse selv blive en Sidekraft, hvis Vinkel
med en anden tilsvarende afgiver en ny Resultant i en ny Retning, saa at
Bevægelsens Spor ved den endelige Forandring bliver en brudt Linie.
Skeer Forandringen continuerlig saaledes, at Sidekræfterne
(Composanterne) uophørlig omsættes i en Resultant, og Resultanten igjen
i en Sidekraft, saa fremkommer der en krum Linie.

Det sande Begreb om Bevægelsen i den krumme Linie, hvis Lov Mechaniken
nøiagtig bestemmer, hviler paa en Opløsning af to Lag Misforstaaelse. De
Gamle meente, at det hørte til visse Legemers eiendommelige Væsen at
bevæge sig i Curver, at navnlig de himmelske Legemer, ifølge deres
høiere Natur, bevægede sig i Cirkelbaner. Det var en Misforstaaelse. Ved
at opløse denne Misforstaaelse blev Physiken en exact Videnskab. Men for
at tydeliggjøre den sammensatte Bevægelse, maatte man tye til
Udstykning. Under Udstykningen bliver den krumme Linie til en brudt
Linie: Brudstykkerne være saa smaa, de ville, det bliver
% --- p. 42 ---
dog altid Brudstykker. Herved fremkom en ny Misforstaaelse.
Udstykningsmethoden har i den Grad forarget den nyere
Naturphilosophie, at Speculationen aabenlyst har brudt med den
mechaniske Physik og søgt Sandheden i en modificeret Opfattelse af den
antike Forestilling. Hegel er saa vred over Udstykningen, at han ikke
blot forbyder Physiken --- hvad han dog i rigeligt Maal indrømmer
Logiken --- en egen Brug af Abstractionen, men han gaaer endog saa vidt,
at han vil have Forestillingen om Kræfter ganske udryddet. „Newton“,
siger han, „forsikkrer, at en Blykugle \textit{in coelos abiret et motu
abeundi pergeret in infinitum}, dersom (rigtignok: dersom) man blot
kunde meddele den den fornødne Hastighed. En saadan Adskillen af den
udvortes og den væsenlige Bevægelse tilhører hverken Erfaringen eller
Begrebet, men alene den abstraherende Reflexion.“ Ved at indføre
„Forestillingen om Kræfter“ skal Newton trods alle sine Reservationer,
have gjort den sande Naturerkjendelse et betydeligt Afbræk. Det er
speculative Misforstaaelser. Disse Misforstaaelser opløse sig og
forsvinde, naar Forholdet imellem Continuiteten, Udstykningen og den
uendelige Discretion tilstrækkelig opklares.

Kastes et Legeme i skraa Retning op efter, da kan den til Stødet
svarende Kraft betragtes som en Sammensætning af to Kræfter, af hvilke
den ene vilde føre Legemet lodret opad, efter Y-Axen, den anden drive
det vandret fremad, efter X-Axen. Dersom nu enhver anden Paavirkning
tænkes borte, vil Legemet gaae efter Diagonalen i det Rectangel, som
Coordinaterne x og y danne med hinanden, og den i Tiden t gjennemløbne
Bane s være ligestor med $\sqrt{x^2 + y^2}$. Den af s og x dannede
Vinkel $\alpha$ vil være constant, og dens Tangens $\frac{y}{x}$
selvfølgelig ogsaa. I et saadant Tilfælde er det da ligegyldigt, om man
fremhæver Continuiteten eller Udstykningen eller Discretionens Uende\-%
% --- p. 43 ---
lighed. Hvert endeligt Stykke af s er et Continuum, og de vilkaarlig
satte Grændser gjøre intet Afbræk; men s kan ogsaa opløses i Elementer
svarende til Elementer af x og y:
\[
ds = \sqrt{dx^2 + dy^2};\quad \frac{y}{x} = \frac{dy}{dx}.
\]

Anderledes, naar Legemet under Bevægelsen forestilles at være en lodret
nedadgaaende Paavirkning underkastet; her viser der sig en væsenlig
Forskjel mellem Udstykning og Continuitet.

Forestilles Paavirkningen at indtræde efter visse Mellemrum, saa vil s
adskille sig i endelige Stykker, $s_1$, $s_2$, $s_3$ o. s. fremdl.;
Vinkelen $\alpha$ vil hver Gang forandre sig, idet hvert Stykke af s
faaer en egen Heldning mod Abscisseaxen, og $\frac{y}{x}$ erholde
forskjellige Værdier. Ethvert Stykke er Diagonal i et Kræfternes
Parallelogram, hvis ene Side er en Fortsættelse af den forrige
Resultant, og hvis anden Side dannes af den tiltrædende Paavirkning:
Legemets Bane bliver en brudt Linie, hvis Dele kunne maales og nøiagtig
bestemmes.

Ombyttes den deelvise Paavirkning med en continuerlig som f. Ex. Tyngden
g, saa forandrer s sin Retning continuerlig, og Banen bliver en krum
Linie. Ligefrem i Continuitetens Form kan en saadan Bane ikke bestemmes,
thi Forandringens Uendelighed for sig kan ikke bestemmes. Ligefrem i
Udstykningens Form kan Banen heller ikke bestemmes, thi Udstykningen er
en endelig Discretion, der bryder Continuiteten. Kun den
Vexelbestemmelse af Endeligt og Uendeligt, der betinges af uendelig
Discretion, hvorved en Operation med Elementer kan blive ligesaa exact,
som en Regning med endelige Størrelser, kan løse Udstykningen op i en
endelig Bestemmen af det Uendelige.

% Bogen skriver „tg“, ikke „tang“ eller „tan“; gjengivet som trykt.
% Punktummet efter α er Bogens eget Multiplicationstegn, ogsaa som trykt.
Ligningen: $y = f\,(x) = \operatorname{tg} \alpha .\, x -
\frac{g}{2\,a^2 \cos^2 \alpha}\, x^2$
% --- p. 44 ---
viser, hvorledes den mindste Forandring i x continuerlig medfører en
tilsvarende Forandring i y; hvoraf igjen følger en continuerlig
Forandring i s. Den Maade, hvorpaa y afhænger af x, udtrykker Loven. Vi
have i Loven det for Phænomenernes Uendelighed --- den uendelige Mængde
af enkelte Tilfælde --- bestemmende Almeenvæsen. Men i denne Uendelighed
have vi tillige Endeligheden. Paa væsenlig Maade er Endeligheden
bestemt ved den bestemte Function; i Tilsyneladelsernes udvortes, paa
Deling og Udstykning anlagte Form kommer det Endeliges Uendelighed
tilsyne i Enkeltværdiernes Vrimmel. Overgangen fra Endeligt til
Uendeligt og fra Uendeligt til Endeligt bestemmes ved Elementernes
indbyrdes Forhold. Hvert Element af Curven falder sammen med et
tilsvarende Element af Curvens Tangent; hvorledes, d. e. efter hvilken
Lov denne Tangent i det Uendelige forandrer sin Retning, sees af
Ligningen: $\frac{dy}{dx} = \operatorname{tg} \alpha -
\frac{g}{a^2 \cos^2 \alpha}\, x = \varphi\,(x)$.

Hvad der gjælder om Loven, gjælder ogsaa om Kræfterne. Kræfternes
Opløsning er ingen ligefrem Udstykning; det er en Analyse, hvorunder
hver Sidekraft ved en for Videnskaben nødvendig Abstraction bliver
fremhævet for sig. Fra Væsenets Side er enhver af Sidekræfterne et
Moment af Resultanten. Hvor en Udstykning virkelig finder Sted, er det
Naturens egen Udstykning, Existensen selv, som i Andethedens Form gjør
den ene Kraft uafhængig af den anden. Heller ikke er Kræfternes
Sammensætning en simpel Addition; den i Resultanten virkende Kraft er
Sidekræfternes Eenhed, en Eenhed, hvori deres Selvstændighed ganske er
ophævet. Her er ikke, hvad dog Hegel paastaaer, mindste Grund til, paa
Begrebets Vegne, at gjøre nogensomhelst Indsigelse mod Physikens
Analyser.

Men er det nu, strengt taget, Lovene, der bestemme
% --- p. 45 ---
Kræfterne, eller kunde man ikke maaskee vende Forholdet om og sige, at
det egenlig er Kræfterne, der bestemme Lovene? En overfladisk
Opfattelse af den Maade, hvorpaa den rationelle Mechanik fremstiller
sine Opgaver, synes nærmest at tale for det Sidste. Et Exempel. Der
meddeles et Legeme en vis Hastighed i en hvilkensomhelst, kun ikke
lodret, Retning, hvorpaa det overlades til Tyngdens Paavirkning: af
hvilken Beskaffenhed vil Banen blive? Svaret er, at Banen, naar
Bevægelsen antages at foregaae i det lufttomme Rum, maa blive en
Parabel. Loven for Parablen er altsaa Loven for Bevægelsen. De
forestillede Kræfter og de opstillede Betingelser synes her at være det
Første, og Loven det Andet, som deraf afledes; altsaa: Kræfterne
bestemmes ved Forestilling, Loven ved Tænkning, og det Tænkte afhænger
af det Forestillede. Hvor eensidig denne Synsmaade er, fremlyser af
Forholdet mellem Forestilling og Tænkning. Den blot forestillende
Bevidsthed erkjender slet ikke Loven. Den kan ikke engang uden Tankens
Bistand stille det Forestillede op som en Forudsætning; thi
Forudsætningen er kun i Reflexion af sin Conseqvens, men Conseqvensen
er kun for Tanken. En Bestemmen af Kræfterne ved den blotte Forestilling
er en Indbildningsact, en flygtig, foreløbig Bestemmen, der taber al
Betydning, naar Tanken mangler. Den bestemmende Forestilling er en
Function af Tanken; den forudskikkes af Tanken selv og staaer fra
Begyndelsen af i Tænkningens Tjeneste.

Fra Mechanikens Synspunkt faaer Sagen let et Udseende, som om Tankens
Opgave udelukkende var at opdage Lovene, som om Lovene i og ved sig selv
vare objective uden at trænge til Objectivering. Objectiveringen skulde
altsaa blot være for Subjectivitetens, ikke for Lovenes Skyld. Det er en
Illusion, som den strenge Videnskab øieblikkelig vil forkaste, naar den
bliver opmærksom paa den. Dersom den lovopdagende Subjectivitet ikke
selv var lovgivende, kunde
% --- p. 46 ---
den slet ikke opdage Lovene, den kunde jo ikke dømme om Betingelserne,
ikke adskille det Mulige fra det Umulige, ikke bestemme det hypothetisk
Nødvendige, ikke forvisse sig om de foregivne Loves aprioriske
Gyldighed. Og hvad bliver der af de objective Love, naar de skilles fra
den lovgivende, objectiverende Subjectivitet? Døde, meningsløse Formler!
En Ligevægtslov for to ved en Linie forbundne Punkter, der paavirkes af
Kræfterne $P_i$ og $P_m$, hvis Vinkler med X-Axen ere $\alpha_i$ og
$\alpha_m$, idet den ubekjendte Kraft i Forbindelseslinien er $R$, dens
Vinkel med X-Axen $\delta$, $\varphi_i$ og $\varphi_m$ vilkaarlige
Vinkler, præformeres for et System af flere indbyrdes forbundne Punkter
ved Ligningerne:
% Anden Ligning har ÷ hvor den første har +. Efterprøvet paa
% KB-Exemplaret ved høi Forstørrelse: Tegnet er utvivlsomt ÷, og begge
% Exemplarer stemme. ÷ var i dansk Sats paa denne Tid det gængse
% Minustegn, saa dette er formodenlig Nielsens egen Notation og ikke en
% Sætterfeil; Trykfeilbladet nævner ikke Stedet. Gjengivet som trykt.
% Prikkerækkerne staae som fire spatierede Punkter i begge Ligninger.
\begin{align*}
P_i \cos (\alpha_i - \varphi_i) + R \cos (\delta - \varphi_i) + \ .\ .\ .\ . &= 0\\
P_m \cos (\alpha_m - \varphi_m) \div R \cos (\delta - \varphi_m) + \ .\ .\ .\ . &= 0
\end{align*}
Men saa lidt som disse Ligninger i sig selv vilde have Betydning, dersom
Ingen forstod dem, ligesaa lidt vilde den ved Ligningerne præformerede
Lov i sig selv have mindste Gyldighed, dersom der ingen, hverken
menneskelig eller guddommelig Subjectivitet var, der kunde tænke Loven,
med selvlovgivende Forvisning opdage Loven, objectiverende hævde Lovens
Objectivitet.

\runin{b) Virkelige Kræfter: Lovenes Opfyldelse.} Fra det subjective
Synspunkt, forsaavidt der blot handles om forestillede Kræfter, er
Overeensstemmelsen mellem den guddommelige og menneskelige Fornuft
gjennemgaaende; thi Naturlovene ere Tankelove, og Tankelovene paa een
Gang baade menneskelige og guddommelige. Men uagtet alle Naturlove ere
Tankelove, ere Tankelovene i deres frie Idealitet dog endnu ikke
Naturlove. Opfattede for sig, have de tænkte Love kun en ideel, ved
Forestillingerne articuleret Objectivitet; i Naturen selv have derimod
Lovene en til Kræfterne svarende materiel Virkelighed; denne Virkelighed
er deres Opfyldelse.

Fra det objective Synspunkt, forsaavidt der udtrykkelig
% --- p. 47 ---
sees paa Naturvirkeligheden, kommer Dualismen, den principielle Forskjel
mellem guddommelig og menneskelig Objectivering igjen tilsyne. Fra den
guddommelige Side maae vistnok alle Kjendsgjerninger kunne afledes af de
bestemmende Love; men fra menneskelig Side maae de virkelige Love
udledes af givne Kjendsgjerninger.

Vel kunne de i factiske Natursætninger som i et guddommeligt Magtsprog
udtalte Love, hvor eiendommelige de end ere, ikke tænkes at staae i
Strid med de almindelige Fornuftlove; men det gjør dog en væsenlig
Forskjel, om Lovene behandles som abstract-almene Formbestemmelser,
altsaa \emph{mathematisk}, eller de paavises i anskuelig Eenhed med det
Virkelighedsindhold, hvori de have deres Opfyldelse,
% Trykfeil: „alsaa“ for „altsaa“. Bogstavet t mangler. Efterprøvet paa
% KB-Exemplaret, hvor Ordet ligeledes staaer „alsaa“, medens det
% samme Ord tre Linier ovenfor staaer „altsaa“; Feilen er altsaa
% Sætterens og ikke dette Exemplars. Bogens eget Trykfeilblad retter
% den ikke. Gjengivet som trykt.
alsaa \emph{physisk}.

„Ligesom Mathematiken“, siger H. C. Ørsted, „maa behandle
Naturgjenstandene paa sin egen Viis, saaledes ogsaa Physiken. Idealet
for hiin er at udrette Alt uden at laane Noget af Erfaringen, og naar
den, ifølge Umuligheden af ganske at naae Idealet, gjør et Laan af
Erfaringen, maa dette være noget saa almindeligt, saa lidet
individualiseret som muligt. Det hører til dens Kunstfuldkommenhed,
hellere at gaae de vidtløftigste Omveie, end at optage Erfaringens
fremmede Stof i sig. Ganske anderledes forholder det sig med
Naturlæren. Naturen er dens videnskabelige Eiendom; den laaner ikke
deraf, men øser deraf som af sin egentlige Kilde. Dens Ideal er det at
fremstille Gjenstandene for Anskuelsen“ --- altsaa Lovene i deres
Opfyldelse --- „og med Hensyn herpaa bearbeider den Former og
Størrelser.“

Den videnskabelige Betydning, der tilkjendes Physiken som
Erfaringslære, er naturligviis afhængig af den Maade, hvorpaa
Erfaringens Væsen i det Hele bliver opfattet. Sand Erfaring beroer paa
en ligelig Samvirken af Sandsning og Tænkning; men denne Samvirken kan
forvanskes, eftersom
%  MATHEMATICAL — work from the image; the OCR is not a witness here.
% --- p. 48 ---
det enten er Sandsningen, der overvælder Tænkningen, eller
omvendt. Er Sandsningen eensidig fremherskende, saa fæster
Blikket sig udelukkende paa Kjendsgjerningerne i deres
Enkelthed. Vadende i Dynger af ophobede Kjendsgjerninger,
faaer den iagttagende Bevidsthed saa travlt med at sondre,
forbinde, addere og subtrahere de mange Enkeltheder, at den
neppe kan kjende Tanken i de Almeenforestillinger, den selv
har abstraheret. Er den aprioriske Tænkning eensidig
fremherskende, saa blive Kjendsgjerningerne uvilkaarlig nedsatte
til flygtige, forbigaaende Tilsyneladelser af det Almene. Uden
at gaae ind paa nogen udtømmende Undersøgelse af Stoffet i
dets Enkeltheder, vil Betragteren, fristet af en eller anden
Yndlingshypothese, kun altfor let udspinde en Theorie, og
udgive den for en paa Erfaring begrundet Sandhed. Mod begge
Udskeielser maa Physiken, den sande Erfaringslære, føre en
stadig Kamp.

Det naturvidenskabelige Grundbegreb, hvorpaa Erfaringslæren
bygger, er Begrebet: \emph{Induction}.

At Inductionen maa blive ufuldstændig, naar Iagttagelserne
ere mangelfulde, forstaaes af sig selv. Men at en udvortes
Sammenhoben af blot sandselige Iagttagelser, om de end
mangfoldiggjøres efter en nok saa uhyre Maalestok, heller ikke
kan give nogen sand Naturerkjendelse, er da ogsaa indlysende.

Et af de grundigste Beviser paa overfladisk Opfattelse af
Inductionens logiske Gyldighed have vi i Stuart Mills inductive
% Citatet af Stuart Mill løber over Sidebruddet og lukkes først
% paa S. 49 („Aarsag og Virkning.“). Ingen Sætterfeil --- kun et
% Citat over to Sider; deraf een „ for meget paa S. 48 og een “
% for meget paa S. 49.
Logik. „Vi lære“, hos denne Philosoph, „af Erfaringen, at der
i Naturen er en uforanderlig Følgeorden, og at enhver
Kjendsgjerning altid foregaaes af en anden Kjendsgjerning. Hvad vi
kalde Aarsag, er det uforanderligt Foregaaende, hvad vi kalde
Virkning, er det uforanderligt Følgende. Begrebet
Nødvendighed, hvoraf man gjør saa meget Væsen, finder her sin
Forklaring. Nødvendigt er det, som vil indtræde, hvad vi saa end
antage med Hensyn til alle
% --- p. 49 ---
andre Ting. Nødvendigheden er tilstede, naar det Forudgaaende
er tilstrækkeligt og fuldstændigt, naar det indeholder alle de
søgte Betingelser, og ingen ny Betingelse udfordres. Ubetinget
at ville følge, det er hele Begrebet om Aarsag og Virkning.“
Forsaavidt en saadan Forklaringsmaade kan antages at have
noget Blændende ved sig, da ligger det i den Ugeneerthed,
hvormed Tanken som Tanke i et og samme Aandedræt paakaldes
og afvises. Seet i Tankens Lys, er det ubestridelig sandt, at
Forholdet mellem Aarsag og Virkning indeslutter Nødvendigheden,
at Følgen ubetinget maa indtræde, naar alle Betingelser ere
tilstede. Men Tanken faaer ikke Lov til at gjælde som Tanke;
det er ikke Tankens aprioriske Forvisning, men Sandsningernes,
de gjentagne Sandsningers Vidnesbyrd, som her i Erfaringens
Navn faaer noget at betyde. At nu en Sum af sandselige
Iagttagelser, om de end sigtes og drøftes nok saa sindrigt, umulig
kan sikkre os et „ubetinget at ville følge“, er i den Grad
indlysende, at det end ikke har kunnet undgaae Forfatterens egen
logiske Opmærksomhed. Stuart Mill „er overbeviist om, at en
Mand, som er vant til Abstraction og Analyse“, meget vel kan
„tænke sig, at paa visse Steder, f. Ex. paa et af hine
Firmamenter, i hvilke Astronomien nu inddeler Universet, kunde
Begivenhederne følge hverandre efter reen Tilfældighed uden
nogen fast Lov, og der er Intet, hverken i vor Erfaring eller i
vor Sindsbygning, der giver os tilstrækkelig Grund til at troe,
at dette ingensteds finder Sted.“ Ifølge en saadan Logik høre
Lovene altsaa ikke til Tingenes Væsen, men ere i visse Naturens
Egne tilfældig klistrede uden paa Tingene, omtrent som man
klistrer en Devise paa en Flaske. Hvis Hume stod op af sin
Grav, vilde han snart overbevise Mill om, at han slet ikke
behøver at uleilige sig med den lange Reise til „et af hine
Firmamenter“ for at faae Lovene bort og Nødvendigheden opløst,
% --- p. 50 ---
at det skeer langt lettere, naar han bliver, hvor han er, og
besinder sig paa, hvad han selv har sagt.

Der er Fornuft i Naturen; denne Hovedtanke, som Ørsted i
vor Litteratur saa smukt har udtalt og hævdet, giver Inductionen
en anden og dybere videnskabelig Betydning, end den Mill'ske
Forklaring; her gaaes da ud fra, at Lovene ligge i Tingenes
Væsen. Denne Hovedtanke maa have staaet klart for Newton,
da Ideen om den almindelige Tyngde gik op for ham. Det Øie,
der saae i det faldende Æble en Yttring af den Kraft, der ogsaa
drager Maanen ned mod Jorden, tilhører ikke Sandsningen, men
Tænkningen. Det var Tankens Anelse om en almindelig, alle
Legemer iboende Væsenhed, der gav Impulsen til den store
Opdagelse. Men med Tænkningen kom ogsaa Sandsningen til sin
Ret. Den almene Væsenhed, der var antydet i en Kjendsgjerning,
maatte igjen nærmere bestemmes ved Kjendsgjerninger. Kun
saaledes kunde den tænkte Lov bekræfte sig som en Naturlov, en
almeengyldig Erfaringslov; med andre Ord: Tyngdeloven maatte
sandses og sees i sin Opfyldelse.

Vexelvirkningen mellem empirisk Physik og rationel Mechanik
viser sig her i det skjønneste Lys. For at Iagttagelserne kunde
begynde, og Erfaringen indledes, maatte Opgaven sættes i
Ligning; men den Ligning, der skulde give Iagttagelserne Betydning,
maatte selv være bygget paa Iagttagelser.

At Maanens Fald maa forholde sig til Æblets, som den Kraft,
der drager Maanen, til den Kraft, der drager Æblet, er en
Selvfølge; men i hvilket Forhold staae nu disse Kræfter til
% Brøken staaer i Satsen som en stablet Brøk midt i Linien;
% Tallet 15 5/8 er Faldrummet i Fod.
hinanden? Den Kraft, der drager Æblet, er bekjendt ved
Faldloven --- $15\tfrac{5}{8}$ Fod i 1ste Secund; den Kraft, der drager
Maanen --- Maanens Fald i 1 Secund --- kan, da Maanens
Middelafstand og Middelomløbstid ere bekjendte Størrelser, ligeledes
bestemmes. Betegnes Maanens Afstand
% --- p. 51 ---
% Θ er stort Theta, ϱ det krøllede rho; begge læste paa
% KB-Exemplaret. De trykte Variable staae opret antiqua; efter
% Husreglen (se Hovedet) gjengives det ikke, men den almindelige
% mathematiske Kursiv beholdes.
fra Jordens Midtpunkt ved R, den til Buen (i et Tidssecund)
svarende Vinkel ved $\Theta$, saa er det Stykke, som betegner
Maanens Fald mod Jorden, $\varrho = 2 R \sin^{2} \tfrac{1}{2} \Theta$.
Kræfternes Forhold er altsaa bestemt ved Qvotienten
$\dfrac{2 R \sin^{2} \tfrac{1}{2} \Theta}{15\tfrac{5}{8}}$;
for at R kan bestemmes, maa tillige Jordradien r være bekjendt.
Her er en Kreds af forudsatte Iagttagelser, og dog kan den ikke
uden særlig Forberedelse beriges ved nye Iagttagelser. Vel laae
det nær at søge Grunden til Tiltrækningens Forskjellighed i
Afstandens Forskjellighed, men med denne almindelige Formodning
kunde endnu ingen Erfaring gjøres; først ved den Antagelse, at
de tiltrækkende Kræfter staae i omvendt Forhold af Afstandenes
Qvadrater, blev Ligningen bestemt:
\[
\frac{2 R \sin^{2} \tfrac{1}{2} \Theta}{15\tfrac{5}{8}}
  = \frac{\dfrac{\mu}{R^{2}}}{\dfrac{\mu}{r^{2}}}
  = \frac{r^{2}}{R^{2}}.
\]
Den saaledes bestemte Ligning var et Spørgsmaal, rettet til
Virkeligheden i en saa exact Form, at Kjendsgjerningen selv
nødvendig maatte svare enten Ja eller Nei.

Det første Svar, Newton fik paa sit Spørgsmaal, var som
bekjendt ugunstigt. Da den virkelige Natur ikke kunde feile, saa
uddrog han mismodig den Slutning, at Feilen laa i hans Tanke;
det var, mærkværdigt nok, hans egenlige Feil. I Sagen selv var
hans oprindelige Tanke lige saa feilfri som Naturen. Feilen laa
i den dagjældende Værdi for Jordradien. Da denne Feil siden
ved en ny Opmaaling blev rettet, fandt Newton sin Tanke
bekræftet; og nu først var det Vendepunkt indtraadt, da den for et
særligt Tilfælde fundne Lov kunde ved en til fortsatte
Iagttagelser svarende Induction udvides til almindelig Verdenslov.

Den Newtonske Lov har i sin Tid afgivet et storartet Exempel
paa, hvorledes Lovopfyldelsen kan formedelst de
% --- p. 52 ---
indtrædende Betingelser forandre en Lovs Udseende, saa at den
næsten bliver ukjendelig. Ved at beskjæftige sig med
Perturbationerne, særlig med de eiendommelige Forandringer,
Maanebanen er underkastet, (Bevægelsen af Apogæum), kom Clairaut
til at opkaste Tvivl om Newtonslovens almene Gyldighed. Han
bemærker i Anledning af en Strid med Buffon, som vilde forsvare
Loven med aprioriske Grunde, at de Maader, hvorpaa man hidtil
havde støttet den Newtonske Lov, vare saa belemrede med
Undtagelser og stykkeviis Tagen Hensyn til forskjellige
Indvirkninger, at en Lov, forskjellig fra $\dfrac{a}{x^{2}}$ kunde give
lige saa god Overeensstemmelse mellem Beregnings- og
Observationsresultaterne. Ved at løse „de tre Legemers Problem“ kom
Clairaut foreløbig til det Resultat, at Udtrykket for Loven maatte
blive, ikke $\dfrac{a}{x^{2}}$, men $\dfrac{a}{x^{2}} + \dfrac{b}{x^{4}}$.
Videre gaaende Undersøgelser førte imidlertid til det Resultat,
at Problemets endelige Løsning netop stadfæstede Newtons Formel.
Vexelvirkningen mellem de tre Legemer, Centrallegemet, det
perturberede og det perturberende Legeme, foregaaer nemlig i Et
og Alt efter samme Lov, saa at de vexlende Betingelser kun
medføre Ændringer: det Specielle er i dette Exempel en
% Fodnotetegnet er trykt *) --- Bogens sædvanlige Mærke.
Modification af det i sig uforandrede Almene.\footnote{At en
Ændring af Grundformlen dog maaskee er mulig, maa vel
indrømmes, siden udmærkede Astronomer, som f. Ex. Bessel, endnu
nære nogen Tvivl.}

Men de Virkelighedsbetingelser, hvorunder en Lov opfyldes,
kunne ogsaa være af den Beskaffenhed, at de indføre een eller
flere nye, med den givne Almeenform ueensartede Lovbestemmelser
og derved bevirke en væsenlig Forandring i Loven. Man tænke
paa den Forandring, der foregaaer med Faldloven, naar Faldet
skeer paa et Skraaplan og de modificerende Omstændigheder tages
i Betragtning;
% --- p. 53 ---
% Fodnotetegnet er trykt *), sat efter „Pendul“ og foran Semikolon.
% I Noten er Bogstavregningen sat i mathematisk Kursiv; de trykte
% Variable staae opret antiqua, hvilket efter Husreglen ikke
% gjengives. Hvor et Bogstav bærer Mærke eller Exponent, er hele
% Rækken sat i Mathematiksats, saa m, m', m'' ikke skifter Snit
% midt i Opregningen. Prikkerækkerne staae som spatierede Punkter.
paa den Forandring, der skeer med Loven for Pendulsvingninger,
eftersom man betragter enten „det enkelte“ eller „det
sammensatte“ Pendul\footnote{Ved lige Fjerningsvinkel og Tyngdekraft
forholde Svingningstiderne for to enkelte Penduler sig som
Qvadratrødderne af Pendulernes Længde. Ved det sammensatte
Pendul gjælder den anførte Lov derimod kun for Pendulets
reducerede Længde, d. e. for Afstanden fra Ophængningspunktet til
Svingningernes Midtpunkt. Ved Bestemmelsen af Svingningernes
Midtpunkt indføres Inertiemomentet. Ere $m$, $m'$, $m''$
$\ .\ .\ .\ .$ Massedelene, og deres Afstand fra
Ophængningspunktet $x$, $x'$, $x''$ $\ .\ .\ .\ .$, saa er Inertiemomentet
$mx^{2} + m'x'^{2} + m''x''^{2} + \ .\ .\ .\ .$ Kaldes hele Massen
M, og Tyngdepunktets Afstand fra Ophængningspunktet a, saa er
den reducerede Længde
$l = \dfrac{mx^{2} + m'x'^{2} + m''x''^{2} + \ .\ .\ .\ .}{Ma}$
Betegnelsen for en egen Lov. At de i Opfyldelsen sammensmeltede
Love ikke modsige hinanden, forstaaes af sig selv.}; paa
Forskjellen mellem Parablen og „den balistiske Curve“ o. s. v.
Afvigelsen er navnlig i sidste Tilfælde saa gjennemgaaende, at
der i Virkeligheden viser sig en ganske ny Lov.

Opfyldelsen af en Lov skeer altsaa paa væsenlig forskjellig
Maade, eftersom de indtrædende Betingelser enten blot ændre det
phænomenale Udtryk eller medføre en Omdannelse af Loven eller
endog forvandle den saaledes, at den optages og forsvinder i en
anden Lov. Det forholder sig med Lovene som med Kræfterne.
Hvorledes de særskilte Love i indbyrdes Vexelvirkning deels
stadfæste, deels ændre, deels ophæve hinanden gjensidig, kan først
tilfulde oplyses, naar vi betragte de tilsvarende Kræfter fra
Totalitetens Synspunkt.

% Indløbende Overskrift, læst paa Siden selv (S. 53) og ikke i
% Indholdsfortegnelsen; Kolon og Ordlyd stemmer med Indholdet.
\runin{c) Totalitet af Kræfter: Lovenes Rige.} Naturlovene ere
uforanderlige; denne Grundsætning, paa hvilken Physiken staaer
fast, og uden hvilken den ikke vilde være den strenge, ophøiede
Videnskab den er, synes, naar den
% --- p. 54 ---
udtales i al Ubetingethed, med eet Slag at nedslaae Alt, hvad
ovenfor er fremsat om Forandringer i Kræfter og Love. Og dog
veed Naturphilosophien sig, hvad dette Punkt angaaer, i fuld
Overeensstemmelse med Physiken. Physiken har nemlig kun Ret
til at opstille sin Sætning: at Naturlovene ere uforanderlige,
fordi Logiken har Ret til at opstille sin: at Tankelovene ere
uforanderlige. Naar da Logiken beviser, at Tankelovenes
Uforanderlighed netop bestaaer i deres Foranderlighed, saa har den
\textit{eo ipso} godtgjort, at det Samme maa gjælde om
Naturlovene. Det Uforanderlige kan ikke fastholdes for sig i endelig
Adskillelse fra det Foranderlige, saa lidt som Væsenet kan stilles
endelig op mod Skin og Phænomen. Væsenet for sig er Skin og
Intet, naar Skin og Phænomen tænkes borte. Væsenet er selv med
i Forandringerne; dets uforanderlige Bestaaen viser sig just deri,
at det under alle sine Forandringer bevarer sig uforandret.
Eenheden af Uforanderligt og Foranderligt udsletter ikke Forskjellen;
nei, den stadfæster Forskjellen: her er væsenlig Forskjel, fordi
her er væsenlig Eenhed.

Naturkræfterne ere uforanderlige; og dog forandres de, idet de
virke og indgaae i Vexelvirkning: i deres Væsen ere de
uforanderlige, men i deres Yttringer ere de høist foranderlige.
Saaledes ogsaa med Lovene: de ere i deres Væsen uforanderlige,
men i deres Opfyldelse foranderlige. Naar Sætningen om
Naturlovenes Uforanderlighed ikke forstaaes paa denne Maade, men
forkyndes slet og ret som et Dogme, en Troesartikel, hvori Væsen
og Phænomen, Mulighed og Virkelighed løbe ud i Eet, saa maa
der protesteres. Trods det katholske Ufeilbarhedsskin, hvormed
Dogmet prædikes, vil den sunde Sands dog let mærke Uraad og
spørge: hvad betyder Tilfældigheden? hvad bliver der af Friheden?

At Nødvendighed og Tilfældighed enes overalt i det store
Naturhele, at det Nødvendige i Kjendsgjerningernes Verden er
tilfældigt, og det Tilfældige nødvendigt, har sin
% --- p. 55 ---
Grund i, at det Uforanderlige bestaaer i Forandringer. Enhver
virkelig Lov maa have virkelige Betingelser og vise sig
bestemmende for virksomme Kræfter. En Naturlov være saa almindelig,
den vil, den kan dog aldrig opfyldes i det blot Almindelige; en
almindelig Opfyldelse er en Abstraction: Opfyldelsens Virkelighed
er i de virkelige Tilfælde. Tilfældene gaae alle ind under Loven;
men derfor tabe de ikke deres Tilfældighed, ophøre ikke at være
Tilfælde. Var der ingen Legemer, som kunde falde, vilde
Faldloven ikke blive opfyldt. Dersom der ingen andre Legemer var
i Himmelrummet, kunde den Kraft, hvormed Jordlegemet virker
paa andre Legemer, ikke yttre sig, ikke blive virkelig. I Væsen
og Mulighed er Faldloven den samme uforanderlige Lov, hvad
enten den saa i givne Tilfælde opfyldes eller ikke opfyldes; men
i Virkeligheden er den ikke den samme, thi naar den ikke
opfyldes, er den uvirkelig. I Væsen og Mulighed er Jordens
Tiltrækningskraft uforanderlig den samme, enten den virker eller den
ikke virker. Men Virkeligheden indfører en Forandring, thi kun
den virkende Kraft er virkelig, den ikke virkende er kun mulig.

Det vil nu let kunne sees, hvorledes Loven og Tilfældet
indgaae i mangesidigt Vexelforhold. At en Steen falder til Jorden,
er vel med Hensyn paa den enkelte Steen kun et enkelt Tilfælde,
men det udtrykker dog tillige den Nødvendighed, hvormed alle
Stene i alle lignende Tilfælde maae falde. Legemerne ere, hvad
Tyngden angaaer, under samme Lov, fordi de ere under samme
Nødvendighed, en Nødvendighed, der udtrykker deres iboende,
almene Væsen. Men, mærkværdigt nok, gives der ogsaa Love, som
tillade Tilfældigheden i hvert enkelt Tilfælde at unddrage sig den
% Fodnotetegnet er trykt *). Noten begynder paa S. 55 og løber
% over paa S. 56, hvor den fortsætter uden nyt Mærke.
almene Nødvendighed.\footnote{Thi fordi at En gifter sig i en vis
Alder, er det ikke nødvendigt, at Andre gifte sig i samme Alder;
fordi En hænger sig, er det slet ikke nødvendigt, at Andre ogsaa
hænge sig; fordi En i Distraction sender et Brev uden Udskrift
paa Posthuset, derfor er det ikke nødvendigt, at Andre maae gjøre
ligesaa. Og dog viser Sandsynlighedsregningen, at store Grupper
af slige Tilfælde, fra hvilken Tilværelsessphære de end hentes,
lade sig henføre under tilsvarende Love.}
% --- p. 56 ---
Der er altsaa i Tilfældenes Verden to væsenlig forskjellige
Kredse af Love, \emph{Væsenslove} og \emph{Tilfældighedslove}.
I den første Kreds er Tilfældigheden nødvendig, i den sidste er
Nødvendigheden tilfældig. De to Kredse danne i deres skarpt
prægede Modsætning to Extremer, som ikke maae forvexles. Men
Extremerne staae ikke overfor hinanden, adskilte ved en gabende
Kløft; utallige Overgangsbestemmelser danne et Væv, hvori
Nødvendighedens Traade under mangehaande Krydsninger ere som
indvirkede i Tilfældighedernes brogede Tæppe.

Der er en gjennemgaaende Sammenhæng i den hele Natur,
intet Spring, ingen Afbrydelse, ingen nyskabt Materie, intet
Indskud af nye Kræfter, ingen væsenlig nye Love: allevegne gribe
Begivenhederne ind i hverandre som Led i en sluttet Kjæde; det
Nye, der dannes, dannes af det Gamle, af det, som i Urtiden
allerede var. Hellenerne, siger Anaxagoras, antage med Uret
Vorden og Forgaaen, thi ingen Ting vorder eller forgaaer, men
Tingene sammensættes af bestaaende Ting, og saaledes vilde de
rigtigere kalde Vorden Sammensætning og Forgaaen Adskillen.
Hvad Anaxagoras anede, har den mechaniske Naturlære fra sit
Synspunkt tilstrækkelig godkjendt. Ligesom der i Universet er et
givet System af „evige Kræfter“, saaledes er der ogsaa et i sig
sluttet System af „evige Love“; dette System er
\emph{Lovenes Rige}.

At have opdaget og begrebet Naturens uendelige
Sammenhængseenhed og derved aabnet sig et storartet Indblik i
% --- p. 57 ---
Lovenes Rige, er „den moderne Videnskabs“ Stolthed.
Naturphilosophien erkjender det Sande og Berettigede heri. Men i de
Hymner, som klinge til Ære for den forgudede Nødvendighed,
høres ikke sjældent uopløste Dissonanser, skurrende Toner, som
minde om, at den moderne Videnskab af overvættes Glæde over
Nødvendigheden er i Færd med at fornegte Friheden.

Kan da, spørger man, Frihed virkelig bestaae med
Nødvendighed? Dette Spørgsmaal frembyder, hvad Naturerkjendelsen
angaaer, to Hovedsider, eftersom der nærmest reflecteres enten paa
Tilfældigheden eller paa den gjennemgaaende Sammenhængseenhed.

Ligesom al materiel Virkelighed har et Moment af
Tilfældighed, saaledes har Frihedens ideelle Virkelighed et Moment af
Vilkaarlighed. Hvad Tilfældigheden er i det Ydre, det factisk
Givne, er Vilkaarligheden i det Indre, det oprindelig Bestemmende.
Kan Nødvendigheden forliges med det Tilfældige udefter, saa maa
den vel ogsaa kunne forliges med det Vilkaarlige indefter. Men
Vilkaarlighed i Naturen? Den blotte Forestilling herom er
Physiken en Forargelse. Forargelsen vil imidlertid fortage sig, naar
Physikeren kommer til Besindelse. Det er jo netop Physiken selv,
som uden i nogen Maade at have med det Vilkaarlige at gjøre,
leverer Philosophien de vigtigste Forudsætninger og paalideligste
Støttepunkter for Vilkaarlighedsmomentets Anerkjendelse.

Physiken har brudt med Speculationen, for ubetinget at holde
sig til Kjendsgjerningen. Det er et afgjørende Vendepunkt i
Videnskabens Historie. At holde sig til Kjendsgjerningen er at
tage den, som den er, uden at raisonnere over, hvorfor den nu
just er saaledes, om den ikke maaskee ligesaa godt kunde have
været anderledes. Kjendsgjerningen er, fordi den er; men hvad
der kun er, fordi det er, er i sin givne Nødvendighed umiddelbart
at opfatte som
% --- p. 58 ---
et Tilfældigt. Iagttagelsen seer i de mange eiendommelige
Kjendsgjerninger lutter eiendommelige Tilfælde; hvorfra have de
givne Tilfælde da deres oprindelige Bestemthed? Fra andre
Tilfælde. Og de andre Tilfælde? Fra andre Tilfælde. Og disse
igjen? Fra andre Tilfælde. Det endeløse Raisonnement, der
uophørlig ombytter Tilfælde med Tilfælde, drukner i Tilfældigheden,
uden at naae saa vidt, at det endog blot kan skimte den
begrundende Nødvendighed.

Kun ved at holde sig til Fornuften og gaae ud fra den
aprioriske Forudsætning, at den i sig begrundede Naturens Orden
hviler paa en Alt omfattende, Alt indenfra bestemmende
Grundeenhed, kan den iagttagende Forsker opdage Nødvendigheden.
Ethvert særskilt Tilfælde maa jo have sin særskilte Grund:
saamange Tilfælde i det Udvortes, saamange særlige Grunde i det
Indvortes. Men Grundeenhedens Adskillelse i særlige Grunde kan
umulig være tilfældig; at gjøre det Begrundende tilfældigt er at
tilintetgjøre al Begrundelse. Grundeenheden maa altsaa selv
% Trykfeil: „Forkjellighed“ for „Forskjellighed“. Bogstavet s
% mangler. Efterprøvet paa KB-Exemplaret, hvor Ordet ligeledes
% staaer „Forkjellighed“, medens „forskjellige“ een Linie ovenfor
% staaer fuldt ud; Feilen er altsaa Sætterens og ikke dette
% Exemplars. Bogens eget Trykfeilblad retter den ikke.
% Gjengivet som trykt.
indenfra bestemme sig til en Mangfoldighed af forskjellige, i deres
Forkjellighed sammenhængende Grunde; Grundeenheden maa, for
Sammenhængens Skyld, nødvendig være bestemmende Selveenhed.

Vil Iagttageren nu fra dette Synspunkt tage Sigte paa
Tilfældigheden, skal han visselig opdage Vilkaarligheden. For ikke
at blændes af uvedkommende Hensyn, kan man jo vælge saadanne
almindelige, usammensatte Kjendsgjerninger som:
Sammenhængskraft, Tyngdekraft, magnetiske Kræfter, elektriske
Tiltrækninger o. s. v. Ingen af disse Kræfter lade sig aflede af de andre;
heller ikke har Physiken Udsigter til at kunne deducere Tingenes
tilsvarende Egenskaber enten af deres andre Egenskaber eller af
Materiens almindelige Væsen. Grundkræfternes Forskjellighed maa
have en bestemmende Grund. Men hvad er den bestemmende
Grund for Tilfældigheden? Det er Vilkaarligheden. Som det Be\-%
% --- p. 59 ---
stemte er, maa det Bestemmende være. Jo større Eftertryk der
falder paa den bestemte Kraft i dens Enkelthed og umiddelbare
Tilfældighed, desto større Eftertryk maa der fra Begrundelsens
Side falde paa den udelukkende Villiesyttring, den
Vilkaarlighedsact, hvorved den er blevet bestemt. Tilfældigheden siger, at
den givne Kraft er, fordi den er; Vilkaarligheden siger, at den er,
fordi den skal være. Fæster Blikket sig paa de enkelte
Kjendsgjerninger i deres Enkelthed, saa er Alt indadtil vilkaarligt,
ligesom Alt udadtil er tilfældigt. Hvad Tilfældigheden er for
Kjendsgjerningernes System, det er Vilkaarligheden for Grundenes
System.

Den rationelle Mechanik er fra Først til Sidst systematisk
betinget af vilkaarlige og tilfældige Forudsætninger og dog en
Articulation af Nødvendigheden, saa sikker og tankefast, at det i
Sandhed gjør den menneskelige Fornuft Ære. Uden
Vexelbestemmelse af Tilfældighed og Vilkaarlighed er det umuligt at stille
sig nogensomhelst mechanisk Opgave. Mathematikeren vælge,
hvilken Opgave, han vil, han seer øieblikkelig, at han enten selv
maa sætte Noget, altsaa begynde vilkaarligt, eller forudsætte
% Fodnotetegnet er trykt *). Denne Note er Bogens hidtil længste:
% den begynder paa S. 59, fylder hele S. 60 under Tekstens to
% første Linier, og slutter paa S. 61 --- overalt uden nyt Mærke.
% Her er hele Noten samlet paa Mærkestedet.
% Prikkerækkerne staae som spatierede Punkter, snart tre, snart
% fire; Antallet er gjengivet som trykt.
Noget som givet, altsaa begynde tilfældigt.\footnote{Hvorledes
Vilkaarlighed og Tilfældighed principielt forholde sig til
Sammenhængsnødvendighed, kan, hvad den exacte Analyse angaaer, sees
f. Ex. af \emph{Interpolationen}. Et Antal spredte Punkter,
beliggende i en vis bestemt Curve, er givet. Om den givne Curve
svarer til en virkelig Kjendsgjerning og forsaavidt falder ind
under det Tilfældige, eller den udtrykkelig er sat ved en fri
Selvbestemmelse og hører ind under det Vilkaarlige, er ligegyldigt.
Det Eneste, som fra Begyndelsen af kan vides, fordi det ligefrem
ligger i Sammenhængsnødvendigheden, er: at Ordinaten y og
Abscissen x maae for hvert Interval nøiagtig svare til hinanden;
$F(x,y) = 0$ eller $y = f(x)$. Det Problem, som ved
Interpolationsregningen skal løses, er da følgende: af en Functions
\emph{givne} Værdier, svarende til \emph{givne} Værdier af den
uafhængige Variable, Argumentet, at beregne den Værdi af
Functionen, der svarer til en hvilkensomhelst Værdi af Argumentet,
beliggende mellem de givne.

Kan det nu forudsættes, at den til den givne Function svarende
Curve har en continuerlig og jævn Gang, saa vil Problemet kunne
løses ved en Tilnærmelse, der er desto større, jo nærmere de
givne Værdier af Argumentet ligge hinanden; idet
Interpolationsformlen bliver:
\[
y = a_{0} + a_{1} x + a_{2} x^{2} + \ .\ .\ .\ a_{n} x^{n}. \qquad (1).
\]
For at bestemme de heri indgaaende (n + 1) Constanter maae
(n + 1) sammenhørende Værdier af Argumentet og Functionen
være givne. Ere disse henholdsviis:
\[
x_{0},\ x_{1},\ x_{2},\ \ .\ .\ .\ .\ x_{n}\ \text{og}\ y_{0},\ y_{1},\ y_{2}\ \ .\ .\ .\ y_{n},
\]
kan man opskrive (n + 1) Ligninger af Formen:
\[
y_{r} = a_{0} + a_{1} x_{r} + a_{2} x_{r}^{\,2} + \ .\ .\ .\ a_{n} x_{r}^{\,n}.
\]
Ved Opløsningen af disse (n + 1) Ligninger af første Grad kunne
Constanterne i Ligningen (1) bestemmes. Geometrisk seet kan
Opgaven altsaa løses ved igjennem de givne Punkter at drage en
Parabel af n\textsuperscript{te} Grad.

Opløsningen af de (n + 1) Ligninger af første Grad undgaaes
ved at bringe Ligningen (1) paa Formen:
$y = X_{0} (x-x_{1}) (x-x_{2}) \ .\ .\ .\ (x-x_{n}) + X_{1} (x-x_{0}) (x-x_{2}) \ .\ .\ .\ (x-x_{n}) + \ .\ .\ .\ X_{n} (x-x_{0}) (x-x_{1}) \ .\ .\ .\ .\ (x-x_{n-1})$,
hvori hvert Led indeholder alle Factorerne $(x-x_{0})$ o. s. v. paa
een nær, idet nemlig det første Led mangler den første Factor,
det andet den anden \ .\ .\ .\ . det n\textsuperscript{te} den
n\textsuperscript{te}. Constanterne $X_{0}$, $X_{1}$,
$\ .\ .\ .\ .\ X_{n}$ bestemmes af de sammenhørende givne Værdier
for x og y; saaledes $X_{r}$ af
$y_{r} = X_{r} (x_{r}-x_{0}) (x_{r}-x_{1}) \ .\ .\ .\ .\ (x_{r}-x_{r-1}) (x_{r}-x_{r+1}) \ .\ .\ .\ (x_{r}-x_{n})$.
Den søgte Interpolationsformel (af Lagrange) bliver altsaa:
% Trykfeilbladet retter S. 60, Linie 1 f. n.: „sidste Led i
% Tælleren (x-x_n) --- læs (x-x_{n-1})“. Efter Husreglen staaer
% Teksten her som trykt, med (x-x_n) i sidste Brøks Tæller.
\begin{align*}
y = {}&+ \frac{(x-x_{1}) (x-x_{2}) \ .\ .\ .\ .\ (x-x_{n})}%
          {(x_{0}-x_{1}) (x_{0}-x_{2}) \ .\ .\ .\ (x_{0}-x_{n})}\ y_{0}\\
      &+ \frac{(x-x_{0}) (x-x_{2}) \ .\ .\ .\ (x-x_{n})}%
          {(x_{1}-x_{0}) (x_{1}-x_{2}) \ .\ .\ .\ .\ (x_{1}-x_{n})}\ y_{1}
          + \ .\ .\ .\ .\ .\\
\ .\ .\ .\ .\ .\ {}&+ \frac{(x-x_{0}) (x-x_{1}) \ .\ .\ .\ (x-x_{n})}%
          {(x_{n}-x_{0}) (x_{n}-x_{1}) \ .\ .\ .\ .\ (x_{n}-x_{n-1})}\ y_{n}.
\end{align*}

Denne saa rigt articulerede Sammenhængsnødvendighed kan, som
sagt, lige saa vel have Vilkaarligheden som Tilfældigheden til
Grundlag.

Lade vi nu den givne Curve med den jævne, continuerlige Gang
være et Billede paa den givne Natur, da vil det kunne sees,
hvorledes Speculationen og Empirien, hver fra sin Side, af
Sammenhængsnødvendigheden fristes til at opfatte denne virkelige
Verden, denne nærværende Tingenes Orden, som den eneste mulige.
Speculationen holder sig til det Almene, bliver staaende ved
$y = f(x)$ og kan i Mangel af mathematisk Indsigt ikke komme
videre. Empirien gaaer ud fra det Specielle, de givne Enkeltled.
Jo mere spredte de iagttagne Enkeltled ere, desto løsere synes
den hele Natursammenhæng at være. Men efterhaanden som
Iagttagelsernes Sum forøges, rykke Enkeltleddene hinanden
nærmere: de utallige mulige Curver nærme sig til at falde sammen
med een eneste virkelig. Under sin rastløse Arbeiden med at
udfinde og beregne de nye Mellemværdier faaer Empirikeren et
stedse stærkere Indtryk af, at denne Naturens Orden, denne Curve
med den jævne Gang, ikke blot under de givne Forudsætninger,
men i ubetinget Almindelighed maa være den eneste mulige. For
at hæve Blændværket, behøver Empirikeren imidlertid blot at gaae
tilbage til den rene Mathematik og besinde sig paa
Interpolationernes almene Grundlag. \emph{Constanterne} i Ligningen (1)
maae for hver \emph{forskjellig} Curve være \emph{forskjellige};
men en uendelig Mængde forskjellige Curver, svarende til en
uendelig Mængde forskjellige Rækker af (n + 1) sammenhørende
Constanter, ere jo mulige, og ligesaa ere uendelig mange
forskjellige Verdener mulige.} Der er i Betingelsernes Kreds
% --- p. 60 ---
% Kun to Linier Brødtekst staaer paa S. 60; Resten af Siden er
% optaget af Noten fra S. 59, som her er sat paa Mærkestedet.
samme Forhold mellem det Nødvendige og det Vilkaarlige, som
imellem det Nødvendige og det Tilfældige: Conseqven\-%
% --- p. 61 ---
sen er lige nødvendig, enten Forudsætningen er vilkaarlig eller
tilfældig.

Ligesom Tilfældighedsmomentet gaaer op i Sammenhængseenheden,
saaledes gaaer Vilkaarlighedsmomentet op i Selveenheden.
Vexelbestemmelsen af Selveenhed og Sammenhængseenhed viser sig
ved Løsningen af en hvilkensomhelst mechanisk Opgave.

% Kapitlet sluttes med en kort, centreret DOBBELT Streg. Efter
% Husreglen gjengives Ornamentets Form ikke: \streg tjener overalt.
\streg


% ======================================================================
%  Første Deel — Andet Kapitel: Mechaniske Aarsager   (printed pp. 62–124)
% ======================================================================

\kapitel{Andet Kapitel.}{Mechaniske Aarsager.}

%  --- A. Mechanisk Idealitet  (pp. 62–86) ---

%  MATHEMATICAL — work from the image; the OCR is not a witness here.
% --- p. 62 ---
\afdeling{A.}{Mechanisk Idealitet.}

At Naturnødvendigheden i sig forener det Vilkaarlige og Tilfældige som
ophævede Bestemmelser har sin Grund i, at den mechaniske Aarsag paa een
Gang er at opfatte baade som Begreb og som Gjenstand. Som Gjenstand er
Aarsagen reel, som Begreb er den ideel. For at gjøre Aarsagsidealiteten
kjendelig, skulle vi nu særlig paavise, hvorledes det i Mechaniken
forholder sig med
% Spærret Sats: „Kraft, Betingelse“ og „Aarsag“ ere spærrede, men det
% mellemstaaende „og“ er det ikke; derfor to Fremhævelser og ikke een.
\emph{Kraft, Betingelse} og \emph{Aarsag}, og hvorpaa det væsenlig
beroer, at Aarsagens Eenhed magter det udvortes Mangfoldige og i sig
optager en Fleerhed af \emph{forskjelligartede Momenter}.

\runin{a) Kraft og Aarsag.} Naar Mechaniken gjør „det, vi kalde Kraft“
til en „ubekjendt“ Aarsag, saa antyder den herved baade Forskjel og
Eenhed af Kraft og Aarsag. Forskjellen er, at det Ubekjendte ligger i
Kraften, ikke i Aarsagen. At Aarsagsbegrebet har Gyldighed, at enhver
Bevægelse af Materien nødvendig maa have sin Aarsag, kan ikke betvivles;
i modsat Fald maatte den rationelle Mechanik jo ophøre at være en
Videnskab. Derimod er man i Tvivl om, hvad det i Aarsagen Virksomme
egenlig er, hvorledes
% --- p. 63 ---
det bedst skal betegnes, og gjør da ligesom en Slags Undskyldning ved at
sige: „vi kalde det Kraft“. Det er i vort Foregaaende viist, at denne
Undskyldning er fra Væsenets Synspunkt ganske overflødig, at Begrebet
Kraft er i sin Art et lige saa tilforladeligt Naturbegreb, som Begrebet
Aarsag. Og dog har Undskyldningen en antagelig Grund, naar vi see paa
Virkeligheden. Netop derved, at Aarsagen bestemmes som Kraft, gaaer
Begrebet ud over sig selv; det ombyttes ikke med et andet Begreb, men det
sættes over i Virksomhed og udtrykkes i Kjendsgjerning. Det Virksomme er
et Begreb og dog mere end et Begreb; det er Kraft. Det er netop gjennem
Kraften at det mechaniske Aarsagsbegreb gaaer over til Virkelighed og
bliver en virkende Aarsag.

Hvorledes Begriben og Virken i denne Henseende forholde sig til hinanden,
vil kunne udfindes ved at søge Oprindelsen til de eiendommelige
Betegnelser, i hvilke Mechaniken, skjøndt den iøvrigt forholder sig
ligegyldig til Spørgsmaalet om Kræfternes „physiske Oprindelse“,
fastholder en qvalitativ Forskjel mellem Kræfterne indbyrdes, imellem
accelererende Kraft, bevægende Kraft og levende Kraft. Det er i
Virkeligheden den samme Kraft --- det Virksomme er Kraft, og Kraften det
Virksomme --- og dog paa Grund af de Indholdsbestemmelser, hvormed
Aarsagsbegrebet under en fremskridende Udvikling beriges, qvalitativt
forskjellige Kræfter; det er i Væsen samme Aarsag, og dog paa Grund af de
virksomme Kræfters Identitet forskjellige Aarsager.

% Det græske Bogstav er læst paa KB-Exemplaret: α, ikke a.
% Arbeidsexemplarets Blyantsunderstregninger under „instantan“ og
% „accelererende“ ere en tidligere Læsers Mærker; paa KB-Exemplaret staaer
% begge Ord i almindelig Antiqua, altsaa hverken kursiveret eller spærret.
\subrunin{α) Accelererende Kraft.} Som Bevægelsens Aarsag kan Kraften
enten virke pludselig, f. Ex. i Stødet, og faaer da Navn af instantan
Kraft, eller vedblivende forandre Hastigheden og kaldes da --- positivt
eller negativt --- accelererende. I Begrebet selv er Forskjellen
forsvindende. Den instantane Kraft udkræver en, om end nok saa lille,
endelig Tid for at bevirke en endelig Hastighed og er for\-%
% --- p. 64 ---
saavidt væsenlig accelererende; omvendt kan den accelererende Kraft
opløses i uendelig mange Stød og er i hvert Stød at betragte som
instantan. Væsenligt for Aarsagsbegrebet er derimod det Eiendommelige, at
Kraften i og for sig som Kraft ikke er Accelerationens Aarsag, at den kun
er det formedelst Inertien. Tænke vi Inertien borte, da maatte Kraften
jo, for at accelerere eller retardere en Bevægelse continuerlig, selv
enten formeres eller formindskes continuerlig; men saaledes kom da
Accelerationens Aarsag ikke til at ligge i Kraften, men i det, som
foraarsagede Forandring i Kraften. Betragtes Kraftforandringens Aarsag
paany som en Kraft, saa maa denne Kraft igjen have en Forandringens
Aarsag, og saa fremdeles i det Uendelige.

% MATHEMATISK SATS, S. 64--70. Efter Husreglen (batch 4) staaer enkelt
% staaende Bogstaver i den løbende Text som almindelig Text (k, g, t, v,
% M, m, P, A, B), medens Ligninger og sammensatte Udtryk sættes i
% Mathematiksats; hvor et enkelt Bogstav i samme Led stilles ligefrem
% overfor et sammensat Udtryk (g og mg, m og mV), sættes begge i
% Mathematiksats, saa at Skriftsnittet ikke skifter midt i Modstillingen.
% Trykkens opreiste Antiqua i Formlerne gjengives ikke; Mathematiksatsens
% kursive Bogstaver staae.
% Arbeidsexemplaret har her Blyantsunderstregninger under „Inertien“ og
% „Kraften“; spacing.py finder ingen Spærring paa Siden, og KB-Exemplaret
% har intet Mærke. Det er en Læsers Streger, ikke Sats.
Kun ved at virke i Eenhed med Inertien er Kraften accelererende. I
Accelerationen have vi da den søgte Eenhed af Reflecteren og Virken, et
kategorisk Udtryk for Aarsagens mechaniske Idealitet. Analysen viser
nøiagtig, hvorledes det forholder sig med denne Eenhed. En constant Kraft
k virker paa et materielt Punkt; i det relative Øieblik dt er
$kdt = dv$. Det materielle Punkt modtager i hvert Øieblik en uendelig
lille, men constant Hastighedstilvæxt; men det, som gjør, at Hastigheden,
der i første Øieblik kun er $dv$, bliver i næste Øieblik $2dv$, i næste
igjen $3dv$ o. s. fremdl., saa at den i det $n^{\text{te}}$ Øieblik er
voxet op til $ndv = v$, for $n = \infty$, er netop Inertien. Det er
Inertien, som optager, og det er Inertien, som bevarer enhver
Hastighedstilvæxt og derved gjør det muligt, at Hastigheden kan voxe med
Tiden, være en Function af Tiden: $v = f(t) = kt$. Det sees heraf, at det
Passive lige saa vel er bestemmende for Aarsagen, som det Active. Det
Active er passivt, thi Kraftens Indvirkning paa sit Angrebspunkt er fra
Først til Sidst afhængig af den Indvirkningen modtagende og bevarende
Inertie; det Passive er activt, thi Inertien er fra sin Side ligesaa
medvirkende under Accelerationen som Kraften.
% --- p. 65 ---
Saa gjennemgaaende er Reflexionen: den accelererende Aarsag er det
Actives og det Passives reflecterede, virksomme Eenhed.

I Reflexionens Form er Aarsagen Begreb; men den Reflexion, hvorom her
tales, er Kjendsgjerningernes egen bestemmende Reflexion, en Reflexion i
Magt. Enhver Reflex i den logiske Begriben har paa Virkelighedssiden sit
antilogisk existentielle Moment i den begribende Magt; den uendelige
Reflecteren har Virksomheden til virkeligt Udtryk: Punkt for Punkt gaaer
denne Reflecteren i Eet med en tilsvarende Virken. I Naturen selv er
Aarsagen Begreb og dog mere end Begreb. Begrebet virker ikke; det er den
i Inertien reflecterede Kraft, den physiske Aarsag, der virker. Men den
physiske Aarsag bestemmes til Virksomhed ved sit Begreb; Alt, hvad
Mechaniken kan oplyse om Kraft og Inertie, maa i enhver Henseende laane
sin Gyldighed fra Aarsagsbegrebet. Tage vi det i Aarsagen Virksomme bort,
saa er der kun et logisk Begreb, ingen virkelig Aarsag; men tage vi
Aarsagsbegrebet bort, saa ere Kraft, Inertie, det Virksomme, og hvad man
ellers vil kalde det, kun ørkesløse Forestillinger; der er da slet ingen
enten virkelig eller uvirkelig Aarsag.

% Det græske Bogstav er læst paa Siden selv: β, ikke B eller ß.
\subrunin{β) Bevægende Kraft.} Da det netop er den bevægende Aarsag, der
kaldes Kraft, kunde det ved første Øiekast synes overflødigt at indføre
en særegen Kategorie og stille den bevægende Kraft op ved Siden af den
accelererende. Den accelererende Kraft er bevægende, og den bevægende
accelererende; hvorpaa beroer da Forskjellen? Paa en særlig Udvikling af
Aarsagsbegrebet. Ved Bestemmelsen af Accelerationen og den accelererende
Aarsag blev det Materielle i Sagen stiltiende underforstaaet, men ikke
udtrykkelig fremdraget. Vel sandt, Kraften udgaaer fra Materien, Inertien
er i Materien, Reflexionen af Inertie og Kraft har Materien til Grundlag;
men det materielle Moment kommer først til sin Ret, naar Massens Begreb
indgaaer i Aarsags\-%
% --- p. 66 ---
begrebet, og Materiens Qvantitet bliver afgjørende for Aarsagen. Vælge vi
til Exempel Tyngden ved Jordens Overflade som accelererende Kraft g, da
er det, hvad Accelerationen angaaer, aldeles ligegyldigt, enten g virker
paa en lille eller paa en stor Masse; hver mindste Deel af Legemet vil
--- i det tomme Rum --- falde med samme Hastighed, enten Delene ere
forbundne eller adskilte: $gdm$ og $\int_{0}^{m} gdm = mg$ give i saa
Henseende selvsamme Resultat.

Anderledes med den bevægende Kraft, der er Aarsag paa Grund af Massen. Er
Massen M større end Massen m, da har den bevægende Kraft i $Mg$ unegtelig
en større Styrke end i $mg$; men medens der imellem $Mg$ og $mg$ kun er
en qvantitativ Forskjel, eftersom de begge høre ind under samme
Aarsagsbegreb, er der derimod en qvalitativ Forskjel mellem $g$ og $mg$,
thi her er Aarsagsbegrebet forskjelligt. For Bevægekraften som Aarsag er
det, strengt taget, ikke nødvendigt, at den foraarsager en virkelig
Bevægelse. Vel er Aarsagen ikke uden Virkning, men Virkningen kan være
ophævet ved en Modvirkning, et Tryk, og Aarsagen har ligefuldt sin
Realitet. Saaledes med Vægten. Vægten er et Product af Masse og
Acceleration uden forudgaaende Bevægelse: $Mg = P$. Vægten P er
reflecteret Eenhed af det Extensive i Massen og det Intensive i Tyngden,
en umiddelbar og dog reflecteret Eenhed af Realitet og Begreb; det er
Vægtens mechaniske Idealitet som Aarsag. Massen i sin umiddelbare
Tilværen er intet Begreb; det kan let forekomme en overfladisk Betragter,
som kunde Massen endda nok være Aarsag uden at indgaae i
Aarsagsbegrebet, men det er en Modsigelse. Af Vægten sees det, at
Massebegrebet er et væsenligt Naturbegreb, thi af Vægten sees det, at
Massen kun er i sig, hvad den i sin Reflexion med Tyngden er for Andet.
Forandrer Tyngden sin Intensitet, saa faaer Massen en anden Vægt og, som
det i Ligningen: $M = \frac{P}{g}$ udtalte Begreb
% --- p. 67 ---
viser, en anden Realitet. De mechaniske Ligninger ere i denne Henseende
mere end Formler med Lighedstegn, de ere Domme. Ligningen
$M = \frac{P}{g}$ kan vel i abstract qvantitativ Forstand opløses i en
Tautologie: $M = \frac{Mg}{g} = M$, men, opfattet som en Dom, betyder den
noget ganske Andet. At $Mg$ udtrykkes ved $P$, betyder nemlig, at
Factorerne M og g ikke maae fastholdes i deres Product som ueensartede
Eenheder, men opfattes som Momenter i Realbegrebet P, det concrete
Begreb, der giver den existerende Masse eiendommelig Værdi som virkende
Aarsag.

Som Aarsag er den bevægende Kraft en Masseaarsag. Ophæves nu Virkningens
Modstand, saa indtræder Bevægelsen, Tiden faaer sin Betydning, og
Bevægekraften lægger sig for Dagen i Productet af Masse og Hastighed,
d. e. i \emph{Be\-%
vægelsesmængden}: igjen en Forandring i Aarsagsbegrebet. Kalde vi
Bevægekraften $mk$, Tiden $t$ og Hastigheden $v$, saa udviser Ligningen
$mkt = mv$, altsaa $mk = \frac{mv}{t}$, at en hvilkensomhelst constant
Kraft kan maales med den Bevægelsesmængde, den frembringer i en
Tidseenhed. At den bevægende Kraft $mk$, virkende i Tiden $t$,
frembringer $mv$, vil da sige, at $mkt$ er Aarsag, $mv$ Virkning.
Forsaavidt den vundne Bevægelsesmængde udtrykkelig sættes som Virkning,
har den virkende Aarsag optaget Tidsbestemmelsen i sit Begreb og er
blevet en Function af Tiden.

Her opstaaer nu det Spørgsmaal: er Begrebet af virkende Aarsag
fuldstændig udtrykt i det Begreb af Tid og Kraft, der har
Bevægelsesmængden til Maal? Cartesius svarer Ja, men Leibnitz Nei. For at
vurdere en Kraft, maa man nøiagtig vide, hvad den kan udrette. Det er,
siger Leibnitz, ikke nok, at Virkningen er proportional med Aarsagen; den
fuldstændige Virkning maa være sin
% --- p. 68 ---
Aarsag fuldkommen lig. For at løfte et Legeme A paa 1 Pund til en Høide
af 4 Alen, udfordres der samme Kraft som for at løfte et Legeme B paa 4
Pund til en Høide af 1 Alen; det er just den Kraft, de to Legemer opnaae
ved at falde, hvert fra sin Høide. Men den Hastighed, A opnaaer ved
Faldet, forholder sig til den, B opnaaer, som 2 til 1.
% Gangetegnet er trykt som et Kryds; sat med \times.
Bevægelsesmængden i A vil da være $1 \times 2$, medens den i B er
$4 \times 1$. Bevægelsesmængden i A er kun halvt saa stor som
Bevægelsesmængden i B, uagtet Kraften er den samme; følgelig kan Kraften
ikke maales med Bevægelsesmængden. Abbé de Conti indvendte, at den af
Cartesius opstillede Lov krævede Hensyn til ligetidige Kræfter eller til
Bevægelser, som to Legemer opnaae til samme Tid; men --- spørger Leibnitz
--- „naar man seer et Legeme af en bestemt Størrelse bevæge sig med en
bestemt Hastighed, skulde man saa ikke kunne vurdere dets Kraft, uden at
vide, hvor lang Tid det har været om at opnaae denne Hastighed? Og naar
to ligestore Legemer have samme Hastighed, men det ene har erholdt den
ved et pludseligt Stød, det andet ved at falde i længere Tid, vil man saa
sige, at deres Kræfter ere forskjellige? Det vilde da være det Samme som
at sige, at en Mand var rigere, fordi han havde været længere Tid om at
vinde sine Penge.“

% Det græske Bogstav er læst paa Siden selv: γ, ikke y eller 7.
\subrunin{γ) Levende Kraft.} Grunden til, at Bevægelsesmængden ikke kan
være Bevægekraftens, den bevægende Aarsags, fuldstændig udtalte Virkning,
ligger i, at Modsætningen af Aarsag og Virkning endnu er uvirkelig.
Ligningen $mkt = mV = Mv$, udsiger, at den bevægende Kraft $mk$ kan i
Tiden $t$ vække en Bevægelsesmængde i $m$, der er $mV$ og i $M$, der er
$Mv$. Spørge vi nu, hvor stor Aarsagen er, saa henvises vi til
Virkningen. Spørge vi, hvor stor da Virkningen er, saa henvises vi til
Aarsagen. Aarsagen er ligestor med Virkningen; det er en blot formel Dom,
en tautologisk Sætning, saalænge det endnu ikke er oplyst,
% --- p. 69 ---
hvori denne Storhed virkelig bestaaer. Spørgsmaalet er, om den bevægende
Aarsag i Form af $mV$ virkelig kan udrette det Samme, som i Form af $Mv$;
thi hvis den ikke kan det, saa er den jo ved at antage to forskjellige
Former gaaet ind under to forskjellige Begreber og er saaledes blevet til
to forskjellige Aarsager.

For at prøve den virkende Aarsag er det nødvendigt at see den i
Virksomhed. Virksomheden kan da ikke være standset af en umiddelbar
Modstand; Modstanden maa overvindes, og Virksomheden hver Gang bestemmes
ved den Maade, hvorpaa den overvinder sin Modstand. Hertil svarer den
Adskillelse, Leibnitz har indført mellem den
% Spærringen er efterprøvet paa KB-Exemplaret: „døde“ er spærret, men det
% foranstaaende „den“ er det ikke, medens det følgende „den levende“ er
% spærret i sin Heelhed. Gjengivet som trykt.
\emph{døde} og \emph{den levende} Kraft. Til døde, elementære Kræfter
regner han dem, hvis Virkning umiddelbart er ophævet af en Modvirkning;
den levende Kraft derimod er „altid forbundet med den virkelige
Bevægelse“. Under Bevægelsen giver Modstanden efter, medens Kraften
efterhaanden tilsættes.

Sættes $mkt = mV$, vil Virksomheden i det første Tidselement være
$mkdt = mdV$. Hvad udrettes nu i den løbende Tid? Kraften bevæger Byrden
fremad i Rummet, men idet den slæber paa Byrden, aftager Hastigheden, og
den aftagende Hastighed er igjen et Udtryk for den aftagende Kraft.
Hvorlænge Kraften kan holde ud, maa da sees paa den Veilængde,
hvorigjennem den formaaer at føre sin Byrde. Ligningen $mdV = mkdt$ maa
altsaa forvandles til $mdV = mk\frac{ds}{V}$,
% TRYKFEIL: der staaer „kvad der udrettes i Tiden“ for „hvad“. Læsningen er
% efterprøvet paa KB-Exemplaret, som har samme Bogstav; Feilen er altsaa
% Sætterens og ikke dette Exemplars. Trykfeilbladet retter den ikke.
% Gjengivet som trykt.
saa at det kjendes paa Rummet, kvad der udrettes i Tiden. Det gamle Ord,
at Tiden er lige lang, men ikke lige nyttig, finder her en egen
Bekræftelse; dt er constant, men V varierer, hvoraf følger, at $ds$ ogsaa
% Punktrækkerne ere trykte med tre Punkter efter ds_2 og fire efter V_2;
% Forskjellen gjengives ikke, \ldots staaer begge Steder.
varierer: $ds > ds_{1} > ds_{2} \ldots$, efterdi $V > V_{1} > V_{2}
\ldots$. Af $mdV = mk\frac{ds}{V}$ følger $mVdV = mkds$; ved at integrere
denne Ligning erholde vi $mV^{2} = 2mks$ og med det Samme den reale
Modsætning af Aarsag og Virkning
% SIDEN SLUTTER HER, uden Punktum efter „og Virkning“ og med henved en
% Trediedeel af Satsfladen staaende tom; S. 70 begynder med et nyt
% Afsnit („Men naar der til Bevægelsesmængden ...“). Baade Arbeids- og
% KB-Exemplaret ere eens paa dette Punkt, saa Manglen er Trykkens og ikke
% dette Exemplars; efter al Sandsynlighed er en udhævet Ligning faldet ud
% af Satsen. Trykfeilbladet retter det ikke. Gjengivet som trykt.
% --- p. 70 ---

Men naar der til Bevægelsesmængden $mV$ svarer en levende Kraft $mV^{2}$,
saa svarer der til Bevægelsesmængden $Mv$ en levende Kraft $Mv^{2}$. Da
nu $mV^{2}$, ifølge Forudsætningen $mV = Mv$, har en anden Værdi end
$Mv^{2}$, saa sees heraf, at Kraften i $mV$, naar den sættes i Arbeide,
udretter noget Andet, end Kraften i $Mv$. Den levende Kraft, der maales
med Arbeidet, er altsaa qvalitativt forskjellig fra den i Tid bevægende
Kraft, der maales med Bevægelsesmængden. Her er hverken i Naturen eller i
Tanken noget Spring; det er Kraften, det i Aarsagen Virksomme, som under
Begrebsudviklingen trinviis optager og samler de eensidig fremtrædende
Bestemmelser og derved naaer sit fuldstændige Udtryk i Form af levende
Kraft.

\runin{b) Betingelse og Aarsag.} Naar de til Aarsagsbegrebet hørende
Existensbestemmelser fastholdes i Forestillingen som tilværende, hver for
sig, blive de betingende, og den virkende Aarsag opløser sig i en Sum af
Betingelser. Men det ligger i Betingelsens Natur at vise ud over sig selv
til et Betinget og til andre Betingelser; idet Betingelserne reflectere
sig i hverandre gjensidig, tabe de deres tilsyneladende Selvstændighed,
og den i sig udeelte Aarsagseenhed kommer igjen tilsyne som det
Ubetingede. Dette Ubetingede, der altsaa ikke falder udenfor, men gaaer i
Eet med sine Betingelser, er den betingede Aarsag.

% Det græske Bogstav er læst paa Siden selv: α.
\subrunin{α) Betingelsen.} Aarsagen er et Begreb, men et
Virkelighedsbegreb, og et Virkelighedsbegreb har altid Reflexer af
virkelige Ting i sit Indhold. De for Virkeligheden nødvendige
Tingsreflexer ere paa Causalitetens Vegne \emph{Be\-%
tingelsen} og \emph{det Betingede}. End ikke Kraften kan i denne
Sammenhæng holdes fast imod Betingelserne, thi Kraften for sig kan blot
forestilles; det Forestillede, „vi kalde Kraft“, synker ned til at være
en Betingelse. Betingelsen er en Forudsætning, som den forestillende
Bevidsthed
%  MATHEMATICAL — work from the image; the OCR is not a witness here.
% --- p. 71 ---
% SØMMEN MOD S. 70: S. 71 begynder UDEN Indrykning, midt i Sætningen fra
% S. 70 („... som den forestillende Bevidsthed | i Tankens Medfør deels
% lægger ind ...“). Der maa altsaa IKKE staae et Afsnitsskifte mellem
% Batch 6 og dette Fragment; joints.py skal føie dem sammen.
i Tankens Medfør deels lægger ind, deels forefinder i Virkeligheden, en
Fordring paa et Uafhængigt, som for det Afhængiges Skyld maa
fyldestgjøres. Betingelsen, det Uafhængige, indeslutter det Betingede,
det Afhængige, og dog stille de sig mod hinanden som Led imod Led.
Betingelsesreflexionen er en Reflexion af forestillede Led, en udvortes
mellem Væsen og Existens, Begreb og Virkelighed svævende, for den
mechaniske Virksomhed betegnende Reflexion. Saaledes ere Kraften og
Inertien Betingelser for Accelerationen; de forestilles hver for sig, og
vise dog uophørlig hen til hinanden. Massen og Accelerationen ere igjen
Betingelser for Vægten; den samme Masse kan undergives en anden
Acceleration, og den samme Acceleration appliceres paa en anden Masse,
forsaavidt maae Betingelserne adskilles; og dog kunne de ikke adskilles,
thi det forandrede Product af Masse og Acceleration forandrer Vægten.
Drages Tiden med ind i Betingelsernes Kreds, saa kommer Vægten i
Bevægelse, viser sig som bevægende Kraft og bliver at bestemme som
Bevægelsesmængde. Føies endelig den Betingelse til, at den givne
Bevægelsesmængde skal anvendes til successivt at overvinde en Modstand og
maales med det Rum, Modstanden gjennemløber, saa have vi den levende
Kraft, dog ikke som Aarsag; thi nu er Aarsagen opløst, den er forandret
til en Sum af Betingelser.

Af Betingelsens kategoriske Eiendommelighed følger, at flere særskilte
Betingelser kunne indesluttes i een Grundbetingelse, og denne igjen
opløses i en Mængde Enkeltbetingelser. Den rationelle Mechanik er lagt an
paa at vise, hvorledes Opløsning og Sammensætning skeer systematisk. For
at Kræfter, der virke paa et Punkt, kunne være i Ligevægt, maa deres
Resultant være Nul; dette er Ligevægtens almene Betingelse, dens
Grundbetingelse. Men hvilken Vrimmel af særskilte Betingelser og
varierende Enkeltbetingelser kunne ikke heri indgaae!
% Arbejdsscanningens Blad har to sorte Klatter paa denne Side (foran
% „Ligevægtens“ og efter „indgaae!“). KB-Exemplaret er reent begge
% Steder; det er altsaa Blæk eller Snavs i dette Exemplar og ingen
% Sætterfeil. Intet gjengivet. „Ligevægtens“ er heller ikke spatieret.

% --- p. 72 ---
% NOTATION: Bogstaverne P, Q, R, A, B, C, O, a, b, c ere paa denne Side
% Led i de opstillede Ligninger og sættes derfor gjennemgaaende i
% Mathematiksats, saa at Snittet ikke skifter midt i Rækken (Husreglen
% fra Batch 5). Bogens egne oprette Antiqua-Variabler gjengives ikke.
Virke to Kræfter i parallele Retninger mod samme Side paa to ved en
ubøielig ret Linie forbundne Punkter, saa kan man bestemme Størrelsen og
Stillingen af deres Resultant, og derved af en denne modsat Kraft, der
holder dem i Ligevægt. Man seer strax, hvorledes den opstillede
Betingelse sammenfatter flere Betingelser: at der er to Kræfter, at de
ere parallele, forbundne o. s. v. Kaldes Kræfterne $P$ og $Q$,
Resultanten $R$, saa have vi det Betingede: $R = P + Q$. Gjøres det
endvidere til en Betingelse, at $A$, Angrebspunktet for $P$, og $B$,
Angrebspunktet for $Q$, skulle fra et vilkaarligt Punkt $O$ i
Forbindelsesliniens Forlængelse have Afstandene $a$ og $b$, da er
Afstanden mellem det vilkaarlige Punkt $O$ og Resultantens
Applicationspunkt $C$, en Afstand, vi kunne betegne ved $c$, udtrykkelig
heraf betinget. Det Betingedes Sammensætning, Relationen mellem $Rc$,
$Pa$ og $Qb$, svarer exact til den sammensatte Betingelse, og dette
Betingede kan ved en Vending i Reflexionen igjen stilles op som en
Betingelse. Da nemlig $Rc = Pa + Qb$, saa er Ligevægten mellem de tre
Kræfter $P$, $Q$ og $R$ jo netop betinget af, at $R$ anbringes i Punktet
$C$ og i modsat Retning af $P$ og $Q$.

Fremskridt i Analysen skeer i det Hele saaledes, at det, der i en
foregaaende Sammenhæng var det Betingede, gaaer i en følgende over til at
blive Betingelse, og at der idelig indvindes Nyt ved deels at
generalisere, deels at specificere en given Grundbetingelse. Til
% Minustegnet er her et ægte Minus, ikke Nielsens ÷ (efterprøvet paa
% KB-Exemplaret ved Forstørrelse); ÷ staaer i denne Bog kun, hvor han
% udtrykkelig trykker det, jvf. S. 46.
ovenstaaende Betingelseskreds høre endvidere Ligningerne $Q = R - P$,
$b = \dfrac{Rc - Pa}{Q}$; specificeres nu Betingelserne saaledes, at
$R = P$, men derimod ikke $c = a$, saa bliver $Q = 0$, og
$b = \infty$. Heraf sees, at to ligestore, i modsat Retning virkende,
parallele Kræfter ($R$ og $P$) umulig kunne holdes i Ligevægt af en
eneste Kraft ($Q$), thi en uendelig lille Kraft ($Q = 0$) i en uendelig
stor Afstand ($b = \infty$) ere Betingelser, som ikke lade sig
tilveiebringe. Saaledes føres vi
% --- p. 73 ---
ad Betingelsernes Vei fra Læren om parallele Kræfter til Læren om
Svingkræfter, et Exempel i det Smaa, som viser, hvad det i det Store vil
sige, at Betingelsernes Vei er Mechanikens Vei, at den exacte Videnskab
med sine Betingelses- og Endeligninger har Grund til at løse Aarsagen op
i Betingelser. Men Naturphilosophie er ikke Mechanik; den gaaer den
omvendte Vei og løser Betingelserne op i

% Det græske Bogstav er læst paa Siden selv (KB-Exemplaret): β, ikke
% „B“, „ß“ eller „P“, som OCR'en giver. Overskriften er spatieret og
% indløbende: Sætningen ovenfor ender uden Punktum og fortsættes i
% Overskriftens Ord („løser Betingelserne op i β) Aarsagen“).
% Bogstaverne A, B og C nedenfor ere løse Bogstaver i Brødteksten paa en
% Side uden Ligninger og staae derfor som almindelig Tekst.
\subrunin{β) Aarsagen.} Betingelserne skulle ikke udslettes, og dog
skulle de forsvinde. Det for Betingelsen Characteristiske er, at den i
sin forestillede Selvstændighed er saa uselvstændig. De hypothetiske
Domme: dersom A er, saa er B, dersom en Fixstjernernes Aberration finder
Sted, maa Jorden bevæge sig om Solen; dersom B er, saa er C, dersom
Jorden bevæger sig, maa Fixstjernerne have Parallaxe, vise os det
Selvstændiges Uselvstændighed, det Uafhængiges Afhængighed, naar de
anførte Domme afgive Præmisser for tilsvarende Slutninger. Der sluttes: A
er, altsaa er B, Fixstjernerne udvise en Aberration, altsaa maa Jorden
bevæge sig, endskjøndt vi ikke umiddelbart kunne sandse det; B er, altsaa
er C, Jorden bevæger sig, altsaa maae Fixstjernerne have Parallaxe, om vi
end paa Grund af dens Lidenhed ikke ere istand til at maale den. Der
sluttes i begge Tilfælde ved \textit{modus ponens}, og Slutningen er et
Beviis for det Betingedes Afhængighed af sin Betingelse. Men i
Virkeligheden kan Relationen ogsaa være omvendt; det er netop
Betingelsen, der er afhængig af sit Betingede; hvad der udtrykkelig sees,
naar vi slutte ved \textit{modus tollens}: B er ikke, altsaa er A ikke; C
er ikke, altsaa er B ikke. Hvilken af de to Slutningsmaader, der i et
givet Tilfælde har virkelig Gyldighed, maa naturligviis beroe paa
Kjendsgjerninger og bestemmes ved Iagttagelser; men at de begge have
logisk Gyldighed, beviser, at Betingelsen, uselvstændig, staaer og falder
med sit Betingede.

Hvor en Grundbetingelse opløses i særskilte Betingelser,
% --- p. 74 ---
viser sig den samme Uselvstændighed. Bevægelsen i en ret Linie har en
anden Grundbetingelse end Bevægelsen i en krum Linie. Analysere vi den
sidste, saa vil Betingelsen for den rette Linie her gjentage sig og
afsætte Enkeltbetingelser, thi Bevægelsen efter Tangenten giver en ret
Linie, og Bevægelsen i en Retning, som danner en Vinkel med Tangenten,
ligeledes. Indføres som særlig Betingelse et fast Centrum, saa vil
Stillingen af de to Retninger --- efter Tangenten og imod Centrum ---
ligeledes gaae ind under særlige Betingelser og igjen være betingende. Er
den ene Retning stadig vinkelret paa den anden, vil Curven blive en
Cirkel, og Hastigheden være constant; varierer Stillingen derimod
saaledes, at de to Retninger danne forskjellige Vinkler, snart en ret,
snart en spids, snart en stump, vil Hastigheden variere, og Curven kan
ikke være en Cirkel.

Det vil nu heraf kunne sees, hvorledes Betingelserne ophæve hinanden og
forsvinde i et Ubetinget. Ingen af de anførte Enkeltbetingelser bliver
særlig opfyldt; den ene forhindrer uafladelig den andens Opfyldelse.
Bevægelsen efter Tangenten bliver i hvert Øieblik forhindret af
Bevægelsen mod Centrum, og denne igjen af hiin. Bevægelsen efter
Tangenten forsvinder i to betingede Resultater; det ene er en Flugt fra
Centrum --- for Forestillingen den saakaldte Centrifugalkraft --- der
ophæves af en modsvarende Bestræbelse, Normalkraften; det andet er
Fremskridtet i Banen, det eneste Resultat, der ikke bliver ophævet.

Den saaledes betingede Eenhed, Bevægelsen i Banen, afhænger nu af en
tilsvarende Betingelsernes Eenhed, der ikke lader sig opløse, efterdi den
ikke indeholder Betingelserne som en Sum, men som ophævede Betingelser.
Denne de ophævede Betingelsers virkelige Eenhed er det Ubetingede, Sagen
i sig, den oprindelige Sag, \emph{Aarsagen}. Men naar det Ubetingede er
bestemt som Aarsag, saa er det deraf
% --- p. 75 ---
Betingede mere end et Betinget, det er et Foraarsaget, en Virkning.

At der er et andet Grundforhold mellem Aarsagen og Virkningen end imellem
Betingelsen og det Betingede, er let at indsee. Dersom det Betingede ikke
er, saa er Betingelsen ikke heller, thi det Betingede er et Betingelsens
Indhold, forestillet som Led for sig; Betingelsen med sit Betingede
tilhører alene den udvortes Reflexion. Men Aarsagen omfatter
Virkeligheden, og derved igjen baade den indvortes og udvortes Reflexion.
Aarsagen har en Væren i Mulighed, i indvortes Reflexion, selv da, naar
den ikke virker, ligesom Kraften har en Væren som et Indre, selv da, naar
den ikke yttrer sig. Betingelsen maa træde til for at en Kraft kan yttre
sig, thi Kraft og Betingelse ere sideordnede Kategorier. Men Aarsagen er
over dem begge; thi den er i dem begge. Aarsagen hutler sig ikke
igjennem ved at laane her en Kraft, der en Betingelse; Aarsagen eier selv
Kræfterne, den magter Betingelserne; dens Væsen er Magt, Betingelsernes
Ubetingede.
% „hutler“ (af „at hutle sig igjennem“) staaer saaledes i begge
% Exemplarer. I Arbejdsscanningen er Ordet understreget med Blyant af en
% tidligere Læser; Blyanten er ikke Tekst.

% Det græske Bogstav er læst paa Siden selv (KB-Exemplaret): γ, ikke
% „y“ eller „7“, som OCR'en giver. Overskriften er spatieret og
% indløbende.
\subrunin{γ) Den betingede Aarsag.} Som Betingelsernes gjennemgaaende
Eenhed er Aarsagen et Ubetinget; men da denne Eenhed selv er bestemt ved
de deri optagne Betingelser, er den i sig bestemte Aarsag betinget.
Reflecteres der paa Forskjellen mellem Aarsag og Betingelse, saa kan den
samme Aarsag ved at optage forskjellige Betingelser have forskjellige
Virkninger; reflecteres der paa Eenheden af Betingelse og Aarsag, saa
bliver den forskjelligt betingede Aarsag til forskjellige Aarsager, og
hver Virkning svarer nøiagtig til sin Aarsag. At denne dobbeltsidige
Reflexion ingenlunde er ørkesløs, men tvertimod afgjørende for den
mechaniske Idealitet, lære vi af den rationelle Mechanik.

Hvilens Aarsag er Kræfternes Ligevægt; Bevægelsens Aarsag er Ligevægtens
Ophævelse: det er i Grunden samme Aarsag og dog i Virkeligheden en
forskjellig Aarsag. Re\-%
% --- p. 76 ---
flectere vi paa Forskjellen, saa have Statiken og Dynamiken hver sin
Kreds af Opgaver og falde forsaavidt ud fra hinanden; reflectere vi paa
Eenheden, saa bliver den statiske Opgave dynamisk, og den dynamiske
statisk. Vexelbestemmelsen grunder sig paa to til hinanden svarende
Principer, de virtuelle Hastigheders Princip og D'Alemberts Princip for
de tabte Kræfters Ligevægt.

% „Tale“ er sat med et løst Mellemrum foran sidste Bogstav („Tal e“),
% men Bogstaverne indbyrdes ere ikke spatierede. Det er altsaa løs Sats
% og ikke Fremhævelse; sat som eet Ord uden \emph. Efterprøvet paa
% KB-Exemplaret.
Overalt, hvor der er Tale om et Princip, maa der, hvorledes Sagen end
iøvrigt skal forstaaes, nødvendig tænkes paa et Ubetinget. Ligevægtens
almindelige Betingelsesligning gjør i denne Henseende ingen Undtagelse;
thi vel erholder den en saadan Udvidelse, at den indgaaer paa alle
særlige Betingelser og finder Anvendelse paa et hvilketsomhelst System af
materielle Punkter; men netop den Maade, hvorpaa de særskilte Betingelser
gaae sammen til Eenhed, giver Ligningen Præg af et Ubetinget. Den i
Ligevægtens Charakteerligning, de virtuelle Hastigheders
% Notens Mærke er trykt *) og staaer efter Kommaet. Noten begynder paa
% S. 76 og fortsættes uden nyt Mærke øverst i Notefeltet paa S. 77
% („Totalvirkningen af alle Kræfterne ...“); hele Noten er efter
% Husreglen sat her ved Mærkestedet.
Formel,\footnote{% Trykfeil i Noten: „af Punktet. som kan tænkes“ ---
% Punktum hvor Sammenhængen kræver Komma. Begge Exemplarer (Google og
% KB) vise Punktum, saa Feilen er Sætterens og ikke dette Exemplars;
% Bogens Trykfeilblad retter den ikke. Gjengivet som trykt.
„Ved et Punkts virtuelle Hastighed forstaaes enhver uendelig lille
Forflyttelse af Punktet. som kan tænkes at finde Sted, uden Hensyn til de
paavirkende Kræfter, alene i Overeensstemmelse med de Betingelser,
Punktet er underkastet“. Kaldes Kraften P, den virtuelle Hastigheds
Projection paa Kraftens Retning $\delta p$, og angives
Vexelbestemmelsen af det Almene og det Enkelte --- det Logiske og det
Antilogiske --- ved en tilføiet Index k, saa er Betingelsesligningen for
Ligevægten af et enkelt Punkt $\Sigma\, P_{k}\, \delta p_{k} = 0$.
Kræfterne regnes positivt, men Projectionerne positive eller negative,
eftersom de falde enten paa Kraftens Retning eller paa dens modsatte
Forlængelse. Den for et System af Punkter gjældende Formel, hvori alle
Betingelser ere optagne, kan udtrykkes ved
\[
P_{1}\, \delta p_{1} + P_{2}\, \delta p_{2} \ .\ .\ .\ .\ + P_{n}\,
\delta p_{n} + \lambda_{1}\, \delta u_{1} + \lambda_{2}\, \delta u_{2}
\ .\ .\ .\ \lambda_{r}\, \delta u_{r} = 0.
\]
Leddene $P_{k}\, \delta p_{k}$ kaldes Kræfternes virtuelle Momenter,
Leddene $\lambda_{s}\, \delta u_{s}$ Modstandens virtuelle Momenter, og
Formlen fremstiller Principet. Summen af de virtuelle Momenter
$P_{1}\, \delta p_{1} + P_{2}\, \delta p_{2} + P_{3}\, \delta p_{3} +
\ .\ .\ .$ angiver i hvert Øieblik
% Her gaaer Noten over paa S. 77.
Totalvirkningen af alle Kræfterne. At denne Sum maa være Nul i det
Øieblik Ligevægtsstillingen passeres, er en ligefrem Følge af Begrebet om
Ligevægt, idet Kræfterne da hverken fremme eller modvirke Bevægelsen.}
ud\-%
% --- p. 77 ---
talte Grundtanke, en Grundtanke, der sammenfatter alle Betingelser i den
bestemte Aarsags betingede Eenhed, det Ubetingede, kan med fuld Ret
kaldes et Princip.

Ifølge dette Princip bliver jo vistnok Ligevægten prøvet paa en
Bevægelse, som dog ingen Bevægelse er, thi en virtuel Hastighed er en
Mulighed, som ikke bliver til Virkelighed. Men denne Mulighed,
Punkternes blot forestillede Forflyttelse, tjener imidlertid til at lade
Bevægelsen skinne igjennem i Ligevægten, lade Ligevægten vise sig som
ophævet Bevægelse og at gjøre det statiske Problem dynamisk.

I Formlen for de virtuelle Hastigheders Princip haves en Sum af virtuelle
Momenter, hvori hvert Led er et Product af Kraften $P_{k}$ og den paa
Kraftens Retning $p_{k}$ projicerede virtuelle Hastighed $\delta p_{k}$,
altsaa $P_{1}\, \delta p_{1} + P_{2}\, \delta p_{2} + P_{3}\,
\delta p_{3} \ .\ .\ .$ Det Eiendommelige ved denne Sum er, at den i
hvert Øieblik angiver Totalvirkningen af alle Kræfterne, hvor forskjellig
end dens Værdi kan være. Hvad der altsaa for Statiken er det Almene, at
nemlig Summen af de virtuelle Momenter udtrykkelig skal være Nul, kan da
med Hensyn paa Dynamiken betragtes som et specielt Tilfælde, en betinget
Aarsag, der gaaer sammen med en anden Aarsagsbetingelse og danner et
andet, mere omfattende Ubetinget, et Princip, som paa een Gang bestemmer
Aarsagen baade til Ligevægt og til Bevægelse --- D'Alemberts
% Ogsaa denne Note er mærket *) --- Siden bærer altsaa Slutningen af
% Noten fra S. 76 (uden Mærke) og dernæst denne, S. 77's egen Note.
% Noten fortsættes uden nyt Mærke i hele Notefeltet paa S. 78 og er
% efter Husreglen sat her ved Mærkestedet.
% Den anden Differentialkvotient er trykt med Exponenten 2 staaende OVER
% Bogstavet (d med 2 over x); Sætningsmaaden gjengives ikke, men Leddet
% sættes som d^{2}x_{k}/dt^{2}.
Princip.\footnote{Et System af Punkter $m_{k}$ være bestemt ved
retvinklede Coordinater og paavirket af Kræfterne $P_{k}$, hvis
Composanter efter Axerne ere $X_{k}$, $Y_{k}$, $Z_{k}$. Den virkelige
Bevægelse af $m_{k}$ i Tidselementet $dt$ kan antages at være den samme,
som hvis dette Punkt var frit og paavirket af Slutningshastigheden
$v_{k}$ og af en accelererende Kraft $Q_{k}$, som i Størrelse
% Her gaaer Noten over paa S. 78.
og Retning er forskjellig fra $P_{k}$. Antages nu ethvert Punkt $m_{k}$
underkastet foruden den givne Kraft $P_{k}$ desuden Kraften $Q_{k}$ og en
anden lig og modsat denne, nemlig $Q'_{k}$, saa er det givne System af
Kræfter ikke derved forandret. Det hele System bevæger sig, som om
Punkterne vare frie og alene underkastede Kraften $Q_{k}$. Den Deel af
Kræfterne, som medgaaer til at overvinde Tryk og Spændinger i Systemet
selv, faaer ingen Indflydelse paa Bevægelsen og maa i denne Henseende
betragtes som tabte Kræfter. De to paa $m_{k}$ virkende Kræfter $P_{k}$
og $Q'_{k}$ have en Resultant $p_{k}$; denne Resultant er \emph{den tabte
Kraft}, der ved Sammensætning med $Q_{k}$ giver som Resultant $P_{k}$.
Betegnes nu Composanterne af
\[
Q'_{k}\ \text{ved}\ -\frac{d^{2}x_{k}}{dt^{2}},\
-\frac{d^{2}y_{k}}{dt^{2}},\ -\frac{d^{2}z_{k}}{dt^{2}},
\]
idet $u_{1} = c_{1}$, $u_{2} = c_{2} \ .\ .\ .\ .$ angive
Betingelsesligninger mellem Coordinaterne, saa haves Charakteerligningen:
\[
\begin{aligned}
\Sigma\, m_{k} \biggl[ \biggl( X_{k} &- \frac{d^{2}x_{k}}{dt^{2}}
   \biggr)\, \delta x_{k}
   + \biggl( Y_{k} - \frac{d^{2}y_{k}}{dt^{2}} \biggr)\, \delta y_{k}\\
   &+ \biggl( Z_{k} - \frac{d^{2}z_{k}}{dt^{2}} \biggr)\, \delta z_{k}
   \biggr] + \lambda_{1}\, \delta u_{1} + \lambda_{2}\, \delta u_{2}
   \ .\ .\ .\ .\ \lambda_{r}\, \delta u_{r} = 0,
\end{aligned}
\]
som viser, hvorledes Dynamikens Problemer reduceres til Problemer af
Statiken. Jvnf. C. Ramus. Anal. Mech. S. 210.}
% --- p. 78 ---
% NOTATION: Aarsags- og Kraftbogstaverne paa S. 78--79 (m_k, P_k, Q_k,
% Q'_k, p_k og Aarsagerne A, A', A'' med Virkningerne B og C) høre til
% een og samme Række; efter Husreglen fra Batch 5 sættes hele Rækken i
% Mathematiksats, saa at Snittet ikke skifter midt i den.
Et System af materielle Punkter $m_{k}$, paavirket af Kræfterne $P_{k}$,
være givet. Da nu Kræfterne med deres Betingelser, ifølge vort
Foregaaende, falde sammen i Aarsagen, have vi altsaa i dette System een
eneste betinget Aarsag, Aarsagen $A$. Men Analysen viser, at en Deel af
Kræfterne $P_{k}$ ophæves af de i Systemet liggende Reactioner, medens en
anden Deel udelukkende anvendes paa Bevægelsen. Aarsagen $A$ afgiver
altsaa to indbyrdes forskjellige Virkninger, af hvilke den ene, $B$, sees
i Systemets Vedligeholdelse, den anden, $C$, i dets Bevægelse. Her er
samme Aarsag og dog to forskjellige Aarsager. Reflectere vi paa $P_{k}$,
saa er det samme Aarsag; thi det er $P_{k}$, der decomponeres, saa at en
% --- p. 79 ---
% Trykfeilbladet retter S. 79, Linie 1 f. o.: „Q_k læs Q'_k“ --- og
% Bladet trykker i netop denne Post intet Tankestreg foran „læs“.
% Rettelsen gjælder det ANDET $Q_{k}$ i Linien, det som er „lig og
% modsat $Q_{k}$“ og med $P_{k}$ danner Resultanten $p_{k}$; jvf. Noten
% S. 77--78, hvor $Q'_{k}$ staaer fuldt ud. Efterprøvet paa
% KB-Exemplaret ved Forstørrelse: der staaer ganske vist $Q_{k}$ uden
% Mærke (Arbejdsscanningen viser her en Klat, der ligner et Mærke).
% Efter Husreglen staaer Teksten her som trykt og Rettelsen er ikke
% indført.
Deel af $P_{k}$, nemlig $Q_{k}$, afgiver Bevægelsen, medens $Q_{k}$, der
er lig og modsat $Q_{k}$, danner med $P_{k}$ en Resultant $p_{k}$, der er
tabt for Bevægelsen; $P_{k}$ er altsaa decomponeret i $p_{k}$ og
$Q_{k}$. Reflectere vi paa Modsætningen mellem $p_{k}$ og $Q_{k}$, saa
have vi to forskjellige Aarsager, nemlig i $p_{k}$ Aarsagen til Systemets
Vedligeholdelse ($A'$) og i $Q_{k}$ Aarsagen til Bevægelsen ($A''$). Hvor
inderlig de saaledes adskilte Aarsager gribe ind i hinanden, sees deraf,
at en Forandring i den ene Punkt for Punkt medfører en Forandring i den
anden. Under Systemets Bevægelse varierer nemlig $Q_{k}$ og følgelig
ogsaa $p_{k}$. Summen af de tabte Kræfters og Modstandenes virtuelle
Momenter maa naturligviis i hvert givet Øieblik være Nul; men idet
$p_{k}$ og $Q_{k}$ bestemmes som Functioner af Tiden og da med det Samme
den til Virkningen $B$ svarende Aarsag $A'$, bliver $Q_{k}$ med den til
Virkningen $C$ svarende Aarsag $A''$ i hvert Øieblik bestemt, og omvendt;
Opgaven er paa een Gang baade statisk og dynamisk, fordi den af Systemet
betingede Aarsag selv er baade en Ligevægtens og en Bevægelsens Aarsag.

% Den indløbende Overskrift c) er trykt i HALVFED Antiqua, ikke
% spatieret som de græske Underoverskrifter α) β) γ). Efterprøvet paa
% KB-Exemplaret. Husmakroen \runin tjener ligefuldt; Snitforskjellen
% gjengives ikke.
\runin{c) Aarsagsmomenternes Eenhed.} Naar der i et mechanisk Phænomen
udtrykkelig skal sees paa det umiddelbart Virksomme, saa nævner man
Kraften og, forsaavidt denne er i Materien, da ogsaa Massen. Vil man
derimod særlig fremhæve alt det, som uden selv at være i og for sig
Virksomt, faaer Indflydelse paa Virksomheden og derved bliver et
Medvirkende, saa taler man om Betingelser. Har nu, som vi have seet,
Aarsagen alle Kræfter og Betingelser i sin Magt, saa er det heraf
forstaaeligt, at den mechaniske Aarsag optager og forener høist
uligeartede Momenter. I sit Væsen er Aarsagen Begreb; heri ligger, at
alle fra Virkeligheden hentede Aarsagsmomenter, hvad enten de iøvrigt
kaldes Kræfter eller Betingelser, maae være \emph{ideelle}. Men i sin
Virken er Aarsagen existerende; heri ligger, at de
% --- p. 80 ---
% Siden begynder plant (uden Indrykning) og fortsætter Sætningen fra
% S. 79, som ender „... heri ligger, at de“. Ingen Orddeling over
% Sideskiftet: Ordgrændsen er reen, saa der behøves hverken nyt Afsnit
% eller \-.
virksomme Aarsagsmomenter alle maae være \emph{reelle}. I det
naturbestemte Grundforhold mellem den mechaniske Aarsags ideelle og
reelle Momenter maa Forskjelseenheden af ideel og reel Aarsag altsaa
blive tydelig udtalt.

% Den græske Underoverskrift er indløbende og spærret; Bogstavet er
% læst paa KB-Exemplaret (α, ikke a).
\subrunin{α) Momenternes Idealitet.} Ligesom Kræfter og Betingelser
gaae op i de derved bestemte Aarsager, saaledes gaae Aarsagerne op i de
existerende Legemer; Legemerne ere fra dette Synspunkt at betragte som
de eneste i Rum og Tid virkende Aarsager: at decomponere Legemerne i
legemlige Elementer er at decomponere Aarsagerne i Aarsagsmomenter.

Det er lærerigt at see, hvorledes den rationelle Mechanik, skjøndt dens
Blik allernærmest er fæstet paa Kræfter, Betingelser og Love, netop for
Lovenes Skyld forvandler de legemlige Aarsager til Aarsagsbegreber og
uden ret at vide deraf, stiller en heel Verden af mechaniske Aarsager
frem i Idealitetens Lys som et System af legemliggjorte
Aarsagsbegreber.

For ret at trænge ind i de legemlige Aarsagsforhold, begynder
Mathematikeren med at løse Legemet op i dets Elementer: $M = \int dm$.
Men han bliver ikke staaende ved den uendelig tomme Forklaring, at
ethvert Legeme bestaaer af uendelig mange smaa Dele; han mærker sig
Begrændsningsmaaden; han tager Sigte paa Legemlighedens Begreb og
sættes derved i Stand til punktlig at bestemme saavel et Legemes indre
Bygning som dets ydre Form, dets Figur.
% Symbolrækken ϱ, V, dm, M holdes i een Skrift (Husreglen fra Batch 5),
% skjøndt V og M staae som opreiste Bogstaver i Trykken.
Er $\varrho$ Tætheden, $V$ Rumfanget, saa er $dm = \varrho\,dv$ og $M =
\int \varrho\,dv$; er Legemet af ulige Tæthed, altsaa $\varrho$
variabel, saa kommer det ham an paa at bestemme det til ethvert
Volumenelement svarende $\varrho$ som en Function af dette Elements
Coordinater. Men hvorledes faaer han Functionen bestemt? Ved at
forudsætte bestemte legemlige Forhold i det paagjældende Legeme, altsaa
ved at gaae ud fra det forestillede Legems Begreb. Han har ikke forud
et givet
% --- p. 81 ---
Legeme, hvorefter han retter Begrebet --- Forestillingen herom er fra
dette Synspunkt kun en forstyrrende Indbildning; nei, han har
formedelst den intellectuelle Anskuelse Legemet fuldstændig i sin Magt;
han bestemmer det i Kraft af sit Begreb og udleder Loven af Begrebet.
Enten han siger: „dersom Legemet er af ulige Tæthed“, eller han siger:
„et Legeme, som skal være af ulige Tæthed“, er med Hensyn til
Legemsbegrebet det Selvsamme. Det er netop det Rationelle i den
rationelle Mechanik, at dens Functioner og Love ikke støtte sig til
empiriske Inductioner, men udvikles af Begreber. De legemlige
Elementer ere ikke Erfaringsgjenstande; de ere Forestillingsreflexer,
som den begribende Forstand modtager fra Anskuelsen, ideelle Momenter i
Billeder af legemlig Realitet.
% „ideelle“: begge OCR-Vidner læste her „ideella“/„ideell&“; KB-Copien
% viser tydelig „ideelle“.

I ubestemt Almindelighed er det ene $dm$ som det andet, et Element, der
igjen kan opløses i Elementer; men i Kraft af Begrebet kan det ene $dm$
blive aldeles uligeartet med det andet. Et Element af en Kugle og et
Element af et Parallelepipedum kunne, hvad den uendelige Forsvinden i
uendelig Lidenhed angaaer, ganske vist gjælde for ligeartede; hvad
enten de uendelig smaa Begrændsningslinier opfattes som Elementer af
rette Linier, $dxdydz$, svarende til et retvinklet Coordinatsystem,
eller der ved Anvendelse af polære Coordinater i Begrændsningen optages
krumliniede Elementer,
% Elementet er her trykt med θ, ω, θ; paa S. 84 skriver Nielsen det
% samme Element med α, α, θ. Uoverensstemmelsen er hans egen og staaer
% i begge Exemplarer; gjengivet som trykt.
$r^{2}\sin\theta\,d\omega\,d\theta\,dr$, gjør i saa Henseende Intet til
Sagen. Men naar det Endelige skal frembringes af det Uendelige, naar
en Integration, $\int\int\int dxdydz =
\int\int\int r^{2}\sin\theta\,d\omega\,d\theta\,dr$, skal udføres,
deels saaledes, at Legemet bliver et Parallelepipedum, deels saaledes,
at det bliver en Kugle, da erkjendes det, at de to Legemer have
uligeartede Elementer: Begrebsmomenterne ere forskjellige, fordi
Begreberne selv ere forskjellige.

Ligesom ethvert Frø har i sig Anlæg til en Plante af en bestemt Art,
svarende til en bestemt Slægt, saaledes har
% --- p. 82 ---
ogsaa ethvert af sit Begreb bestemt Legemselement i sig Anlæg til ved
Integrationen at voxe en eiendommelig Væxt og blive til et bestemt
Legeme.
% Brøkerne staae i Trykken som opsatte Brøker midt i Linien; \tfrac
% følger Husreglen fra Batch 4. Grændsen „r = o“ er sat med Bogstavet
% o, ikke med Nullet — som i Noten til S. 21; saaledes gjengivet.
At Massen af en Kugleskal $dm = d\tfrac{4}{3}\pi r^{3} =
4\pi r^{2} dr$, naar der integreres mellem Grændserne $r = o$ og $r =
a$, maa give en Kugle $= \tfrac{4}{3}\pi a^{3}$, er lige saa
nødvendigt, som at Frøet af en Rose maa, naar det kommer til Væxt og
Udvikling, frembringe en Rose. Kategorierne \emph{Slægt}, \emph{Art}
og \emph{Individ} have lige saa vel Gyldighed i Mechaniken som i
Botaniken.

At en homogen Kugle under Forudsætning af, at Tiltrækningen forholder
sig omvendt som Qvadratet af Afstanden, maa tiltrække et udvendigt
Punkt, som om hele Kuglemassen var samlet i Centrum, dette virkelige
Aarsagsforhold bestemmes ikke ved Iagttagelse af nogen udvortes
Kjendsgjerning, men alene ved Indsigt i Kuglens og Tiltrækningens
Væsen.
% Arbeidsscanningen viser et lille høit Mærke efter „Afstanden,“, som
% tesseract læste som Anførselstegn. KB-Exemplaret har her et reent
% Komma og Intet andet: Mærket er en Plet i den ene Copi, ikke Typen.
Selv da, naar empiriske Bestemmelser indgaae i Beregningen, er det dog
altid et System af apriorisk-legemlige Aarsagsbegreber, der afgiver det
rationelle Grundlag. Physiken har veiet Kloderne, vistnok ifølge
empiriske Iagttagelser, men paa Idealitetens Vægtskaal.

\subrunin{β) Momenternes Realitet.} Man skulde nu mene, at den
analyserende Mathematiker maatte være fuldt forvisset om sine
Legemsbegrebers physiske Gyldighed, da han jo er fuldt forvisset om de
fundne Loves physiske Gyldighed; men nei! Om Lovenes Naturgyldighed
nærer han ingen Tvivl, de stadfæste sig for ham gjennem alle de
Ændringer, som deres Opfyldelse under særlige Betingelser i
Virkeligheden medfører. Men med Forholdet mellem Legemsbegreberne og
de virkelige Legemer er det en anden Sag; her er den mathematiske
Physiker tilbøielig til at ansee de for Analysen nødvendige
Grundbestemmelser for vilkaarlige Betegnelser, subjective
Forestillingsmaader, der Intet have med Legemerne i Naturen,
% --- p. 83 ---
følgelig heller Intet med de legemlige Aarsager selv, at gjøre. At
stole paa Lovene og dog tvivle om de Begreber, af hvilke Lovene
udledes, er tilvisse en høist besynderlig Inconseqvens; men hvorfra
stammer den? Fra en uklar Opfattelse af Momenternes Realitet.
% „høist“: Google-Laget læser „heist“ (det kjendte ø→e-Feillæsning);
% KB-Copien viser Skraastregen i ø tydelig.

Et Element som $dm$ er ikke til i Naturen: \textit{un infiniment petit
est une grandeur moindre que toute grandeur donnée de la même nature
(Poisson)}; det er Delingsproblemet, som her paatrænger
sig.\footnote{% trykt Mærke: *)
% Titlen er her spærret; i den tilsvarende Note til S. 30 staaer den
% uspærret. Gjengivet som trykt paa hver sin Side.
Jvnfr. \emph{Mathematik og Dialektik}. Kbhvn. 1859; Delingsproblem
S. 116--131.}
Det gaaer med Legemselementet $dm$ som med Rumselementet $ds$ og med
Tidselementet $dt$; de mathematiske Elementer ere ikke til og ere dog
til. De ere ikke til som særlig bestaaende Dele, men de ere Momenter i
den virkelige Tilværelse. Det materielle Legeme er lige saa vist et
Integral af materielle $dm$ som det mathematiske Legeme af mathematiske
$dm$; den virkelige Begivenhed er lige saa vist et Integral af
Begivenhedsmomenter, som Tiden, hvori den foregaaer, er et Integral af
Tidsmomenter. Spørgsmaalet dreier sig slet ikke om Forholdet mellem
det absolut Enkelte og det deraf Sammensatte, thi Elementet er lige saa
lidt i reel som ideel Forstand noget absolut Enkelt. Grundspørgsmaalet
er: hvorledes kan det, som er uden Størrelse, frembringe Størrelse?
hvorledes kan det, som selv er uden Udstrækning, frembringe noget i
Virkeligheden Udstrakt?

For at aflede det Reelle af det Ideelle indførte Leibnitz Monaderne og
kaldte dem metaphysiske Punkter:
% Sætteren har ikke Plads til det sidste t i „soient“ og har sat det
% nedsænket ved Linieenden. Begge Exemplarer viser det samme, saa det
% er Sætterens Nødhjælp og ikke en Skade i den ene Copi; Ordet er
% gjengivet helt.
\textit{il n'y a que les points metaphysiques ou de substance, qui
soient exacts et réels, et sans eux il n'y auroit rien de réel,
puisque sans les véritables unités il n'y auroit point de multitude.}
Men af de monadiske Punkter, de sjælelige Eenheder, var
% --- p. 84 ---
han dog ikke i Stand til at deducere den legemlige Virkelighed; det
blev, hvad Materien angaaer, ved „dunkle, forvirrede Forestillinger“.
I sin Fremstilling af „Reflexionsbegrebernes Amphibolie“ paaviste Kant
meget rigtig, at der er en Forskjel mellem Materien som et paa
„Forestillingskraften“ bygget Forstandsbegreb og Naturphænomenet
Materie, der er Gjenstand for Sandsning, og hvis Opfattelse er betinget
af vor Sandseanskuelses aprioriske Form. Fra Intellectualphilosophiens
Synspunkt er Materien kun et Complex af Forestillinger, et Phænomen af
det sig forestillende Væsen, et Umiddelbarhedens Skin af det
Intellectuelle; for den kritiske Philosophie er Kløften mellem det
subjectivt Ideelle og det objectivt Reelle saa ubegribelig stor, at
Materien hverken kan frembringes af Forestillingen, eiheller vedkomme
„Tingen i sig“. Begge Synsmaader ere eensidige. Den virkelige Eenhed
af det Ideelle og det Reelle er en Eenhed i Modsætning, en Magteenhed.
Det Ideelle i $dm$ frembringer ikke det Materielle, saa lidt som
omvendt; men begge hinanden modsatte Grundbestemmelser frembringes i og
ved den objective, almene Magt, der i sin Fordobling er baade ideel og
reel.

% Bindestregen i „Ideel-reel“ staaer virkelig i den trykte Overskrift
% (kontrolleret paa KB-Exemplaret); jvnfr. „den ideel-reelle Aarsag“ i
% Brødteksten S. 87.
\subrunin{γ) Ideel-reel Aarsag.} For end yderligere at bestemme,
hvorledes det i mechanisk Forstand forholder sig med Aarsagsbegrebet og
den virkende Aarsag, vælge vi med Henblik paa det Foregaaende et
bestemt Exempel: Dannelsen af en Kugle.

Ved at gaae ud fra det prismatiske Element
\[ r^{2}\sin\alpha\, d\alpha\, d\theta\, dr \]
% Integralernes nedre Grændser ere sat med Bogstavet o, ikke med
% Nullet, ligesom „r = o“ paa S. 82; saaledes gjengivet.
og foretage Integrationerne $\int_{o}^{\pi}\int_{o}^{2\pi}\int_{o}^{r}
r^{2}\sin\alpha\, d\alpha\, d\theta\, dr$ indseer man, hvorledes det
ene Element opstaaer af det andet, og den hele Kugle af Elementernes
System.\footnote{% trykt Mærke: *)
Jvnfr. Forelæsninger over Phil. Propæd. 1861--62. S. 31--33.}
Kuglens
% --- p. 85 ---
Begreb i Eenhed med den frembringende Anskuelse er Kuglens ideelle
Aarsag. Men det saaledes Frembragte er ogsaa kun noget Ideelt,
kun Kuglens Form, ikke den virkelige Kugle. Det Indbegreb af
Formbetingelser, hvortil Integrationerne svare, er vilkaarlig valgt af
særligt Hensyn til bekjendte Functioner og staaer ikke i udelukkende
nødvendig Forbindelse med det Indbegreb af materielle Bestanddele og
udvortes Betingelser, hvori Dannelsen af en virkelig Kugle maa have sin
reelle Aarsag. Begrebets Momenter og Virkelighedens Elementer falde i
to fra hinanden adskilte Kredse: her er en gabende Kløft imellem det
Ideelle og det Reelle.

Denne Kløft er i Sagen udfyldt af det guddommelige Magtbegreb, hvis
Energie ligelig omfatter baade det Ideelle og det Reelle. Men en
menneskelig Indsigt i det guddommelige Begreb opnaaes kun paa
menneskelige Vilkaar; derfor er Exemplet valgt.

Fra den ideelle Side kan Aarsagsbegrebets mechaniske Væsen
tydeliggjøres ved Hjælp af de fordrede Integrationer. Af det rene Intet
kan Intet dannes. Hver Dannelse maa begynde med Noget; det Noget,
hvormed der begyndes, er selv en begyndende Dannelse, et
Dannelseselement. Elementets Mulighed ligger i den begribende Anskuen;
dets definitive Bestemthed sættes vilkaarlig som Dannelsens første
Forudsætning: det prismatiske Element med sine tre uendelig smaa
Dimensioner er kun en almindelig Dannelsesbetingelse. Hvad kan der nu
dannes paa en saadan Betingelse? Ifølge den opstillede Formel ere de
successive Dannelser særlig bestemte ved det særlige Begreb. Integrere
vi først med Hensyn paa r, saa fremkommer der en elementær Pyramide, et
Element af een endelig og to uendelig smaa Dimensioner; integrere vi
dernæst med Hensyn paa $\theta$, saa fremkommer der en Kegleskal, et
Element med to endelige og een uendelig lille Dimension; integrere vi
endelig med Hensyn paa $\alpha$,
% --- p. 86 ---
erholde vi den søgte Form, Kugleformen med tre endelige Dimensioner.
Altsaa vil, ifølge det herskende Begreb, der er Kugledannelsens ideelle
Aarsag, Prismeelementet betinge Pyramideelementet, Pyramideelementet
Kegleskallen, og Kegleskallen Kuglen.\footnote{% trykt Mærke: *)
Ved at foretage Integrationerne i sex forskjellige Følgerækker ville
nye Elementer vel indtræde og Betingelsernes Gang ændres, men da det
herskende Begreb er det samme, saa ere disse Ændringer uvæsenlige.
Først naar Kuglebegrebet ombyttes med et andet Begreb, maae Elementerne
benyttes paa en anden Maade, og Dannelsen bliver en anden.}

Denne Betingelsesrække viser, hvad det mechaniske Begreb, bortseet fra
endelige Forstandsoperationer, er i sig selv som Aarsagsbegreb. Det er
ikke abstract logiske Almeenformer, men punktuelle Betingelser, der
udgjøre Begrebets Indhold. Hvad der enkeltviis sees som vilkaarligt,
bliver i Sammenhængen nødvendigt. Ingen Enkeltbetingelse er Noget for
sig, men alle ere de Momenter for hverandre indbyrdes. Udenfor
Begrebet ere de indvortes ideelle Betingelser intetsigende
Forestillinger; kun idet de gjennemtrænges og beaandes af det herskende
Begreb, have de i Begribningens Eenhed hver sin eiendommelige
Betydning.

Til Betingelserne i det Ideelle svare Betingelserne i det Reelle Punkt
for Punkt, naar Begrebet er et Magtbegreb. Magtens Dialektik er, som
vi have seet, dette Dobbeltsidige, at Momenterne ere tilværende Led, og
de tilværende Led Momenter. Udenfor Magtbegrebet ere de udvortes,
reelle Betingelser intetsigende Forestillinger. Enhver Betingelse maa
nødvendig være Betingelse for Noget, men det Noget, hvis Betingelse den
er, kan ikke umiddelbart bestemmes ved en anden Betingelse. For at
Betingelserne kunne bestemmes ved hverandre, maae de i Sagen selv
% S. 86 slutter midt i Sætningen; Teksten løber ned til Sidefoden (med
% Noten under Stregen) og fortsætter paa S. 87 med „reflecteres i
% hverandre; ...“. Der staaer INGEN Streg paa S. 86: Afdeling A løber
% videre over paa S. 87 og sluttes først der, hvor den korte,
% centrerede Streg staaer foran Overskriften B. Mechanisk Realitet.
% Slutningen af Afdeling A (to Afsnit) hører altsaa til næste Batch.
%  --- B. Mechanisk Realitet  (pp. 87–105) ---
%  MATHEMATICAL — work from the image; the OCR is not a witness here.
% --- p. 87 ---
% SØMMEN: Sætningen fra S. 86 („... maae de i Sagen selv“) fortsætter
% umiddelbart her. Der maa IKKE staae en tom Linie foran „reflecteres“,
% naar dette Stykke splejses ind — det vilde skabe et falsk Afsnitsskifte.
reflecteres i hverandre; den Sag, hvori de gjensidig reflecteres, er
Aarsagen, Begrebet i Magt, den ideel-reelle Aarsag.

Vende vi nu med disse Forudsætninger tilbage til vort oprindelige
Spørgsmaal: \emph{hvorledes kan det, som selv er uden Udstrækning,
frembringe noget i Virkeligheden Udstrakt}? vil en punktlig Besvarelse
kunne gives. Jo større Kuglemassen er, desto større er den Kraft,
hvormed den tiltrækker. Opløses Kuglen i sine materielle Elementer, saa
opløses Kraften i lutter tilsvarende Momenter: hvert Kraftmoment er
ligesaa afhængigt af sit eiendommelig bestemte Masseelement som den
samlede Kraft af den hele Masse. At Kraftmomenterne kunne integreres til
endelig Kraft er betinget af, at Masseelementerne kunne integreres til
endelig Masse. Integrationen forudsætter, at der i hvert Masseelement er
et lige saa oprindeligt Anlæg til virkelig Udstrækning, umiddelbar
Realitet i det Extensive, som der i hvert Kraftmoment er Anlæg til
virkelig Styrke, reflecteret Realitet i det Intensive. Den ideel-reelle
Aarsag saavel til Materien som Kraften er Magten, Aandens Magt, den
guddommelige Gjenstandsmagt.

% Kort, centreret Streg midt paa S. 87: her — og først her — sluttes
% Afdeling A. Mechanisk Idealitet, som løber to Stykker ind paa denne Side.
\streg

\afdeling{B.}{Mechanisk Realitet.}

I mechanisk Realitet er Aarsagen at opfatte som umiddelbar, naturlig
Gjenstandsmagt. Det kommer da til at see ud som var Begrebet Aarsag kun
en Benævnelse, thi Gjenstanden træder i Aarsagens Sted: den hertil
svarende Undersøgelse bevæger sig om \emph{Magt og Aarsag}. Endvidere er
Gjenstandsmagten Magt i Andet. Det kommer da til at see ud, som om de
existerende Aarsager faldt udenfor hverandre
% --- p. 88 ---
i al Tilfældighed: Undersøgelsen bevæger sig om \emph{Udstykning af
Aarsager}. Endelig forudsætter den udvortes Aarsagernes Mangfoldighed et
Grundlag for mangesidig Vexelvirkning, saa at Leddene i deres indbyrdes
Sammenhæng optræde som foraarsagede Aarsager, og Realitetsspørgsmaalet
bliver, hvorledes det oprindelig forholder sig med de uendelige
\emph{Aarsagsrækker}.

\runin{a) Magt og Aarsag.} En nøiagtig bestemt Adskillelse af, hvad der
tilhører Gjenstanden, det udvortes Virkelige, og hvad den betragtende
Bevidsthed selv medbringer til Gjenstandens Opfattelse, er ikke saa let,
som det ved første Øiekast kunde synes. Erklæringer som denne: „der
gives hverken Legemer eller Aarsager, men kun Systemer af
Kjendsgjerninger, Grupper af nærværende eller mulige Bevægelser og af
nærværende eller mulige Tanker“, vise en Bestræbelse for at fastsætte
Adskillelsen, men afgive tillige Beviser paa den Overfladiskhed, hvormed
en umoden Reflexion kan bilde sig ind at være hævet over dybere gaaende
Undersøgelser. Siger man: Gjenstanden \emph{Steen} er ikke et Legeme;
\emph{Gjenstanden}, Subjectet i Dommen, er en \emph{Kjendsgjerning}, der
ligger udenfor Bevidstheden, men \emph{Legeme}, Dommens Prædicat, er et
Ord, et Almeenbegreb, der falder indenfor Bevidstheden: saa have vi
Modsigelsen. I den ene Dom er det Ordet Legeme, der er Prædicat; men i
den anden er det Ordet Kjendsgjerning. Skal Prædicatet Legeme ikke
vedkomme Gjenstanden som Gjenstand, men alene den prædicerende
Bevidsthed, da maa det Samme nødvendig gjælde om Prædicatet
Kjendsgjerning. Vil man indskærpe, at det udelukkende er
Kjendsgjerningen, der er Gjenstand, saa bliver det endnu værre, thi i
Dommen: Kjendsgjerningen er Gjenstand, er Gjenstanden selv Prædicat og
gaaer ind i Bevidstheden. Siger man: Kjendsgjerningen er ikke Aarsag;
thi Aarsagen er et Begreb, men Kjendsgjerningen er kun Kjendsgjerning:
da er det aabenbart, at
% --- p. 89 ---
man driver Spil med Ord. I det ene Tilfælde lader man Ordet vise ud over
sig selv og pege paa det, som ikke er et Ord, men en existerende Ting; i
det andet Tilfælde lader man Tingen falde sammen med Ordet og gjør
Tingen til et Ord. For at ende dette Ordspil, er det nødvendigt at
besinde sig paa Forholdet mellem det Logiske og det
Antilogiske.\footnote{% trykt Mærke: *) — eneste Note paa S. 89, og den
% staaer helt paa denne Side.
Jvnf. Grundideernes Logik I. Side 146--166. Grundideernes Logik i kort
Begreb. Side 28 og flgd.}

I antilogisk Forstand er der slet ingen Tanker eller Begreber i Naturen:
Alt er Virkelighed, Alt er Gjenstand, Ting, Kjendsgjerning i det Mindste
som i det Største. Ordene: Gjenstand, Ting, Kjendsgjerning ville ikke
forstaaes som Ord; de vise ud over sig selv, de tale til den Sandsende,
de sige til Betragteren: luk Øinene op, kom og see! Men i logisk
Forstand er Kjendsgjerningen Billede, Forestilling, Begreb. Den
forskende Bevidsthed forefinder i den hele Natur intet Andet end lutter
Billeder, lutter Forestillinger, Tanker og Begreber; og det af den gode
Grund, at Bevidstheden, hvor travlt den end har med at opdage
Gjenstande, Ting og Kjendsgjerninger, møder i Alt sine egne Opfattelser
og umulig paa noget Punkt kan springe ud over sig selv. Jorden
tiltrækker Maanen. I antilogisk Forstand er Jorden Jord, Maanen Maane og
Bevægelsen en Kjendsgjerning, der iagttages; men Aarsagen iagttages ikke
som Aarsag, Virkningen ikke som Virkning, thi Begreber ere ikke
Gjenstande for umiddelbar Iagttagelse. I logisk Forstand derimod er
Jorden Aarsag, og Maanebevægelsen Virkning, thi det virkelige Forhold
mellem Jord og Maane er umuligt uden et Forholdsbegreb af almægtig
Begriben.

Erfaringsforholdet mellem Gjenstand og Bevidsthed egner sig ret til at
opklare Sagforholdet mellem Magt og Aarsag. Gjenstandens Realitet er
dens Magt; Magten kjendes
% --- p. 90 ---
paa Realiteten, og omvendt. At Realiteten er Magt, bringes til
Bevidsthed gjennem Erfaringens umiddelbare Indtryk. Kunde Indtrykket
ikke med Vished antages at udgaae fra Gjenstanden som Virkning fra
Aarsag, saa vilde en Erfaring af Gjenstanden være umulig; man kunde jo
ikke vide, at der udenfor Bevidstheden virkelig var nogen Gjenstand. I
Erfaringsforholdet er Gjenstanden Aarsag og dens Opfattelse Virkning.
Gjenstanden er Aarsag, fordi den er Magt, og den er for Bevidstheden
Magt, fordi den er Indtrykkets Aarsag. Ligesaa er Gjenstandens
Opfattelse, dens Billede og tilsvarende Begreb, en ved Indtrykket
bestemt Virkning: Virkningen bestemmes ved at henføres til Aarsagen. At
tvivle om Aarsagen og stole paa Erfaringen er altsaa meningsløst, da
Erfaringsforholdet er et Forhold af Aarsag og Virkning. I Erkjendelsen
af den reelle Aarsag gaaer Tanken sammen med Sandsningen, det Logiske
med det Antilogiske. Fra logisk Side er Gjenstanden Aarsag, fra
antilogisk Side er Aarsagen Gjenstand: Erfaringens Sprog er i
Kjendsgjerningens Stiil et Magtsprog, der taler for Aarsagens Realitet.

Men hvoraf veed nu Iagttageren, at ogsaa Forholdet mellem Legemerne
indbyrdes er et virkeligt Aarsagsforhold? Alene deraf, at Erfaringens
Grundforhold selv er et Aarsagsforhold. At Sandseorganerne paavirkes af
den sandselige Gjenstand, er et umiskjendeligt Erfaringsbeviis for, at
Legeme paavirker Legeme, at, med andre Ord, Aarsagsforholdet har
legemlig Gyldighed. Og dog er, hvad der ligefrem følger af Extremerne i
Magtens Væsen, Modsætningen af det Logiske og det Antilogiske saa brat,
at en Iagttager kan være fuldt forvisset om Aarsagsbegrebets
Anvendelighed som Erfaringsbegreb og endda nære Tvivl om dets
aprioriske Værdi som tilforladeligt Udtryk for Tingenes Væsen.

Det mechaniske Phænomen, hvori Aarsagsforholdet paa
% --- p. 91 ---
en ret umiddelbar og anskuelig Maade udtaler sig som Magtforhold, er
Stødet. Aarhundreder igjennem har Stødet været anerkjendt som virkelig
Aarsag, før den menneskelige Viden naaede saa vidt, at den kunde
fastsætte Lovene for dets Virkning eller, hvad der er det Samme, faae
det endnu ubestemte Aarsagsbegreb nøiagtig bestemt. Af de sex
Sætninger, Cartesius har opstillet, sees det, at han forfeiler Opgaven,
idet han mere retter sig efter sit selvdannede Begreb end efter
Kjendsgjerningerne. Hans Grundtanke, at Bevægelsesmængden i hele den
legemlige Verden altid maa blive sig selv lig, forleder ham nemlig til
den Antagelse, at to i Retning modsatte Bevægelser ikke ophæve hinanden
som positive og negative Størrelser, men at de paa denne Maade
sammenstødende Legemer maae gaae tilbage, hvert med samme Hastighed, som
% Det latinske Citat er sat med kursiv mod Sidens antikva. „æqvalia“ og
% „æqve“ er Bogens egen qv-stavemaade, efterseet paa KB-Exemplaret.
det havde før Stødet --- \textit{si duo illa corpora essent plane
æqvalia et æqve velociter moverentur, B a dextra versus sinistram,
C a sinistra versus dextram,
cum sibi mutuo occurrerent, reflecterentur et postea pergerent moveri:
B versus dextram, C versus sinistram, nulla parte suæ celeritatis
amissa.} Den for uelastiske Legemer gjældende Formel
% Tegnet i Tælleren er ÷, Nielsens Minus (jvnf. S. 46), efterseet paa
% begge Exemplarer.
$x = \dfrac{MH \div mh}{M + m}$, som for Massen $M = m$ og Hastigheden
$H = h$ giver Hastigheden efter Stødet $x = 0$, viser nu, at Cartesius
slet ikke er blevet opmærksom paa Forskjellen mellem elastiske og
uelastiske Legemer.

Feiltagelser af denne Natur tyde paa, at den bestemmende Magt ikke
ligger i det menneskelige Begreb, men i Kjendsgjerningen. Forskerens
Erkjendelse betinges af Forsøg; han maa experimentere. Og dog kan
Kjendsgjerningen, umiddelbart seet, ikke lære ham Noget; den kan opløses
i en Mangfoldighed af Kjendsgjerninger, men hver af Delene er som det
Hele --- en Kjendsgjerning, intet Begreb. For at lære Noget af
Kjendsgjerningen, maa han selv tænke og dømme, altsaa medbringe
Begreber. Men naar det er
% --- p. 92 ---
Forskeren, som af egen Bevidsthed maa lægge Begreberne til, de specielle
lige saa vel som de almene, saa ere alle hans Kjendsgjerninger jo, hvad
Begribningen angaaer, kun Anledninger, ikke virkende Aarsager: saa lærer
han kun anledningsviis af Kjendsgjerningen, hvad han oprindelig lærer af
sig selv. Saa kommer altsaa dog den bestemmende Magt til at ligge i
Begrebet.

Det elastiske Legeme, f. Ex., er i antilogisk Bestemthed en
Kjendsgjerning, men Elasticiteten er et Begreb. Hvorledes kommer nu
Iagttageren fra Kjendsgjerning til Begreb og fra Begreb til
Kjendsgjerning? Ere, siger han, B og C elastiske, ville de ved at støde
mod hinanden sammentrykkes til et vist Maximum. I dette Maximum af
Sammentrykning ere de, om end kun for et Øieblik, uelastiske, og den
fælles Hastighed bliver, ligesom for uelastiske Legemer, at udtrykke ved
% Tegnet i Tælleren er her ±, ikke det ÷, som staaer i samme Formel paa
% S. 91 og i de følgende Parenteser paa denne Side: et Plus over en
% Streg, uden Punkter. Efterseet paa begge Exemplarer, som stemme
% overeens; gjengivet som trykt.
$x = \dfrac{MH \pm mh}{M + m}$. B har altsaa tabt Bevægelsesmængden
$M(H \div x)$, og C Bevægelsesmængden $m (h \div x)$. Men nu ligger det
i elastiske Legemers Væsen, at de efter Sammentrykningen igjen stræbe at
vende tilbage til den oprindelige Form; lignende Indvirkninger, som de,
der fandt Sted under Sammentrykningen, finde da ogsaa Sted under
Udvidelsen, hvorved B og C atter tabe samme Bevægelsesmængde som før.
Hele Tabet i Bevægelsesmængde er altsaa for B's Vedkommende
$2 M (H \div x)$ og for C's $2 m (h \div x)$. Dersom Iagttageren kunde
gjøre punktlig Rede for, hvormeget han i denne Forklaring har fra
Kjendsgjerningen, og hvormeget han har fra sit eget Begreb, vilde en for
Erfaringslæren betydningsfuld Opgave være løst.

Den givne Forklaring er bygget paa Iagttagelse af Kjendsgjerninger; og
dog gaaer den langt ud over de iagttagne Kjendsgjerninger, thi den er
bygget paa Begreber og Domme. Hver Kjendsgjerning er antilogisk et
\emph{Dette}; kun
% --- p. 93 ---
Dette, det Enkelte, kan hver Gang iagttages, det Almene kan ikke
umiddelbart iagttages. \emph{Dette} Legeme kan iagttages, men et Legeme
i Almindelighed eller det Almindelige i Legemet kan ikke iagttages;
\emph{denne} Sammentrykning og \emph{denne} Spændighed kunne iagttages,
men Spændigheden i Almindelighed eller det Almindelige i Spændingen kan
ikke iagttages. Ikke desto mindre tør Iagttageren ubetinget stole paa
sit ved Kjendsgjerningen bestemte Begreb. At der i Naturen ikke findes
enten absolut spændige eller absolut uspændige Legemer, afholder ham
ingensinde fra at lægge Begreberne Spændighed og Uspændighed til Grund
for en mechanisk Analyse og trods alle særlige Betingelser lade
Begrebsconseqvenserne være gjældende for Kjendsgjerningernes almene
Væsen. Ledet af sit Begreb iagttager han endog i Kjendsgjerningen selv
ikke blot det Synlige, men ogsaa det Usynlige.

Hvorledes det gaaer til, naar to Legemer, de være nu mere eller mindre
elastiske, under Sammenstødet øiebliksviis forandre hinandens Form, kan
ikke sees med sandselige Øine, og dog seer Iagttageren mathematisk, hvad
der skeer. Før Stødet havde Legemet B med Massen M og Hastigheden H en
levende Kraft $MH^2$; efter Stødet, d. v. s. i Maximum af
Sammentrykning har det en levende Kraft $Mx^2$. Legemet har altsaa i
levende Kraft tabt $M (H^2 \div x^2)$, og hvorfor? Fordi det ved at
sammentrykke Legemet C har udført et Arbeide. Kaldes det Tryk P, som
svarer til Indtrængningen s, saa er $P = f(s)$, og
$M (H^2 \div x^2) = 2 \int\limits_{s_2}^{s_1} Pds$. Det ligger i Sagens
% Bemærk Fortegnene paa denne Side: ÷ (Nielsens Minus) i Parenteserne
% $M (H^2 \div x^2)$, men en almindelig Minusstreg i „$s_1 - s_2$“.
% Begge Exemplarer stemme overeens; gjengivet som trykt.
Natur, at Intervallet $s_1 - s_2$ maa være større, naar Indtrykket under
forøvrigt samme Vilkaar gjøres i et blødere end i et haardere Legeme.
Skal i begge Tilfælde venstre Side af Ligningen være eens, altsaa
$M (H^2 \div x^2)$ have samme Værdi, saa maa høire Side, betragtet som
Sum af $Pds$, indeholde, midt i Uendeligheden, et større
% --- p. 94 ---
% S. 94 aabner med sex Linier sat med MINDRE Skrift end Texten;
% Skriftskiftet falder netop ved Bladskiftet, midt i Sætningen, og den
% mindre Skrift løber til Afsnittets Slutning („... ved det blødere.“).
% Efterseet paa begge Exemplarer, som stemme overeens. Gjengivet med
% \small efter Husreglen fra Trykfeil-Bladet.
{\small Antal Addender for det blødere end for det haardere. Summen er
nemlig i begge Tilfælde lige stor, thi, som Ligningen viser, afhænger
den kun af Stødets Begyndelses- og Slutningshastighed; men de enkelte
Addender maae i ligestore Summer nødvendig være større, hvor der er et
mindre Antal; med andre Ord: Trykket ved det haardere Legeme er større
end ved det blødere.}

Naar Forskeren i Kraft af Iagttagelser paa denne Maade gaaer ud over
det, som lader sig iagttage, saa argumenterer han ikke af
Kjendsgjerninger, men af Begreber. Ved at lægge Begreber til Grund for
Analysen sættes han i Stand til at aflede Lovene for elastiske Legemer
af Lovene for uelastiske, og indsee Nødvendigheden af, at disse ved
Sammenstød altid maae tabe i levende Kraft, hine derimod ikke.

Men endnu staaer Besvarelsen af det Spørgsmaal tilbage: hvad berettiger
Iagttageren til saaledes at ombytte det Enkelte, han seer i
Kjendsgjerningen, med det Almene, han ikke seer? Han abstraherer, hedder
det, Begrebet af Kjendsgjerningen; vi maae da mærke os, hvad en saadan
Abstraction væsenlig forudsætter. Det Enkelte er ikke det Almene; men
dersom det Enkelte i Kjendsgjerningen ikke kunde fremstille det Almene,
saa var det umuligt at abstrahere det Almene af det Enkelte.
Gjenstanden er ikke et Begreb; men dersom der intet Begreb var udtalt i
den givne Gjenstand, vilde Forskeren bedrage sig selv, naar han meente
at kunne abstrahere et Begreb af Gjenstanden: at abstrahere et Begreb af
Noget, hvori der intet Begreb er, er en Modsigelse; at opdage og
forfølge Conseqvenser af de i Naturtingene udtalte Tanker, naar der i
Tingene ingen Tanker er, er en Modsigelse; at gjøre Tingen til Begreb og
Begrebet til Ting er en Modsigelse. Som Gjenstand for videnskabelig
Erfaring er Tingen paa een Gang Begrebets Virkning og dets Aarsag. Den
er Virkning; thi det guddommelige Magtbegreb er dens ideelle Aarsag.
Den er
% --- p. 95 ---
Aarsag, nemlig Erfaringsaarsag i Vexelvirkning med den iagttagende
Bevidsthed, som derved kommer til Begreb om Tingen og ved Tingens Begreb
til Begreb om Tingenes Vexelvirkning, altsaa til Begreb om Aarsagens
Realitet. Som Gjenstandens Begreb er Magten Aarsag, som Begrebets
Gjenstand er Aarsagen Magt.

\runin{b) Udstykning af Aarsager.} Den reelle Aarsag lader sig
udstykke; hele den ydre Natur er et storartet Exempel paa Udstykning af
Aarsager. Er nu enhver reel Aarsag et legemliggjort Begreb, saa skulde
det jo synes, som maatte de legemliggjorte Begreber have samme Skjæbne
som Legemerne; altsaa: naar legemlige Aarsager udstykkes, maae
Aarsagsbegreberne udstykkes; og naar legemlige Aarsager sammenstykkes,
maae Begreberne sammenstykkes. Men Begreber kunne hverken udstykkes
eller sammenstykkes; Begrebet er som Begreb en uopløselig Totalitet af
sammenhørende Bestemmelser: hvorledes forholder det sig da med den
mechaniske Udstykning?

Den alle Legemer iboende Tyngde virker som Aarsag; Tyngdebegrebet er et
i Universet gjennemgaaende Magtbegreb; og dog er her en Deling.
Kloderne i Rummet have i Virkeligheden hver sin Deel af den almindelige
Tyngde, Jorden sin, og Maanen sin. Men Deelbarheden i det Reelle er
netop begrundet paa Udeelbarheden i det Ideelle. Hvis Tyngdebegrebet og
dermed Tyngdens mægtige Væsen selv var deelt og ved Delingen gaaet i
Stykker: saa vilde det Stykke Tyngde, der boer i Jorden, jo befinde sig
i en Middelafstand paa 51032 danske Miil fra det Stykke, der boer i
Maanen; men hermed vilde Tyngdens Baand være brudt: Tyngden kunde ikke
være Tyngde. Aarsagens Hemmelighed er netop den, at det Reelle kan
deles, medens det Ideelle er og forbliver en udeelbar Magt. I reel
Forstand er Jorden Aarsag, og Maanen er Aarsag, og de to
% S. 95 slutter midt i Sætningen; den fortsætter paa S. 96 med
% „reelle Aarsager befinde sig virkelig i 51032 Miles Afstand ...“.
%  MATHEMATICAL — work from the image; the OCR is not a witness here.
% --- p. 96 ---
% S. 96 begynder UDEN Indrykning, midt i Sætningen fra S. 95: „...  de to
% reelle Aarsager befinde sig virkelig i 51032 Miles Afstand ...“. Der maa
% altsaa ikke aabnes et nyt Afsnit her. Overskriften b) Udstykning af
% Aarsager staaer IKKE paa S. 96; den hører til den foregaaende Batch.
reelle Aarsager befinde sig virkelig i 51032 Miles Afstand fra hinanden;
men i ideel Forstand er Tyngden Aarsag, og Tyngdens ideelle Baand kan
ikke briste. Imellem de to for Jord og Maane ideelle Aarsager er der
ingen Afstand: de ere, hver for sig, den hele Aarsag, og dog Momenter af
een Aarsag.

At Udeelbarheden af det Ideelle kan bestaae med en Udstykning af det
Reelle, har sin Grund i, at den ideelle Udeelbarhed selv bestaaer i en
Grunddeling, men en Grunddeling er ingen
Udstykning\footnote{% trykt Mærke: *)
Grundideernes Logik. II. Side 347--458.}. De samme Magtbegreber kunne i
Naturen individualiseres paa forskjellige Maader, ligesom de samme Love
jo kunne opfyldes paa utallige Maader. Det maa bestandig fastholdes, at
Naturlovene i alle Kjendsgjerninger lige saa vel som Kjendsgjerningerne
selv beroe paa Naturbegreberne. Vistnok opdager Physiken mangfoldige
Love, hvis Gyldighed den sikkrer sig ved almindelig Induction, uden endnu
at kunne begrunde dem paa tydelig Indsigt i de reelle Aarsagers Natur;
men Inductionen selv vilde mangle sand videnskabelig Grund, dersom
Forskeren ikke turde forudsætte, at det overalt er Begreber, om end
skjulte, endnu uopdagede Begreber, der afgive Lovenes Grundlag. Vistnok
kunne Lovene for Electricitetens
% SATSFEIL: der staaer „Virkniuger“ med omvendt n (u) --- eftersect paa
% KB-Exemplaret, hvor det tredie Bogstav i Stavelsen tydelig er et u med
% aaben Overkant, medens n'et i samme Ord er lukket. Begge Exemplarer
% vise det, altsaa Sætterens Feil. Gjengivet som trykt; læs „Virkninger“.
Virkniuger endnu ikke udledes af Electricitetens Begreb paa samme
indtrængende Maade, som Lovene for elastiske Legemers Sammenstød kunne
udledes af Elasticitetens Begreb; men derfor er der dog vel ingen
Physiker, der tvivler om, at Grundforholdet mellem Begreber og Love
bestandig maa være det samme.

Hvorledes Lovene forholde sig til Aarsagsbegreberne og disse til de
reelle Aarsager, kan allerede sees af de saakaldte mechaniske Potenser.
Om saadanne simple Maskiner som Vægtstangen, Skraaplanet o. s. v. kan man
jo vistnok
% --- p. 97 ---
sige, at enhver af dem er et Stykke Aarsag, et vilkaarlig dannet Redskab,
men at hele Aarsagsbegrebet lige saa vel er realiseret i Vægtstangen som
i Skraaplanet, at ethvert af disse Redskaber fremstiller et heelt Begreb
for sig, er indlysende. At disse indskrænkede Begreber i al deres
Eensidighed ikke ere Brudstykker, men Momenter af det store, Universet
omfattende mechaniske Aarsagsbegreb, fremgaaer jo deraf, at de Love for
Ligevægt og Bevægelse, som ved disse Redskaber individualiseres, ere
Naturlove, almindelige Verdenslove.

Loven som Lov er en Side af Begrebet og selv det hele Begreb, nemlig i
Almindelighedens Form, medens de legemlige Aarsager i deres indbyrdes
Forskjellighed fremstille Begrebet i det Særskilte, og hvert Legeme for
sig det antilogisk existerende Enkeltbegreb. Fra Realitetssiden er
Enkeltbegrebet imidlertid ikke at see som Begreb, men som Materie, som
Legeme. Iagttageren mærker ikke, at Legemerne i deres Væsen ere Begreber,
lige saa vel som Loven, og denne desaarsag sandselig i sin Opfyldelse,
lige saa vel som Legemerne. Han sandser Legemerne og abstraherer Loven;
den saaledes abstraherede Lov bliver ham derfor ikke et Begreb, men en
kategorisk Form for en uforklaret Nødvendighed. Standpunktets Eensidighed
kan imidlertid siges at være retfærdiggjort ved Erfaringslærens
eiendommelige Charakteer; det ligger ikke i Forskerens Interesse at
udvikle Aarsagsbegreberne som reale Naturbegreber: han holder sig til
Kjendsgjerninger og Love.

Men hvad har Immanensphilosophien udrettet? Hegel ivrer mod Physikerne og
beklager sig over deres Udstykninger; men Physikerne ere jo undskyldte,
naar Naturens „Væren, som den er, ikke svarer til sit Begreb“. Hvorledes
skulde Naturforskerne kunne holde paa „Ideen“, naar Naturen ikke engang
kan det, men maa i Hegelsk Forstand betragtes som „Ideens Affald fra sig
selv“. Under Be\-%
% --- p. 98 ---
stræbelsen for at bringe den til „Udvorteshed“ forfaldne Natur paa Fode
ved Hjælp af „Begrebet“ eller rettere: gjøre det begribeligt, hvorledes
„Ideen“ Trin for Trin arbeider sig op af sin Fornedrelse, har den
speculative Naturphilosophie i sørgelig Maade maattet overanstrenge den
logiske Nødvendighed.

Saasandt det er „den absolute Idee“, som nedenfra begynder sin absolut
nødvendige Vandring opad, maa Naturphilosophien ganske rigtig føre
Beviset for, at den trinvise Opadstigen skeer med immanent Nødvendighed,
og at denne nærværende Verdensorden er den eneste tænkelige, i hvilken
Ideen har kunnet komme til udvortes, virkelig Selvaabenbarelse. Dersom et
saadant Beviis nu virkelig var ført eller kunde føres, da vare vi ganske
vist pligtige at anerkjende Beviset og bøie os for Nødvendigheden; vi
maatte jo, som H. C. Ørsted siger, „skamme os i vort Inderste, naar vi
grebe os selv i at ville have en anden Sandhed end den virkelige“. Men
naar et saadant Beviis hverken er ført eller nogensinde kan føres, saa
krænke vi jo den virkelige Sandhed, naar vi bøie os for en Fixidee.

Det er sandt, at Alt hvad der skeer maa have sin Aarsag og forsaavidt sin
Nødvendighed; men det er usandt, at nogensomhelst relativ Virkelighed
skulde være absolut nødvendig. I abstract Almindelighed er
Nødvendigheden en tom Tanke; for at give Tanken et Indhold, maae vi give
den en Form og optage den i et bestemt Begreb. Først ved Begrebet kan der
komme noget Nødvendigt ud af Nødvendigheden. Nødvendigt er det, hvis
Modsatte er umuligt. Begrebsnødvendigheden betegner overalt Grændsen for
det Mulige og det Umulige, men naaer aldrig Grændsen for det Virkelige.

Nødvendighedsgrændsen for et virkeligt Skraaplan f. Ex. er en Grændse for
dets Mulighed. Alle Bestemmelser af logisk Nødvendighed vedrøre kun
Virkeligheden forsaavidt
% --- p. 99 ---
de ligge i Muligheden. Det er nødvendigt, at ethvert Skraaplan maa have
en Grundlinie, en Høide og en Længde; hvorfor? Fordi et Skraaplan uden
Høide, Længde og Grundlinie er en Umulighed. Men indenfor
% SATSFEIL: „Mulighedeu“ med omvendt n (u) i Ordets sidste Bogstav.
% Eftersect paa KB-Exemplaret ved Siden af „kan der“ i samme Linie, hvor
% n'et er lukket foroven; her staaer et aabent u. Begge Exemplarer vise
% det. Gjengivet som trykt; læs „Muligheden“.
Mulighedeu kan der gives en Uendelighed af virkelige, i deres
Virkelighed, forskjellige Skraaplaner; mod de virkelige Forskjelligheder
er Nødvendigheden ligegyldig: den er i Virkeligheden ligesaa tilfreds med
det ene Exemplar som med det andet. Det Samme gjælder om Lovene.
„Faldbestræbelsen paa Skraaplanet maa forholde sig til Trykket mod
samme, som Skraaplanets Høide til dets Grundlinie.“ Det er en Lov, thi
det er nødvendigt; dets Modsatte er umuligt. Men denne Nødvendighed
afgjør ikke det Mindste angaaende den virkelige Høide, den virkelige
Grundlinie, det virkelige Tryk. Begrebsnødvendigheden er haard og streng
mod det Umulige, men smidig, føielig, yderst liberal mod det Virkelige.
Speculationen kan udvide sin Synskreds saa meget, den vil, dens logiske
Nødvendighed naaer aldrig ud over Muligheden.

Maaskee kunde det ved første Øiekast forekomme En og Anden, som havde
Læren om denne Verdens absolute Nødvendighed et paalideligt Grundlag i
den rationelle Mechanik; men ikke at tale om, at Mechaniken i sin
Fremstilling af Lovene for Himmellegemernes Bevægelse lægger
Newtonsloven til Grund uden apriorisk Beviis for denne Lovs udelukkende,
absolute Nødvendighed, saa gaaer det jo i det store Hele med
Nødvendigheden for Verdenskloderne, som i det Particulære med
Nødvendigheden for Skraaplanet. Forestille vi os en Planet med Massen m,
bevægende sig om et Centrallegeme med Massen M, og decomponere de
accelererende Kræfter efter tre paa hverandre lodrette Axer, saa at
% Der staaer i Satsen en kort Bindestreg uden Mellemrum: „og for m-x, y og
% z:“. Meningen kræver et Skilletegn (Kolon eller Tankestreg); begge
% Exemplarer vise den samme korte Streg, saa den er gjengivet som trykt.
Coordinaterne for M blive X, Y og Z og for m-x, y og z: da vil Afstanden
r mellem M og m være
% ÷ er Nielsens Minustegn (jvnfr. S. 46 og Husreglen i Hovedet); det staaer
% saaledes i begge Exemplarer under Rodtegnet, hvis Vinkel dækker hele
% Udtrykket.
$r = \sqrt{(x \div X)^{2} + (y \div Y)^{2} + (z \div Z)^{2}}$; det er
nødvendigt, thi dets Modsatte er umuligt. Reflectere vi alene paa
% --- p. 100 ---
X-Axen og kalde den Vinkel, som r danner med denne Axe, α, saa at
$\cos\alpha = \frac{x \div X}{r}$: da vil den Tiltrækning,
Centrallegemet udøver paa Planeten, altsaa dennes Acceleration, være
% Efter Lighedstegnet staaer der et selvstændigt ÷ foran Brøken: det er
% Nielsens Minustegn, ikke et Divisionstegn. Anden Ligning har intet
% Fortegn. Begge Exemplarer stemme.
$\frac{d^{2}x}{dt^{2}} = \div\frac{M(x \div X)}{r^{3}}$, og den
Tiltrækning, Planeten udøver paa Centrallegemet, altsaa dettes
Acceleration, være $\frac{d^{2}X}{dt^{2}} = \frac{m(x \div X)}{r^{3}}$;
det er nødvendigt, thi dets Modsatte er umuligt. Lade vi Centrallegemet
være ubevægeligt, saa maa dets Acceleration lægges negativt til
Planetens, og Ligningen for dennes Bevægelse bliver da
$\frac{d^{2}(x \div X)}{dt^{2}} + \frac{(M + m)(x \div X)}{r^{3}} = 0$;
det er nødvendigt, thi dets Modsatte er umuligt. At der under slige
Analyser kan foresvæve Mathematikeren Billeder af virkelige Kloder, et
virkeligt Centrallegeme, virkelige Planeter og virkelige Afstande,
betyder jo Intet med Hensyn paa Virkeligheden selv og den virkelige
Udstykning. Hvor virkelige Bevægelser finde Sted, er det ganske vist
nødvendigt, at de almindelige Love maae opfyldes; men fra denne logiske
Nødvendighed til et antilogisk bestemt Planetsystems absolute
Nødvendighed er der et uendeligt Spring.

Den absolute Idee, siger man, har alle relative Momenter, alle væsenlige
Betingelser for sin Virkeliggjørelse i sig selv. Den hele udad kastede
Endelighed er altsaa en Udstykning, hvori det Absolute absolut maa være
allestedsnærværende. Jo mere man udvider Heelhedsbegrebet, jo flere
Betingelser man efterhaanden optager i Analysernes almindelige
Forudsætning, desto mere vil man ad den indre Nødvendigheds krydsende
Veie være i Stand til at nærme sig det udvortes Virkelige. Vi have i
denne Synsmaade en af de mange immanente Tvetydigheder, hvormed den
„moderne“ Philosophie er bleven saa overlæsset, at kun en Henviisning
til exacte Formbestemmelser kan antages at ville udøve en
% --- p. 101 ---
% NOTATION: Rækken m, m', m'' bærer Mærker, og efter Husreglen (batch 5)
% sættes hele Opregningen i Mathematik, for at Skriftsnittet ikke skifter
% midt i Rækken. De enkeltstaaende Bogstaver R, x, X, r i det følgende
% staae derimod som trykt, i almindelig Sats.
rensende Virkning. Lad os da til de angivne Forudsætninger endnu føie
den, at der foruden Planeten $m$ bevæge sig andre Planeter $m'$, $m''$
o. s. v. om Centrallegemet. Hvad fører nu dette til? Ifølge den sidste
Forudsætning kan Summen
$\frac{d^{2}(x \div X)}{dt^{2}} + \frac{(M + m)(x \div X)}{r^{3}}$
eller, naar $x \div X$, m's Abscisse med Hensyn til M's Midtpunkt,
betegnes ved $\xi$, Summen
$\frac{d^{2}\xi}{dt^{2}} + \frac{(M + m)\xi}{r^{3}}$ rigtignok ikke være
Nul; her fremkommer en perturberende Function R, og Ligningen bliver
$\frac{d^{2}\xi}{dt^{2}} + \frac{(M + m)\xi}{r^{3}} =
\left(\frac{dR}{d\xi}\right)$. En saadan Ligning giver ganske vist
Nødvendighedsudtrykket en mere articuleret Form; men Grændsen er kun i
Muligheden en Grændse for det Umulige. Om vi endog, veiledede af Methoden
for „de arbitrære Constanters Variation“, gjennemvandrede den hele Urskov
af Ligninger, Rækker og Formler, som ligger bagved ovenstaaende
Antydninger: vi naaede dog ikke de antilogisk bestemte Kjendsgjerninger,
fik ikke Udstykningens materielle Virkelighed presset ind i den absolute
Nødvendighed.

% Overskriften er trykt HALVFED, ikke spærret som de øvrige a)/b)/c)-Hoveder
% (samme Forhold som ved c) Aarsagsmomenternes Eenhed, S. 79); „c)“ selv
% staaer i almindelig Antikva. Udgaven gjengiver ikke Forskjellen, saa
% \runin bruges ogsaa her.
\runin{c) Aarsagsrækker.} Udstykningen af de mechaniske Aarsager maa
nødvendig have en egen Tingenes Sammenhæng til Grundlag. Fra
Almindelighedens Side er Sammenhængen realiseret i Tyngden, der varierer
under Særskilthedens Form med de legemlige Tilstande. Physiken seer i den
Modstand, Legemerne gjøre mod Delenes Adskillelse, en fra Tyngdekraften
forskjellig Sammenhængskraft, og vistnok med
% NOTEN LØBER OVER TO BLADE: den begynder under Stregen paa S. 101 og
% fylder hele Notefeltet paa S. 102; Mærket gjentages ikke paa
% Fortsættelsen. Hele Noten er her sat ved sit Mærke paa S. 101.
% NOTATION: Bogstavrækken T, r, μ, ϱ, ε blandes af latinske og græske
% Tegn; efter Husreglen om Opregninger er hele Rækken sat i Mathematik,
% for at Skriftsnittet ikke skifter midt i Opregningen.
Rette.\footnote{% trykt Mærke: *)
Cohæsionstiltrækningen gjælder nemlig kun for Smaadelene indbyrdes;
skulde denne Tiltrækning være en Yttring af Tyngdekraften, maatte den paa
Grund af Delenes Lidenhed og store indbyrdes Nærhed blive overordentlig
ringe. Ifølge Tyngdeloven $T = \frac{\mu}{r^{2}}$ --- idet $T$ er
Tiltrækningen, $r$ Afstanden og $\mu$ Intensiteten --- vil, for en
Kuglemasse med Radius $r$ og Tætheden $\varrho$, Intensiteten være
$\mu = \tfrac{4}{3}\pi\varrho r^{3}$. Indsvinder nu Kuglen til et Atom,
vil Intensiteten for Atomet være
$\varepsilon = \tfrac{4}{3}\pi\varrho(dr)^{3}$ og følgelig
$\frac{\varepsilon}{(dr)^{2}} = \tfrac{4}{3}\pi\varrho\, dr$, en
forsvindende Størrelse. Den Indvending, at Smaadelene ikke ere uendelig
smaa, har i denne Sammenhæng Intet at betyde; thi Formlen viser, at
Sammenhængskraften maa blive desto svagere, jo inderligere
Sammenføiningen er, Noget, som baade er i Strid med Erfaringen og i sig
selv en Modsigelse.}

% --- p. 102 ---
% Selve Brødteksten paa S. 102 begynder her; hele Notefeltet paa Bladet er
% optaget af Fortsættelsen af Noten fra S. 101, som staaer ovenfor.
Men ihvorvel Sammenhængskraften ikke tør betragtes som en blot Aflægger
af Tyngdekraften, vil det dog kunne forstaaes, at de to uligeartede
Kræfter maae henføres, vel ikke til een Grundkraft, men dog til eet
Grundbegreb, det omfattende Magtbegreb, hvis indre Articulation er alle
Kræfters ideelle Aarsag. Ligesom Tyngdekraften, hvad der paa sit Sted er
viist, beroer paa Reflexion i Magt, saaledes ogsaa Sammenhængskraften. De
særlige Begreber i denne Reflexion ere \emph{Hele} og \emph{Dele}. At
Legemet er et i Udstykning reelt Hele, er eenstydigt med, at det bestaaer
af Dele. Disse Dele ere i mechanisk Forstand uselvstændig selvstændige
Led. Deres Uselvstændighed er et Udtryk for Realbegrebets Magt i det
givne Hele; men deres Selvstændighed viser det Heles Sammenstykning, thi
Realbegrebet er kun Magt i det Hele ved særlig at være Magt i Delene: det
i Existensen reflecterede Aarsagsudtryk for de selvstændige Deles
Uselvstændighed er Sammenhængskraften. Af de paa særegne Forhold mellem
„Udvidekraft“ og „Sammenhængskraft“ beroende Tilstandsformer sees det,
hvorledes Delenes Egenskaber reflectere sig i det Hele, og heraf igjen,
hvorledes det paa Udstykningens Grundlag forholder sig med
Legemsbegreberne og de existerende Legemer.

Det er en Grundsætning i Mechaniken, at Virken og Modvirken, Action og
Reaction, nøiagtig maae svare til
% --- p. 103 ---
% NOTATION: Rækken k, M, v, M', v' bærer Mærker, saa hele Opregningen er
% efter Husreglen sat i Mathematik. ÷ foran kt er atter Nielsens Minustegn;
% de øvrige Fortegn i Afsnittet ere ægte Plus og Lighedstegn.
hinanden. Grundsætningen gjælder for Tyngden; thi den Tiltrækning, en
Masse udøver paa et Punkt, er altid lige stor med den, som Punktet udøver
paa Massen. Et faldende Sandskorn tiltrækker hele Jorden med en Kraft lig
den, hvormed hele Jorden tiltrækker Sandskornet --- at Accelerationen er
høist forskjellig er noget ganske Andet. --- Grundsætningen gjælder for
Bevægelsesmængder. Virker en constant Kraft $k$ en vis Tid paa Massen
$M$, saa er Accelerationen $\frac{k}{M}$, og Bevægelsesmængden i en vis
bestemt Retning $Mv = kt$. Men den Kraft, som paavirker Massen $M$, maa
være udgaaet fra en anden Masse, f. Ex. $M'$, og den af $k$ paavirkede
Masse $M$ virker da tilbage paa $M'$ med en Kraft, der ogsaa er $k$; kun
at Virksomheden gaaer i modsat Retning, og Bevægelsesmængden bliver
$M'v' = \div kt$. Heraf følger da $Mv + M'v' = 0$, en Ligning, som viser,
at Bevægelsesmængdernes Sum for begge Legemer vedbliver at være
uforandret. --- Grundsætningen gjælder for den levende Kraft. Vel kan det
ikke siges, at det Arbeide, som afgives af et Legeme, altid modtages af
et andet, thi Virksomheden kan antage forskjellige Former; Varme f. Ex.
kan afgive Arbeide og Arbeide igjen opspares. Men betragtes Virksomheden
under alle Former som samme Kraftfunction, og maales den med samme Maal,
da vil ogsaa denne
% SATSFEIL: „Functiou“ med omvendt n (u). Eftersect paa KB-Exemplaret mod
% „Kraftfunction“ to Linier ovenfor, hvor n'et er lukket foroven. Begge
% Exemplarer vise det. Gjengivet som trykt; læs „Function“.
Functiou ved Overgang fra Legeme til Legeme forblive uforandret: Virken
og Modvirken svare til hinanden.

Grundsætningen lærer os, at Kræfterne altid fremtræde to og to parallele,
ligestore og modsatte: Tryk svarer til Modtryk, Stød til Modstød, A's
Tiltrækning af B til B's Tiltrækning af A. Den mechaniske Aarsag er, med
andre Ord, aldrig nogen umiddelbar Enkeltaarsag, men altid en
Vexelvirkning af to Aarsager. At enhver udvortes Aarsag, væsenlig seet,
er Vexelvirkning af to udvortes Aarsager, udbreder et eget Lys over alle
mechaniske Aarsagsforhold.

% --- p. 104 ---
For at begynde med Tyngden, den mest omfattende af alle udvortes
Aarsager, saa er den ikke blot udvortes i Forholdet mellem de Masser,
hvorpaa den virker, men den er i sig selv udvortes. Massen $m$ tiltrækker
$m'$, forsaavidt $m'$ tiltrækker $m$, at Tiltrækningen fra begge Sider er
lige stor, beviser, at Aarsagen i begge Masser er lige oprindelig og lige
afledet; begge Aarsager ere i deres Virksomhed \emph{foraarsagede}: de
foraarsage hinanden gjensidig. Hvad der gjælder om to Led, gjælder om
alle; her danner sig en uendelig Række af foraarsagede Aarsager; intet
Led betegner den første Aarsag. Den Modsigelse, at Tyngden overalt er
lige oprindelig og lige afledet, løses i Naturen saaledes, at Tyngden,
som i det Uendelige forudsætter sig selv, kommer udenfor sig selv og
concentrerer sine Virkninger i andre, fra den forskjellige Aarsager:
\emph{Tryk}, \emph{Stød}, \emph{Bevægelsesmængde}.

Hvor Afstanden mellem de tunge Masser er ophævet, virker Tyngden som
\emph{Tryk}. Ifølge Grundsætningen for Ligestorheden af Tryk og Modtryk
er et Legemes Vægt at bestemme ved Tyngden; og dog er Trykket som Aarsag
forskjellig fra Tyngden. Ifølge Tyngden stræbe Massedelene at forene sig
i eet Punkt, Tyngdepunktet, men Trykket modvirker denne Bestræbelse og
er, skjøndt selv en Virkning, dog Aarsag, Aarsag nemlig til at Delene
holde sig udenfor hverandre, saa at Masserne udfylde Rummet. Kun det
Trykkede trykker; Trykket som Aarsag er ikke blot udvortes imod sin
Gjenstand; men det er i sig selv en udvortes i det Uendelige foraarsaget
Aarsag. Aarsagsrækker, dannede af Tryk og Modtryk, kunne umulig have
noget første Led, hvori Trykket som Tryk var oprindelig Aarsag. Den
Modsigelse, at ethvert særskilt Tryk i dets Umiddelbarhed er en afledet,
foraarsaget Aarsag, løser sig op i en Omdannelse af Trykket. Ved denne
Omdannelse bliver det Moment i en
% --- p. 105 ---
anden Aarsag, \emph{Stødet}. Stød forudsætter Bevægelse, og Bevægelsen
maa have sin Aarsag.

Hvor en Afstand mellem Legemerne er givet, kan Tyngden vel ophæve
Afstanden og blive Aarsag til Stødet; men Tyngden, der ikke oprindelig er
Aarsag til Afstanden, kan ikke være Stødets oprindelige Aarsag. Ligesom
Trykket foraarsages ved en Sammenstøden, skulde Legemernes Afstande
altsaa være foraarsagede ved en Frastøden; Stødet vilde da fra dette
Synspunkt være at opfatte som den concrete og forsaavidt oprindelige
Aarsag, der havde Tyngden til sit almene og Trykket til sit særskilte
Moment. Den ældre mechaniske Naturphilosophie (Cartesius) har ikke ladet
det mangle paa Henviisninger til et første, oprindeligt Stød, et
Grundstød, hvoraf alle Bevægelser og Verdensphænomener skulde være at
forklare. Men at Forestillingen om et første oprindeligt Stød maa løse
sig op i et begrebløst Begreb, er klart af Grundsætningen: kun fra Stød
kommer Stød. Men naar kun det, som forud er stødt, kan støde, saa er jo
ethvert Stød en af et Stød foraarsaget Aarsag. Den uendelige
Aarsagsrække, hvori Stød uafbrudt forudsætter Stød, kan da umulig begynde
med et første Stød. Begyndes der med et første Stød, d. e. med et Stød,
hvis Aarsag ikke er et Stød, saa falder Rækkens Begyndelse udenfor
Rækken: der begyndes da med en Aarsag, der vel indeholder Stødets
Mulighed, men dog er forskjellig fra Stødet; der begyndes med en
\emph{Bevægelsesmængde}.

Saa skulde Bevægelsesmængden altsaa være den indholdsrige Aarsag, som ved
at optage Tyngden, Trykket og Stødet i Kredsen af sine reale Momenter,
kunde trods Udstykningen være Bevægelsernes Princip og fremtræde som
Grundaarsag; men ogsaa dette er umuligt. Bevægelsesmængder kunne, siger
Physiken, endog uden Hensyn til Kraftens Uforanderlighed forplantes fra
Legeme til Legeme uden at forandre Størrelse eller Retning, saaledes at
hvad det ene
% S. 105 slutter midt i Sætningen: „... saaledes at hvad det ene“ ---
% Fortsættelsen staaer paa S. 106 og hører til næste Batch.
%  --- C. Grundaarsagen: det første Bevægende  (pp. 106–124) ---
%  MATHEMATICAL — work from the image; the OCR is not a witness here.
% --- p. 106 ---
% S. 106 aabner IKKE med Afdeling C. Afdeling B (Mechanisk Realitet) løber
% videre hertil med TO Afsnit: det første fortsætter Sætningen fra S. 105
% („... hvad det ene / Legeme modtager ...“), det andet sluttes med
% „mechanisk Realitet.“ Derefter staaer den korte, centrerede Streg, og
% først da Overskriften C. Samme Form som ved Overgangen A -> B paa S. 87.
Legeme modtager, det taber det andet. For en Verden af Legemer ville
altsaa Bevægelsesmængderne i enhver Retning være uforanderlige
Størrelser, hvilke Kræfter der end virke mellem Legemerne eller de
legemlige Dele indbyrdes; men ogsaa i Form af Bevægelsesmængde er
Aarsagen sig selv udvortes.
% Tegnet foran M'v' er Nielsens MINUS (÷), læst paa KB-Exemplaret;
% jvnfr. Husreglen fra Batch 4 (S. 46).
Aarsagen til et $Mv$ er i et $\div M'v'$;
den sig meddelende Bevægelsesmængde maa overalt være meddeelt:
Aarsagernes uendelige Række er en Række af lutter meddeelte,
\emph{foraarsagede}, Aarsager. Causalitetsforholdet forandrer da heller
ikke sin Natur derved, at Bevægelsesmængden viser ud over sig selv,
optages i en endnu concretere Aarsag og omsættes i \emph{levende
Kraft}.

% „levende Kraft“ er spærret ved Afsnittets Slutning ovenfor, men IKKE i
% dette Afsnit, hvor Ordene staae tre Gange i almindelig Antiqva
% (kontrolleret paa KB-Exemplaret).
Et Legemes Forøgelse i levende Kraft er, hedder det, lig med det
Dobbelte af det modtagne Arbeide, og et Legemes Tab i levende Kraft lig
med det Dobbelte af det afgivne Arbeide; Ligningen
$M(v^{2} \div c^{2}) = \pm\, 2\,MGs$ viser begge Dele: Omsætningen af
Arbeide i levende Kraft og af levende Kraft i Arbeide kan fortsættes i
det Uendelige. Under Medvirkning af de forskjellige faste, flydende,
luftformige Tilstande, i hvilke Legemerne befinde sig, kunne de
udvortes Aarsager paa utallige Maader omdannes og indgaae i mangfoldig
forviklede Vexelvirkninger; men at disse udvortes Aarsager overalt, i
det store Hele som i de mindste Dele, nødvendig forudsætte en
Grundaarsag og dog holde Grundaarsagen skjult, er og bliver
Charakteermærket for al mechanisk Realitet.

% Kort, centreret Streg mellem Afdeling B og Afdeling C, midt paa S. 106.
% Efter Husreglen gjengives Ornamentets Længde ikke: \streg tjener overalt.
\streg

%  --- C. Grundaarsagen: det første Bevægende  (pp. 106–125) ---
% Overskriften er læst paa Siden selv (KB-Exemplaret), ikke i Indholdet:
% „C.“ paa egen Linie, derunder „Grundaarsagen: det første Bevægende.“
% med Kolon, sat med fed, spærret Skrift.
\afdeling{C.}{Grundaarsagen: det første Bevægende.}

Ifølge den mechaniske Idealitet ere de virkende Aarsager reflecterede i
tilsvarende Aarsagsbegreber, og disse
% --- p. 107 ---
Causalitetens Magthavere vise sig da at være de ideelle, indenfra
bestemmende Aarsager. Ifølge den mechaniske Realitet ere
Aarsagsbegreberne derimod formedelst Gjenstandsmagten skjulte under den
udvortes materielle Existens, hvori de have deres Virkelighed;
Grundaarsagen er forsvundet, og Betragteren af de uendelige Rækker seer
sig forgjæves omkring efter et første Bevægende. Skal det være muligt i
Bevægelsernes Rækker at opdage Grundaarsagen, da maa Forskjellen mellem
de to uligeartede Bevægelser, den \emph{ideelle} i Begrebet og den
\emph{reelle} i Kjendsgjerningen, udtrykkelig paavises. Det er under
den ideelle Bevægelse, at Aanden, den Alt foraarsagende Selvmagt,
bestemmer de særlige Magtbegreber til Gjenstandsmagtens materielle
Aarsager. Men den ideelle Bevægelse forudsætter netop, hvad den selv
bestemmer, nemlig den reelle. Ikke umiddelbart, men igjennem virkelige
Naturaarsager og reelle Bevægelser er Selvmagten \emph{den hele
mechaniske Udstyknings oprindelige Aarsag}. Udstykningen beroer lige
saa lidt paa en absolut Nødvendighed som paa en blind Tilfældighed:
Grundaarsagen selv, den almægtige Frihed, er \emph{det første
Bevægende}.

% Indløbende Overskrift, læst paa Siden selv (KB-Exemplaret): sat med fed
% Antiqva, ikke spærret; „er“ foran den er almindelig Skrift.
\runin{a) Ideel og reel Bevægelse.} Immanenslæren har gjort sit
Yderste, for at faae den ideelle og den reelle Bevægelse til at tage
sig ud som Eet og det Samme. I Hegels Logik skal Begrebsbevægelsen
gjælde for Sagbevægelse; i hans Naturphilosophie skal til Gjengjæld
Sagbevægelsen gjælde for Begrebsbevægelse. Man behøver blot at besinde
sig paa den Grundmodsigelse, som heri ligger, for at komme til Vished
om, hvor huult det Absolute er, paa hvilket „Immanensen“ fremdeles maa
bygge sine „pantheistiske“, „naturalistiske“, „atheistiske“ Luftslotte.

Skal det for Alvor gjælde, at den logiske Idee har en evig, i sig selv
hvilende, af Naturen uafhængig Væren, da maa der jo nødvendig
forudsættes en over Rum, Tid og
% --- p. 108 ---
Materie ophøiet Tænken, Reflecteren og Begriben; thi uden Tænken,
Reflecteren og Begriben gives der ikke en eneste logisk Bestemmelse,
end sige en absolut logisk Idee. Ved at erklære sig for den logiske
Idees Oprindelighed og fastholde dens overnaturlige Væren gaaer det
Absolutes Philosophie ud over sin egen anpriste Immanens og gjør et
betænkeligt Laan hos Transcendensen. Fra dette Synspunkt er den ideelle
Bevægelse i Begrebet da ogsaa væsenlig adskilt fra Naturudviklingens
reelle Bevægelse, og dog skulde de to Bevægelser være Eet.

For at undvige det Modsigende i den anførte Conseqvens, griber man til
en anden Forklaring. At det Absolute i sit Væsen er logisk, og Begrebet
dets oprindelige Form, maa da ikke forstaaes saaledes, at Logikens Idee
virkelig skulde have gjennemtænkt sine kategoriske
Indholdsbestemmelser, før den forandrede sig, slog over i det Reelle og
begyndte forfra som Naturidee. Det Absolutes evige Væren er tvertimod
en oprindelig Naturværen; dets immanente Væsensbestemmelser forvandle
sig da som ved et Trylleslag; de ere ikke paa absolut Maade i det
Absolute; deres Tilværelse bliver at opfatte som Naturens,
Udvorteshedens sandselige, materielle, for Ideen uadæqvate Former:
Bevægelsen er reel. For at det Absolute midt i Adsplittelsen ikke desto
mindre kan vedblive at være absolut, maa det Tankeløse selv være
tænkende. Det er det Almene, det Alttænkende, det store tænkende
Neutrum, hvorpaa Immanensen hviler. Dette Almene er rigtignok en i
Tingene fortabt, naturbunden Form, der først i Menneskeaanden kan
frigjøres til bevidst Tænkning; men den ufrie Naturtilstand gjør intet
Skaar i det Almenes uendelige Tænken; Bevidsthed er ganske overflødig;
den bevidste Tænkning er kun en lavere Art af Tænkning. Den absolute
Tænkning er en Tænkning uden Bevidsthed, ligesom den absolute Viden er
„en ubevidst Viden“, en Klarhed uden Lys, en Begriben uden Indsigt,
% --- p. 109 ---
en Forstand uden Forstaaelse. Paa disse Vilkaar er det slet ikke
urimeligt, at Tiden kan begribe Rummet, og Rummet Tiden, og at den ved
Rummets og Tidens gjensidige Begriben frembragte Materie igjen kan
begribe og bevæge sig selv: den reelle Bevægelse er ideel, thi den
ideelle er reel.

Det Sande i alt dette er, at Naturbegreberne ere Naturtingene iboende.
Det er Hegels udødelige Fortjeneste at have udviklet Spinozas golde
Immanenslære til en saa frugtbar Alsidighed, at „Begrebet“ sees at
rumme et Verdensindhold. Hegels Naturphilosophie fortjener i denne
Henseende lige saa fuldt Paaskjønnelse som hans Logik. Men
Grundvildfarelsen er, at Immanensen, hvor meget der end tales om det
„overgribende“ --- altsaa transcendente --- Begreb, overalt opløser og
fornegter den oprindelige Forudsætning, hvorpaa den selv beroer:
Immanens er umulig uden Transcendens. Naturbegreberne have et dobbelt
Medium; de ere Tingene iboende: den virkelige Natur er deres reelle
Medium; de fremgaae af uendelig Begriben: den Alt vidende, selvmægtige
Aand er deres ideelle Medium. Eenheden af de to Medier er en
Modsætningseenhed, der begrunder den væsenlige Forskjel mellem ideel og
reel Bevægelse.

% Det græske Bogstav er læst paa KB-Exemplaret: α, ikke „a“ eller „&“,
% som de to OCR-Vidner give. Overskriften er indløbende og spærret.
\subrunin{α) Ideel Bevægelse.} For at den mechaniske Aarsag kan være
Begreb, maa den begribes. Hvorledes Begribningen foregaaer i
menneskelig Viden, er ovenfor viist. Tænken og Forestilling ere
Functioner af hinanden gjensidig. Aarsagen kan ikke tænkes uden et
Noget, der forestiller det Virksomme, nemlig Kraften, og det Virksomme
kan ikke tænkes uden et Noget, hvorpaa der virkes, nemlig Materien,
hvis Inertie netop bestemmes derved, at den modtager og bevarer
Paavirkningen. Kraft og Inertie reflectere hinanden gjensidig: det
hertil svarende Phænomen i Tid og Rum er Accelerationen, og den i
Begrebet forestillede Aarsag er den
% --- p. 110 ---
accelererende Kraft. I Accelerationen er Massen, den begrændsede
Materie, underforstaaet, men endnu ikke optaget i Begrebet selv.
Aarsagsbegrebet fordrer, at Acceleration og Masse udtrykkelig indgaae i
reflecteret Eenhed af Bevægekraften, den i Tid begrebne Aarsags
bestemtere Udtryk, hvis Maal er Bevægelsesmængde. Men Bevægelsesmængde
i exact Form er et eensidigt Maal. Causalitetsforholdet fordrer, at
Virkningen modsættes Aarsagen, nemlig saaledes, at den i sin Virkning
reflecterede Aarsag nøiagtig svarer til det, som bevirkes. Den
fuldstændig bestemte Virkning er Arbeide; Aarsagen sættes i Arbeide og
antager Form af levende Kraft. Hvormange Bestemmelser, Forestillingen
under denne Begrebsudvikling end har tilført Tænkningen, saa er det dog
indlysende, at ingen af disse Forestillingsbestemmelser enten
vilkaarlig, eller tilfældig drages ind under Aarsagen udenfra, men at
de alle udvikle sig indenfra af den mechaniske Aarsags eget Begreb.
Begyndes der i det Abstracte med Accelerationen, saa kan Bevægelsen
ikke standse før det Concrete er naaet i levende Kraft; begyndes der i
det Concrete med levende Kraft, saa vil Analysen deri opdage Rummet,
Tiden, Bevægelsesmængden og i Bevægelsesmængden Accelerationen: det er
ideel Bevægelse.

Og dog er Begrebsudviklingen endnu ikke fuldendt. Den mechaniske Aarsag
er af Charakteer en Andethedsaarsag; dens Virksomhed er i det Udvortes;
Udvorteshedens Begreb hører oprindelig med i dens Væsen. Den begribende
Reflexion bestemmes ved Reflexer af udvortes Ting, og det Foraarsagende
spreder sig i en Kreds af Betingelser. Men Betingelsernes
Selvstændighed er kun tilsyneladende; de hverandre ophævende
Betingelser forsvinde som Momenter i Aarsagens Eenhed, det Ubetingede.
Ved at optage Betingelserne i Ubetingethedens Form er Aarsagen i sig
selv betinget, og det saaledes, at den indenfra af sit Begrebsindhold
kan udvikle sin egen Forskjellighed: een Aarsag kan
% --- p. 111 ---
paa Betingelsernes Grundlag blive til flere, og flere igjen til een.
Det er ingen reel, men en ideel Bevægelse, der sees i disse Overgange.
Aarsagsbegrebet er endnu kun et Begreb; men Begribningen viser, at
Momenternes Idealitet Punkt for Punkt reflecterer sig i tilsvarende
Realitet.

% spacing.py meldte her „Aand“ og „vanvittigt“ som mulig Spærring; begge
% ere forkastede efter Eftersyn paa KB-Exemplaret --- almindelig Antiqva,
% løs Udslutning i Arbeidsscanningen.
Som Naturbegreb har Aarsagen sit ideelle Medium i den Alt begribende
Aand; til den uendelige Begriben maa der da ogsaa svare en ideel
Bevægelse. Saa urimeligt, ja vanvittigt det end vilde være at forestille
sig den absolute Aand i Lighed med den menneskelige --- iagttagende,
overveiende, abstraherende, beregnende, differentierende, integrerende,
ligesaa urimeligt, ja vanvittigt er det at ville tænke sig en Aarsag
uden Begreb, et Begreb uden Begrebsudvikling, en Begrebsudvikling uden
indre, fri, ideel Bevægelse. Immanensmethaphysiken med det selvløse,
% „Grund“: begge OCR-Vidner læse „Gruud“ (den kjendte n/u-Forvexling).
% KB-Exemplaret afgjør det ikke positivt ved 300 ppi, og efter
% Playbook §2 (Sagens Læsemaade, naar Trykbilledet ikke positivt viser
% andet) staaer her „Grund“. Den forkastede Læsemaade er „Gruud“.
tænkende Neutrum burde dog idetmindste kunne lære Saameget af den
Uendelighedsanalyse, der ligger til Grund for den rationelle Mechanik,
at den omsider kom til Bevidsthed om, hvorledes det i Begribningen
selv, den være menneskelig eller guddommelig, absolut maa forholde sig
med Articulationen af det Uendelige ved det Endelige.

% Det græske Bogstav er læst paa KB-Exemplaret: β, ikke „B“, „ß“, „6“
% eller „ØB“, som OCR-Vidnerne give.
\subrunin{β) Reel Bevægelse.} Den reelle Bevægelse staaer i skarp
Modsætning til den ideelle, fordi den existerende Aarsag staaer i skarp
Modsætning til Aarsagsbegrebet. Et Legemes Bevægelse fra Sted til Sted
har, udvortes seet, ingen Lighed med en Begrebsbevægelse. Det er et
Magtsprog, der dicterer Modsætningen; men dette Magtsprog er
Realbegrebernes medfødte Natursprog. Begrebernes Overgang fra det
ideelle i det reelle Medium er ingen endelig naturhistorisk Tildragelse,
det er en Grundbevægelse, en uendelig Begivenhed, der ligelig tilhører
Tid og Evighed. Overgangen skeer ikke saaledes, at Realbegreberne,
uberørte af alle Realiteter, opbevares i Selvmagtens Evighed, for at
Gjenstandsmagten siden kan træde til i Tiden og fylde de tomme Former
med
% --- p. 112 ---
stofagtigt Indhold. Selvmagt og Gjenstandsmagt ere Extremer af samme
Almagt i evig-timelig Fordoblingseenhed. Det er ved Analyser
tydeliggjort, at det samme Differential, $dm$, der er ideelt i Legemets
Naturbegreb, er reelt i det existerende Legeme, at, med andre Ord,
Gjenstandsmagten, der sætter Existensen, er uendelig gjennemtrængt og
inderlig bevæget af Existensens Begreb. Men netop ved at udtrykke sit
Begreb kommer Gjenstanden til at skjule Begrebet; Kjendsgjerningen
tilhyller Idealiteten, og det intellectuelle Skin er saa svagt, at
Iagttageren neppe kan skjønne, om det er Tingens eller hans eget
Begreb, han har med at gjøre. Denne fra Momentet, $dm$, stammende
Modsætning af det Ideelle og det Reelle er Objectivitetens Modsætning i
Magt; men Iagttageren objectiverer i Afmagt.

Realbegrebernes Overgang fra det ideelle i det reelle Medium er
begrundet i en almægtig Objectivering. Objectiveringen er dobbelt,
indvortes og udvortes. Forsaavidt et bestemt System af Verdensbegreber
gjennemtænkes i Selvmagten, er Bevægelsen logisk, og Objectiveringsacten
ideel; forsaavidt det samme Begrebssystem oversættes i Gjenstandsmagten,
er Bevægelsen --- nu \emph{her}, nu \emph{der} --- antilogisk, og
Objectiveringsacten reel: den reelle Bevægelse har den ideelle til
Forudsætning.

% Det græske Bogstav er læst paa KB-Exemplaret: γ, ikke „y“ eller „7“.
\subrunin{γ) Vexelbestemmelse af ideel og reel Bevægelse.} For at
Realbegreber kunne dannes, maa et Udvortes uafbrudt reflectere sig i
sit Indvortes; Reflexer af Rum og Tid, Materie og Kraft med alle dertil
hørende Bestemmelser af Udvorteshed høre væsenlig med i
Begrebsdannelsen. Der tænkes ikke i Rummet, thi Rummet er ingen Form
for abstract Tænkning; og dog maa der, naar den materielle Virkelighed
skal begribes, nødvendig falde et Skjær af Rumsforhold ind i
Begribningens ideelle Medium. Realbegrebernes Momenter ere opløste
% TRYKFEIL, ikke rettet i Bogens eget Trykfeilblad: der staaer
% „Forstillingerne“ for „Forestillingerne“ --- et bortfaldet e. Ordet
% „Forestillinger“ staaer korrekt to Ord tidligere paa SAMME Linie, saa
% Forskjellen er umiddelbart maalelig; begge Exemplarer have Feilen,
% altsaa Sætterens. Gjengivet som trykt.
Forestillinger, og Forstillingerne opløste Billeder; der er i hvert
Billede en, om end nok saa
% --- p. 113 ---
% TRYKFEIL: der staaer „Rummct“ for „Rummet“ --- et c sat for et e.
% Læst paa KB-Exemplaret ved høi Forstørrelse og calibreret mod de sikre
% e'er i „Materien“, „er“, „ingen“ paa samme Linie, der alle vise den
% vandrette Tværstreg, som dette Bogstav mangler. Begge Exemplarer have
% Feilen. Trykfeilbladet retter den ikke. Gjengivet som trykt.
flygtig Anskuen af Rummct. Materien er ingen Tanke, men Tankens
Gjenstand, dens Modstand; det i Materien Materielle tænkes ikke, det
begribes som Tankens Andet i Reflexion af Anskuen og Tænken:
Sandsebilledet, hvis Betydning er, at Tankens Andet har Tilværelse
udenfor Tanken, bringer Reflexen med af det Materielle. Der er altsaa i
Realbegrebernes ideelle Bevægelse overalt Antydninger af reel
Bevægelse.

Den reelle Bevægelse har nu, hvad den ideelle forudsætter, en virkelig
Materie i virkelige Masser, et virkeligt Rum med bestemte Steder, en
virkelig Tid inden givne Grændser. Vistnok kunne de to uligeartede
Bevægelser ikke være parallele, den intellectuelle Vei fra Begreb til
Begreb er incomparabel med den udvortes Vei fra Sted til Sted; men at
den reelle Bevægelse er umulig uden den ideelle, er Naturphilosophiens
pulserende Hovedtanke.

Den ideelle Bevægelse tilhører Evigheden; den reelle foregaaer i Tiden.
Det store Verdenssystem af mechaniske Aarsagsbegreber bliver ikke til i
Tiden; det er evig færdigt. Den uendelige Begriben er simultan, og den
Succession, som er uadskillelig fra ideel Bevægelse, er kun den i
Evigheden selv reflecterede Antydning af Tid. De reelle Aarsager, hvori
Begreberne virkeliggjøres, tilhøre derimod Tiden som Tid, og
Bevægelserne foregaae med forskjellig Hastighed i Rummet. I evig
Forstand ere alle Bevægelser forud bestemte, thi det Ideelle maa i sig
% Bogens eget Trykfeilblad retter S. 113, Linie 8 f. n.:
% „blive til --- læs blive til.“ ---: Nielsen tilføier det manglende
% Punktum. Efter Husreglen staaer her Teksten SOM TRYKT, altsaa uden
% Punktum; Rettelsen er ikke udført. Linietællingen f. n. udelader
% Arkesignaturen „8“ nederst paa Siden.
være fuldendt, før det Reelle kan blive til
Bevægelsen efter Diagonalen i Hastighedernes Parallelogram maa i
Begrebet være fuldendt, thi den reelle Bevægelse kan kun paa
eftergjørende Maade udtrykke, hvad der oprindelig er i Begrebet.
Lovene, Selvmagtens Tanker, som evig tænkes og ere tænkte, bestemmes
ved \emph{ideel}, men Lovenes Opfyldelse ved \emph{reel} Bevægelse.

% Indløbende Overskrift NEDERST paa S. 113, sat med fed Antiqva.
% Arkesignaturen „8“ staaer under den og er ikke Text.
\runin{b) Udstykningens Aarsag.} I Vexelbestemmelsen af
% --- p. 114 ---
ideel og reel Bevægelse ligger, at den mechaniske Fordeling af Masser
og Kræfter i Rummet maa henføres til en gjennem udvortes Aarsager
virkende Grundaarsag. Fra den ideelle Side er Grundaarsagen at søge i
den Alt bestemmende Selvmagt, fra den reelle Side derimod i
Gjenstandsmagten, den til Udvorteshed svarende Tilværelsesmagt. Falder
nu Eftertrykket særlig paa Selvmagten, saa vise Grundbestemmelserne hen
paa lutter Villiesacter, og den i Sammenhængen udtalte Nødvendighed
hviler paa \emph{Vilkaarlighed}. Falder omvendt Eftertrykket
udelukkende paa Gjenstandsmagten, saa fremgaaer Delingen af særskilt
tilværende Aarsager, og Sammenhængsnødvendigheden hviler paa
\emph{Tilfældighed}. Kun derved, at den absolute Selvmagt erkjendes at
have al Magt i Andet som Realitet af sin ubetingede Overmagt, kan
Vilkaarlighedsmomentet fra den indre og Tilfældighedsmomentet fra den
ydre Causalitet enes med Nødvendigheden, og Viisdommens Plan i den
givne Verdensorden fremstiller sig da som et Værk af guddommelig
\emph{Frihed}.

% Det græske Bogstav er læst paa KB-Exemplaret: α, ikke „a“ eller „&“.
\subrunin{α) Vilkaarlighed.} Enhver oprindelig Bestemmelse, den være
forøvrigt af hvad Art den vil, er en Bestemmelse af Selvet, en
Selvbestemmelse. Opfattes Selvbestemmelserne deelviis, enkeltviis,
antage de Præg af Vilkaarlighed. Dog indsees det let, at en Flerhed af
Vilkaarlighedsbestemmelser enten maa falde ud fra hverandre, modsige
hverandre og tabe sig i det Meningsløse, eller ved indbyrdes
Overeensstemmelse forudsætte hverandre, ligge i en Plan og danne en
Sammenhængseenhed.

% S. 114 slutter midt i et Ord: Linien ender „hin-“ og S. 115 begynder
% med „anden, som hørte det Vilkaarlige til en Deel, ...“. Orddelingen
% opløses med \-% ; næste Batch (S. 115--124) fortsætter umiddelbart.
Formedelst Sammenhængen gaae Vilkaarlighedsbestemmelserne ind under
Nødvendigheden: det Vilkaarlige bliver nødvendigt, og det Nødvendige
vilkaarligt --- i den uendelige Reflexion. Men i den endelige Reflexion
seer det ud, som om Vilkaarlighed og Nødvendighed faldt ud fra hin\-%
%  MATHEMATICAL — work from the image; the OCR is not a witness here.
% --- p. 115 ---
% S. 115 begynder MIDT I ET ORD: S. 114 slutter „…faldt ud fra hin-“ og
% Orddelingen opløses med \-% paa den foregaaende Side. Ingen ny Periode,
% intet Afsnitsskifte, ingen gjentagen Bindestreg.
anden, som hørte det Vilkaarlige til en Deel, det Nødvendige til en
anden Deel af Planen.

Lad det f. Ex. ligge i Almagtens Plan, at to Kloder i Verdensrummet
skulle forholde sig til hinanden saaledes, at den ene afgiver et fast
Centrum, hvorom den anden bevæger sig i en vis Afstand: da kan jo
vistnok, naar det er Almagten, der absolut vil det, enhver af Kloderne
vilkaarlig sættes for sig, den ene meddeles en Bevægelse efter
Tangenten, Inertien lægges til, og en Tyngdelov fastsættes for at
bevare Forholdet; her er i denne Udstykning ikke mindre end fem
almægtige Vilkaarlighedsacter. Men nu kommer Nødvendigheden.

Vilkaarlighedsacterne forudsætte i Planen selv hverandre indbyrdes.
Den rationelle Mechanik gjør med fuld Ret gjældende, at der til Løsning
af en saadan Opgave høre
Betingelsesligninger,\footnote{% trykt Mærke: *)
% Denne Note begynder paa S. 115 og fortsætter under Texten paa S. 116;
% Mærket gjentages IKKE paa Fortsættelsen. Hele Noten staaer her ved sit
% Mærke, saaledes som Huusreglen foreskriver.
% De trykte Variable ere opreiste Antiqua; det gjengives ikke.
Opløses Kræfterne efter tre vinkelrette Axer i X, Y og Z, og betegnes de
løbende Coordinater ved x, y og z, da angiver Ligningen: $X : Y : Z = x :
y : z$, en Betingelse, der nødvendig maa opfyldes, dersom Kræfternes
Resultant skal gaae igjennem Coordinaternes Begyndelsespunkt, og deri
have et fast Centrum. Da $\dfrac{d^2x}{dt^2} = X$,
$\dfrac{d^2y}{dt^2} = Y$ og $\dfrac{d^2z}{dt^2} = Z$, saa vise
Ligningerne: $\dfrac{xdy - ydx}{dt} = a$, $\dfrac{zdx - xdz}{dt} = b$
og $\dfrac{ydz - zdy}{dt} = c$, at alle Constanterne a, b og c umulig
kunne være Nul, naar det bevægede Legeme ikke skal styrte lige imod
Centrum, men dets Rad. v. beskrive et Areal. Af Ligningerne følger
endvidere med Nødvendighed, at de paa Planerne XY, XZ og YZ projicerede
Arealer maae forholde sig som de anvendte Tider, og at Banen maa være en
plan Curve, efterdi den af Forudsætningen afledede Ligning:
$az + by + cx = 0$, netop er Ligningen for et Plan, der gaaer igjennem
% „0“ er her Nullet, ikke Bogstavet o; efterseet paa KB-Exemplaret.
Coordinaternes Begyndelsespunkt. Endelig, da der af de opstillede
Ligninger ligeledes kan udledes, at
$\dfrac{dx^2 + dy^2 + dz^2}{dt^2} = \dfrac{ds^2}{dt^2} = 2F(x, y, z) + C$
(idet s betegner et Banestykke, og $F(x, y, z)$ Integralet af
$Xdx + Ydy + Zdz$) og at Hastigheden $\dfrac{ds}{dt}$ --- saasandt
$Xdx + Ydy + Zdz$ er et fuldstændigt Differential, der hverken
indeholder Tiden eller Hastigheden, --- alene afhænger af
Coordinaterne, følger heraf, at de levende Kræfter vedligeholdes, at det
bevægede Legeme, hver Gang det kommer til samme Sted i Banen, har samme
Hastighed, at Bevægelsen altsaa fortsættes i det Uendelige i samme
Bane.} som ufravigelig maae opfyldes, at
% --- p. 116 ---
% Nederst paa S. 115 staaer Arksignaturen „8*“; den er ikke Text.
Principerne for Arealernes og de levende Kræfters Vedligeholdelse under
den angivne Forudsætning have en Nødvendighed, som end ikke Almagten
selv kan unddrage sig.

Den, der vil Øiemedet, maa ville Midlerne; det betyder i
Uendelighedens Sprog, at Øiemedet villes i Midlerne og Midlerne i
Øiemedet. Den begribende Selvmagt, som i ufordunklet Klarhed seer det
Hele bestemt ved Delene og omvendt, vælger lige saa vel det Betingede
for Betingelsernes som Betingelserne for det Betingedes Skyld. Paa
intet Punkt i den valgte Sammenhæng falder det Vilkaarlige udenfor det
Nødvendige. Her kan i den uendelige Begriben altsaa ikke være Tanke om,
at Vilkaarligheden skulde enten svække eller bryde Nødvendighedens
Love; men vel kan her være Fare for, at de reelle Aarsager fra dette
Synspunkt skulle blive tomme Magtphænomener af de ideelle, at
Vilkaarlighedsbestemmelsernes abstracte Idealitet skal undergrave den
naturlige Realitet, at den altbestemmende Selvmagt i en saadan
Eensidighed skal fortære Gjenstandsmagten.

% Underoverskriften er indløbende og SPÆRRET (de græske α) β) γ) er
% spærrede, medens a) b) c) er halvfede); det græske Bogstav er læst paa
% KB-Exemplaret: β, ikke „B“, „ß“ eller „P“.
\subrunin{β) Tilfældighed.} Den reelle Aarsag falder sammen med
Gjenstandsmagten i legemlig Existens. Seer man nu bort fra
Spørgsmaalet om den ideelle Causalitet og tager Tingene som de vise
sig, saa ere Lovene Naturlove, og
% --- p. 117 ---
den betingede Nødvendighed lider intet Afbræk. Ja det kunde næsten
synes, som om Vexelvirkningen mellem de naturlige Aarsager afgav en
sikkrere Garantie for lovbestemt Nødvendighed end en formel Reflexion
af guddommelige Vilkaarlighedsacter. For at blive ved vort Exempel, saa
ville nu de to Legemer ikke indtage deres Plads og tiltrække hinanden
ifølge guddommelig Ordre, men i Kraft af deres egen materielle Natur.
Enhver Biforestilling om, at det maaskee engang kunde behage den
absolute Selvmagt at omdanne Materien, meddele den nye Egenskaber,
indføre nye Love, en ny Sammenhæng med en ny Nødvendighed --- alle
slige Vilkaarlighedsforestillinger falde bort, og den factisk givne
Nødvendighed er som grundmuret.

Ved denne Opfattelse er der ganske vist Noget vundet; Aarsagens
mechaniske Realitet er hævdet, og Naturen forsaavidt kommen til sin
Ret. Men de, der mene, at Nødvendigheden selv, den rationelle
Nødvendighed, skulde være bedre sikkret paa det reelle end paa det
ideelle Grundlag, tage mærkelig feil. Om Nødvendighedsbegrebet gjælder
det vel, at det i Charakteer af Selvhedsform har Betingelserne som
Momenter i sit Væsen; men om den mechanisk reelle Nødvendighed gjælder
netop det Modsatte; den har som Andethedsform Betingelserne udenfor
sig: de existerende Led, hvoraf den betinges, frembyde overalt en
Tilværelsesside, der underligger Tilfældigheden. Vi lære det af vort
Exempel. De to Kloder, af hvilke den ene bevæger sig om den anden, maae
vel erkjendes at være i et reelt Nødvendighedsforhold og forudsætte
hinanden i udvortes Reflexion; men den ene er jo dog ikke Aarsag til
den anden: deres Samværen er i Udstykningens Form en reen
Tilfældighed. Forsøger man nu at hæve Tilfældigheden ved at føre begge
Masser tilbage til en fælles oprindelig Masse, saa er Tilfældigheden
der igjen; den er i Massens Tilværelse. Hæver man Tilfældigheden af
denne Tilværelse ved at lade den oprindelig rote\-%
% --- p. 118 ---
rende Masse fremkomme som et nødvendigt Resultat af Atomer, spredte i
Rummet, saa er Tilfældigheden der igjen; den er i Atomernes
Tilværelse.\footnote{% trykt Mærke: *)
% S. 118 bærer TO Fodnoter, mærkede *) og **); den anden sættes efter
% Mønstret fra S. 12 med en lokal \renewcommand.
% Bogtitlerne staae i almindelig Antiqua, ikke i Cursiv; „Über“ i den
% første Titel og „Ueber“ i den anden ere trykte netop saaledes.
I sit Skrift: Über die Natur der Cometen, Beiträge zur Geschichte und
Theorie der Erkenntniss; 2\textsuperscript{te} Auflage, Leipzig 1872,
har J. C. F. Zöllner --- under Overskrift: Ueber die Endlichkeit der
Materie im unendlichen Raume --- ved Raisonnement og Beregning søgt at
vise, at, efterdi Antagelsen af en Skaberact er en \emph{vilkaarlig}
Begrændsning af Causalitetsrækken, saa er det --- under Forudsætning
af, at Materiens Elementer ere den Newtonske og den Mariotte'ske Lov
underlagte --- nødvendigt at antage et endeligt, fra Evighed af
tilværende Qvantum Materie i det uendelige Rum. --- Saa interessante
disse Undersøgelser end i deres Enkeltheder ere, og saa villigt vi end
indrømme, at ogsaa Physiken kan have Gavn af „at gaae ud fra almene
Begreber“, at Erkjendelsens Fremskridt „ikke skal hæmmes ved
overleverede Fordomme“: saa staaer det dog fast, at en physisk
Aarsagsrække umulig kan begynde med et udvortes Givet, et første Led,
der ikke er foraarsaget. For at undgaae „Vilkaarligheden“, kaster man
sig i Armene paa Tilfældigheden; den endelige Materie i det uendelige
Rum bliver fra det Zöllnerske Synspunkt at tage som den første,
uforaarsagede Aarsag, det store Tilfælde i evig grundløs
Tilfældighed.}

% Det græske Bogstav er læst paa KB-Exemplaret: γ, ikke „y“ eller „7“.
% Overskriften er indløbende og spærret.
\subrunin{γ) Frihed.} Nødvendigheden i og for sig er et afmægtigt
Væsen, en rugende Nat i Mulighedernes Skjød, for tung og formløs til at
hjælpe sig selv frem. For at vinde Skikkelse og komme til Reisning, maa
Nødvendigheden støttes indenfra af Vilkaarligheden, udenfra af
Tilfældigheden. Det er Aanden, Selvmagten i evig Vexelvirkning med
Gjenstandsmagten, som bryder Mulighedernes Segl og bringer den dunkle
Nødvendighed for Lyset.%
\renewcommand{\thefootnote}{**)}%
\footnote{% trykt Mærke: **)
Om Nødvendighedsbegrebet, jvnfr. Grundideernes Logik. II. S.
253--275. Grundid. Logik. i kort Begreb. S. 78--83.}%
\renewcommand{\thefootnote}{*)}
Enhver fuldstændig Virkelighed, ethvert mechanisk, relativt sluttet
Verdenshele maa oprindelig bestemmes indenfra som en Begrebstotalitet,
% --- p. 119 ---
der fyldestgjør Selvhedens Fordring. Dette er Planen, det vilkaarlig
formede, men indholdsrige, omfattende og gjennemgaaende Almeenbegreb,
hvori de særlige Begreber reflectere og bestemme hverandre med rationel
Nødvendighed; dette er Udstykningens ideelle Aarsag.

Og dog er her endnu ingen Udstykning. Den virkelige Udstykning viser i
det Hele som i det Enkelte overalt hen til reelle Aarsager. Med de
reelle Aarsager fremkomme Kjendsgjerningerne, med Kjendsgjerningerne
Tilfældigheden; Tilfældigheden væves ind i Nødvendigheden, og den
udvortes betingede Nødvendighed, urolig og dog rolig, hviler ud paa
tilfældige Forudsætninger.

Dersom der nu i det hele Universum kunde opdages en reel Aarsag, saa
forudsætningsfri, saa factisk, at den i given Tilværelse faldt udenfor
ethvert bestemmende Aarsagsbegreb, da havde vi her en ægte Repræsentant
for Facticitetens ubetingede Magt, den evig blinde Tilfældighed. Men
saalænge en saadan Opdagelse endnu ikke er gjort, kunne vi frit lade
„den moderne“ Videnskabs „frie Tanke“ lege Nødvendighedsleg med den
% De to Anførsler ere efterseete paa KB: almindelig Antiqua, IKKE spærrede.
blinde Tilfældighed, idet vi ufravigelig fastholde den Grundsætning, at
enhver Kjendsgjerning som udvortes Aarsag indenfra maa være bestemt ved
sit Aarsagsbegreb, sin ideelle Aarsag.

Paa den oprindelige Eenhed af ideel og reel Aarsag maa al
Naturerkjendelse begrundes. Den philosophiske Kritik har tydelig
paaviist, at Materien saa lidt som dens Egenskaber og Kræfter lader sig
af nogen udenfor Naturen værende Magt meddele naturlig Tilværelse.
Dette er en for Naturphilosophien afgjørende Tanke; men Tanken mistydes
og forvrænges, naar den anvendes til at forsvare Naturens ubetingede
Autonomie, dens absolut-nødvendige Selvbegrundelse. Sandheden er, at al
Naturudvikling beroer paa naturlige Forudsætninger: Natur begynder med
Natur, hvor langt vi end gaae tilbage. Men naar denne Sandhed ikke
% --- p. 120 ---
skal forvanskes, maa det nødvendig underforstaaes, at Naturaarsagerne
ere reelle, og at enhver reel Aarsag er i Væsen bestemt af en ideel.
Udstykningens Aarsag er lige saa lidt at henføre til en absolut
Nødvendighed som til Facticitetens Tilfældighed; den naturlige Aarsag
er i Gjenstandsmagten --- et altomfattende System af udvortes
Aarsager.

Naturaarsagernes System er en Verdensplan, der udtaler sig i et
kategorisk Magtsprog. Vilkaarlighed, Tilfældighed, Nødvendighed ere
eensidige Begreber. Tilfældigheden hæves, naar det insees, at det
% TRYKFEIL: „insees“ for „indsees“ — d'et mangler. Samme Side sætter tre
% Linier længere nede „indsees“ rigtigt. Begge Exemplarer ere eens, saa
% Feilen er Sætterens; Trykfeilsbladet berører ikke S. 120. Gjengivet som trykt.
Tilfældige er, fordi det \emph{skal} være; altsaa: Tilfældigheden
hæves, naar den reflecteres i Vilkaarligheden. Vilkaarligheden hæves,
naar det indsees, at det Vilkaarlige er, fordi det \emph{maa} være, at
det sluttede System er nødvendigt for Planen; altsaa: Vilkaarligheden
hæves, naar den reflecteres i Nødvendigheden. Nødvendigheden hæves,
naar den i Leddenes Sammenhængseenhed reflecteres som Form for den
absolute Causalitets Indhold. Den absolute Causalitet, ideel i
Selvmagten, reel i Gjenstandsmagten, er overgribende Selvbestemmelse,
er Frihed; altsaa: Nødvendigheden hæves, naar den reflecteres i
Friheden.

% Den indløbende Overskrift c) er trykt HALVFED (ikke spærret), efter
% samme Mønster som a) og b) tidligere i Bogen. Det gjengives ikke.
\runin{c) Det første Bevægende.} Forholdet mellem Tankebevægelse og
Bevægelse i Rummet er paa eiendommelig Maade blevet udviklet af
Aristoteles. Bevægelse i Naturen er Overgang fra Mulighed til
Virkelighed, og omvendt; Bevægelse forudsætter to Led, et Bevægende og
et Bevæget, en actuel Væren i Formen og en potentiel i Stoffet:
Bevægelsen er evig som Form og Stof. Dersom Bevægelsen havde en
Begyndelse, saa maatte der forud for denne Begyndelse enten have været
et Bevægende og et Bevæget eller ikke. Have de ikke forud været, saa
maae de jo være blevet til, og der er da gaaet en Bevægelse forud for
den første Bevægelse. Have de forud været til, saa er det utænkeligt,
% --- p. 121 ---
at de ikke skulle have frembragt Bevægelse; der maatte i modsat Fald
være indtraadt en Virksomhed, hvorved de kom til at bevæge, og vi fik
da ogsaa i dette Tilfælde en Bevægelse forud for al Bevægelse. Ligesom
Bevægelsen paa endelig Maade forudsætter sig selv i det Uendelige,
saaledes forudsætter den i absolut Forstand et første Bevægende, der
ikke bevæges (τὸ πρῶτον κινοῦν ἀκίνητον ὄν); dette første Bevægende er
% Græsken er læst paa KB-Exemplaret. Trykken bruger ϰ-formen af kappa og
% ϑ-formen af theta; det er Skriftsnittets Varianter og gjengives med κ og θ.
Guddommen. Overalt i Tingenes Orden er det Bevægende selv et Bevæget;
men hvad der bevæges, det forandres og kan ikke virke nogen stadig,
eensformig Bevægelse. Det første Bevægende maa være uforanderligt; men
uforanderligt kan det kun være derved, at det efter sit Væsen udelukker
Muligheden af Andetværen; men at udelukke Mulighed af Andetværen er at
udelukke Stoffet. Stofløst, ulegemligt, udeleligt, udenfor Rummet, uden
Bevægelse, Liden eller Forandring (ἀπαθὲς καὶ ἀναλλοίωτον) er det
første Bevægende den fuldkomne Form, den evige Virkelighed, den rene
Energie, med eet Ord: Gud.

Med denne metaphysiske Slutningstanke bliver Aristoteles staaende paa
Grændsen mellem to indbyrdes uforligelige Verdensanskuelser,
Pantheismen og Theismen.

I theistisk Retning hævder Aristoteles Aandens absolute Oprindelighed
og lægger en afgjørende Vægt paa Modsætningen mellem Gud og Verden.
Guddommens Begreb falder i den fuldkomne Tænkning sammen med det
fuldkomne Gode. Som det evige, fuldkomne Gode bevæger Guddommen Alt,
idet den attraaes af Alt; thi Alt i Naturen er anlagt paa det Gode som
paa et høieste fælles Maal (πρὸς μὲν γὰρ ἕν ἅπαντα συντέτακται). Som
% „ἕν“ er trykt med acut, ikke med gravis; Mærket er det samme som over
% α'et i „ἅπαντα“. Den forkastede Læsning er „ἓν“, som Udgaverne har.
det i sig fuldkomne Gode kan Guddommen bevæge Alt, uden selv at
bevæges; thi det, der attraaes og tænkes, bevæger, uden selv at
bevæges. Hvad der \emph{viser sig} skjønt, attraaes, hvad der
\emph{er} skjønt, villes. Vi attraae det Skjønne, fordi det viser sig,
% „viser sig“ og det tostavelses „er“ ere begge spærrede paa KB;
% spacing.py fandt det første som RUN og oversaae som ventet det andet.
men det viser sig ikke, fordi vi attraae det. For\-%
% --- p. 122 ---
maalet bevæger, fordi det elskes, og hvad der bevæges af Formaalet,
bevæger igjen det Øvrige. Da det Guddommelige bestaaer i evig fuldendt
Energie, bestaaer det i evig Leven, thi Livet er Energie (ὥστε ζωὴ καὶ
αἰὼν συνεχὴς καὶ ἀΐδιος ὑπάρχει τῷ θεῷ). Dette Liv er saligt; en saadan
Glæde er os kun i korte Øieblikke forundt; men al Glæde (ἡδονή) er et
Drag af det guddommelige Princip.

At den Aristoteliske Theisme desuagtet staaer paa svage Fødder, har
Kritiken let ved at godtgjøre. En Guddom, der ikke gjør Andet end evig
at tænke Tænkningens Tænken (αὑτὸν ἄρα νοεῖ \ldots\ καὶ ἔστιν ἡ νόησις
νοήσεως νόησις) og som, for ikke at plette sin Fuldkommenhed med noget
Ufuldkomment, umulig kan have en Tanke tilovers for nogen Ting i
Verden, er et tomt, naturløst Væsen, hvis absolut nødvendige
Tankeindhold opløser sig i rene Tautologier.

Hvor den for sig værende Guds Indvirken paa Verden omhandles, ere
Vendingerne tvetydige. Der var, hedder det, ikke enten evigt Chaos
eller evig Nat, da Bevægelsen begyndte; men det Samme har altid været,
% Disse Linier ere blot løst udsluttede, ikke spærrede; efterseet paa KB.
enten i stadig Afvexling eller paa anden Maade, saasandt Virkeligheden
er forud for Muligheden. Skal nu det Samme altid virke i stadig
Afvexling, saa maa der være Noget, der altid virker paa samme Maade;
skal der være Opstaaen og Forgaaen, saa maa et Andet uophørlig virke
anderledes og atter anderledes. Det første Bevægende, der altid er sig
selv ligt, kommer da herved i et betænkeligt Samkvem med det anderledes
Virkende, der altid er sig selv uligt; og dog skulde den evige
Tænkningens Tænkning være uberørt af alt Andet. Hvad der siges om den
Trang og Attraa, hvormed Tingene bevæge sig efter det fuldkomne Ene, er
kun et teleologisk billedligt Udtryk for den alt Endeligt iboende
Negativitet, som viser, at det Relative er Moment for det Absolute.
% --- p. 123 ---
Paa en langt mere umiddelbar Maade have Tænkere som Cartesius og
Gassendi gjort Gud til det første Bevægende. Hvor uenige om
Naturaarsagerne disse to Philosopher end ellers ere, saa gaae de dog
begge ud fra en dogmatisk Adskillelse mellem det absolute Princip, den
første Aarsag til Alt, Skaberen, og de secundære Aarsager, hvortil
regnes „Alt, hvad der foruden Gud har en Kraft til at bevirke og
udrette Noget i Verden“.

Bevægelsens Aarsag, siger Cartesius, „er en dobbelt; deels en
% Anførselen er efterseet paa KB: almindelig Antiqua, ikke spærret.
almindelig og oprindelig, der foraarsager al Bevægelse, som overhovedet
er i Verden; deels en særegen, der bevirker, at Materiens enkelte Dele
erholde Bevægelser, som de forhen ikke havde. Den almindelige Aarsag er
Gud, som skabte Materien baade med Bevægelse og Hvile og derhos ved sin
naturlige Medvirken opholder al den Bevægelse og Hvile i Materien, som
han fra først af lagde ind i den“.

Det Charakteristiske ved denne Synsmaade er, at der ud fra en brat
Modsætning af det Overnaturlige og det Naturlige lægges an paa at
forklare Naturens Tilblivelse ved en endelig Skaberact, altsaa ved en
Act, som ikke er naturlig. Først skabes Materien; Skaberen indplanter
den visse Egenskaber, en vis Bevægelsesmængde o. s. v. Med en skabt
Materie, med skabte Kræfter og Egenskaber skal saa den skabte Natur
begynde at virke paa naturlig Maade. En saadan Begyndelse er under den
opstillede Forudsætning umulig. Som overnaturlig Virkning af en
overnaturlig Aarsag faaer den skabte Materie med sine Egenskaber og
Kræfter ved Begyndelsen selv en overnaturlig Tilværelse; dens første
Bevægelse er altsaa overnaturlig: hvorfra skulle de naturlige Aarsager
med de naturlige Bevægelser nu komme? De kunne ikke ligge i Materiens
Væsen, thi Materiens skabte Væren er jo overnaturlig; de kunne lige saa
lidt tilskrives Kræfterne, da Kræfterne selv ligeledes ere
overnaturlige. Som Aarsagen er, maa Virkningen være; Aarsag og Virk\-%
% --- p. 124 ---
ning maae trods Adskillelsen i det Ydre absolut have samme Væsen. At
den skabte Materie med sine Egenskaber og Kræfter viselig fordeles af
Skaberen og bliver ved den overnaturlige Fordelingsact til „secundære
Aarsager“, vil da sige, at de secundære Aarsager heller ikke ere
naturlige.

Den dogmatiske Forsikkring, at her nødvendig maa forudsættes en
overnaturlig Begyndelse, skjøndt vi Mennesker ikke ere i Stand til at
aflede det Naturlige af det Overnaturlige, er kun en passende Udflugt
for Theologer, men har Intet at betyde i Videnskaben. Et af To! Enten
% „Videnskaben.“ med Punktum; begge OCR-Vidner læste Komma, KB afgjør.
har Naturen som Natur en overnaturlig Begyndelse, og da maa den
nødvendig fortsætte som den har begyndt: alle saakaldte naturlige
Aarsager, Virkninger og Vexelvirkninger maae i Væren og Virken være
overnaturlige; eller, dersom der virkelig gives naturlige Aarsager og
naturnødvendige Bevægelser, da maa Begyndelsen svare til
Fortsættelsen: den i Væren og Virken naturlige Natur maa have en
naturlig Begyndelse.

Det er en over Almagtsbegrebet hvilende Uklarhed, som fra alle Sider
afføder Misforstaaelserne. Medens Aristoteles opfatter det Absolute i
Tænkningens Form som et metaphysisk Selvvæsen uden virkelig Almagt,
bliver det Absolute hos Gassendi og Cartesius forestillet som et
almægtigt Selv, men uden Gjenstandsmagt, uden Magt i Andet; medens
Aristoteles faaer et første Bevægende, der i Virkeligheden Intet kan
bevæge og ved Siden heraf maa lade Naturen være Natur og bevæge sig
selv paa Grundlag af en i negativ Form tilsnegen Bevægelse, gjøre
Gassendi og Cartesius den absolut ideelle Selvmagt til umiddelbar
Realitetsmagt og forudsætte en Almagt, som conseqvent maa fortære den
Natur, den selv frembringer.

Grundmodsigelserne finde, som vi have seet, deres Løsning i Magtens
oprindelige Selvfordobling. De mechaniske Naturaarsager tilhøre
Gjenstandsmagten og bevæge med
% S. 124 slutter MIDT I EN SÆTNING, med det hele Ord „med“ --- ingen
% Orddeling. Afdeling C er ikke sluttet: den fortsætter paa S. 125, som
% hører til næste Batch. Derfor hverken \streg eller Kapiteloverskrift her.
% ======================================================================
%  Første Deel — Tredie Kapitel: Himlens Mechanik   (printed pp. 125–196)
% ======================================================================
%  NOTE: this chapter head does NOT open a page. It stands
%  part-way down printed p. 125, below the close of the
%  preceding chapter. The head is therefore NOT pre-placed here:
%  the agent transcribing p. 125 writes it, in its printed
%  position, as  \kapitel{Tredie Kapitel.}{Himlens Mechanik.}
%  --- A. Verdensbygningen  (pp. 125–146) ---
% --- p. 125 ---
% S. 125 aabner med de tre sidste Linier af ANDET Kapitel, Afdeling C
% (Grundaarsagen: det første Bevægende). De fortsætter uden nyt Afsnit
% fra S. 124: Linien begynder ved Margenen uden Indrykning, og „Nødvendighed“ er
% et helt Ord (intet Orddelingsmærke paa S. 124's sidste Linie).
% Spatieringen af „Nødvendighed“ og „Frihed“ er bekræftet paa
% KB-Exemplaret.
\emph{Nødvendighed}; men den ideelle Causalitet, de naturlige Aarsagers
indre Bestemmende, tilhører Selvmagten, det første Bevægende, der
bevæger med \emph{Frihed}.

% Kort centreret Streg mellem Andet og Tredie Kapitel. Den er trykt som
% en DOBBELT Linie (to tynde Streger tæt over hinanden) i begge
% Exemplarer; hverken Længden eller Dobbeltheden gjengives --- \streg
% tjener overalt, efter Husreglen.
\streg

% Kapiteloverskriften AABNER IKKE EN SIDE: den staaer midt paa S. 125,
% under Slutningen af Andet Kapitel og over Afdeling A. Læst paa Siden
% selv (KB-Exemplaret), ikke i Indholdet.
\kapitel{Tredie Kapitel.}{Himlens Mechanik.}

% Afdelingsoverskriften følger umiddelbart efter, endnu paa S. 125:
% „A.“ centreret, derunder „Verdensbygningen.“ sat med fed, spatieret
% Skrift. Den Forskjel gjengives ikke.
\afdeling{A.}{Verdensbygningen.}

Der er i Verden Intet objectivt uden det, som objectiveres: Rummet og
Tiden, Materien og dens Kræfter, Alt i den ydre Natur, i Magtens Rige,
har Objectivitet --- ikke umiddelbart af sig selv, men af Aandens
Almagt. Fra Endelighedens Synspunkt gjælder det jo vel, at de synlige
Ting i deres ydre Virkelighed ere uafhængige af Øiet, hvormed de sees;
% „sees“ uden Accent, læst paa KB-Exemplaret; begge OCR-Vidner give
% „sées“, hvilket er Arbeidsscanningens Prikker og ikke Typen.
men spørges der om Erkjendelsens Sandhed, saa lærer Objectiveringens
Lov, at den Gyldighed, der kan tilkomme Gjenstanden som tilværende
Object, er afhængig af den Gyldighed, som ifølge videnskabelig Indsigt
kan tilkomme Gjenstandens Objectivering. Verdensbygningens Objectivitet
er herpaa et storartet, talende Exempel. Vel er Stjernehimlen Gjenstand
for umiddelbar Betragtning, og det kunde forsaavidt synes, som om
Enhver, der seer op mod den høie Hvælving, uvilkaarlig maatte see, i
hvilken Stiil denne Bygning er opført. Og dog er det en Illusion; den
uoverskuelige Bygning med de mange Værelser staaer ikke aaben for naivt
% --- p. 126 ---
nysgjerrige Tilskuere. Ved jordiske Bygninger er det den virkelige
Opførelse, der koster Arbeide, Betragtningen kommer af sig selv. At
faae Himmelbygningen reist koster intet Arbeide, det Værk er fuldført;
men at see ind i Bygningen og, om end kun tilnærmelsesviis gjøre sig
bekjendt med Planen i dens eiendommelige Construction, er betinget af
Arbeide. Det har kostet Menneskeaanden Aartusinders Arbeide at faae
Verdensbygningen objectiveret i alle til Virkeligheden svarende
Hovedtræk; dette ideelle Bygningsarbeide er af uendelig større Værd end
alt det materielle Slid, der er anvendt paa Ægyptens Pyramider. Vi
følge Objectiveringens Fremgang fra \emph{Illusionens Trin} gjennem
\emph{Reflexionen til den endelige Prøvelse}.
% Spatieringen bekræftet paa KB-Exemplaret: „Illusionens Trin“ og
% „Reflexionen til den endelige Prøvelse“ ere spatierede, „gjennem“
% derimod ikke. Punktet efter „Prøvelse“ sættes udenfor \emph.

% Indløbende Overskrift, læst paa Siden selv (KB-Exemplaret) og ikke i
% Indholdet: sat med halvfed Skrift, paa samme Linie som det Afsnit den
% aabner. Skjelnen mellem halvfed og spatieret gjengives ikke.
\runin{a) Illusionen: den Ptolemæiske Verdensbygning.} Oldtidens
Naturbevidsthed er væsenlig mythisk. Hvad der tilføres Sandserne
udenfra og maa henføres til den ydre Natur, smelter uvilkaarlig sammen
med hvad der under en stemningsrig Medvirkning af Følelse, Phantasie og
Tænkning kommer indenfra, fra Sindets Grund og omdanner det Udvortes;
hvorledes det gaaer til at Verdensbilledet og den virkelige Verden
trods deres indbyrdes Forskjellighed mødes og sammensmelte i
Bevidstheden selv, denne Objectiveringens Hemmelighed er ikke til for
en naiv Betragtning. Synets Sands er objectiv, men paa Grund af en
Illusion; i Opfattelsen af den synlige Himmel hersker Illusionen saa
meget mere som det jordiske Standpunkt ved sig selv er fastsat, og
ingen af de øvrige Sandser rækker med sit Vidnesbyrd ud i det uendelige
Rum. Det er da en Selvfølge, at det første Billede af
Verdensbygningen maa være illusorisk. Men Illusionen er ikke et simpelt
Bedrag. Himmelkuglens Omdreining om Jorden, Solens Op- og Nedgang,
Maanens særlige Bevægelse og vexlende Udseende, alt Sligt har fra et
jordisk Stade lige saa vist sin Sandhed som
% --- p. 127 ---
Adskillelsen af Lys og Mørke, Morgen og Aften, Dag og Nat.

Den Ptolemæiske Verdensbygning hviler paa gamle Grundstene; det første
Materiale blev samlet paa en Tid, da en vedholdende Beskuelse af
Phænomenerne paa Himlen mindre var et Videnskabens end et Religionens,
et Folketroens Anliggende. De syv Planeter: Maanen, Mercur, Venus,
Solen o. s. v. bleve naturligviis ikke betragtede som livløse Masser af
sammenrullet Materie, men som høie Magter i straalende Glands, som
Guder. I denne Illusion blev Aandens Eenhed med Naturen fastholdt. De
henover Himlen vandrende Guder gave ikke blot ved deres afmaalte
Bevægelser Menneskene en Maalestok til Beregning af den borgerlige Tid;
de vare selv Tidens Beherskere. For at overskue Begivenhederne og styre
de jordiske Anliggender deelte disse Tidens Herskere Tiden imellem sig;
det var da ikke tilfældigt, at enhver af dem fik sin egen Dag i
Ugen.\footnote{% trykt Mærke: *) --- eneste Note paa S. 127; den staaer
% helt paa denne Side og gaaer ikke over paa næste Blad.
Ved at ordne Planeterne i følgende Række: Saturn, Jupiter, Mars, Solen,
Venus, Mercur og Maanen, og gaae ud fra den Forudsætning, at hver af
dem vel har sine Timer i Døgnet, men dog saaledes, at den, der regjerer
i den første Time, er Overherre og giver Dagen Navn, vil man i ---
Illusionens Forstand --- forstaae, hvorledes de syv Dage i Ugen kunne
undergives de syv Planetguder. Begynder Dagens første Time f. Ex. med
Saturn, saa have vi en Løverdag, \textit{Dies Saturni}. Ved at dele
Herredømmet med de sex Medregenter vil Saturn fremdeles i egen Person
komme til at beherske den ottende, den femtende og den to og tyvende
Time; derpaa vil Jupiter faae den tre og tyvende, Mars den fire og
tyvende, og Regentskabet i den følgende Dags første Time tilfalde
Solen, saa at Søndagen bliver Solens Dag, \textit{Dies Solis}. Som det
gik med Saturn, gaaer det nu med Solen; naar den sidst har hersket i
den to og tyvende Time, hersker Venus i den tre og tyvende, Mercur i
den fire og tyvende, saa at den følgende Dags første Time tilfalder
Maanen, og bliver en Mandag, en Maanens Dag, \textit{Dies Lunæ}; o. s.
frmdl.}

% --- p. 128 ---
% S. 128 aabner et NYT Afsnit (Linien er indrykket paa Siden).
Er Illusionen den naturnødvendige Form, hvorunder Menneske-Aanden paa
den naive Anskuelses Trin formaaer ved Gudernes Bistand at see en
Afglands af sit objective Væsen, saa kan det ikke undre os, at Oldtiden
ikke blot var uvilkaarlig blændet af Illusionen, men vi forstaae, at
den tillige paa vilkaarlig og dog instinktagtig Maade, halvt bevidst,
halvt ubevidst, holdt Illusionen fast som den sandeste Form for sit
rigeste og dybeste Indhold. Saaledes ere Stjernebillederne ingenlunde
dannede ved blot vilkaarlig Gruppering med vilkaarlige Navne; og lige
saa lidt kan det antages, at de Gamle virkelig have bildt sig ind at
see Fortidens Heroer mellem Slanger og Uhyrer lyslevende paa Himlen.
Det Vilkaarlige og det Uvilkaarlige smelter sammen i en Sindets Trang
til ophøiet Symbolik, til illusorisk Form for Aandens
Selvobjectivering i Naturen. Stjernebillederne ere, som J. L. Heiberg
saa træffende har udtrykt det, „sande Billeder, som Oldtiden har
ophængt i det store Gallerie, hvor alle Slægter skulle betragte dem.
Her viser det sig ret, at Stjernehimlen, langt fra at være fremmed for
Jorden, er Reflex af de jordiske Skjæbner, de jordiske Tanker og
Drømme.“

Grundtegningen til den Ptolemæiske Verdensbygning har den græske
Philosophie leveret; det er en Tegning, der ligger for Illusionen fra
Først til Sidst. Den af Verdensideen, det skjønt afrundede kosmiske
Hele opfyldte Phantasie forstod at forme det stofagtige Indhold til
Tankens Behag saa sindrigt, at Illusionens Phænomener fik en saare dyb
rationel Begrundelse.

I mythisk ideale Billeder beskriver Platos Timæus, hvorledes Guden, den
øverste Bygmester, fra Begyndelsen af indrettede Alt paa det Bedste,
idet han dannede Verdenssjælen til det Mellemvæsen, hvori Fornuft og
Sandselighed, Idee og Legemlighed alene lade sig forene. Alt Legemligt
bevæges af et Andet; kun Sjælen, Selvbevægelsens egen
% --- p. 129 ---
% Det græske er læst paa KB-Exemplaret ved stærk Forstørrelse: ἡ med
% spiritus asper, αὐτὴ med spiritus lenis, αὑτὴν med spiritus asper.
% Circumflex'en i „κινεῖν“ er i Trykket sat over ε, ikke over ι; det er
% Skriftens Vane igjennem hele Bogens Græsk og lader sig ikke gjengives
% i Unicode, hvor Tegnet maa sidde paa ι.
Kraft (ἡ δυναμένη αὐτὴ αὑτὴν κινεῖν κίνησις), bevæger sig selv og Alt.
Endnu før de legemlige Elementer vare dannede, blandede Guden den
udeelbare, sig selv lige Substans og det legemlig Deelbare saaledes, at
der fremkom et tredie Væsen, Sjælesubstansen nemlig, i hvilken det
Selvlige og det Andet ere sammenføiede. Dette Hele deeltes nu efter
Grundtallene for det harmoniske og astronomiske System paa en saadan
Maade, at Fixstjernehimlens og Planeternes Kredse fremkom ved en egen
Længdedeling; thi alle Tal- og Maalforhold indeholdes i Sjælen, dens
Væsen er lutter Tal og Harmonie. Plato lader Verdenssjælens syv Dele
forholde sig til hverandre som Tallene: 1, 2, 3, $2^2$, $3^2$, $2^3$,
$3^3$; paa denne Deling grunder sig Planeternes Afstand fra Verdens
Midte; Solen er 2 Gange, Venus 3 Gange, Mercur 4 Gange, Mars 8 Gange,
Jupiter 9 Gange og Saturn 27 Gange saa langt borte fra Jorden som
Maanen. Idet Sjælen under Kredsbevægelsen vender tilbage til sig selv,
siger den gjennem sit hele Væsen overalt til sig selv, hvad det er, den
i sit Omløb berører; om det er noget Deelbart eller Udeelbart, hvormed
det er eensartet, hvorfra det er forskjelligt og hvorledes det
overhovedet forholder sig i sin Vorden og Væren. Denne Tale forplanter
sig lydløst i det sig selv Bevægende og avler Erkjendelsen. Fra Kredsen
af det Andet, Ekliptiken, udgaaer de rigtige Forestillinger og
Meninger; fra Fixstjernekredsen, Æqvator, Kredsen for det Selvlige,
stammer Fornufterkjendelse og forstandig Viden. Den illusorisk
dybsindige Grundtanke er, at det harmonisk bevægede Verdenshele er et
levende, besjælet Væsen.

Aristoteles gjør et Skridt fremad. Han forsmaaer den mythiske
Indklædning; han underkaster Platos Lære om Verdenssjælen en streng
Kritik og holder sig til Virkeligheden. Man skulde nu mene, at den
aristoteliske Bygning, befriet fra Mythens og Illusionens blændende
Lystaager, maatte staae fast paa Fornuftens Grund, og de virkelige
% --- p. 130 ---
Omrids holde sig i Videnskabens klare Dag. Og dog er den hele
Objectivering en fuldstændig Illusion. Hvor lidt den menneskelige
Fornuft ved nogle faa almindelige Iagttagelser er i Stand til at trænge
saaledes objectivt ind i Naturen, at den i Kraft af sin egen
Verdensidee kan construere den virkelige Verden, derpaa er den
aristoteliske Physik et storstilet Exempel. Alt fra Først til Sidst har
denne Philosoph gjennemraisonneret, begrundet, defineret, demonstreret;
det med udholdende Skarpsindighed gjennemarbeidede System gjør Indtryk
af en omfattende uimodsigelig Tankenødvendighed; men Nødvendigheden
hviler paa en Illusion.

Hvorledes Aristoteles har tænkt sig det første Bevægende, er viist i
foregaaende Kapitel; her skal nu i Korthed berøres, hvorledes han har
opfattet Bevægelserne i Legemverdenen. Al Bevægelse i Rummet er enten
retliniet eller kredsformig eller sammensat af begge. Den linierette
Bevægelse er ifølge Elementernes Natur igjen enten \emph{mod} eller
\emph{fra} Midtpunktet.
% Spatieringen af „mod“ og „fra“ er bekræftet paa KB-Exemplaret; det
% senere „mod Midtpunktet“ i samme Afsnit er derimod ikke spatieret.
De fire Elementer bevæge sig efter den rette Linie, Jorden, der er
absolut tung, mod Midtpunktet, Ilden, der er absolut let, fra
Midtpunktet; Vand og Luft danne Overgange; men Kredsbevægelsen tilhører
Ætheren. Da nu ethvert Element har sin naturlige Bevægelse og sit
naturlige Sted, saa er Verdensbygningens Grundskikkelse med
Nødvendighed bestemt. Ubevægelig, i Midten af det Hele, hviler den
kugleformige Jord. Det er umuligt, at Jorden kan bevæge sig; det
strider mod dens Natur. Det Jordiske kan kun bevæge sig, forsaavidt det
bevæger sig mod Midten; men Jordkuglen er jo Midten; ethvert Legeme
kommer til Hvile, naar det har naaet det Sted, hvortil dets Bevægelse
fører det: alt Jordisk har altsaa sin Hvile i Midten. Staaer det fast,
at Jorden med dens Luft- og Ildkreds er kugleformig, saa følger heraf,
at Himlen, der omgiver Jorden, ogsaa maa være kugleformig.
% --- p. 131 ---
Men det indsees desuden af andre Grunde. Himlen maa nemlig som det
Fuldkomneste have den fuldkomneste Form; men den fuldkomneste legemlige
Form er Kuglen. Som Maalet for alle Bevægelser maa Himlen bevæge sig
med den største Hastighed, altsaa følge den korteste Vei; men den
korteste Vei fra det Samme til det Samme er Cirkelen. Alle
Himmellegemer bevæge sig daglig fra Øst til Vest; men syv iblandt dem
bevæge sig desuden i kortere eller længere Tidsmaal fra Vest til Øst.
Da det nu er afgjort, at Jorden hviler, saa maa enten hele Himlen
bevæge sig eller blot Stjernerne eller baade Himlen og Stjernerne. Men
at baade Himlen og Stjernerne skulde bevæge sig, er utænkeligt; thi
hvorledes skulde det vel gaae til, at Stjernerne i Hastighed som af sig
selv holdt Skridt med Hastigheden i deres egne Kredse? At Stjernerne
skulde bevæge sig og deres Kredse hvile, er ogsaa utænkeligt; det
maatte jo være en uforklarlig Tilfældighed, om en Stjernes Bevægelse
skulde nøiagtig svare til Kredsens Størrelse. Der staaer da kun tilbage
at antage, at Kredsene bevæge sig, medens de paa dem befæstede Stjerner
hvile (τοῦς μὲν κύκλους κενεῖσθαι τὰ δὲ ἄστρα ῆρεμεῖν).
% Det græske Citat (Aristoteles, De caelo II, 8) er læst paa
% KB-Exemplaret ved stærk Forstørrelse og gjengives som TRYKT. Det
% afviger paa tre Punkter fra den rigtige Form
% (τοὺς μὲν κύκλους κινεῖσθαι, τὰ δ' ἄστρα ἠρεμεῖν):
%   1) τοῦς med circumflex over υ, ikke τοὺς med gravis;
%   2) κενεῖσθαι med ε i anden Stavelse, ikke κινεῖσθαι --- ε'et er
%      utvetydigt og lig ε'et i fjerde Stavelse;
%   3) ῆρεμεῖν med circumflex over η, ikke ἠ med spiritus lenis: Tegnet
%      er en flad, bølget Streg, tydelig forskjellig fra de tynde
%      Aandetegn i „ἄστρα“ og fra gravis i „δὲ“ paa samme Linie.
% Der er intet Komma trykt mellem κενεῖσθαι og τὰ. Circumflex'en i
% κενεῖσθαι og ῆρεμεῖν sidder paa ε, ikke paa ι, ganske som paa S. 129.
Den berømte Astronom Eudoxus fra Knidos havde allerede udkastet et
System med 27 Sphærer, af hvilke de 26 gik paa Planeterne. Hertil
føiede Kallippus endnu 7 Sphærer, saa at Sol og Maane fik hver to,
Mercur, Venus og Mars hver een. Af den Maade, hvorpaa Aristoteles
optager og begrunder de indviklede Sphæretheorier, bliver det ret
indlysende, at han, langt fra at ophæve Illusionen, tvertimod arbeider
den saa dybt ind i den tænkende Bevidsthed, at henved to Aartusinder
maatte hengaae inden Videnskaben fra Grunden af fik den udrenset.

Paa et illusorisk Grundlag kan en forstandig Objectivering imidlertid
foregaae paa det Besindighedens Vilkaar, at Iagttageren ikke vil
paatage sig at give Grunde for Alt, men holde sig til Kjendsgjerninger
og ved empirisk Induc\-%
% --- p. 132 ---
tion lægge an paa at udfinde Bevægelseslovene. Hermed begynder den
videnskabelige Astronomie.

Den forstandige Objectivering gjennemløber nu tre Hovedtrin. Det første
Nødvendige er at drage Kredse og afsætte Punkter, hvorved det bliver
muligt at orientere sig paa Himmelkuglen. Hertil svare de bekjendte
Storcirkler, Horizonten med Zenith og Nadir, Æqvator med
Verdenspolerne, Ekliptika, kort: hvad der ud fra forskjellige
Synspunkter hører til, for at angive en Stjernes Coordinater, d. v. s.
bestemme dens sphæriske eller angulære Sted. Uagtet slige Punkter,
Poler og Cirkler ikke have materiel Virkelighed, saa have de dog en for
Anskuelse og Tænkning uafviselig Gyldighed; de have en væsenlig
Objectivitet.

Efterhaanden som Indbegrebet af Objectiveringens første Betingelser
udvides, og den fortsatte Iagttagelse af Bevægelsesphænomenerne leder
til stedse nøiagtigere og speciellere Resultater, vil det gaae op for
Bevidstheden, at Verdensharmonien, det Heles vidunderlige Sammenhæng,
er langt mere indviklet end man fra Begyndelsen af har forestillet sig
den. Her indtræder en Modsigelse mellem den oprindelige Forudsætning af
Jorden som Midtpunktet for alle Kredse, af Himmellegemernes jævne,
eensformige Bevægelse o. s. v. og de øiensynlige, idetmindste
tilsyneladende Uregelmæssigheder, som Observationerne udvise;
Bestræbelsen for at løse dette store, ved Modsigelsen angivne, Problem
er betegnende for det næste Trin. Er Solbanen en Cirkel og Bevægelsen
jævn, hvorledes kan det da gaae til, at Aarstiderne ikke ere lige
lange, at Solen bruger en meget forskjellig Tid til at gjennemløbe
enhver af Ekliptikens fire Qvadranter? Til Phænomenets Forklaring
opstillede Hipparch to Hypotheser, som dog --- hvad Ptolemæus siden
indsaae --- kun i Formen, ikke i Virkeligheden, ere forskjellige.
Ifølge den ene af disse Hypotheser beskriver Solen med jævn Hastighed
en Cirkel, hvis Centrum ligger udenfor Jorden;
% --- p. 133 ---
ifølge den anden tillægges der Solen en jævn Bevægelse paa en Cirkel
--- Epicyklen --- hvis Centrum selv bevæger sig paa en Cirkel --- den
defererende Cirkel --- der har Jorden til Centrum.
% Trykfeil: mellem „Centrum“ og „selv“ staaer i BEGGE Exemplarer et
% lodret indfarvet Mærke, der hverken har Kommaets Krog eller dets
% Plads under Linien --- rimeligviis en opstegen Udslutning, der har
% taget Farve. Gjengivet som intet Tegn; tesseract læser Komma.
Hipparch er Epicykeltheoriens egenlige Begrunder. Endelig, og dette er,
hvad Methoden angaaer, Objectiveringens sidste, fuldendende Skridt,
bliver den hele Planettheorie ved fortsat Anvendelse og Udvidelse af
Epicykelhypothesen bragt til relativ Afslutning. Præcessionen,
Jævndøgnspunktets tilbageskridende Bevægelse, blev paaviist af
Hipparch; den Ujævnhed i Maanebanen, som siden fik Navn af Evectionen,
blev opdaget af Ptolemæus, hvis \textit{Μεγάλη σύνταξις} har sat
Krandsen paa den gamle Lærebygning. I illusorisk Objectivitet med nogle
arabiske Reparationer blev Bygningen staaende urokket i fjorten
hundrede Aar.

% Indløbende Overskrift, læst paa Siden selv (KB-Exemplaret): halvfed.
\runin{b) Reflexionen: den Copernicanske Verdensbygning.} Phantasien
kan aande i Tænken, sætte Tanker i Bevægelse og reflectere; omvendt kan
Tænkningen bygge paa en Anskuelse og, naar Anskuelsen er i en Svæven,
phantasere. Illusion og Reflexion danne i og for sig ikke nogen skarp
Modsætning. Der kan, som vi nylig have seet, opbydes megen Reflexion
for at uddanne og begrunde en illusorisk Theorie, og ligesaa kunne,
hvad vi i det Følgende skulle see, forskjelligartede Illusioner
indsnige sig i en for øvrigt gjennemreflecteret, til Virkeligheden
svarende Grundanskuelse. Men naar Talen er om, hvad der i det Hele
characteriserer et Objectiveringens Standpunkt, bliver Forskjellen
kjendelig.

Der er en dybt gaaende Forskjel mellem en Verdensbetragtning, hvis
Udgangspunkt ubevidst og ligefrem tilhører Illusionen, og en
Undersøgelse af Phænomenernes Sammenhæng, som begynder med at
underkaste Betragtningens Grundlag, den første, almindelige
Forudsætning, en streng videnskabelig Prøvelse. I denne Forstand kunne
vi, hvad
% --- p. 134 ---
Læren om Verdensbygningen angaaer, med Grund sige, at Copernicus staaer
i Astronomiens Historie som en heroisk Skikkelse, forsaavidt han med
opladt Syn paa en høiere Verdensharmonie, er den Første, der har havt
den Genialitetens Overlegenhed og det intellectuelle Mod at bryde med
Vanens Magt, stille sig paa Reflexionens Standpunkt og bekæmpe
Illusionen. „Forfatterne“, siger han, „ere sædvanligviis saa enige om,
at Jorden hviler i Verdens Midte, at de ansee det for meningsløst, ja
endog latterligt, at antage det Modsatte. Men Spørgsmaalet er endnu
ikke afgjort og kan ikke afvises med
% Bogens Vane: Anførselen gjenaabnes med „ efter Indskuddet „siger
% han“. Her lukkes der imidlertid ogsaa efter „Forfatterne“, saa
% Tegnene gaae op paa Siden (to „ og to “). Gjengivet som trykt.
Foragt.“\footnote{% trykt Mærke: *) --- første af TRE Noter paa S. 134.
De Revolutionibus orbium cœlestium. Norimbergæ MDXLIII. Lib. I. Cap.
V.}

At vælge sit Standpunkt i Solen, naar alle Traditionens Tilhængere ere
enige om at holde sig ved Jorden, maa vel i Manges Øine tage sig ud som
en Grille, et forvovent Indfald af en høitflyvende Phantasie. Men
Copernicus vil ikke phantasere, han vil i Alvor reflectere. Det er det
virkelige Forhold mellem Jord og Himmel, et Forhold, som kun en
gjennemført Reflexion, ikke et umiddelbart Skjøn, kan bedømme, han fra
Grunden af vil have undersøgt.%
\renewcommand{\thefootnote}{**)}%
\footnote{% trykt Mærke: **) --- anden af tre Noter paa S. 134.
Qvam ob causam ante omnia puto necessarium, ut diligenter
animadvertamus, qvæ sit ad cœlum terræ habitudo, ne dum excelsissima
scrutari volumus, qvæ nobis proxima sunt, ignoremus, ac eodem errore
qvæ telluris sunt attribuamus cælestibus. De Revol. Lib. I. Cap. IV.}%
% Sætteren blander Ligaturerne: „cœlum“ med œ, men „cælestibus“ med æ i
% een og samme Note. Læst paa KB-Exemplaret. Gjengivet som trykt.
\renewcommand{\thefootnote}{*)}
Det er Forskjellen mellem tilsyneladende og virkelig Bevægelse, altsaa
Bevægelsens Objectivitet, paa hvis Bestemmelse det kommer an.%
\renewcommand{\thefootnote}{***)}%
\footnote{% trykt Mærke: ***) --- tredie af tre Noter paa S. 134.
% DENNE NOTE GAAER OVER BLADET: den begynder under Texten paa S. 134
% („Omnis enim … Si igitur motus ali-“) og fortsættes øverst i
% Notefeltet paa S. 135 („qvis terræ deputetur … Lib. I. cap. V.“) uden
% at Mærket gjentages. Efter Husreglen sættes hele Noten her ved sit
% Mærke; Orddelingen „ali-/qvis“ er opløst.
Omnis enim, qvæ videtur secundum locum mutatio, aut est propter
spectatæ rei motum, aut videntis, aut certe disparem utriusqve
mutationem. Nam inter mota æqvaliter ad eadem non percipitur motus,
inter rem visam dico et videntem. Terra autem est, unde cœlestis ille
circuitus aspicitur, et visui reproducitur nostro. Si igitur motus
aliqvis terræ deputetur, ipse in universis, qvæ extrinsecus sunt, idem
apparebit, sed ad partem oppositam, tanqvam prætereuntibus, qvalis est
revolutio cotidiana imprimis. Lib. I. cap. V.}%
\renewcommand{\thefootnote}{*)}

% --- p. 135 ---
% S. 135 aabner et NYT Afsnit (Linien er indrykket paa Siden).
At Himlens daglige Bevægelse fra Øst til Vest kan forklares som en
skuffende Gjengivelse af Jordens virkelige Bevægelse fra Vest til Øst,
er Copernicus saa langt fra at ville udgive for nogen ny Tanke, at han
tvertimod minder om, at blandt de Gamle have f. Ex. Pythagoræerne
Heraclides og Ecphantos og Syracusaneren
Nicetas\footnote{% trykt Mærke: *) --- eneste EGEN Note paa S. 135; den
% staaer under Slutningen af den fra S. 134 fortsatte ***)-Note.
Skal nok være Hicetas.} ifølge Cicero været af denne Mening. Men,
tilføier han, naar dette indrømmes, saa følger en anden ikke mindre
Tvivl om Jordens Sted, endskjøndt det nu næsten af alle antages, at
Jorden er Verdens Midte. Ud fra denne Tvivl begynder Reflexionen da
saaledes, at den vragede Synsmaade stilles imod den autoriserede, og
Begrundelsen skeer ved Raisonnement mod Raisonnement. Ptolemæus finder
det urimeligt, at Jorden skulde bevæge sig. Ved den hurtige Omdreining
maatte jo nemlig Tingenes Sammenhæng opløses, alle jordiske Baand
briste, ja, hvad der ret er en latterlig Forestilling, Himlen selv
maatte falde ned. Men, spørger Copernicus, hvorfor anvender man ikke
hellere dette paa Verden, hvis Bevægelse maa være saa meget hurtigere,
som Verden er større end Jorden? Mon Himlen maaskee er blevet saa
uendelig stor, fordi den med ubeskrivelig Hastighed
(\textit{ineffabili motus vehementia}) drives bort fra Midten? I
Sandhed, dersom dette Raisonnement holder Stik, maa Himlens Storhed
være i uendelig Tiltagen; Hastigheden maa frembringe Storheden og
Storheden igjen Hastigheden uden Maal og Grændse. Men ifølge det
physiske Axiom, at det, som er uendeligt, hverken kan overskrides eller
bevæges, maa Himlen nødvendig staae stille (\textit{stabit necessario
cœlum}). Men naar Himlen staaer
% --- p. 136 ---
stille, maa det være Jorden, der bevæger sig. Skulde Jorden være
Verdens stillestaaende Midte, da maatte den jo være fælles Centrum for
alle Baner; men at den ikke er det faste fælles Centrum, fremlyser
saavel af den øiensynlige Ujævnhed i Planeternes Bevægelse som af deres
foranderlige Afstande fra Jorden. Da der nu vitterlig existerer flere
Centra, saa er der al Grund til at see sig om efter et andet
Verdenscentrum, et andet
Tyngdepunkt.\footnote{% trykt Mærke: *) --- første af TO Noter paa S. 136.
Pluribus ergo existentibus centris, de centro qvoqve mundi non temere
qvis dubitabit, an videlicet fuerit istud gravitatis terrenæ, an aliud.
Eqvidem existimo, gravitatem non aliud esse, qvam appetentiam qvandam
naturalem partibus inditam a divina providentia opificis universorum,
ut in unitatem integritatemque suam sese conferant in formam globi
coëuntes. Qvam affectionem credibile est etiam Soli, Lunæ, cæterisque
errantium fulgoribus inesse, ut ejus efficacia in ea qua se
repræsentant rotunditate permaneant, qvæ nihilominus multis modis suos
efficiunt circuitus. Lib. I. cap. IX.}
% Noten blander qv og qu: „qvoqve, qvis, qvam, qvandam, qvæ“ mod
% „integritatemque, cæterisque, qua“. Læst paa KB-Exemplaret og
% gjengivet som trykt.
Alle Planeternes Vidnesbyrd (\textit{tot tantaqve errantium syderum
testimonia}) stemme overeens i (\textit{consentiunt}), at Jorden
bevæger sig fremad i Rummet. Copernicus tillægger den en tredobbelt
Bevægelse: en \emph{daglig} omkring dens Axe, en \emph{aarlig} omkring
Solen igjennem Ekliptiken og --- til Overflødighed%
% Spatieringen af „daglig“ og „aarlig“ er bekræftet paa KB-Exemplaret;
% det følgende „til“ er derimod ikke spatieret (spacing.py's Udslag her
% er Justeringen).
\renewcommand{\thefootnote}{**)}%
\footnote{% trykt Mærke: **) --- anden af to Noter paa S. 136.
Hvad Rothmann allerede blev opmærksom paa.}%
\renewcommand{\thefootnote}{*)}
--- endnu en Bevægelse, hvorved Axens Hældning stedse bliver parallel
% Trykfeilbladet retter S. 136, Linie 14 f. o.: „Hældning --- læs
% Heldning“. Efter Husreglen staaer Texten som TRYKT: Hældning.
med sig selv (\textit{tertius declinationis motus}).

Vi see de gamle Philosopher gaae ud fra den Grundbetragtning, at
Legemer, som bevæge sig med lige Hastighed, maae synes at bevæge sig
langsommere, i Forhold til som de ere i større Afstand, at Maanen med
den mindste Bane desaarsag har den korteste Omløbstid, Saturn med den
største Bane den længste. Men om Venus og Mercur ere Meningerne deelte.
Enten man med Plato sætter dem ovenfor Solen eller med Ptolemæus
nedenfor, er Sagen lige for\-%
% --- p. 137 ---
virret. Som en Grund, hvorfor Venus og Mercur bør sættes nedenfor
Solen, har man anført, at der ellers vilde blive et altfor uhyre stort
Rum mellem Solen og Maanen. Det er en høist upaalidelig Grund; thi
netop de Afstandsmaal, hvorpaa Antagelsen skal begrundes, udvise, at
Venusepicyklen, naar Jorden skal være Deferentens Centrum, omgiver et
Rum saa stort, at dette tomme Rum kunde mere end rumme baade Jorden,
Luften, Ætheren, Maanen og Mercur. Et andet Argument er endnu
uheldigere. Ptolemæus anseer det for passende, at Solen faaer en
Mellemplads mellem de Planeter, som fjerne sig fra den paa alle Sider
(\textit{omnifariam}), og de Planeter, som ikke gjøre det; men hvor
falsk Beviset er, kan sees paa Maanen, som i dette Tilfælde maatte have
havt sin Plads ovenfor Solen. Altsaa, slutter Copernicus, enten er
Jorden ikke Centrum, eller den Orden, hvori man stiller Planeterne, er
uden Betydning.\footnote{% trykt Mærke: *) --- første af TO Noter paa S. 137.
Oportebit igitur, vel terram non esse centrum, ad qvod ordo syderum
orbinmqve referatur: aut certe rationem ordinis non esse, nec apparere
cur magis Saturno qvam Jovi seu alii cuivis superior debeatur locus.
Lib. I. cap. X.}
% Trykfeil i Noten: „orbinmqve“ for „orbiumqve“. n'et er utvetydigt paa
% KB-Exemplaret ved stærk Forstørrelse (Bue foroven, modsat u'et i
% „syderum“ paa samme Linie), saa Feilen er Sætterens og ikke
% Læsningens. Gjengivet som trykt; læs „orbiumqve“.

At forandre det geocentriske Standpunkt til et heliocentrisk er kun
muligt for den, der er blevet sig Forskjellen mellem Tilsyneladelse og
Virkelighed ret tydelig bevidst. Fra Betragterens virkelige Standpunkt
ere alle Bevægelser tilsyneladende, kun fra et uvirkeligt,
tilreflecteret Standpunkt ere Bevægelserne virkelige; det er en
Omvendthed, som kun en overlegen Reflexion har været i Stand til, paa
første Haand, at fatte og gjennemføre. Om enten Pythagoræeren
Philolaus%
\renewcommand{\thefootnote}{**)}%
\footnote{% trykt Mærke: **) --- anden af to Noter paa S. 137.
Jvnfr. De Revol. Lib. I. cap. V.}%
\renewcommand{\thefootnote}{*)}
eller Aristarch fra Samos have været af den Mening, at Jorden var en
Planet, gjør i denne Sammenhæng Intet til Sagen; thi det, hvorpaa det
kommer an, er en astronomisk-geometrisk Reflexion, som med behørig
Indsigt i
% --- p. 138 ---
de forviklede Relationers Mangfoldighed formaaer at føre
Tilsyneladelserne tilbage til Virkeligheden og saa igjen ud fra
Virkeligheden faae det Tilsyneladende fra alle Sider i alle
Enkeltheder nøiagtig forklaret og begrundet. Det er denne Opgave,
Copernicus i Forhold til sin Tidsalders astronomiske Viden fuldstændig
har løst. En opheiet Grundtanke har ledet ham, en Idee, hvorom hans
objectiverende Phantasie fra Først til Sidst har været enig med hans
Reflexion, Ideen af en „beundringsværdig Verdens Symmetrie“
(\textit{admiranda mundi symmetria}) og „en harmonisk Sammenhæng mellem
Bevægelserne og Banernes Størrelse“. Fixstjernehimlen, „det Høieste af
Alt“, er ubevægelig, thi den omslutter Alt; men Planeterne Saturn,
Jupiter, Mars, Jorden med sin Maane, Venus og Mercur bevæge sig alle
omkring det fælles Centrum, Solen. „Hvo kunde vel i dette skjønne
Tempel sætte denne Lampe paa et andet eller bedre Sted end der, hvorfra
den paa een Gang kan bestraale Alt. Nogle kalde den ret træffende
Verdens Lygte, Andre dens Sjæl, Andre dens Styrer, Trimegistus, en
synlig Gud, Sophocles's Alt skuende Elektra. Ja, saaledes sidder Solen
som paa en kongelig Throne, styrende den hele Familie af omkredsende
Himmellegemer.“\footnote{% trykt Mærke: *) --- eneste Note paa S. 138.
De Revol. Lib. I. cap. X.}

Copernicus har ikke beriget Astronomien med nogen ny
Iagttagelseskonst eller med Opdagelse af nye Enkeltheder; hans
Opdagelse tilhører ikke saameget Kjendsgjerningerne som Reflexionen;
det er en Objectiveringens Hemmelighed, han har opdaget; ved denne
Opdagelse har han forvandlet ikke Verden, men vor Bevidsthed om
Verden: han har for vor Viden omdannet Himmel og Jord.

% Indløbende Overskrift, læst paa Siden selv (KB-Exemplaret): halvfed,
% og her virkelig indløbende --- Texten fortsætter paa samme Linie.
% Overskriften staaer nederst paa S. 138, ikke paa S. 139.
\runin{c) Discussionen: tre Hypotheser.} Det Ptolemæiske System var
dogmatisk; Dogmatismen med sin naive Stolen
% S. 138 slutter midt i en Sætning: sidste Linie ender „med sin naive
% Stolen“, uden Orddeling; næste Batch (S. 139--146) fortsætter
% umiddelbart.
% --- p. 139 ---
% S. 139 begynder UDEN Indrykning: Linien staaer lige ud til Marginen og
% fortsætter Sætningen fra S. 138, der ender „... med sin naive Stolen“.
% Ingen ny Paragraf her. --- Overskriften „c) Discussionen: tre Hypotheser.“
% staaer paa S. 138, ikke her; den skrives af Batch 13.
paa den menneskelige Objectiverings objective Gyldighed er i god
Forstaaelse med Illusionen. Om Enkeltheder kunde der opkomme Tvivl, men
Grundlaget i det Hele stod fast, thi det hvilede paa umiddelbar
Anskuelse. Det Copernicanske System er thetisk i Reflexionens Form og
forsaavidt ogsaa dogmatisk, skjøndt paa en egen Maade. Vel blive
Sandsernes umiddelbare Vidnesbyrd i Et og Alt underordnede Forstandens
Prøvelse; men det Standpunkt, Reflexionen har valgt, maa indtil videre
holde sig paa Harmoniens Dogme. Om end Copernicus nu og da udtaler sig
med Varsomhed, som var det ham nok, naar Læseren blot fandt hans
Anskuelse sandsynlig eller dog ikke urimelig, saa er det dog af hele
Indtrykket klart, at han selv anseer den for uigjendrivelig; den
Verdensbygning, han har reist, er ham ikke en af flere mulige, men den
eneste mulige, fordi den er den eneste fornuftige.

Nu ligger det i Sagens Natur, at to hinanden modstridende Theorier, der
begge optræde med lige Adkomst, eftersom de begge, hver paa sin Viis,
omtrent i lige Grad tilfredsstille Observationerne, maae gjøre hinanden
problematiske. Men naar Grundlagene ere problematiske, saa maa enten
den ene af Theorierne slaaes fast som en Troesartikel, eller begge
blive staaende som Hypotheser.

% „charachteristisk“: saaledes trykt, læst paa KB-Exemplaret (ch...cht).
% Bogen skriver ellers „charakteristisk“ (S. 142, S. 145); Trykfeilbladet
% retter det ikke. Gjengivet som trykt.
Den ubenævnte Forfatter af Forordet til Værket \textit{De
Revolutionibus} har i denne Henseende udtalt sig paa en ret
charachteristisk Maade. Da det er umuligt, siger han, at opdage de
sande Aarsager, saa har Astronomen kun at udtænke og digte saadanne
Forudsætninger eller Hypotheser, som kunne give ham et beqvemt og
tilfredsstillende Grundlag for Beregningerne. Under flere Vendinger
bliver det da paa det Eftertrykkeligste stillet Læseren for Øie, at den
Copernicanske Theorie hverken er eller vil være Andet end en slet og
ret Hypothese. Det er sikkerlig mere end Klog\-%
% --- p. 140 ---
skabshensyn, der præger denne Fremstilling; det er ærlig
% Fodnotens Latin er sat med opreist Antiqua, IKKE med Cursiv --- kontrolleret
% paa KB-Exemplaret. Noten begynder og ender paa S. 140.
% Trykfeilbladet: S. 140, Linie 2 f. n.: „quiqvam --- læs quidqvam“.
% Efter Husreglen staaer Texten som trykt (quiqvam); Rettelsen anføres kun her.
% Ellipserne staae i Trykket som FIRE spatierede Punkter, saaledes gjengivne.
Overbeviisning.\footnote{% trykt Mærke: *)
Est enim Astronomi proprium, historiam motuum cælestium diligenti et
artificiosa observatione colligere. Deinde causas earundem, seu
hypotheses, cum veras assequi nulla ratione possit, qualesumque
excogitare et confingere, quibus suppositis, iidem motus ex Geometriæ
principiis, tam in futurum quam in præteritum recte possint calculari.\par
Neque enim necesse est, eas hypotheses esse veras, immo ne verisimiles
quidem, sed sufficit hoc unum, si calculum observationibus congruentem
exhibeant .\,.\,.\,. Sinamus igitur et has novas hypotheses, inter
veteres, nihilo verisimiliores, innotescere .\,.\,.\,. Neque quisquam,
quodad Hypotheses attinet, quiqvam certi ab Astronomia expectet, cum
ipsa nihil tale præstare queat.}

De to Hypotheser stode imidlertid paa en saa spændt Fod med hinanden,
at Tilhængerne af den ene viste sig uforsonlige mod Tilhængerne af den
anden; det kunde til Sindenes Beroligelse da være ret hensigtssvarende
at forsøge en Mægling ved at opstille en tredie Hypothese, en
Hypothese, som borttog det Forargelige ved den nye, og for en Deel
afhjalp Ufuldkommenhederne ved den gamle: denne tredie Hypothese blev
opstillet af Tycho Brahe. Illusionens gamle Standpunkt hævdes,
forsaavidt som Jorden vedbliver at være det ubevægelige Midtpunkt: om
Jorden bevæge sig Maanen og Solen; men Solen --- det er den
Indrømmelse, der gjøres den nye Theorie --- er Centrum for samtlige
% „Ændring“: tesseract læser „Æudring“ (den kjendte n/u-Forvexling);
% ordet er læst paa KB-Exemplaret.
Planeter. Man kan fra Astronomiens nuværende Standpunkt opfatte den
Tychoniske Hypothese som en Ændring af den Copernicanske, fremkommen
ved, at Coordinataxernes Skæringspunkt henlægges i Jordens Centrum; men
man kan med lige saa megen Grund betragte den som et specielt Tilfælde
af den Ptolemæiske, det nemlig, som erholdes ved at tage Solbanen til
fælles defererende Cirkel for alle Planeterne. Ifølge denne Hypothese
kjender man da Forholdet mellem
% --- p. 141 ---
samtlige Planeters Middelafstande, men Afstandens Værdi i
Solbaneradier er, ligesom Forholdet mellem Maanens og Solens Afstand,
endnu ubekjendt.

% „Discussion“ er spærret --- bekræftet paa KB-Exemplaret.
Med de tre Hypotheser er Verdensbygningens Objectivitet gaaet tabt;
Virkeligheden selv er blevet borte for Bevidstheden. Men at den
virkelige Sammenhæng mellem Jord og Himmel ikke skulde lade sig
erkjende, er en utaalelig Forestilling, lige unaturlig for den
videnskabelige Forskning og den naive Betragtning. Har den
reflecterende Forstand i et sindrigt, combineret Spil af Hypotheser sat
sin oprindelige Medgift, Visheden om Erkjendelsens første umiddelbare
Realitet, overstyr, da finder den heller ingen Fred, før den paa Ny er
blevet mættet med virkeligt Indhold og har erobret den tabte
Objectivitet tilbage. Tvivlen og Famlen ledte til omfattende
\emph{Discussion}; modsigende Betragtninger krydsede hverandre;
Synspunkterne, hvorfra de udgik, ere betegnende for Datidens
videnskabelige Tilstand.

% Første af de tre Synspunkter. Spærringen omfatter paa denne Side KUN
% „kirkelig dogmatiske“ --- „Det“ og „Synspunkt“ ere sat med almindelig
% Antiqua. Kontrolleret paa KB-Exemplaret ved at sammenholde „Synspunkt“
% med „Synspunkterne“ to Linier ovenfor. Paa S. 142 omfatter den
% tilsvarende Spærring derimod ogsaa „Synspunkt“; Uligheden er Sætterens
% og gjengives som trykt. Ikke en Overskrift: der staaer intet a)/b)/c)
% og intet græsk Bogstav foran.
Det \emph{kirkelig dogmatiske} Synspunkt var i Grunden subjectivt, men
tog, som sædvanligt, feil af sig selv, og bildte sig ind, at det i al
Ufeilbarhed var objectivt. Det varede nogen Tid, inden Theologerne ret
bleve opmærksomme paa den Fare, hvormed den Copernicanske Theorie
truede saavel den protestantiske som den catholske Rettroenhed; men
neppe havde de faaet Øinene op, før de ud fra begge Kirker bøde Faren
Trods og gik i Kampen med Dødsforagt; Historien giver dem det
Vidnesbyrd, at de i denne „for Saligheden vigtige“ Sag
samvittighedsfuldt have benyttet enhver Anledning, hvorved de paa
passende Maade kunde komme til at prostituere sig ret videnskabeligt.
% Det græske staaer med Cursiv og med Akut paa ό; læst paa KB-Exemplaret
% (ματαιόλογοι, ikke „avaióhoyo“, som Google-Laget giver).
% Citatet „tomme Ordgydere ... detortum)“ løber over Bladskiftet og
% sluttes først paa S. 142. Det er ikke en Trykfeil og faaer ingen Note.
Hvad Copernicus i Dedicationen til Pave Paul III forudsaae, at „tomme
Ordgydere (\textit{ματαιόλογοι}), blottede for al mathematisk Viden,
vilde anmasse sig en Dom over Værket ved med Villie at fordreie et
eller andet Sted i den hellige
% --- p. 142 ---
Skrift (\textit{propter aliqvem locum scripturæ male ad suum propositum
detortum})“, gik rigelig i Opfyldelse.

% Andet Synspunkt. Her omfatter Spærringen ogsaa „Synspunkt“ ---
% bekræftet paa KB-Exemplaret; smlg. Noten til S. 141.
% „spiritus“ er ligeledes spærret og sat med opreist Antiqua, ikke Cursiv;
% bekræftet paa KB-Exemplaret. Paa Arbeidsscanningen er Ordet tillige
% understreget med Blyant af en tidligere Læser --- Blyanten findes ikke
% paa KB-Exemplaret og er ikke Text.
Det \emph{empirisk philosophiske Synspunkt}, repræsenteret af Baco,
egnede sig unegtelig bedre til en besindig, paa Inductionsmethoden
begrundet Drøftelse af de stridige Hypotheser; men, mærkeligt nok,
havde Baco, i aabenbar Modsigelse med sin egen Theorie om Inductionen,
sat sig saa grundig fast i en vis Opfattelse af Naturens
Grundqvaliteter og Legemernes iboende \emph{spiritus}, at han ikke
alene holdt paa den ubevægelige Jord uden dog enten at slaae sig til Ro
hos Ptolemæus eller antage Tycho Brahes Ændring, men ikke engang var i
Stand til at gaae ind paa exacte Undersøgelser. Han seer i Copernicus
en Mand, som optager Indfald af enhver Art og indfører dem i Naturen,
naar de blot paa nogen Maade lade sig bringe i Overeensstemmelse med
hans Beregninger. Ved en saadan Bedømmelse har Baco imidlertid ikke
charakteriseret Copernicus, men hvad de astronomiske Problemer angaaer,
% „Thema Coeli“ staaer med Cursiv OG spærret (T h e m a  C o e l i);
% bekræftet paa KB-Exemplaret. Denne Udgave gjengiver Cursiven og noterer
% Spærringen her, da \emph inden i \textit vilde slaae Skriften opreist.
snarere sig selv; thi en Paastand i \textit{Thema Coeli} som den, at
Planeternes saakaldte retrograde Bevægelse fra Øst til Vest i det Hele
slet ikke existerer, men er et tomt Skin, som kun opstaaer derved, at
den faste Himmelhvælving rykker mere frem paa den ene Side, saa at
Planeten paa den anden Side synes at blive tilbage, maa dog vel regnes
til Indfald af den Art, som ingen Astronom finder det Umagen værd at
% De fire sidste Linier paa Arbeidsscanningen bære en tidligere Læsers
% Blyantstreger under „Ved kjækt at bryde ... et vældigt Stød fremad“.
% KB-Exemplaret har intet af det, og Skriften er der almindelig Antiqua:
% ingen Fremhævelse. spacing.py's RUN paa „henvise til“, „Baco sin“ og
% „mellem den“ er forkastet af samme Grund.
gjendrive. Ved kjækt at bryde med det scholastiske Formelvæsen, blotte
Syllogistikens Ufrugtbarhed, hævde Inductionen og henvise til
Erfaringen gav Baco sin Tidsalders Tænkning et vældigt Stød fremad. Men
Grundforholdet mellem den empiriske og den philosophiske Methode har
han ikke formaaet at gjennemskue. Til en indtrængende Analyse af den
empiriske og den aprioriske Tænknings eiendommelige Vexelvirkning
mangler han paa den ene Side Indsigt i det Mathematisk-Geometriske, paa
den anden Side Blik for det
% --- p. 143 ---
% Siden bærer kun 7 Linier Text; Resten er Fodnoten, som begynder og
% ender her.
Kritisk-Dialektiske. Han har i Aandrighedens Form de almindelige
Sætninger; men ude af Stand til at lade Ideen selv gjennemtrængende
opløse og stadfæste den concrete Virkeligheds eiendommelige
Kjendsgjerninger bliver han staaende ved det blot Aandrige og naaer
ikke Aanden: den Baconiske Empirisme er kun en noget forbedret Udgave
% „o. s. v.;“ --- Semikolon, læst paa KB-Exemplaret (tesseract giver „v.7“).
% De to Tankestreger om „saa lidt som Engellænderne endnu den Dag idag“
% ere begge lange Tankestreger paa KB-Exemplaret; Arbeidsscanningens
% kortere første Streg er en Tærskelvirkning.
af den Aristoteliske.\footnote{% trykt Mærke: *)
I heftig Polemik mod den Copernicanske Hypothese anfører Baco: at
Copernicus betynger Jorden med en tredobbelt Bevægelse og løsriver
Solen fra Planeterne, med hvilke den dog er saa nær beslægtet; at hans
System indfører for meget Ubevægeligt i Naturen, idet det berøver Solen
og Fixstjernerne, altsaa netop de Legemer, der lyse og straale
allerklarest, Bevægelse o. s. v.; disse Argumenter ere i
naturvidenskabelig Værdi aldeles Aristoteliske. Ogsaa den i Sandhed
betydningsfulde Grundtanke, at Formen er Sagens inderste Væsen,
\textit{ipsissima res, natura naturans, fons emanationis}, er
Aristotelisk. Vel gjør Baco det i Legemerne Formgivende, Principet for
Bevægelse og Forandring, den hvert Legeme iboende Spiritus, til en
materiel Substans og vil paa ingen Maade have dette materielle
Formvæsen forvexlet med en Aristotelisk Energie, men ved denne uklare
Opfattelse af Legemernes Natur er han saa langt fra at hæve sig i
Begrebet op over Aristoteles, at han tvertimod synker ned til en Slags
aandrig Leflen med Magie og Troldom. At Baco med al sin Ivren for
Eenhed af Erfaring og Tænkning ikke har været i Stand til at forsone
den empirisk forskende Naturlære med den egenlige Naturphilosophie, er
forklarligt deraf, at han ikke for sin Tid --- saa lidt som
Engellænderne endnu den Dag idag --- har været i Stand til at blive sig
Modsætningen mellem de to qvalitativt forskjellige Videnskaber bestemt
og tydelig bevidst. Hos Baco er Interessen for det Empiriske
overveiende philosophisk; derfor har han ingen Sands for Mechanik og
seer fornemt ned paa Mathematiken; hos Newton hviler Interessen
afgjørende paa det Mathematisk-Mechaniske, men desuagtet anseer Newton
sine physisk exacte Principer for naturphilosophiske.} De, der i
Erfaringsphilosophiens Navn
% --- p. 144 ---
endnu ansee Baco for Non-plus-ultra, skulle ved nærmere Prøvelse
visselig gjøre den Opdagelse, at der, for Principet selv, i hans
Philosophie er for lidt Erfaring, og i hans Erfaring for lidt
Philosophie.

% Tredie Synspunkt. „exact empiriske“ er spærret --- begge Ord, over
% Linieskiftet („det e x a c t  e m p i- / r i s k e.“) --- bekræftet paa
% KB-Exemplaret. Samme Udtryk to Gange paa S. 145 er derimod IKKE spærret.
% „Gilbert og Galilei“ ser fremhævet ud paa Arbeidsscanningen, men det er
% en tidligere Læsers Blyantstreg: KB-Exemplaret har almindelig Antiqua.
% Kaldet forkastet. Ligesaa Blyantklammerne om „Det er ikke det endelige
% Resultat ... kommer an.“, som tesseract gjengiver som „[“ og „]“.
Saa bliver der altsaa kun et tredie Synspunkt tilbage, som i denne
Discussion kunde faae og da ogsaa virkelig fik en sand videnskabelig
Betydning, nemlig det \emph{exact empiriske}. De store Repræsentanter
ere Gilbert og Galilei. Det er ikke det endelige Resultat, men Veien
dertil, ikke Dommen om en Hypothese, hvorvidt den bør antages,
forkastes eller betvivles, men Dommens Begrundelse, hvorpaa det ved
Discussionen først og sidst kommer an. Hvad Tvivl og Betænkeligheder
angaaer, staaer Gilbert nær ved Baco; men i begrundende Reflexion ere
de vidt adskilte. Gilbert speculerer ikke, han forsker; naar han ikke
kan overbevise sig om, at Jorden har en aarlig Bevægelse, men vakler
bestandig mellem Tycho og Copernicus, da er det ikke Ideerne, men
Kjendsgjerningerne, der volde ham Bryderier. I sit Værk De Magnete
% Titlen „De Magnete“ staaer her med almindelig Antiqua, ikke med Cursiv;
% de latinske Indskud i Parenthes derimod med Cursiv. Saaledes trykt.
% Noten begynder og ender paa S. 144.
antager han vel, at Jorden dreier sig om sin Axe,\footnote{% trykt Mærke: *)
Guilielmi Gilberti Colcestrensis, Medici Londinensis. De Magnete
Magneticisque corporibus et de Magno magnete tellure, Physiologia nova.
Londini MDC. --- De terrestris globi diurna revolutione magnetica. Lib.
VI. cap. III. Terram circulariter moveri. cap. IV. Terræ motum
negantium rationes et earum confutatio. cap. V.} men sætter Læren
herom i uklar Forbindelse med Læren om Magneten og den elektriske
Tiltrækning (\textit{Globus telluris per se electrice congregatur et
cohæret}), idet han antager, at den elektriske Stræben gaaer ud paa en
bindende Sammenhobning af Materien (\textit{motus electricus est motus
coacervationis materiæ}). Skjøndt hans Beviser i denne Henseende ere
utilfredsstillende, saa er der dog Noget, som viser, at han staaer paa
et naturvidenskabeligt Grundlag i ganske
% --- p. 145 ---
% Denne Side bærer TRE Fodnoter, mærkede *), **) og ***) ved Foden.
% „Ingenium“ er spærret --- bekræftet paa KB-Exemplaret; opreist Antiqua,
% ikke Cursiv.
anden Forstand end Baco. Han afgjør ikke ved nogen speculativ
Forklaring, hvad Naturen kan og hvad den ikke kan, hvad der passer sig
for dens \emph{Ingenium}, og hvad der ikke passer
% TRYKFEIL i Noten: „hnmanum“ for „humanum“ --- et omvendt u. Bogstavet er
% utvetydigt et n paa KB-Exemplaret (smlg. u'erne i „inutilia“ i samme
% Linie), og begge Exemplarer vise det, altsaa Sætterens Feil. Ligeledes
% staaer der „docenque“ for „docentque“. Gjengivet som trykt.
sig:\footnote{% trykt Mærke: *)
Theologi etiam curiosi, mysteria divina ultra hnmanum sensum posita,
per magnetem et succinum illustrant, ut vani metaphysici, cum inutilia
phantasmata fundunt docenque, magnetem habent tanqvam Delphicum
gladium, exemplum semper ad omnia accommodandum. De Magnete. Lib. II
cap. II.} han tager Kjendsgjerningerne som de vise sig for ham; han kan
tage feil i Bestemmelsen af Præmisserne, men deri tager han ikke feil,
at Begrundelsen af physiske Phænomener udelukkende maa afledes af
% Latinen i denne Note veksler mellem qv og qu i samme Ord:
% „æqvinoctiorum ... obliqvitatis“ mod „inæquali ... obliquitatis“, og
% „Qvo magis“, „qvam“ mod „Quare“, „neque“, „aliquas“. Alt læst paa
% KB-Exemplaret og gjengivet som trykt.
physisk bestemte Data;%
\renewcommand{\thefootnote}{**)}%
\footnote{% trykt Mærke: **)
Det sidste Kapitel --- Lib. VI cap. IX. --- De præcessionis
æqvinoctiorum et obliqvitatis Zodiaci anomalia sluttes paa en ret
charakteristisk Maade med følgende Ord: Qvo magis hæc omnia de inæquali
motu tam præcessionis qvam obliquitatis incerta et incognita sunt.
Quare neque nos illius causas aliquas naturales proferre et certo
statuere possumus. Quare etiam et nos magneticis rationibus et
experimentis hic finem et periodum imponimus.}%
\renewcommand{\thefootnote}{*)}
han fastholder selv i sine Feilgreb den exact empiriske Methode.

% TRYKFEIL: Henvisningsmærket i Texten ved „Dialog“ er trykt **), medens
% Noten ved Foden er mærket ***). Begge Exemplarer vise det, altsaa
% Sætterens Feil. Gjengivet som trykt: Mærket i Texten staaer **), Noten
% ***), ved \footnotemark/\footnotetext.
De betydeligste Indlæg i den store Proces leveredes imidlertid af
Galilei. Ikke nok, at han i sin Dialog%
\renewcommand{\thefootnote}{**)}%
\footnotemark%
\renewcommand{\thefootnote}{***)}%
\footnotetext{% trykt Mærke ved Foden: ***)
% „gjornate“ for „giornate“ og „Firense“ for „Firenze“ staae saaledes i
% Trykket; læst paa KB-Exemplaret. Gjengivet som trykt.
Dialogo di Galileo Galilei, dove ne' Congressi di quattro gjornate si
discorre de' due massimi sistemi Tolemaico e Copernicano. Firense per
Landini 1632.}%
\renewcommand{\thefootnote}{*)}
lader En af de Samtalende føre et Forsvar for den Copernicanske
Hypothese, der siden paadrog ham Inqvisitionens Forfølgelse; han sigter
høiere, han væbner Øiet og udvider Synskredsen; med den Kikkert, han
selv construerede, gjør han forbausende Opdagelser i Himmelrummet. Med
Galilei holder den exact empiriske Methode sit triumpherende Indtog i
% Arkesignaturen „10“ staaer nederst paa S. 145 under Noterne; den er
% ikke Text og gjengives ikke.
Viden\-%
% --- p. 146 ---
% Siden bærer TO Fodnoter, mærkede *) og **); begge begynde og ende her.
skabens Historie.\footnote{% trykt Mærke: *)
Det er naturligviis om Galilei som Repræsentant for sin Tidsalder, ikke
om Galilei som enkeltstaaende Forsker, at dette kan siges. At Kikkerten
først blev opfundet i Holland, og at det var Efterretningen om den i
Holland gjorte Opfindelse, der foranledigede Galilei til at construere
sin Kikkert, er bekjendt.} Af alle Sandser er Synet den mest objective;
med det udvidede Syn faaer den sig udvidende Tænkning et udvidet
Grundlag; saa indhentes Virkeligheden, og den bortflygtede Objectivitet
kaldes igjen tilbage. Maanebjergene, Venusphaserne og fremfor Alt
Jupitersmaanerne,%
\renewcommand{\thefootnote}{**)}%
\footnote{% trykt Mærke: **)
„Jupiters Maaner, de første af alle ved Kikkerten fundne Biplaneter,
bleve, som det synes, næsten samtidig opdagede af Simon Marius i
Ansbach den 29. December 1609 og af Galilei i Padua den 7de Januar
1610, uden at disse Mænd stode i mindste Forbindelse med hinanden. Med
Hensyn til Publicationen derimod af denne Opdagelse kom Galilei ved sin
Nuntius Sidereus (1610) Simon Marius's Mundus Jovialis (1614) i
Forkjøbet.“ A. v. Humboldt. Kosmos, oversat af C. A. Schumacher. Andet
Bind. Kbhvn. 1849. S. 289--90.}%
\renewcommand{\thefootnote}{*)}
høre til en Cyclus af Kjendsgjerninger, som ved det Nye og
Overraskende, de frembøde for det skjærpede Syn, nødvendig maatte gjøre
et stærkt og storartet Indtryk paa Samtiden, rokke Tilliden til den
gamle Forestillingsmaade, aabne Øiet for den Copernicanske
Verdensbygning og vække tankerige Anelser om en hidtil ukjendt Himlens
Mechanik.

% AFDELING A SLUTTER HER, paa S. 146, og ikke først paa S. 147.
% Under den sidste Textlinie staaer en kort, CENTRERET Streg (kontrolleret
% paa begge Exemplarer; den lavere, venstrestillede Streg længere nede er
% Fodnotestregen). S. 147 aabner derefter med Overskriften
% „B. Solsystemet og Universet.“, som skrives af næste Batch (S. 147--155).
% Den krøllede Parenthes, tesseract læser efter „Mechanik.“, er en tidligere
% Læsers Blyantklamme, ikke Text.
\streg

%  --- B. Solsystemet og Universet  (pp. 147–171) ---

%  MATHEMATICAL — work from the image; the OCR is not a witness here.
% --- p. 147 ---
% Afdeling B aabner paa selve S. 147: „B.“ centreret paa egen Linie, derunder
% „Solsystemet og Universet.“ i halvfed (kontrolleret paa KB-Exemplaret).
% Stregen ved Slutningen af Afdeling A staaer paa S. 146 og er allerede skrevet.
\afdeling{B.}{Solsystemet og Universet.}

Saavel paa Illusionens som paa Reflexionens Trin viser der sig et eget
Forhold mellem Planetsphærerne og Fixstjernehimlen; men det er umuligt
at danne sig et fornuftigt Begreb om Verdensaltet, saalænge
Forestillingen om Solsystemet endnu er uklar. Den nærmeste Opgave
bliver altsaa: ved skærpet Iagttagelse og tilsvarende Beregning at
tilveiebringe den nøiagtigst mulige Overeensstemmelse mellem Banerne og
Lovene og derved sikkre sig \emph{Objectiviteten af det formelle
Grundlag}. Dernæst, naar Bevægelsesphænomenerne ere saavidt opklarede,
at de vise sig som lovbestemte Virkninger af bestemte Aarsager, kunne
Aarsagerne efterforskes, og Kjendsgjerningernes causale Sammenhæng lede
til en Forvisning om \emph{Objectiviteten af det reelle Grundlag}.
Endelig, naar Grundlaget, det reelle saa vel som det formelle, alsidig
er sikkret, ville Methodens to Elementer: Iagttagelseskonsten og den
mathematiske Analyse, gjensidig drive hinanden fremad, og
Objectiveringen af den mechaniske Verdensbygning saavel i Enkelthederne
som i de store Omrids fuldende Billedet af \emph{et udvortes objectivt,
i Endelighed uendeligt Universum}.

% Løbende Overskrift, trykt HALVFED paa samme Linie som Afsnittet; Spærringen
% i de tre Fremhævninger ovenfor er derimod ægte Sperrsatz — begge Dele
% kontrollerede paa KB-Exemplaret.
\runin{a) Objectiviteten af det formelle Grundlag: de Keplerske Love.}
En væsenlig Hindring for en virkelig og sand Indsigt i Solsystemets Plan
og Bygning stod der endnu tilbage at fjerne; det var Epicykeltheorien.
Hvor resolut selvstændig Copernicus end var, forsaavidt det gjaldt om at
gjøre Solen til Centrum, saa var han dog endnu afhængig af den gamle
Forudsætning om Cirkelbanernes Nødvendighed. „Vi maae“, siger
han,\footnote{% trykt Mærke: *)
De Revol. Lib. I Cap. IV.} „antage, at de himmelske Bevæ\-%
% Arkesignaturen „10*“ staaer nederst paa S. 147; den er ikke Text.
% Citatet „antage \dots\ i Cirkler.“ løber over Bladskiftet; derfor er
% Anførselsregnskabet for S. 147 +1 og for S. 148 tilsvarende ÷1.
% --- p. 148 ---
gelser ere cirkelformige, thi deres Uligheder følge en bestemt Lov og
vende tilbage i uforanderlige Perioder, hvad de kun kunne gjøre, naar de
foregaae i Cirkler.“ Vi skulle fra Naturphilosophiens Synspunkt nu see,
at Epicykelmethoden, ifølge Sagens Natur, maa gjøre en endelig Løsning
af det store Objectiveringsproblem umulig, idet vi særlig mærke os,
hvori det for Objectiveringen Forvirrende egenlig bestaaer.

Det Blændende og Forvirrende kan ikke sættes deri, at Iagttageren
fristes til, at tage de mange Cirkler, store og smaa, for ligesaa mange
virkelige Baner og tillægge enhver af dem en materiel Tilværelse; thi
det Urimelige i en saadan Forestillingsmaade er for iøinefaldende til at
kunne blænde den Kyndige. „Om end dogmatiske Philosopher, som
Aristoteles, have anseet de himmelske Sphærer for virkelig bestaaende,
solide Legemer, saa taler dog Ptolemæus kun om dem, som om blot
imaginære Ting, og allerede af det Beviis, han giver for Identiteten af
en Epicykel med en excentrisk Cirkel, følger klart, at han ikke har
holdt Epicyklen for Andet end en geometrisk Conception, hvorved han
forsøgte at fremstille Himlens tilsyneladende Bevægelser
overeensstemmende med Iagttagelserne.“ (Whewel). Middelalderens og
Reformationstidens mystiske Drømmerier om faste Himle og Himmelsphærer
bleve af Astronomerne paa det Skarpeste adskilte fra den mathematisk
geometriske Epicykeltheorie.\footnote{% trykt Mærke: *)
% Denne Note begynder paa S. 148 og ender paa S. 149; Mærket gjentages
% ikke ved Fortsættelsen. Hele Noten staaer her ved sit Mærke.
\dots\ „i Midten af det sextende Aarhundrede, da den lærde Polyhistor
Girolamo Fracastoro's Theorie om de 77 homocentriske Sphærer fandt saa
meget Bifald, og da senere Copernicus's Modstandere opsøgte alle Midler
til Opretholdelsen af det Ptolemæiske System, var den især af
Kirkefædrene begunstigede Forestilling om Tilværelsen af solide Sphærer,
Cirkler og Epicykler endnu meget udbredt. Tycho Brahe roser sig
udtrykkelig af den Fortjeneste, ved sine Betragtninger over
Cometbanerne at være den Første, som har beviist Umuligheden af solide
Sphærer og sønderknust deres konstige Stillads.“ (A. Humboldt).}

% --- p. 149 ---
Det for Objectiveringen Vildledende kan da heller ikke søges deri, at
Epicykeltheorien ved at stable Cirkel paa Cirkel tilsidst maatte blive
for indviklet; thi for det Første lode en Deel af de overflødige
Hypotheser sig jo efterhaanden reducere. Man havde begyndt med at lægge
de defererende Cirkler og første Epicykler i Eklipticas Plan, og maatte
desaarsag forklare Brederne ved at indføre en Bevægelse paa smaa
Cirkler, der vare vinkelrette paa Ekliptica. Copernicus hjalp sig
imidlertid ved at antage visse Librationer eller Svingninger af
Epicyklerne; og Tycho Brahe undgik de vinkelrette Cirkler ved at give
Planetbanerne en lille Heldning mod Ekliptiken. Dernæst, og dette er
det Afgjørende, hængte de voxende Complicationer paa det Nøieste sammen
med Iagttagelsernes Mangfoldiggjørelse og de mathematiske Methoders
daværende Ufuldkommenhed. Skulde det Indviklede\footnote{% trykt Mærke: *)
% S. 149 bærer TO Noter i Satsen: Slutningen af Noten fra S. 148 (uden
% Mærke) og derefter denne, mærket *).
Naar det Indviklede fremkommer ved Ubehjælpsomhed i at stille og
behandle en Opgave, saa vil en forbedret Methode være i Stand til endog
i en forbausende Grad at hæve alle Vanskeligheder. Saaledes behandler nu
den høiere Mathematik med største Lethed en Mængde Opgaver, hvis Løsning
er vanskelig, vidtløftig, ja umulig, naar den skal skee alene ved de
Hjælpemidler, den elementære Mathematik tilbyder. En anden Sag er det,
naar det Indviklede ligger i Opgavens egen Natur; thi i dette Tilfælde
vil Opgavernes tiltagende Forvikling staae i ligefremt Forhold til
Videnskabens Fremskridt. Indvikling forudsætter Udvikling. I denne
Forstand maa det siges, at Opgaverne i den nyere Astronomie ere langt
mere indviklede end i den ældre.} i og for sig være vildledende, da
maatte den nyere Analyses uendelige Rækker være ganske anderledes
vildledende, thi i Sammenligning med disse kunne Epicykeltheoriens
Complicationer vistnok siges at være meget enfoldige.

% --- p. 150 ---
Nei, det for den objectiverende Bevidsthed Vildledende ligger i, at de
alternerende Hypotheser gjøre en objectiv Skjelnen mellem virkelige og
fingerede Baner umulig. At Solbanen, abstract geometrisk opfattet, maa
blive den samme, enten man antager, at Solen med jævn Hastighed
beskriver en Cirkel, hvis Centrum ligger udenfor Jorden, eller man
tillægger Solen en jævn Bevægelse paa en Cirkel, hvis Centrum bevæger
sig jævnt paa en Cirkel, der har Jorden til Centrum, er
uimodsigeligt.\footnote{% trykt Mærke: *)
% Den ubestemte Størrelse skrives i Trykken y med hævet 1 (et Tal, ikke
% en Mærkeprim: Cifret har baade Fane og Fodserif og er identisk med
% 1-Tallet i „1+e cos v“; kontrolleret paa KB-Exemplaret).  De ellers
% isolerede Bogstaver i den løbende Text sættes her ligeledes i Math,
% for at Skriftbilledet ikke skal skifte midt i Opregningen y, y^1
% (Præcedens: Lagrange-Noten S. 59--61).
% MINUS er her et ægte Minus, ikke Nielsens ÷; læst af Siden.
% TRYKFEIL: anden Gang staaer der „detererende“ for „defererende“ —
% samme Ord staaer korrekt to Linier ovenfor. Begge Exemplarer have det;
% Trykfeilsbladet nævner det ikke. Gjengivet som trykt.
Lægges, ifølge den sidste Hypothese, Coordinataxernes Skjæringspunkt i
Jordens Centrum, medens Ordinataxen falder sammen med Apsidelinien, og
betegnes Ordinaten for den defererende Cirkel ved $y$, for Solen paa
Epicyklen ved $y^{1}$; saa er $x^{2} + y^{2} = r^{2}$ Ligningen for den
detererende Cirkel, og $y = y^{1} - \varrho$; herved erholdes:
$x^{2} + (y^{1} - \varrho)^{2} = r^{2}$, som netop er Ligningen for den
excentriske Cirkel, der svarer til første Hypothese, og Excentriciteten
$= \varrho$.} Men uagtet det mathematiske Resultat er det Samme, saa er
Sagens mechaniske Sammenhæng dog høist forskjellig. Ifølge den første
Hypothese er her i objectiv Forstand kun een Bane, ifølge den sidste er
her derimod to. I første Tilfælde vil man indtil videre kunne forklare
Banen ved at henvise til den alle Himmellegemer iboende Natur at bevæge
sig i Cirkler; den eneste metaphysiske Vanskelighed, der bliver tilbage,
maa da dreie sig om Grunden til, at Jorden ikke er blevet Solbanens
virkelige Centrum. I det andet Tilfælde er det egenlig ikke Solen, men
Solepicyklen, der bevæger sig om Jorden. Her maa da i objectiv Forstand
være en dobbelt Grund til den dobbelte Bevægelse, en særegen Grund til
Epicykelcentrets Bevægelse om Jordens Centrum, og ved Siden heraf en
særegen Grund til Epicyklens Bevægelse med Solen eller, hvad der for
Begrundelsen giver en ny Opgave, til Solens Bevægelse om
Epicykelcentret. Kunne to alternerende Hypotheser, hvor
% --- p. 151 ---
simple de end, hver for sig, ere, allerede give Sagens Sammenhæng et saa
subjectivt og tvivlsomt Udseende, da følger det af sig selv, at man ad
denne Vei med hvert Skridt kommer bestandig længere bort fra det
Objective. Ved Marsbanen f. Ex. er den mechaniske Sammenhæng og dens
Begrundelse endnu mere forviklet, da der maa tages Hensyn til den mellem
Jorden og Mars beliggende Solbane. Lader man Mars bevæge sig paa en
Epicykel, hvis Centrum bevæger sig jævnt med en Omløbstid af 687 Dage
paa en defererende Cirkel om Jorden saaledes, at Mars's Sted paa
Epicyklen er bestemt ved Solens Sted, idet Epicykelradien til Planeten
stedse er parallel med Forbindelseslinien mellem Jorden og Solen: da er
det Objective i Marsbanen for sig allerede tvivlsomt og bliver endnu
mere tvivlsomt derved, at Solbanens objective Beskaffenhed er tvivlsom.
Til Bevægelse i Cirkelbaner kunde den speculative Fornuft dog tænke sig
en Grund; men til excentriske Baner og Bevægelser af Epicykelcentrer har
end ikke den sindrigste Speculation været i Stand til at udfinde Noget,
der kunde ligne en Naturgrund. Der kan tænkes Grunde uden Aarsager, men
der kan ikke tænkes Aarsager uden Grunde; ere Phænomenerne saa
forviklede, at man ikke kan skimte en bestemmende Grund, hvor skulde det
da være muligt at opdage en virkende Aarsag? Skulle Bevægelsens Love
erkjendes, da maa Epicykeltheorien bort: her er det Vendepunkt, hvor
Keplers Genie bryder igjennem.

% Det latinske Indskud er trykt KURSIVT, ikke spærret; Udraabstegnet med.
At følge Opdageren af de store Verdenslove paa hans tornede Vei med de
vildsomme Bugtninger --- \textit{per errores ad veritatem, per ardua ad
astra!} er ikke Naturphilosophiens Opgave; en Antydning af Bevægelsen
maa være nok til en Fremstilling af Objectiveringens Gang.

I Aaret 1600 sluttede Kepler sig til Tycho Brahe, hvem han i Prag fandt
ivrig beskjæftiget med en Theorie over Planeten Mars, svarende til de
Observationer, han selv
% --- p. 152 ---
havde anstillet. Det var disse nye og forbedrede Observationer, der
ledte Kepler ind paa den nye Bane og satte ham i Stand til ni Aar efter
at udgive sit epochegjørende Værk:
% Titlen er trykt KURSIVT OG SPÆRRET paa eengang (samme Tilfælde som
% „Thema Coeli“ S. 142): \textit, og Spærringen kun noteret her, da \emph
% inde i Kursiv vilde vende Skriften opret igjen. Punktummet er sat
% udenfor Kursiven, da et Punktums Skriftsnit ikke kan iagttages.
% Understregningen under Titlen i Arbeidsscanningen er en tidligere
% Læsers Blyant; KB-Exemplaret har den ikke.
\textit{De motibus stellæ Martis}. Havde det nu ved umiddelbare
Iagttagelser været muligt at bestemme et vilkaarligt Antal Steder i
Marsbanen omkring Solen, da vilde den Opgave at udfinde Banens
Beskaffenhed derved unegtelig været betydelig simplificeret; men
Problemet var i den Grad indviklet, at Relationerne i det Objective
maatte gjennemgaae en Mængde Prøvelser i subjectiv Reflexion, før
Løsningen kunde opnaaes. Vexelbestemmelserne af det Bestemte og
Ubestemte med det Bekjendte og det Ubekjendte ere mangfoldige. Vel maa
der i objectiv Forstand være et System af Punkter beliggende i Mars's
heliocentriske Bane; men disse Punkter maatte dog først bestemmes i
Relation til Jordbanen. Først maatte Omløbstiden findes ved Hjælp af
Oppositionerne\footnote{% trykt Mærke: *)
% S. 152 bærer TO Noter, mærkede *) og **); begge begynde og ende her.
I Oppositionens Øieblik er Planetens heliocentriske Længde den samme som
den geocentriske; Planetens og Solens geocentriske Længde differere 180
Grader fra hinanden. Den Tychoniske Observationsrække over Mars omfatter
et Tidsrum af 20 Aar og angiver 12 Oppositioner.}; dernæst maatte eet og
samme Sted, paa hvilket Mars to Gange havde været i sin Bane%
\renewcommand{\thefootnote}{**)}%
\footnote{% trykt Mærke: **)
Vælger man iblandt de geocentrisk angulære Steder, som en
Observationsrække kan give, tvende, som netop adskille sig ved et
Tidsmellemrum af 687 eller nøiagtigere af 686,98 Dage, da veed man, at
Mars begge Gange har været i det samme Punkt af sin Bane.}%
\renewcommand{\thefootnote}{*)}%
, sees i Relation til de to tilsvarende Steder for Jorden i dens Bane:
kun ved to Retningslinier fra to bekjendte Standpunkter imod samme Punkt
blev det muligt at udfinde Planetens Afstand fra Solen og bestemme det
fuldstændig heliocentriske Sted i Marsbanen. Først naar eet Punkt var
fundet, kunde Bestemmelsen gjentages for et vilkaarligt Antal af
Punkter.
% --- p. 153 ---
Da nu Kepler med disse Punkter ad Beregningens Vei gav sig i Færd med at
fastsætte Liniens Natur, fandt han, at Marsbanen ikke, som man hidtil
havde antaget, kunde være nogen Cirkel, men at den maatte have en oval
Form. Ved at antage Linien for en Ellipse fandt han, vel ikke
fuldstændig Overeensstemmelse med Observationerne, men dog en
betydningsfuld Tilnærmelse. Hvorfor kun en Tilnærmelse? Blændet af
Cirkeltheorien\footnote{% trykt Mærke: *)
% S. 153 bærer TRE Noter, mærkede *), **) og ***). Den latinske Kepler-
% Text i denne Note er trykt ANTIQVA (opret), ikke kursiv; ligesaa
% Titelhenvisningen — modsat den kursiverede Titel i Texten paa S. 152.
% TRYKFEIL: „metaphycicæ“ (med c) for metaphysicæ. Bogstavet er positivt
% et c paa KB-Exemplaret, sammenlignet med s'et i „specie“ paa samme
% Linie; begge Exemplarer have det, og Trykfeilsbladet nævner det ikke.
% Gjengivet som trykt.
% Klammerne, tesseract sætter om denne Note, ere en tidligere Læsers
% Blyant, ikke Text.
Primus meus error fuit, viam planetæ perfectum esse circulum, tanto
nocentior temporis fur, qvanto erat ab auctoritate omnium philosophorum
instructior et metaphycicæ in specie convenientior. De motibus stellæ
Martis. Pars III. cap. XI.}, var Kepler gaaet ud fra den Forudsætning,
at Jordbanen, hvis Excentricitet forholdsviis er ubetydelig, maatte være
en Cirkel; men da denne Feil blev rettet, var ogsaa den saakaldte
\emph{første} Lov opdaget: \emph{Planetbanerne ere Ellipser, i hvis ene
Brændpunkt Solen befinder sig.}

Den polære Ligning for Ellipsen viser, at \textit{Radius vector} er en
Function af sin Vinkel med Storaxen%
\renewcommand{\thefootnote}{**)}%
\footnote{% trykt Mærke: **)
% Noten bestaaer alene af den udhævede Formel. Minus er et ægte Minus,
% ikke ÷; Exponenten er et 2, ikke et 3 (kontrolleret paa KB).
$r = f(v) = \dfrac{a\,(1 - e^{2})}{1 + e \cos v}.$}%
\renewcommand{\thefootnote}{*)}%
; dette leder igjen til en nærmere Bestemmelse af Bevægelsen i Banen.
Hidtil gjaldt det som et astronomisk Dogme, at Planeten absolut maatte
bevæge sig med jævn angulær Hastighed, med andre Ord, tilbagelægge lige
store Vinkler i lige store Tider. Til Ære for Dogmet opfandt man
\textit{Punctum æqvans} med tilsvarende epicykliske Kunstlerier. Kepler
indsaae omsider det Urimelige heri. Vistnok var der ved Bevægelsen
Noget, som var jævnt; men det var ikke de tilbagelagte Vinkler, det var
derimod de af \textit{Radius vector} beskrevne elliptiske Sectorer%
\renewcommand{\thefootnote}{***)}%
\footnote{% trykt Mærke: ***)
% Denne Note begynder paa S. 153 og ender paa S. 154; Mærket gjentages
% ikke ved Fortsættelsen. Hele Noten staaer her ved sit Mærke.
Af Loven om Arealerne følger, at Planetens angulære Hastighed i ethvert
Punkt af dens Bane forholder sig omvendt som Qvadratet paa Afstanden.
Herved vindes et nyt Middel til at bestemme Planetens Afstande fra
Solen, naar en af disse Afstande antages for givet.}%
\renewcommand{\thefootnote}{*)}%
. Herved bestemmes den saa\-%
% --- p. 154 ---
% „Radii vectores“ er trykt KURSIVT OG SPÆRRET midt inde i den spærrede
% Lovformel. \emph brydes derfor om det, saa Kursiven bevares: \emph inde
% i \emph vilde vende Skriften opret og tabe det, Siden trykker.
kaldte \emph{anden} Lov: \emph{De Arealer, som} \textit{Radii vectores}
\emph{beskrive, forholde sig som de til Bevægelsen anvendte Tider.}

Ligesom Copernicus har ogsaa Kepler en storartet Opfattelse og levende
Forvisning om en objectivt bestaaende Verdensharmonie. Der var i denne
Henseende ingen Tanke saa dristig, at dette vidunderlig kosmiske Genie
jo vovede at tage den op. I sit
% ogsaa denne Titel er kursiv OG spærret; se Noten ved S. 152.
\textit{Mysterium Cosmographicum} (1596) udtaler han sig om
Planetbanernes Forhold paa følgende Maade: „Jordbanen er en Cirkel.
Beskriver man om den Kugle, hvortil Jordbanen svarer som Storcirkel, et
Dodekaeder, saa giver den omskrevne Kugle Marsbanen. Beskriver man om
denne Bane et Tetraeder, saa fremstiller den tilsvarende Cirkel
Jupitersbanen. Beskriver man om Jupitersbanen en Kubus, vil den til
samme svarende Cirkel give Saturnsbanen“ o. s. v. Disse
Pythagorisk-Platoniske Speculationer fandt omsider en frugtbar
videnskabelig Afklaring, forsaavidt de undergik en Forvandling og efter
27 Aars Anstrengelser ledte til Opdagelsen af den \emph{tredie} Lov, som
siger, at \emph{Qvadraterne af de forskjellige Planeters Omløbstider
forholde sig som Cuberne af deres Middelafstande fra Solen.}

Disse Love, hvis objective Gyldighed Ingen kan drage i Tvivl uden med
det Samme at drage hele Erfaringslærens Gyldighed i Tvivl, afgive det
formelle Grundlag for videnskabelig Objectivering, ikke alene af
Solsystemet, men af Universet. At det Formelle i disse Loves
Objectivitet peger ud over sig selv med en Forjættelse om det Reelle, er
øiensynligt. Istedenfor tomme geometriske Centra, som i
Epicykeltheorien, have de elliptiske Baner et fyldigt legemligt
% --- p. 155 ---
Centrum, et virkeligt Centrallegeme i Brændpunktet; istedenfor den
subjective Reflexions imaginære Kredse forefinde vi her objective
Relationer mellem Centrallegemet og de peripheriske, i Banerne sig
bevægende Legemer. Idet Opmærksomheden fæster sig paa disse
Grundrelationer mellem Legemerne indbyrdes, faaer Spørgsmaalet om
bevægende Aarsager sikkre Støttepunkter og et nyt Liv, der lover et nyt
og rigt Udbytte: den formelle Erkjendelse har forberedt den reelle; den
abstract empiriske Mechanik er i Begreb med at blive til rationel
Physik.

% Ogsaa denne Overskrift er trykt HALVFED og løber ind i Afsnittet;
% den staaer paa S. 155, ikke paa S. 156.
\runin{b) Objectiviteten af det reelle Grundlag: Gravitationen.}
De første Indvendinger af nogen Betydning mod det Copernicanske System
hidrørte deels fra Uklarhed over Afstanden mellem Planetsystemet og
Fixstjernehimlen, idet man fordrede, at Copernicanerne skulde paavise de
Bevægelser, som Fixstjernerne, ifølge den nye Lære, nødvendig maatte
have\footnote{% trykt Mærke: *)
Mærkværdigt nok skulde Spørgsmaalet om Fixstjernernes Afstand siden, da
Loven for Lysets Forplantning var opdaget, lede til et af de stærkeste
Beviser for det Copernicanske System ved Bradley's Opdagelse af
Aberrationen.}, deels, og det fornemmelig, fra Ukyndighed i Mechaniken.
Dersom Jorden bevægede sig fra Vest mod Øst, saa maatte en Steen, der
faldt fra et høit Taarn, nødvendig --- indvendte man --- falde i vestlig
Retning, skjøndt Kjendsgjerningen er, at den falder i østlig: man
kjendte endnu ikke Inertien. Ved Galileis Opdagelse af Faldloven blev
Vexelbestemmelsen af Inertie og Tiltrækning paaviist i exact Form, og
den mechaniske Physik grundlagt. Imidlertid maatte der, hvad
Verdensmechaniken angaaer, endnu gjøres flere indledende Skridt og
famlende Forsøg, inden Newton kom til at udtale Grundsætningerne,
sammenarbeide det Brudstykkeagtige til et rationelt Hele og betegne den
nye
%  MATHEMATICAL — work from the image; the OCR is not a witness here.
% --- p. 156 ---
% S. 156 begynder FLUGTENDE (ingen Indrykning) og fortsætter Sætningen fra
% S. 155; Overskriften b) Objectiviteten af det reelle Grundlag: Gravitationen
% staaer altsaa IKKE paa denne Side --- den falder paa S. 155 (Batch 15).
Epoche ved et Værk som: \textit{Philosophiæ naturalis Principia
mathematica.}

Det er vistnok mere Værkets Overskrift end dets Ind\-%
hold, der i den nyere Tid har foranlediget Hegel og den hegelske
Skole%
\footnote{% trykt Mærke: *) --- første af to Noter paa S. 156.
Jul. Schaller. Geschichte der Naturphilosphie. 1ster Theil.
% „Naturphilosphie“: o'et mangler i „philos[o]phie“. Sætterfeil; begge
% Exemplarer (Arbeidsscanningen og KB-Exemplaret) vise det samme.
% Gjengivet som trykt; læs „Naturphilosophie“.
Leipzig 1841. S. 353--400.}
til en Deel lige saa uklare som ubeføiede Indsigelser imod de
Newtonske Principer. Holder man sig strængt til Benævnelsen:
Naturphilosophie (\textit{philosophia naturalis}), og tager Ordet
som Kategorie i hegelsk Forstand, da kan der unegtelig gjøres
Indvendinger mod den Maade, hvorpaa Newton har defineret og
opstillet sine Grundsætninger; men en saadan Opfatning beroer paa
en kolossal Misforstaaelse af det Newtonske Synspunkt. For at
undgaae unødvendige Kævlerier om dunkle Ord, erklærer Newton
ligefrem, at de af ham valgte Udtryk, Inertie, Til\-%
trækning o. fl., ikke skulle tages i physisk, men kun i
mathematisk Betydning.%
\renewcommand{\thefootnote}{**)}%
\footnote{% trykt Mærke: **) --- anden af to Noter paa S. 156.
Voces autem attractionis, impulsus, vel propensionis cujuscunque in
centrum, indifferenter et pro se mutuo promiscue usurpo; has vires
non physice, sed mathematice tantum considerando.
Princ.Mathem.\ Amstælodami MDCCXXIII. pag.\ 5.}%
% „Amstælodami“ med æ (for Amstelodami) staaer saaledes i begge
% Exemplarer; læst paa KB-Exemplaret. Gjengivet som trykt.
\renewcommand{\thefootnote}{*)}
Deraf har den speculative Philosophie taget Anledning til at vise,
hvor ubegrundet hans Adskillelse mellem det Physiske og det
Mathematiske er i sig selv, og hvor lidt han har været i Stand til
at fastholde Adskillelsen conseqvent. Men hvor fremmed denne hele
Opfattelse af Physik og Mathematik er for Newtons Tanke\-%
gang, vil man ved nogen Eftertanke let kunne indsee. Newton
forvarer sig imod, at man tager hans Udtryk physisk d. v. s.
metaphysisk, og Meningen er, at han aldeles ikke vil indlade sig
paa speculative Undersøgelser
% --- p. 157 ---
angaaende de physiske Virksomheders indre Hvorledes%
\footnote{% trykt Mærke: *) --- eneste Note paa S. 157.
Vocem attractionis hic generaliter usurpo pro corporum conatu
qvocunque accedendi ad invicem; sive conatus iste fiat ab actione
corporum vel se mutuo petentium vel per spiritus emissos se invicem
agitantium, sive is ab actione ætheris aut aëris mediive cujuscunque
seu corporei seu incorporei oriatur corpora innatantia in se invicem
utcunque impellentis \ldots{} l.\ l.\ p.\ 172, confer.\ p.\ 483.
% De to Bogstaver efter Prikkerne ere ens og have ingen Fane øverst til
% venstre (til Forskjel fra Ettallet i „172“): altsaa to smaa l'er, ikke
% „1. 1.“ og ikke „l. I.“ (liber I), som blev overveiet og forkastet.
% Læst paa KB-Exemplaret. Noten blander qv og cu i een og samme Sætning:
% „qvocunque“ ved Siden af „cujuscunque“ og „utcunque“ --- som trykt.
Tydeligere kan Grundtanken for det hele Synspunkt dog vel neppe
udtales.}%
; men at løsrive det Physiske fra det Mathematiske, kunde dog vel
umulig falde den Mand ind, som fra Først til Sidst lægger an paa at
tydeliggjøre, hvorledes det Qvantitative i Virkningen Punkt for
Punkt viser hen til det Qvalitative i Aarsagen.

Mechanik er oprindelig Eenhed af Mathematik og Physik; Rum og Tid,
Materie og Kraft ere i lige Grad physiske og mathematiske Factorer.
For at udvikle det Mathematiske maa Newton lægge Vægt paa det
Physiske, og omvendt. Mechaniken fordrer udtrykkelig, at Forholdet
mellem begge skal opfattes dobbeltsidig, deels fra Syns\-%
punktet af Aarsag og Virkning, deels fra Synspunktet af Lov og
Phænomen. Hvad Lov og Phænomen angaaer, er det forstaaeligt, at
Newton kan ombytte Kategorier som Attraction og Stød med hinanden,
thi hvad Attractionen er i Continuitetens, er Stødet da i
Discretionens Form. Hvad Aarsag og Virkning angaaer, er der derimod
en væsenlig Forskjel mellem Attraction og Stød. Men, at der,
mechanisk forstaaet, ingen Strid er mellem den første og den sidste
Opfattelse, sees ligefrem deraf, at et Legeme ligesaa vel paavirkes
\emph{udenfra}, naar det tiltrækkes, som naar det stødes.
% Spatieringen af „udenfra“ er bekræftet paa KB-Exemplaret.
Den Modsigelse mellem „\textit{vis attractiva}“, physisk opfattet,
% Latinen er sat kursiv indenfor Anførselstegnene; læst paa KB.
og Træghedsloven, hvorpaa Schaller beraaber sig, finder i
Virkeligheden ikke Sted. Hvad Newton frem\-%
% --- p. 158 ---
hæver, at Attractionskraften ikke kan have physisk Betyd\-%
ning allerede af den Grund, at de Centra, hvorfra den udgaaer, kun
ere mathematiske Punkter, maa af Sammenhængen forstaaes saaledes,
at den mechaniske Theorie, der gjør en Begyndelse med at
concentrere de tiltrækkende Aarsager i mathematiske Punkter, just
derved giver tilkjende, at den ikke vil undersøge den indre
Væsenseenhed af Materie og Kraft; men deraf kan paa ingen Maade
udledes, hverken at det Physiske er udelukket, eller at det
Mathematiske i Tiltrækningen skulde høre op, fordi Centret lægges
ind i en legemlig Masse.

(Forholdet mellem Kepler og Newton er paa Mechanikens Trin et
% Parenthesen aabnes her paa S. 158 og lukkes først paa S. 161
% („Tænkning.)“). Det er ikke en Trykfeil og faaer ingen Note.
Forhold mellem Objectivering i empirisk Reflexion og Objectivering i
gjennemgaaende Begriben. Den Keplerske Reflexion abstraherer Lovene
af foreliggende Til\-%
fælde og sikkrer sig dem ved empirisk Induction; den Newtonske
Methode gaaer derimod ud paa at erkjende Aarsagerne af deres
Virkninger, bestemme det Særlige paa Grundlag af det Almene og
begribe Lovene i deres Opfyldelse. Det er paa Objectiveringens Vei
et stort Fremskridt.

Kepler taler kun om de Arealer, som beskrives i Planetbanerne; men
Newton beviser, at Loven for Arealerne maa gjælde for en
hvilkensomhelst Bevægelse af et materielt Punkt om et fast Centrum,
da den i Almindelighed er bestemmende for Centralbevægelsens
objective Begreb.

Kepler holder sig til det Factiske, at Planetbanerne ere Ellipser,
men, bunden til Kjendsgjerningen, er han ikke i Stand til at angive
Aarsagen; af Newtons Principer kan derimod Aarsagen sees i
Virkningen, forsaavidt det kan indsees, at ikke blot Ellipsen, men i
Almindelighed Keglesnittet er en naturnødvendig Conseqvens, en
udtrykkelig Aarsagsconseqvens af Tiltrækningens eiendommelige Lov:
% --- p. 159 ---
$T = \dfrac{\mu}{r^{2}}$, og at omvendt denne Lov
% Formlen staaer inde i Textlinien, ikke udhævet paa egen Linie.
er en nødvendig Conseqvens af Bevægelsen i et Keglesnit.%
\footnote{% trykt Mærke: *) --- eneste Note paa S. 159; den fylder
% hele Bladet under de to Linier Brødtext.
J. Schaller bemærker, at Newton ikke har kunnet „almindelig
bevise“, at Curven maa blive et Keglesnit, naar Tiltrækningen
forholder sig omvendt som Qvadratet af Afstanden. Det maa i denne
Sammenhæng udtrykkelig fremhæves, at det ikke væsenlig kommer an paa
den foreløbige Beviismaade, men paa det, der af Grundtanken selv
lader sig udlede. Enhver der gjør sig bekjendt med den rationelle
Mechaniks Fundamentalsætninger, vil have rig Leilighed til at
overbevise sig om de Newtonske Principers Almeengyldighed. For
Centralbevægelsen i Almindelighed, naar r er Radius vector, c en
vilkaarlig Constant, θ den Vinkel, r danner med en fast Axe, og φ
den accelererende Kraft, finder man (jfr. Duhamel. Cours de
Mécanique. Paris 1846. II. S. 3) Formlen:
\[
\varphi = \frac{c^{2}}{r^{2}} \left( \frac{1}{r} +
    \frac{d^{2}\dfrac{1}{r}}{d \theta^{2}} \right)
\]
Er nu Banen et Keglesnit, Vinklen mellem r og Polaraxen θ, mellem
Storaxen og Polaraxen ω, Ligningen for Banen
\[
\frac{1}{r} = \frac{1 + e \cos (\theta - \omega)}{\dfrac{p}{2}},
\]
saa følger heraf:
\[
\frac{d^{2}\dfrac{1}{r}}{d \theta^{2}} = \div \frac{e \cos (\theta - \omega)}{\dfrac{p}{2}},
\]
% Fortegnet foran Brøken er Nielsens MINUS (÷), læst paa KB-Exemplaret.
% Tegnet mellem θ og ω er derimod et almindeligt Minus i begge Formler
% paa denne Side, medens S. 160 sætter ÷ paa samme Sted. Ikke harmoniseret.
som i Forbindelse med Ligningen for Banen giver:
\[
\frac{1}{r} + \frac{d^{2}\dfrac{1}{r}}{d \theta^{2}} = \frac{1}{\dfrac{p}{2}}
\]
og følgelig:
\[
\varphi = \frac{c^{2}}{\dfrac{p}{2}} \cdot \frac{1}{r^{2}} = \frac{k}{r^{2}}
\]
Heraf sees, at Tiltrækningen forholder sig omvendt som Qvadratet af
Afstanden, naar Banen er et Keglesnit, og Centralkraften udgaaer fra
Brændpunktet.}

% --- p. 160 ---
Kepler indskrænker sig til at opstille de tre Love ved Siden af
hverandre uden nogensomhelst Deduction; men ifølge Newtons Principer
kan den ene Lov, naar man gaaer ud fra det Almindelige og indfører
Betingelser, deduceres af den anden.%
\footnote{% trykt Mærke: *) --- eneste Note paa S. 160. Noten STRÆKKER
% SIG OVER TO BLADE: den fylder Foden af S. 160 og fortsættes øverst i
% Notefeltet paa S. 161 (fra „c = ...“), uden at Mærket gjentages.
% Efter Husreglen staaer hele Noten her ved sit Mærke.
% Formlerne staae i Trykken INDE I Textlinierne, ikke udhævede paa egne
% Linier; saaledes gjengivne her.
Sættes Tiltrækningskraften $\varphi = \dfrac{\mu}{r^{2}}$, saa blive
Ligningerne i Planet (xy):
$\dfrac{d^{2}x}{dt^{2}} = \div \dfrac{\mu x}{r^{3}}$,
$\dfrac{d^{2}y}{dt^{2}} = \div \dfrac{\mu y}{r^{3}}$ (I).
% Tegnet er Nielsens MINUS (÷), læst paa KB-Exemplaret, ikke Division.
% Nummereringen er et stort I med Seriffer (ikke Ettallet) begge Steder:
% baade her og ved Henvisningen „Ligningerne (I)“ nedenfor. Læst paa KB.
Indfører man Polarcoordinater, saa at $x = r \cos \theta $, $y = r \sin \theta $,
og det af Radius vector beskrevne Areal
$\tfrac{1}{2} c\, dt = \tfrac{1}{2} r^{2}\, d\theta $, saa faaer man ved at
multiplicere Ligningerne (I) med $cdt = r^{2}\, d\theta $, Ligningerne
$c\, d \dfrac{dx}{dt} = \div \mu \cos \theta d\theta $ (α) og
$c\, d \dfrac{dy}{dt} = \div \mu \sin \theta d\theta $ (β).
Integreres disse Ligninger saaledes, at Constanten for (α) bliver
$\div \gamma \sin \omega $, for (β) $\gamma \cos \omega $; multipliceres derpaa den første
med $y = r \sin \theta $, den anden med $x = r \cos \theta $ (idet man erindrer,
at $xdy \div ydx = cdt$ og $\sin^{2} \theta + \cos^{2} \theta = 1$), saa
erholder man:
$r = \dfrac{\dfrac{c}{\mu}}{1 + \dfrac{\gamma}{\mu} \cos (\theta \div \omega)}$
% TRYKFEILBLADET retter S. 160, Linie 13 f. n.: „c/μ i Tælleren --- læs
% c²/μ“. Efter Husreglen staaer Texten her SOM TRYKT, med c (uden Exponent)
% i Tælleren. KB-Exemplaret er entydigt: der staaer c alene. Arbeidsscanningen
% viser et hævet 2 ved siden af c'et --- det er en tidligere Læsers BLYANT,
% ikke Trykken, ganske som det forfalskede Punktum paa S. 113.
% Tegnet i (θ ÷ ω) er Nielsens ÷; S. 159 sætter almindeligt Minus samme Sted.
d. v. s. \emph{Banen er et Keglesnit med Tiltrækningens Centrum i
Brændpunktet.}
% Spatieringen bekræftet paa KB-Exemplaret.

Integreres Ligningen $r^{2}\, d\theta = cdt$ saaledes, at $\theta = o$, naar
$t = o$, saa bliver $t = \dfrac{\displaystyle\int_{o}^{\theta} r^{2}\, d\theta}{c}$,
% Grænserne og Værdierne ere sat med BOGSTAVET o, ikke med Nullet;
% læst paa KB-Exemplaret.
hvilket med Ord vil sige, at \emph{de af Radius vector beskrevne
Arealer ere proportionale med Tiderne, hvori de beskrives.}
% Spatieringen bekræftet paa KB-Exemplaret; den begynder ved „de“.

Vælges endelig det særlige Tilfælde, at Keglesnittet er en Ellipse,
og tages Integralet $ct = \int r^{2}\, d\theta $ fra o til $2\pi $, saa bliver
dette Ellipsens dobbelte Areal,
$2 ab\pi = 2a^{2} \sqrt{1 \div e^{2}}\, \pi $, og
$T = \dfrac{2 a^{2} \sqrt{1 \div e^{2}}\, \pi}{c}$ den Tid, hvori hele
Banen tilbagelægges. Kjender man a, e og T, kan man bestemme c, idet
$c = \sqrt{a (1 \div e^{2}) \mu}$, hvor man erholder
$t = \dfrac{\int r^{2}\, d\theta}{\sqrt{a (1 \div e^{2}) \mu}}$, følgelig
$T = \dfrac{2 a^{2} \sqrt{1 \div e^{2}}\, \pi}{\sqrt{a (1 \div e^{2}) \mu}}$;
altsaa er $\dfrac{T^{2}}{a^{3}} = \dfrac{4\pi^{2}}{\mu} = K$.
% Den trykte Rodstreg naaer i de to sidste Formler ikke ud over μ'et;
% Sammenhængen kræver μ under Roden, og saaledes er det sat her.
At Forholdet mellem Omløbstidens Qvadrat og tredie Potens af Storaxen
er en constant Størrelse, viser, at dette Forhold gjælder for
forskjellige materielle Punkter, der bevæge sig i Ellipser om et
fælles Brændpunkt; med andre Ord: \emph{Omløbstidernes Qvadrater
forholde sig som Kuberne af de halve Storaxer eller
Middelafstande.}
% Spatieringen bekræftet paa KB-Exemplaret (spacing.py melder to RUN her).

Paa disse Betingelser kunne de Keplerske Love deduceres af Newtons
Principer.}
At Deductionen bestandig maa knyttes til
% --- p. 161 ---
givne Betingelser, minder om, at Beviserne, hvad det Physiske
angaaer, hverken ere blot aprioriske eller blot empiriske, men at
de, ifølge Sagens Natur, maae beroe paa en exact articuleret
Vexelbestemmelse af Erfaring og Tænkning.)
% Parenthesen, som blev aabnet paa S. 158, lukkes her.

Man har fra „Begrebets“ Standpunkt paastaaet, at „jo almindeligere
Tiltrækningen opfattes, desto mere taber det mathematiske Beviis i
Betydning og viser sig utilstrække\-%
ligt til at træffe Phænomenets Væsen.“%
\footnote{% trykt Mærke: *) --- egen Note paa S. 161. I Trykken staaer
% den under Fortsættelsen af Noten fra S. 160.
Schaller. Anfr. Skr. S. 367.}
% „Anfr.“ (ikke „Anf.“) staaer der; læst paa KB-Exemplaret.
Det Modsatte heraf er netop det Sande. Videnskabens Historie lærer,
at Indsigten i Phænomenernes Væsen stadig er skredet frem med
Lovenes Erkjendelse, og Lovenes Erkjendelse med den mathematisk
mechaniske Bestemmelse af de objective Relationer, af hvilke
Aarsagerne og deres Betingelser alene lade sig udfinde.

Gravitationslæren, Læren om Klodernes gjensidige Tiltrækning,
grundes ligelig paa Begreb og Kjendsgjerning, er i lige Grad
apriorisk og empirisk. Det kan af Erfaring godtgjøres og ved
Induction fastsættes, hvorledes Masserne indbyrdes paavirke
hinanden; men den gjensidige Paavirkning, som lader sig iagttage,
begrundes igjen ved en Virk\-%
% Nederst paa S. 161 staaer Arksignaturen „11“. Ikke Text; ikke gjengivet.
% --- p. 162 ---
somhed, som ikke umiddelbart lader sig iagttage, men maa bestemmes
ved Tænkning, nemlig den Vexelvirkning, hvorved de mindste Dele i
den ene Masse tiltrække Delene i den anden. Delenes Tiltrækning
forklares af Massernes, Massernes igjen af Delenes. Den gjensidige
Tiltrækning mellem to Masser er proportional med deres Product,
medens den Kraft, hvormed den ene Masse accelereres i sin Bevægelse
om den anden, er proportional med Massernes Sum.%
\footnote{% trykt Mærke: *) --- første af to Noter paa S. 162.
Er M Solens, m Planetens Masse, og Afstanden r, saa bliver den
gjensidige Tiltrækning at udtrykke ved $\dfrac{f.\,Mm}{r^{2}}$; den
accelererende Kraft for Planeten ved $\dfrac{f.M}{r^{2}}$, for Solen
ved $\dfrac{f.m}{r^{2}}$; for Planetens \emph{relative} Bevægelse
% Spatieringen af „relative“ er bekræftet paa KB-Exemplaret.
omkring Solen maa den drivende Kraft derimod blive
$\dfrac{f\,(M + m)}{r^{2}}$. Ved f betegnes den gjensidige
% Punktummet efter f er Nielsens Multiplicationstegn (jvnfr. S. 46);
% den sidste Brøk trykker intet Punktum. Gjengivet som trykt.
Tiltrækning af to Masse-Eenheder, reducerede til simple Punkter og
stillede i en Afstand, der angiver en Længde-Eenhed.}
Idet Masserne, de virkelige Kloder, indføres i Stedet for materi\-%
elle Punkter, undergaae de Keplerske Love ganske vist mangfoldige
Ændringer; men det er saa langt fra at disse Ændringer svække deres
Gyldighed, at de tvertimod ere Beviser paa deres virkelige
Opfyldelse.

Dersom Kloderne vare fuldkomment kugleformige, dannede af ligeartede
concentriske Lag, vilde de paavirke hinanden, som om enhver af dem
havde sin hele Masse sammentrængt i sit Centrum; men da de i Form
afvige noget fra Kuglen og, hvad der er væsenligere, befinde sig i
Tiltrækningsforhold ikke blot til Solen men til hverandre indbyrdes,
saa bliver Vexelvirkningen uendelig forviklet, og Perturbationerne,
de periodiske saa vel som de seculære, stille i denne Henseende
store Fordringer%
\renewcommand{\thefootnote}{**)}%
\footnote{% trykt Mærke: **) --- anden af to Noter paa S. 162. Noten
% STRÆKKER SIG OVER TO BLADE: den begynder paa Foden af S. 162 og
% sluttes øverst i Notefeltet paa S. 163, uden at Mærket gjentages.
% Efter Husreglen staaer hele Noten her ved sit Mærke.
Den af Delaunay i connaissance des temps for 1865 fremstillede
Formel for Maanens Længde optager ikke mindre end 45 trykte Sider,
og dog tager denne Formel kun Hensyn til de Perturbationer, der
skyldes Solen.}%
\renewcommand{\thefootnote}{*)}
til den iagt\-%
% --- p. 163 ---
tagende og beregnende Videnskab. Perturbationerne ere dog ikke, hvad
Ordet eensidig antyder, Forstyrrelser af noget Virkeligt; tvertimod!
de høre væsenlig med i Charakteren af mechanisk Totalitet; de ere
Phænomener af mangesidig Vexelvirkning; det er just
Perturbationerne, der give det Virkelige Præg af udvortes materiel
Objectivitet. Til videnskabelig Objectivering af Verdensbygningen er
det ikke tilstrækkeligt at betragte de enkelte Kloder i de enkelte
Baner, hver for sig, eller blive staaende ved en almindelig
Forestilling om de forskjellige tilstødende Paavirkninger, en Planet
i sin Bane kan være underkastet. Virkelighedsbegrebet maa udvikles
saa exact og fuldstændigt, at Vexelbestemmelsen af det Almene og det
Enkelte, det Logiske og det Antilogiske
% TRYKFEILBLADET retter S. 163, Linie 13 f. o.: „det Antilogiske --- læs
% det Antilogiske,“. Efter Husreglen staaer Texten som TRYKT, altsaa UDEN
% Komma. KB-Exemplaret bekræfter, at der intet Komma staaer i Satsen;
% Arbeidsscanningens Komma er en tidligere Læsers BLYANT (med Understregning).
Lovene og deres Ændringer i de særlige Tilfælde, staaer klart for
Tanken.%
\footnote{% trykt Mærke: *) --- eneste Note paa S. 163. Noten STRÆKKER SIG
% OVER TO BLADE: den begynder paa Foden af S. 163 og fylder hele Notefeltet
% paa S. 164, uden at Mærket gjentages. Efter Husreglen staaer hele Noten
% her ved sit Mærke.
% Prikkerækkerne ere i Trykken sat med tre eller fire spatierede Punkter;
% Antallet er ikke gjengivet, men sat med \ldots overalt.
Betegner man Solens Masse ved M, dens Coordinater ved X, Y og Z,
Planeternes Masser ved $m$, $m'$, $m''$ \ldots{}, deres Coordinater
% Enumerationen holdes i een Skriftart: da Leddene bære Mærker, er hele
% Rækken sat i Mathematik (Husreglen fra Batch 5).
ved $x$, $y$, $z$, $x'$, $y'$, $z'$, $x''$, $y''$, $z''$ \ldots{},
lader man endvidere Solens Afstande fra Planeterne være $r$, $r'$,
$r''$ \ldots; den perturberede Planets Afstande fra de perturberende
$s'$, $s''$ \ldots{}, saa ville Ligningerne for Bevægelsen af den
enkelte Planet, m, om Solen som fast Ceutrum være med Hensyn paa
% „Ceutrum“ for „Centrum“: omvendt eller forkert Type. Bogstavet er
% utvetydigt et u paa KB-Exemplaret (aabent foroven, lukket forneden,
% til Forskjel fra n'et i „Solen“ paa samme Linie). Sætterfeil; gjengivet
% som trykt.
% TRYKFEILBLADET retter S. 163, Linie 10 f. n.: „være med --- læs være, med“.
% Efter Husreglen staaer Texten som TRYKT, UDEN Komma efter „være“.
% KB-Exemplaret bekræfter det; Kommaet i Arbeidsscanningen er Blyant.
X-Axen:
\[
\frac{d^{2}\,(x - X)}{dt^{2}} + \frac{(M + m)\,(x - X)}{r^{3}} = 0;
\]
% TRYKFEILBLADET retter S. 163, Linie 8 f. n.: „$=0$; --- læs $=0$, o. s. v.“
% Efter Husreglen staaer Texten som TRYKT: „= 0;“ med Semikolon og uden
% „o. s. v.“. KB-Exemplaret bekræfter det; „o. s. v.“ i Arbeidsscanningen
% er skrevet med Blyant af en tidligere Læser.
men dersom Planeten m perturberes af $m'$, $m''$ \ldots, vil
Ligningen for X-Axen være:
\begin{align*}
\frac{d^{2}\,(x - X)}{dt^{2}} + \frac{(M + m)\,(x - X)}{r^{3}}
  = {}& \frac{m'(x' - x)}{{s'}^{3}} + \frac{m''\,(x'' - x)}{{s''}^{3}} + \\
\ldots {}& - \frac{m'\,(x' - X)}{{r'}^{3}}
        - \frac{m''\,(x'' - X)}{{r''}^{3}} - \ldots \quad (I)
\end{align*}
% Tegnene her ere almindelige Minus, ikke Nielsens ÷; læst paa KB-Exemplaret.
og de tilsvarende for Y- og Z-Axen have samme Form. At vi i slige
Ligninger have et System af objective Relationer, er vist; men
hvorved er Objectiviteten sikkret? Ved Objec\-%
tiveringen! Og hvorved er Objectiveringen sikkret? Ved articulerende
% Notens Text gaaer her over paa S. 164.
Eenhed af Iagttagelse og Tænkning. Det er Tænkningen, der indseer,
at, naar Loven for Bevægelsen af den enkelte Planet, m, er
$\dfrac{d^{2}x}{dt^{2}} = - \dfrac{M\,(x - X)}{r^{3}}$ o. s. v., for
Solens Bevægelse
$\dfrac{d^{2}X}{dt^{2}} = \dfrac{m\,(x - X)}{r^{3}}$ o. s. v., saa
maa Loven for den pertuberede Planet, m, blive:
% TRYKFEILBLADET retter S. 164, Linie 15 f. o.: „pertuberede --- læs
% perturberede“. Efter Husreglen staaer Texten som TRYKT: „pertuberede“.
% Bekræftet paa KB-Exemplaret --- ingen Blyant her. (Linietællingen
% indbefatter Notelinierne, jvnfr. Filhovedet.)
\[
\frac{d^{2}x}{dt^{2}} = - \frac{M\,(x - X)}{r^{3}}
  + \frac{m'\,(x' - x)}{{s'}^{3}}
  + \frac{m''\,(x'' - x)}{{s''}^{3}} + \ldots
\]
og for Solen:
\[
\frac{d^{2}X}{dt^{2}} = \frac{m\,(x - X)}{r^{3}}
  + \frac{m'(x' - X)}{{r'}^{3}}
  + \frac{m''\,(x'' - X)}{{r''}^{3}} + \ldots \quad (II)
\]
Paa venstre Side bestemmes Begrebet af (perturberende) Kraft i exact
logisk Almindelighed; paa høire Side vises det, hvorledes Begrebet
udtrykker sig i existerende Tilfælde. Af Ligningerne (II) for de
absolute Bevægelser udledes saa Ligningerne (I) for de relative. Den
objectiverende Reflexion articulerer det Logiske, det Almene, ved at
tydeliggjøre Loven i dens forskjellige Former og Formler; men
Tegnsproget viser dog ogsaa, at de almindelige Formler kun ere tomme
Symboler, saalænge Virkelighedsmomenterne med Tal og Maal i
tilsvarende antilogiske Bestemmelser endnu mangle.
Kjendsgjerningerne ere objective, men Analyserne subjective. Hvad
lære vi heraf? At ingen, end ikke den mest objective Videnskab kan
tage det Subjective bort og beholde det Objective, at, med andre
Ord, Intet i Verden er objectivt uden det, som er objectiveret.}
Først fra Objectiveringens Standpunkt
% Nederst paa S. 163 staaer Arksignaturen „11*“. Ikke Text; ikke gjengivet.
% --- p. 164 ---
bliver det tilfulde indlysende, hvad det vil sige, at saadanne
realistiske Videnskaber som Mechanik og physisk Astronomie staae i
Aandens Tjeneste. Det er noget langt Høiere end en smaalig Higen
efter at faae de mangfoldige Bevægelses\-%
phænomener punktlig udregnede, der driver Undersøgelserne fremad;
det er den forskende Aand, der theoretisk vil have Virkeligheden i
sin Magt; det er den tænkende Subjectivitet, der attraaer det
Objective og rastløs arbeider med Objecti\-%
veringens Betingelser. Et storartet Exempel, som viser, hvad
%  MATHEMATICAL — work from the image; the OCR is not a witness here.
% --- p. 165 ---
% S. 165 AABNER FLUSH, uden Indrykning: Linien fortsætter den Sætning, der
% løber fra S. 164 (Batch 16). Ingen ny Paragraf her --- Sømmen skal ikke
% have noget Afsnitsskifte.
% „Mécanique céleste“ staaer i Kursiv mod den omgivende Antiqua, med Accent
% baade paa „Mécanique“ og „céleste“; „Laplace's“ er almindelig Antiqua.
% Læst paa KB-Exemplaret (tesseract læste „celeste“ uden Accent).
Videnskaben kan naae ved at bygge paa Newtons Grundlag, har Nutiden
beundret i Laplace's \textit{Mécanique céleste}.

% Indløbende Overskrift, læst paa Siden selv (KB-Exemplaret) og ikke i
% Indholdet, hvis OCR har tabt Bogstavpræfixet: Præfixet er „c)“. Trykt
% HALVFED og paa samme Linie som den Paragraf, den aabner --- Mønstret fra
% S. 79, 101 og 107. Kolon efter „Objectivitet“ som trykt.
\runin{c) Universets mechaniske Objectivitet: et System af Systemer.}
Da Hephaistos fordum blev kastet ud af Olympens Vindue maatte han falde
i 9 Dage og 9 Nætter, inden han naaede Lemnos. Vel kan et Fald fra en
saadan Høide opvække Gysen; men naar man betænker, at det er Veien fra
den øverste Himmel ned til Jorden, som paa denne Fart maatte
gjennemløbes, skal man ikke kunne sige, at Længden er overdreven. I
Grækernes Verdensanskuelse er Begrændsning et saa væsenligt Moment, at
Aristoteles endog paa en Maade kan siges at gjøre Guddommen selv til
% GRÆSKEN er læst paa begge Exemplarer ved høi Forstørrelse. Sidste Ord
% staaer som οὐρανοῦ, ganske regelret. Derimod sætter Trykken i det tredie
% Ord Perispomenen over OMIKRON og ikke over ypsilon --- „τõυ“ --- i begge
% Copier, altsaa Sætterens (eller Skriftsnittets) Form og ikke denne Copis
% Skade. Unicode/LaTeX kan ikke bekvemt bære den Placering, saa Ordet
% gjengives her som τοῦ, og Afvigelsen noteres alene her.
Himlens Grændse (τὸ πέρας τοῦ οὐρανοῦ). Vel maatte den iagttagende
Astronomie fra mange Sider komme ind paa Spørgsmaalet om
Himmellegemernes Afstandsforhold, men Observationerne vare ufuldkomne,
og den Ptolemæiske Epicykeltheorie maatte nøies med en indbildt
Maalestok. Tycho Brahe selv antog endnu Solens Parallaxe for at være
saa stor, at Jorden blev en betydelig Størrelse i Universet. Først
efter Kikkertens Opdagelse og Indførelsen af de nye
Observationsmethoder naaede Videnskaben saa vidt, at den fik Maal og
Maalestok. Ikke nok, at Jordbaneradien blev en bekjendt Størrelse; selv
Fixstjernernes uendelig ringe, i Almindelighed forsvindende, Parallaxe
førte til Indblik i umaalelige Afstandsmaal. Da Ole Rømer ved Hjælp af
Jupitersmaanernes Formørkelse havde opdaget Lysets Hastighed, gik
Tanken med Straalen gjennem Himmelrummet; det Umaalelige blev maalt,
det Ufattelige forstaaet. Phantasien gav tabt, men Indsigten seirede.
For at gjennemløbe Afstanden fra Solen til Jorden, 20 Millioner Miil,
behøver Lyset 8 Minuter og 13 Secunder; for at gjennemløbe Afstanden
fra den 61de Stjerne i Svanen, hvis Parallaxe Bessel har bestemt,
% Brøken staaer i Trykken som opsat Brøk i Linien: 9 med lille ophøiet 1,
% Skraastreg og lille nedsat 4. \tfrac efter Husreglen (jvnfr. „15\tfrac{5}{8}
% Fod“ paa S. 50).
behøver Lyset $9\tfrac{1}{4}$ Aar; en saadan
% --- p. 166 ---
Veilængde er for Anskuelsen ufattelig, og dog bliver den nøiagtig
opfattet af Forstanden, den bliver jo maalt. Men naar man saa fremdeles
„af Teleskopets Beskaffenhed“ maa slutte, at Lyset fra de fjerneste
Stjerner i Melkeveien har været 2 à 3000 Aar underveis, ja, at der er
Stjerner, vi see i præadamitisk Lys: da bringer det Umaalelige endogsaa
Tanken til at svimle.

For den klassiske Phantasieanskuelse var Verden begrændset; for den
moderne, vi kunne sige: romantiske Reflexion, er Universet
ubegrændset. Enhver af disse Synsmaader er eensidig. Strækker det
virkelige Universum sig virkelig saa vidt som Rummet, da maa det jo
være ubegrændset som Rummet. Men det Ubegrændsede er et Uvirkeligt; alt
Virkeligt er begrændset: forsaavidt Universet er virkeligt, er det
begrændset. Den tilsyneladende Modsigelse, at Universet baade er
begrændset og ubegrændset, er imidlertid, som allerede viist i Læren om
Rummet, ingen uopløselig Antinomie. Astronomien selv har antydet
% Humboldt-Citatet: Anførselstegnene ere balancerede paa denne Side ---
% tre „ og tre “. Afbrydelsen „, siger Humboldt (Kosmos I, S. 68),“ lukker
% og aabner paany, men lukker ogsaa til Slut; altsaa IKKE Bogens bekjendte
% Vane (S. 10, 11, 16) med kun een Lukning.
Gaadens Løsning. „Sammenligner man“, siger Humboldt (Kosmos I, S. 68),
„Verdensrummet med et af vor Planets ø-rige Have saa kan man tænke sig
Materien gruppeviis fordeelt: snart som uopløselige Stjernetaager af
forskjellig Alder, fortættede omkring een eller flere Kjærner, snart
% „sammenboldtet“ staaer saaledes paa begge Exemplarer; gjengivet som trykt.
som sammenboldtet i Stjernehobe eller isolerede Sporader. Vor
% Spærring bekræftet paa KB-Exemplaret: „V e r d e n s ø“.
Stjernehob, den \emph{Verdensø}, til hvilken vi høre, danner et
lindseformigt, fladtrykt, til alle Sider afsondret Lag, hvis store Axe
antages til syv- eller ottehundrede, den lille derimod til hundrede og
% Kommaet efter „Universet“ staaer klart paa Arbeidsscanningen; paa
% KB-Exemplaret er Halen slidt bort, saa Tegnet dér ligner et Punktum.
% Typen er et Komma.
halvtredsindstyve Siriusafstande“. Vi ledes herved til at opfatte
Universet, ikke blot som en uoverskuelig Mangfoldighed af Verdener og
% Spærring bekræftet paa KB-Exemplaret: „e t  S y s t e m  a f  S y s t e m e r“.
Verdensøer, men som \emph{et System af Systemer}.

% Spærring bekræftet paa KB-Exemplaret: „p h y s i s k e“.
De \emph{physiske} Dobbeltstjerner, som ved gjensidig Tiltrækning
beskrive Baner om hinanden efter de Keplerske Love, indeholde en
Bekræftelse paa, at Newtons Principer
% --- p. 167 ---
ere gjældende ogsaa udenfor vort Solsystem. Er nu Tiltrækningen en al
Materien iboende Egenskab, saa følger jo heraf, at ikke blot Kloderne i
hvert enkelt System maae tiltrække hverandre, men at Systemerne tillige
maae paavirke hverandre gjensidig og vise sig som Led, der ad
Middelbarhedens uendelige Omvei kunne betragtes som Realisationer af et
eneste, altomfattende Grundsystem. Denne Synsmaade stadfæstes endvidere
derved, at der i Verdensrummet ikke gives noget absolut Stillestaaende;
Alt er i Bevægelse, ikke blot Planeter og Kometer, men Solen selv og
Solene: der gives, egenlig talt, ingen Fixstjerner. Det Spørgsmaal
paatrænger sig: er her for alle disse Bevægelser et fælles Centrum?
Legemligt kan dette Centrum vistnok ikke være. Som bekjendt, bevæger
Maanen sig ikke om Jordens Centrum, men Jord og Maane om et fælles
Centrum; og det Samme gjælder om Solen og Planeterne. Naar da hele vort
Solsystem har en translatorisk Bevægelse, saa tyder det just ikke paa,
at der er et enkelt fast Legeme, hvorom Bevægelsen skeer; det
forudsætter snarere, at Systemet paavirkes af et andet eller rettere af
andre Systemer, saa at det fælles Tyngdepunkt bestemmes ved Systemernes
gjensidige Tiltrækning. I Begrebets Almindelighed bliver altsaa Sagen
at opfatte saaledes: Universet er begrændset, thi ethvert virkeligt
System har sin Grændse; Universet er ubegrændset, thi udenfor det ene
virkelige System er der et andet virkeligt, udenfor dette igjen et
andet: hvert System er et Hele for sig, og dog en Deel, et Led af et
mere omfattende, endnu mere sammensat System: den virkelige Grændse er
allevegne og overskrides allevegne. Det Endelige ender ikke her eller
der, saa lidt som det Uendelige begynder her eller der; det Uendelige
er ikke det Endeliges Ophør, betyder ikke, at det Endelige hører op,
% SPÆRRINGEN LØBER OVER BLADSKIFTET: „d e t  E n d e l i g e  e r“ staaer
% som de sidste Ord paa S. 167 og fortsættes med „o v e r a l t ,  f o r d i
% d e t  U e n d e l i g e  e r  o v e r a l t“ paa S. 168. Bekræftet paa
% KB-Exemplaret paa begge Blade. Kolon efter „overalt“ er sat udenfor
% Spærringen her.
men tvertimod, at det gjentager sig uophørlig; \emph{det Endelige er
% --- p. 168 ---
overalt, fordi det Uendelige er overalt}: dette er i kort Begreb
% Spærring bekræftet paa KB-Exemplaret: „U n i v e r s e t s  U e n d e l i g h e d“.
\emph{Universets Uendelighed}.

Hvorledes maa nu Objectiveringen af et i Endelighed uendeligt
Universum nødvendig foregaae? Som Objectet er, maa Objectiveringen
være, og omvendt. At det Uendelige i Virkeligheden er endeligt, og det
Endelige uendeligt, fremlyser paa det Allerklareste af den nyere
Videnskabs Iagttagelseskonst og analytiske Methoder. Den articulerende
Vexelbestemmelse af Endeligt og Uendeligt, der i lige Grad er
charakteristisk for Iagttagelsen og Beregningen, viser hen til Læren om
% Spærring bekræftet paa KB-Exemplaret: „d e n  u e n d e l i g e
% T i l n æ r m e l s e“.
\emph{den uendelige Tilnærmelse}.

Endskjøndt der med de Keplerske Love indførtes legemlige Centra, saa
have Gravitationslærens Tyngdepunkter dog ikke altid materiel
Tilværelse. De bestemmes ved Samvirken af givne Kræfter, og adskille
sig derved fra Epicykeltheoriens tomme Stedbetegnelser. Men foruden
slige ideelle Realiteter har den nyere Astronomie en Mængde fingerede
Størrelser, især Middelstørrelser, Middelsol, Middelanomalie,
Middelbevægelse, Middeltid o. s. frmdl. For at see, hvad en uendelig
Tilnærmelse i denne Henseende har at betyde, vil det være tilstrækkeligt
at fæste Opmærksomheden paa et enkelt Problem, f. Ex. det saakaldte
Keplerske: Centrets Æqvation.

Det er noget Eiendommeligt ved de mathematisk empiriske
Videnskaber --- og Astronomien er en mathematisk empirisk Videnskab --- at de fordre
det Exacte og dog uophørlig maae kæmpe med det Omtrentlige. Jo
strengere Methoden er, desto tydeligere bliver det, at der er en
Uovereensstemmelse mellem det Virkelige og dets Opfattelse, mellem
Objectet og dets Objectivering, som aldrig ganske lader sig hæve. Den
gamle Astronomie blev i al Naivetet staaende ved det Omtrentlige.
Ifølge den Hypothese, at Solen med jævn Hastighed beskriver en Cirkel,
hvis Midtpunkt ligger udenfor Jorden, var det da en temmelig simpel
Opgave at
% --- p. 169 ---
% S. 169 bærer TO Fodnoter, mærkede *) og **). Selve Textmassen paa Bladet er
% kun fire Linier; Resten af Siden er Noterne, adskilte fra Texten ved den
% sædvanlige korte Fodnotestreg.
% NB. Prikkerækkerne i Formlerne staae i Trykken som spatierede Punkter;
% Antallet er gjengivet som trykt, efter Mønstret fra Interpolationsnoten
% til S. 59--61. Produkterne i Nævnerne trykkes med Nielsens
% Multiplicationspunktum og med Luft omkring det („1 . 2“); Luften gjengives
% ikke, jvnfr. „1.2.3“ paa S. 12.
finde Centrets Æqvation.\footnote{% trykt Mærke: *)
Diametren, som gaaer gjennem Jorden er Apsidelinien, dens Endepunkter,
være P og A, Perigæum og Apogæum; Excentriciteten, Afstanden fra Jorden
J til Cirkelens Centrum C, være betegnet ved $\varepsilon$; den Vinkel,
som Linierne fra J og C til Solen S danne med hinanden, ved E; Vinkelen
PCS, Middelanomalien, ved nt: man har da til Bestemmelse af Centrets
Æqvation
\[
\frac{\sin E}{\sin (nt + E)} = \frac{JC}{CS} = \varepsilon
\]
eller i al Omtrentlighed $E = \varepsilon \sin nt$.}
Gaae vi derimod ud fra de Keplerske Love og gjennemføre Analysen i
eiendommelig Retning, saa faae vi en Sinusrække med multiplicerede
Buer, der vel er stærkt convergerende, men dog uendelig:%
\renewcommand{\thefootnote}{**)}%
\footnote{% trykt Mærke: **)
% DENNE NOTE LØBER OVER TO BLADE. Den begynder under Textstregen paa S. 169,
% fylder Resten af det Blad, og fortsættes under Textstregen paa S. 170 fra
% „erholdes saa Ligningen“ til Enden („en ubevidst Viden.“). Bogen gjentager
% ikke Mærket paa Fortsættelsen. Efter Husreglen staaer hele Noten her, ved
% sit Mærke, og Sidemærket for S. 170 sættes ved Textens, ikke ved Notens,
% Overgang.
Ellipsens Storaxe er Apsidelinien og
% TRYKFEIL, SOM NIELSEN IKKE HAR RETTET: Trykken har „Fxcentriciteten“ med
% stort F, hvor E er Meningen. Læst ved høi Forstørrelse paa KB-Exemplaret og
% bekræftet paa Arbeidsscanningen --- begge Copier vise et F med to Tværbjælker
% og ingen Fod. Altsaa Sætterens forkeerte Type. Gjengivet som trykt; læs
% „Excentriciteten“.
Fxcentriciteten e. Jorden, ikke Solen, antages at være i Brændpunktet. E
er Vinkelafstanden mellem Middelsolen, en fingeret Sol, og den virkelige
Sol; m Middelanomalien. Sammenligne vi Udtrykket fra forrige Løsning:
$E = \varepsilon \sin nt$ med første Led af omstaaende uendelige Række,
sees, at Hypothesen om den excentriske Cirkel, idet
% TRYKFEILBLADET retter S. 169, Linie 19 f. o.: „E = 2e, --- læs ε = 2e,“.
% Efter Husreglen staaer Texten her SOM TRYKT, og Rettelsen anvendes ikke.
% Typen er kontrolleret paa KB-EXEMPLARET, netop fordi den fede tobitede
% Arbeidsscanning slører Forskjellen paa stort E og græsk ε: der staaer et
% stort, tre-streget E --- samme Bogstav som „E“ fire Linier ovenfor og
% tydelig forskjelligt fra det lunede ε, Noten *) bruger to Gange.
% Optællingen af „19 f. o.“ passer: Fodnotelinierne tælles med, og Linie 19
% er netop denne.
$E = 2 e$, giver Centrets Æqvation en Værdi, der stemmer med Antagelsen
af den elliptiske Bane, saa længe man, vel at mærke, nøies med det
Omtrentlige. Men skal Omtrentligheden hæves, og Fordringen paa det
Exacte skee Fyldest, da maa Endelighedsgrændsen indenfor Endeligheden
overskrides i det Uendelige. Det er, hvad Analyserne angaaer, et uhyre
Skridt. Hvilke Forvandlinger, den mathematiske Methode har maattet
undergaae, inden det blev muligt at komme fra første Led i omstaaende
Række til alle de følgende, deraf faaer man et samlet Indtryk ved at
kaste et Blik paa den elegante Løsning af Keplers Problem, der findes
% „Cours de Mécanique“ staaer i Kursiv; „Duhamel“ og „II. S. 63--69“ i
% almindelig Antiqua. Sidetallene læste paa KB-Exemplaret: 63--69 (tesseract
% læste „68--69“).
hos Duhamel (\textit{Cours de Mécanique} II. S. 63--69). Som almindeligt
Grundlag henvises til Lagrange's Formel $z = x + \alpha f(z)$, hvoraf
faaes:
% Tælleren trykkes „d.(f(x))²“ med Argumentet (x) i mindre, nedsat Skrift;
% derfor $f_{(x)}$.
\[
z = x + \alpha f (x) + \frac{\alpha^{2}}{1.2}\ \frac{d.(f_{(x)})^{2}}{dx}
+ \ .\ .\ .\ .\ ;
\]
deraf, med Hensyn til Formlerne for Problemet:
\[
u = nt + e \sin u;
\]
\[
r = a (1 - e \cos u), \quad \operatorname{tg} \tfrac{1}{2} \theta =
\sqrt{\frac{1 + e}{1 - e}}\ \operatorname{tg} \tfrac{1}{2} u,
\]
% Her gaaer Noten over paa S. 170; se Bemærkningen ved Notens Begyndelse.
erholdes saa Ligningen:
\[
u = x + e \sin x + \frac{e^{2}\, d\, .\, \sin^{2} x}{1.2\ dx}
+ \ .\ .\ .\ \frac{e^{m}\, d^{m-1} \sin^{m} x}{1.2 \ldots m\ dx^{m-1}}
+ \ .\ .\ .
\]
For at foretage Differentiationen, maa man imidlertid forud kjende,
hvad der er givet i Ligninger som:
\[
(2 \sin z)^{m} \cos \frac{m \pi}{2} = \cos mz - \frac{m}{1} \cos (m-2) z
+ \frac{m (m-1)}{1.2} \cos (m-4) z \ .\ .\ .
\]
\[
(2 \sin z)^{m} \sin \frac{m \pi}{2} = \sin mz - \frac{m}{1} \sin (m-2) z
+ \frac{m (m-1)}{1.2} \sin (m-4) z \ .\ .\ .\ .
\]
% Kildeangivelsen staaer i Trykken CENTRERET paa egen Linie under de to
% Formler. Placeringen gjengives ikke; Linien løber her med i Texten.
(Ramus. Algeb. og Functionsl. S. 202.)
o. s. fremdl., inden man (Duhamel. S. 69) naaer den definitive Form for
den omspurgte Række. For de objective Relationers Skyld maa den
analyserende Tænkning fordybe sig i opererende, combinerende
% Spærring bekræftet paa KB-Exemplaret: „s u b j e c t i v t“. Derimod er
% „Mørke,“ to Linier længere nede IKKE spærret --- spacing.py's Udslag dér
% er Justeringsluft, forkastet paa KB.
Reflexioner. Tankerne i det Objective maae tænkes \emph{subjectivt}, for
at være Tanker; en ubevidst Tankernes Objectivering er en ligesaa
meningsløs Talemaade, som et skinnende Mørke, en seende Blindhed, en
ubevidst Viden.}%
\renewcommand{\thefootnote}{*)}
% --- p. 170 ---
% S. 170 aabner med denne opsatte Formel, som er Fortsættelsen af Texten fra
% S. 169 („men dog uendelig: **)“). Fodnoten **) fra S. 169 løber videre
% under Textstregen paa dette Blad og er skrevet ud ovenfor, ved sit Mærke.
\[
E = 2e \sin nt + \frac{5}{4}\ e^{2} \sin 2nt + \frac{1}{4}\ e^{3}
\left( \frac{13}{3} \sin 3nt - \sin nt \right) + \ .\ .\ .\ .
\]

Her have vi i Beregningens Form et Exempel paa den uendelige
Tilnærmelse.

Fæste vi Opmærksomheden paa Iagttagelsens Opgave, saa viser sig det
Samme fra en anden Side. Ingen enkelt Observation er exact; man kan
gjentage en Observation saa ofte man vil, man begaaer hver Gang en
Feil, skjøndt ikke den samme. Det er ikke Ubehjælpsomhed fra
Iagttagerens Side, ikke tilfældigt Drilleri af Skæbnens Lune; det
ligger i Betingelsernes Natur, i Tilfældighedens objective
Nødvendighed, det er en Lov. Hvor charakteristisk, at det netop er
% TRYKFEIL, SOM NIELSEN IKKE HAR RETTET: Trykken har „dan exacte Videnskab“,
% hvor „den“ er Meningen. Læst paa KB-Exemplaret ved høi Forstørrelse: et
% reent a med Skaal og Stamme, ikke et e; begge Copier vise det. Gjengivet
% som trykt; læs „den“.
dan exacte Videnskab, som er blevet opmærksom paa, at der gives en
% Spærring bekræftet paa KB-Exemplaret: „F e i l l o v“. „Feilenes Lov“
% umiddelbart foran er derimod almindelig Antiqua.
Feilenes Lov, en \emph{Feillov}. Sandsynlighedsreg\-%
% --- p. 171 ---
ningen med Feilloven og de mindste Qvadraters Methode er et betegnende
Vidnesbyrd om den vedholdende Kamp, som den mathematisk empiriske Viden
i alle Retninger maa føre mod det Omtrentlige ved Hjælp af uendelig
Tilnærmelse.

Den astronomiske Objectivering af Verdensbygningen culminerer i
Tilnærmelse; en saadan Culmination er uendelig dialektisk. Fra den ene
Side seer det ud, som om det Virkelige, det i sig Objective, unddrog
sig Iagttagerens Maal i det Uendelige; thi dersom den exacte,
definitive Bestemthed virkelig kunde naaes, saa hørte jo Tilnærmelsen
op af sig selv. Fra den anden Side kommer Virkeligheden med de bestemte
Grændser bestandig nærmere, thi dersom der var en endelig fast Grændse
for Tilnærmelsen, vilde denne jo ikke være uendelig. Det Exacte kan
ikke naaes ved et indskrænket Antal af enkelte Beregninger og
omhyggelig verificerede Iagttagelser; thi den ene Observationsrække
viser tilbage til den anden og selvfølgelig ogsaa den ene Beregning til
den anden; men heri ligger tillige, at Observationer og Beregninger
kunne corrigeres og sikkres ved hverandre indbyrdes. Jo nøiagtigere de
forskjelligartede Observationsrækker blive, desto flere Midler erholder
den analyserende Forsker til yderligere at corrigere dem og nærme sig
det Exacte; den Rest af Ubestemthed, der endda bliver tilbage,
forstyrrer ham ikke ved nogen taaget Omtrentlighed, thi nu veed han
exact, mellem hvilke Grændser den falder, og under hvilke Betingelser,
% Mærket staaer efter „forsvindende“ og FORAN Punktummet, som trykt.
den kan betragtes som forsvindende\footnote{% trykt Mærke: *)
Jvnfr. Mathematik og Dialektik S. 34--56.}.

% AFDELING B SLUTTER IKKE PAA DETTE BLAD. S. 171 ender midt i et Ord ---
% „Det er nødvendigt for Iagt-“ --- og Texten løber videre paa S. 172, hvor
% Afsnittet først slutter med „at objectivere i Afmagt.“; DERPAA staaer den
% korte, centrerede Streg og Overskriften „C. / Solsystemets Oprindelse.“
% Kontrolleret paa Siden selv. Hverken \streg eller \afdeling{C.} hører
% altsaa hjemme i denne Batch: begge skrives af Batch 18 (S. 172--180).
Den iagttagende Videnskab maa have Instrumenter og sætte Priis paa
deres Fuldkommengjørelse; Kikkerter, Teleskoper, Universalinstrumenter
o. s. v. spille en vigtig Rolle ved de exacte Undersøgelser. Det er
% Orddelingen løber over Bladskiftet: „Iagt-“ paa S. 171, „tageren“ paa
% S. 172 (næste Batch).
nødvendigt for Iagt\-%
%  --- C. Solsystemets Oprindelse  (pp. 172–196) ---
%  MATHEMATICAL — work from the image; the OCR is not a witness here.
% --- p. 172 ---
% Bladet aabner MIDT I ET ORD: „Iagt-“ staaer paa S. 171, „tageren“ her.
% Ingen ny Periode; Afdeling B løber videre og sluttes først nede paa denne
% Side med „at objectivere i Afmagt.“
tageren at kjende sit Instrument og forstaae at bruge det. For
Objectiveringens Skyld har Naturphilosophien dog endnu at gjøre
opmærksom paa et uundværligt Instrument, i Sandhed et
Universalinstrument, hvis Medbestemmen ved Objectiveringen ellers saa
let oversees, og det er Bevidstheden selv. Udenfor Bevidstheden gives
der vel en forudsat Virkelighed, men ingen Viden af dens Objectivitet.
Ikke blot i Bestemmelsen af det Objective, men i Objectivitetens Væsen,
det være nu Jordens eller Solens, Melkeveiens eller Cometernes
Objectivitet, er den objectiverende Bevidsthed en væsenlig Factor. Saa
gigantiske Skridt, den exact empiriske Videnskab end har gjort med
Objectiveringen af den virkelige Verdensbygning, saa meget den end i
Theorien har udvidet sin ideelle Magt ved fortsatte Erobringer fra
Stjernehimlens umaalelige Rige: Bygningen selv har den ikke erobret ---
Naturvirkeligheden, det i sig Objective, er i Smaat som i Stort en
Bevidsthedens Skranke, en Skranke, som ved hvert Skridt minder den om,
at den, uendelig langt fra at være det første Objectiverende, der
objectiverer i Magt, maa være tilfreds med sin Lod: at objectivere i
Afmagt.

% Kort, centreret Streg mellem Afdeling B og Afdeling C, midt paa S. 172.
% Kontrolleret paa KB-Exemplaret. Længden gjengives ikke; \streg bruges.
\streg

% Overskriften staaer paa to Linier: „C.“ og derunder „Solsystemets
% Oprindelse.“ --- den sidste sat med halvfed OG spærret Antiqua.
% Kontrolleret paa KB-Exemplaret. Udgaven gjengiver ikke Forskjellen.
\afdeling{C.}{Solsystemets Oprindelse.}

Som et uendeligt System af Verdener er Universet selv uden endelig
Begyndelse. Vorden, Opstaaen og Forgaaen af de enkelte Verdener finder
Sted i Grundsystemet; men dette
% „ikke.“ med Punktum: tesseract læser Komma, men Punktummet staaer
% tydeligt paa KB-Exemplaret (den kjendte Punktum/Komma-Feillæsning).
er og vorder ikke. For nu at undersøge, hvad Videnskaben paa dens
nuværende Trin formaaer at oplyse om Dannelsen
% --- p. 173 ---
% Spærringen af „Urtaagen“, „den immanente Idee“ og „Natur og Aand“ er
% bekræftet paa KB-Exemplaret. „Urtaagen“ er i Trykken delt over to
% Linier: „Urtaa-“ / „gen,“.
af vort Solsystem, fremhæve vi: fra Kjendsgjerningernes Synspunkt den
physisk astronomiske Hypothese om \emph{Urtaagen}, og fra „Begrebets“
Synspunkt Systemets Selvdannelse i Kraft af \emph{den immanente Idee}.
Ved at stille de to hinanden skarpt modsatte Anskuelser under en fælles
Belysning af objectiverende Almagt, kunne vi da overbevise os om, at de
begge ere utilstrækkelige, medens de dog, hver fra sin Side, afgive
væsenlige Bidrag til Sagens Opklaring, naar der spørges om Systemets
Oprindelse, og Grundforholdet mellem \emph{Natur og Aand} tilfulde
erkjendes.

% Indløbende Overskrift, læst paa KB-Exemplaret: sat med halvfed Antiqua
% (ikke spærret), paa samme Linie som den Periode, den aabner.
\runin{a) Urtaagen: den physisk astronomiske Hypothese.} Grunden, hvorfor der ikke kan lægges synderlig Vægt paa de supranaturale
Forudsætninger af Skabelse og skabende Almagt, hvormed de ældre
Physikere og Philosopher, fra Gassendi og Cartesius til Newton og
Leibnitz, gjerne udstyre Grundlaget for deres Naturbetragtning, er
omhandlet i foregaaende Kapitel. Det eneste Udbytte,
Gravitationslæren skulde kunne antages at yde den supranaturalistiske
Anskuelse, maatte i al Fald bestaae deri, at en dybere gaaende Indsigt i
Verdensbygningens Mechanik gav Betragteren fornyet Anledning til at
efterspore og beundre den guddommelige Viisdom. At Skaberen maatte være
en lige saa stor Mathematiker som Physiker, fremgik da fra dette
Synspunkt tydelig nok af den hele Plan. Hvad kunde det vel hjælpe, at
den Almægtige havde „indplantet“ Materien Inertie og Tiltrækning,
dersom han havde glemt at fordele Masserne, at anbringe Sol, Planeter og
Biplaneter hver paa sit bestemte Sted og beregne det første Stød
saaledes, at Banerne kunde blive Ellipser. Ikke nok hermed. Der er
nemlig ikke alene en gjensidig Tiltrækning af Solen og hver enkelt
Planet for sig, men en saa mangfoldigt krydsende Vexelvirkning mellem
Planeterne indbyrdes, at de Perturbationer, som herved foraarsages,
aabenbart true Sy\-%
% --- p. 174 ---
stemet med Undergang. Dette har Skaberen forudseet og fra Begyndelsen
af taget sine Forholdsregler. Vel har han ikke villet gjøre Kloderne
fuldkomment kugleformige, ligesom han i det Hele ikke har skyet det
Indviklede, at hans mechaniske Viisdom ret maatte blive Naturforskeren
aabenbar. Men han har sørget for, at Banerne kun fik ringe Heldning, at
% Formlen er læst paa KB-Exemplaret: T over t = A^{3/2} over a^{3/2},
% Exponenterne satte som opstillede Brøker. De trykte Bogstaver ere
% opreiste Antiqua; det gjengives ikke (Husreglen, se Hovedet).
Omløbstiderne bleve incommensurable:
$\frac{T}{t} = \frac{A^{\frac{3}{2}}}{a^{\frac{3}{2}}}$,
ja, hvad mere er, at Solen fik en saa stor Masse, at den mere end
opveier alle de øvrige\footnote{% trykt Mærke: *) --- eneste Note paa S. 174.
Solen kan endnu, efter at Neptun er opdaget, antages at være omtrent 720
Gange større end alle de Masser tilsammen, der bevæge sig omkring den.}.
Det Sidste især er af overordenlig Betydning; thi den ene Planet kan nu
ikke indvirke forstyrrende paa den anden, uden at Solen i majestætisk
Overlegenhed, saa at sige, lægger sig imellem og holder Orden. Jo mere
man saaledes fordyber sig i de mathematiske Analyser, i den
Mangfoldighed af uendelige Rækker, hvorigjennem en Laplace har udviklet
Perturbationslovene, desto mere voxer Beundringen af Verdensskaberens
guddommelige Mechanik.

Under slige opbyggelige Betragtninger maa det da i høi Grad overraske at
høre Laplace selv udbryde: \textit{Philosophe, montre moi la main qui a
jeté les Planètes sur le tangent de leur orbite!} Hvad betyder denne
Udfordring? Er Mathematikeren da i al sin Viisdom ikke i Stand til at
opdage Skaberens Haand og erkjende hans Viisdom?

% Spærringen af „Alt i Naturen maa gaae naturligt til“ er bekræftet paa
% KB-Exemplaret (Ordet „Naturen“ er der delt: „Natu-“ / „ren“).
% Punktummet holdes udenfor \emph, som Husreglen er.
Der er en Grundsætning, som siger, at \emph{Alt i Naturen maa gaae
naturligt til}. Ved at afvige fra denne Grundsætning vilde
Naturvidenskaben gjøre sig skyldig i en lige saa livsfarlig Modsigelse,
som Bonden i Træet, da han gjennemsaugede den Green, hvorpaa han selv
sad. Men at
% --- p. 175 ---
% Her er „Alt i Naturen gaaer naturligt til“ derimod IKKE spærret;
% kontrolleret paa Siden. spacing.py melder intet.
Alt i Naturen gaaer naturligt til, vil jo med andre Ord sige, at alle i
Naturen virkende Aarsager, uden Undtagelse, maae være naturlige. Der
maa da lige saa nødvendigt være en naturlig Aarsag til Klodernes
Fordeling i Solsystemet, som til deres Bevægelser den Dag idag; lige saa
vel en naturlig Aarsag til Solmassens Overvægt over de øvrige Masser,
som til Forskjellen mellem Stjernetid og Soltid. Det er ikke vilkaarlig,
anmassende Higen ud over Grændsen for det, der kan erkjendes, men
videnskabelig Trang til en sammenhængende Redegjørelse for givne
Kjendsgjerninger, som ligger til Grund for den Verdenshypothese, ved
hvilken Kant først og udførligt\footnote{% trykt Mærke: *) --- første af to Noter paa S. 175.
Allgemeine Naturgeschichte und Theorie des Himmels. Kant's Werke.
Achter Band. Leipzig 1838. S. 224 flgd.}
og siden Laplace i al Korthed%
\renewcommand{\thefootnote}{**)}%
\footnote{% trykt Mærke: **) --- anden af to Noter paa S. 175.
% Noten er sat med Cursiv indtil „Tome second.“; „S. 293 flgd.“ er Antiqua.
% „francaise“ trykkes UDEN Cedille --- baade Arbeids- og KB-Exemplaret vise
% det. Gjengivet som trykt; læs „française“.
\textit{Exposition du Système du Monde. A Paris. L'an IV de la
republique francaise. Tome second.} S. 293 flgd.}%
\renewcommand{\thefootnote}{*)}
have søgt at forklare sig Planetsystemets naturlige Tilblivelse.

Den physiske Astronomie forudsætter, at Materien med dens Kræfter er
„evig“ som Rummet, og lægger saa an paa: ud fra det første Værende at
begribe den Vorden og trinvise Dannelse, som det Nuværende i Tidernes
Løb maa have gjennemgaaet. Det første Værende er at opfatte som en
Tilstand, i hvilken alt Tilværende fra Begyndelsen kan antages at have
befundet sig. Ved en storartet Reduction føres vi da tilbage til en i
Rummet ligelig udbredt, endnu indifferent Materie, en overordenlig tynd,
men uhyre Dampmasse, hvis Tvermaal af Nogle anslaaes til 12000 Billioner
Mile.

Dannelsen begynder nu med, at Delene tiltrække hverandre, hvorved der
fremkommer Fortætning, Rotation og Varme. Dampmassen bliver
kugleformig, ophøiet i Æqvator; Rotationen tiltager, og
Centrifugalkraften voxer: altsammen naturligt! I det ophøiede Bælte af
Kuglemassen vil der
% --- p. 176 ---
saa efterhaanden danne sig Ringe, som løsne sig, afsondres og igjen
sammenrulles; de sammenrullede Masser bevæge sig om deres Axer og om
Hovedmassen: her have vi et første Udkast til Systemet, i al
Naturlighed. Hvad der skete med den store Masse, kan gjentage sig med
de mindre; nogle iblandt disse kunne igjen afsætte een eller flere
Ringe, og disse Ringe afsondres og blive til Maaner. Saturnsringen, der
i Virkeligheden bestaaer af tre tydelig fra hverandre adskilte Ringe, er
et i Udviklingen hæmmet Element, hvoraf man kan lære, hvorledes det i
det Hele er gaaet til med Dannelserne. Da Solmassen netop er den
resterende Hovedmasse, saa indsee vi let den naturlige Grund til dens
Overvægt; kun angaaende Planetbanernes nærmere Beskaffenhed bliver der
en Tvivl tilbage, som Theorien, hvis den skal antages, maa kunne hæve.
Vel maa det i Almindelighed indrømmes, at vi nu ikke længer behøve en
overnaturlig Aarsag for at forklare det første Stød efter Tangenten, da
Svingkraften jo allerede er i den roterende Masse; endvidere følger det
af Loven for Tiltrækningen, at alle Banerne maae blive Keglesnit. Men
af de tre Keglesnit er kun det ene, nemlig Ellipsen, egnet til at give
et sluttet System af Bevægelser; de to øvrige vilde derimod gjøre
Systemet meningsløst, idet Planeterne, naar de bevægede sig enten i en
Parabel- eller i en Hyperbelgreen, efterhaanden vilde fjerne sig fra
Solen uden nogensinde at vende tilbage. Vil man nu ikke skrive det paa
Tilfældets Regning, at der er Fornuft i Solsystemet, saa synes det
unegtelig, som om netop den særegne Begyndelseshastighed, af hvilken de
elliptiske Baner ere betingede, maatte forudsætte en særegen
Viisdommens Foranstaltning og altsaa dog vise hen til en usynlig Haand,
% Trykfeilbladet retter S. 176, Linie 4 f. n.: „Ingriben --- læs
% Indgriben“. Trykken har „Ingriben“ uden d --- læst paa KB-Exemplaret,
% fordi Arbeids-Exemplaret netop her bærer en Blyantsrettelse. Gjengivet
% som trykt. Jvnfr. „Indgriben“ paa S. 177, Linie 1.
en vilkaarlig Ingriben af en overnaturlig Aarsag.

Forsvarerne af den opstillede Hypothese finde det ingenlunde vanskeligt
at imødegaae slige Betænkeligheder. For at afværge, at Banerne skulde
blive Parabler, behøvedes
% --- p. 177 ---
der i Sandhed ingen guddommelig Indgriben. Blandt uendelig mange mulige
Tilfælde er der for Parablen kun et eneste
% Mærket staaer efter „gunstigt“ og FORAN Kolonet, som trykt.
gunstigt\footnote{% trykt Mærke: *) --- eneste Note paa S. 177; fylder Bladets nedre Halvdeel.
Betegnes Radius vector ved r, den halve Storaxe ved a, Intensiteten ved
$\mu$ og Hastigheden ved v, saa er den for alle tre Keglesnit fælles
Formel:
% Tegnet i Parenthesen er et virkeligt Minus (ikke Nielsens ÷); læst paa
% KB-Exemplaret. De trykte Bogstaver ere opreiste Antiqua, hvad Husreglen
% ikke gjengiver.
\[ v^2 = \mu\left(\frac{2}{r} - \frac{1}{a}\right). \]
Da nu a er positiv, negativ eller $\infty$, eftersom Keglesnittet er en
Ellipse, Parabel eller Hyperbel, saa sees heraf, at det maa beroe paa
Hastigheden (Begyndelseshastigheden), hvilken af de tre Keglesnit, der
skal fremkomme. For Ellipsen er $v^2 < \mu\,\frac{2}{r}$, for Hyperblen
er $v^2 > \mu\,\frac{2}{r}$, for Parablen er $v^2 = \mu\,\frac{2}{r}$.
Af de uendelig mange mulige Tilfælde vil altsaa kun eet eneste være
gunstigt for Parablen, medens Antallet af gunstige Tilfælde er uendelig
stort saavel for Hyperbel som for Ellipse.}:
Parablen er, ifølge de i Hypothesen selv indeholdte Forudsætninger, den
allerusandsynligste Form; men at det Allerusandsynligste ikke skeer, er
jo ganske naturligt. Hvad Hyperblen angaaer, da har denne Curve, naar
man betragter Sagen i al Almindelighed, ganske vist en Sandsynlighed for
sig, lige saa vel som Ellipsen; men bringer man de nærmere
Omstændigheder med i Anslag, da er det for det Første et Spørgsmaal, om
Svingkraften i den roterende Masse nogensinde har været stor nok til at
give en Begyndelseshastighed, der kunde svare til Hyperblen. Men om
ogsaa Muligheden heraf indrømmes, hvad saa? Saa bliver i al Fald Følgen
kun den, at nogle Planeter maae i Urtiden have forladt Systemet, og at
vi kun have dem, der fik en ringere Begyndelseshastighed, tilbage.

% Arksignaturen „12“ staaer nederst paa S. 177. Ikke Text.
Gaaer man i det Enkelte med Hypothesens Anvendelse, da er der vel ikke
faa Kjendsgjerninger --- alle Bevægelser fra Vest til Øst,
Axedreiningerne svarende til Solens Om\-%
% --- p. 178 ---
dreining og til Retningerne i Banerne; smaa Heldninger; lille
Excentricitet; de ydre Planeters hurtigere Omdreining, ringere
% „de ydre Planeters“ her er IKKE spærret --- kontrolleret paa
% KB-Exemplaret; spacing.py's Udslag er Justeringen.
Banehastighed\footnote{% trykt Mærke: *) --- første af to Noter paa S. 178.
Ogsaa Spectralanalysen synes at tale til Hypothesens Fordeel. Vi ville
imidlertid ikke undlade at anføre følgende, skjøndt, som det forekommer
os, noget svage Indvendinger af en udmærket fransk Geognost: Mais
l'hypothèse émise par Laplace, il ne la présentait lui-même qu'avec
„defiance“ . . . ses conjectures n'expliquent point la forme elliptique
des orbites planétaires, l'inclinaison de leur axe; elles paraissent en
outre démenties par le mouvement rétrograde des satellites d'Uranus
% Prikkerne ere satte med Mellemrum, som trykt: tre efter „defiance“,
% fire efter „d'Uranus“. Cursiven begynder ved „Enfin“ --- bekræftet paa
% KB-Exemplaret; alt foran er almindelig Antiqua.
. . . . \textit{Enfin la découverte de l'analyse spectrale, qui sera
l'éternel titre de gloire de M. M. Kirchhoff et Bunsen, autorise à
croire que la composition chimique du soleil diffère
% „très-notablement“: Bindestregen falder ved Linieskiftet, saa det kan
% ikke afgjøres, om den er Sætterens Delingstegn eller en virkelig
% Bindestreg. Den ældre franske Retskrivning med Bindestreg er fulgt.
très-notablement de celle des planètes de son système, car cet astre ne
contient, du moins dans ses couches extérieures, ni silice, ni étain, ni
plomb, ni mercure, ni argent, ni or (!).} La Terre par Élisée Reclus.
I. Paris 1870. pag. 19.}
o. s. v. --- som i store Træk tale for dens Gyldighed; men der er dog
ogsaa adskillige Phænomener, der ikke ret ville føie sig efter
% Trykfeilbladet retter S. 178, Linie 5 f. o.: „dem. --- læs den.“
% Trykken har „dem.“ --- læst paa KB-Exemplaret, fordi Arbeids-Exemplaret
% her bærer en Blyantsrettelse. Gjengivet som trykt.
dem. Alt slaaer ikke saa lige til, som Kant antager.

% Anførselen løber over Bladskiftet og sluttes paa S. 179 („Men“).
% Det er ikke en Trykfeil og faaer ingen Note.
„De teleskopiske Planeter, Vesta, Juno, Ceres, Pallas (og den nye
% Trykfeilbladet retter S. 178, Linie 9 f. o.: „hældende --- læs
% heldende“. Trykken har „hældende“ --- læst paa KB-Exemplaret, fordi
% Arbeids-Exemplaret her bærer en Blyantsrettelse. Gjengivet som trykt.
Astræa) med deres gjensidig sammenslyngede, stærkt hældende og mere
excentriske Baner, har man“, siger Humboldt%
\renewcommand{\thefootnote}{**)}%
\footnote{% trykt Mærke: **) --- anden af to Noter paa S. 178.
Kosmos. I. S. 71 flgd.}%
\renewcommand{\thefootnote}{*)}%
, „forsøgt at betragte som en vort Planetsystem i tvende rumlige
Afdelinger adskillende Zone, ja ligesom en Middelgruppe i dette. Efter
% Spærringen af „indre“ og af de to „ydre“ er bekræftet paa
% KB-Exemplaret; det senere „De indre, Solen nærmere Planeter“ er
% derimod almindelig Antiqua, og „nær,“ ligesaa.
denne Anskuelse frembyder den \emph{indre} Planetgruppe (Mercur, Venus,
Jorden og Mars) i Sammenligning med den \emph{ydre} (Jupiter, Saturn og
Uranus) adskillige paafaldende Contraster. De indre, Solen nærmere
Planeter ere af middelmaadig Størrelse, tættere, temmelig eens og
langsomt roterende (i henved 24 Timers Omløbstid), mindre fladtrykte og
paa een nær, nemlig Jorden, aldeles maaneløse. De \emph{ydre}, fra
Solen fjernere Planeter ere be\-%
% --- p. 179 ---
tydelig større, fem Gange mindre tætte, have to Gange hurtigere
Axeomdreining og ere desaarsag stærkt fladtrykte samt maanerige, i
Forhold som 17 til 1, dersom det stadfæster sig, at Uranus virkelig har
% Mærket staaer EFTER Punktummet, som trykt.
6 Satelliter.\footnote{% trykt Mærke: *) --- første af to Noter paa S. 179.
I den nyeste Tid ere kun 4 Maaner udtrykkelig iagttagne.}
Men“ --- tilføier han --- „denne almindelige Betragtning over visse, for
hele Grupper charakteristiske Egenskaber lader sig ikke med lige Ret
anvende paa de enkelte Planeter i enhver af Grupperne; ei heller paa
noget Forhold mellem de enkelte kredsende Himmellegemers
% Spærringen af „Afstand fra Centrallegemet“ er bekræftet paa
% KB-Exemplaret.
\emph{Afstand fra Centrallegemet} eller paa deres absolute Størrelse,
deres Tæthed, Omdreiningstid, Excentricitet, Baners Heldning og
% ANFØRSELEN LUKKES IKKE. Citatet aabnes med „denne almindelige
% Betragtning“ og faaer intet sluttende “ efter „Axe**).“ --- begge
% Exemplarer vise det, og Trykfeilsbladet nævner det ikke. Gjengivet som
% trykt; Sidens (og Batchens) „/“-Balance staaer derfor +1.
% Mærket staaer FORAN Punktummet, modsat S. 179 Linie 4 f. o.
Axe%
\renewcommand{\thefootnote}{**)}%
\footnote{% trykt Mærke: **) --- anden af to Noter paa S. 179.
Som Exempel paa, at vi ikke kjende nogen mechanisk Naturlov, der, „lig
den skjønne Lov, som binder Qvadraterne af Omløbstiderne til de tredie
Potenser af Banernes store Axe, gjør Planetlegemernes nysnævnte 6
Elementer eller deres Baners Form afhængige enten af hinanden eller af
Middelafstandene“, anføres, at den fra Solen fjernere Mars er mindre
baade end Jorden og Venus, at Saturn er mindre end Jupiter og dog meget
større end Uranus; at Uranus selv synes at være tættere end Saturn; at
vi, hvad den indre Gruppe angaaer, til begge Sider af Jorden finde baade
Venus og Mars mindre tætte end denne selv. Vel aftager Rotationstiden
(om Axen) i det Hele med Afstanden fra Solen; men dog er den ved Mars
større end ved Jorden, ved Saturn større end ved Jupiter. Saturn
befinder sig midt imellem Jupiter, hvis Rotationsaxe næsten staaer
lodret, og Uranus, hvis Axe derimod næsten falder sammen med Banen;
o. s. v.}%
\renewcommand{\thefootnote}{*)}%
.

% TRYKFEIL: „antage ig“ for „antagelig“ --- l'et mangler, og der staaer
% et rent Ordmellemrum i dets Sted. Baade Arbeids- og KB-Exemplaret vise
% det, altsaa Sætterens Feil (udfaldet Type); Trykfeilsbladet nævner den
% ikke. Gjengivet som trykt; læs „antagelig“.
% Arksignaturen „12*“ staaer nederst paa S. 179. Ikke Text.
Hypothesen kan i det Hele være antage ig, om de end ikke lykkes at gjøre
fuldstændigt Rede for alle Enkeltheder og fjerne alle Betænkeligheder.
Det Eiendommelige i den Samvirken af Betingelser og Kræfter, hvorunder
hine første Dannelser ere foregaaede, taber sig i Urtidens Mørke,
utilgængeligt for specielle Undersøgelser. Hypothesen har
Sandsynligheden for sig, saalænge der ikke kan paavises
% --- p. 180 ---
Kjendsgjerninger, der staae i aabenbar Strid med dens Antagelse. Men
sæt endog, at uigjendrivelige Indvendinger omsider førte til en ligefrem
Forkastelse af Hypothesen, vilde saa deraf følge, at Naturvidenskaben
blev nødt til at tye til supranaturale Forudsætninger og gjøre Laan hos
Theologien? Ingenlunde! Kan Hypothesen ikke forsvares, saa skeer der Et
af To: enten bliver Spørgsmaalet om Solsystemets Dannelse stillet i Bero
som ubesvarligt, eller der danner sig en ny og mere tilfredsstillende
Hypothese. I begge Tilfælde vil Videnskaben med lige urokkelig Vished
holde sin Grundsætning fast, den uafviselige Grundsætning, at Alt i
Naturen maa gaae naturligt til. Dette er det Afgjørende;
Forklaringsmaaden selv er ikke det Afgjørende. Om ogsaa Solsystemets
trinvise Tilblivelse ved Overgange fra Tilstand til Tilstand ud fra en
Urtilstand lod sig fastsætte ved en Theorie, der var lige saa
tilforladelig som Gravitationslæren: saa vilde Theorien om en Række
Tilstandsforandringer endnu ikke være nogen Forklaring af det Heles
Oprindelse; men vistnok vilde en saadan Theorie ved at klare vort Blik
for de naturlige Aarsagers indre Sammenkjædning afgive et stort reelt
% Spærringen af „Begrebet: Natur“ er bekræftet paa KB-Exemplaret.
Bidrag til Bestemmelsen af \emph{Begrebet: Natur}.

% INDLØBENDE OVERSKRIFT MED FORKERT BOGSTAV. Trykfeilbladet retter
% S. 180, Linie 13 f. n.: „c) Systemets --- læs b) Systemets“; ogsaa
% Indholdet og Rækkefølgen (a) paa S. 173) kræve b). Trykken har „c)“,
% læst paa KB-Exemplaret --- Arbeids-Exemplaret bærer her en
% Blyantsrettelse af Bogstavet. Gjengivet som trykt, ikke omnummereret.
% Overskriften er sat med halvfed Antiqua, ogsaa Bogstavet „c)“.
\runin{c) Systemets Selvdannelse: den immanente Idee.} Medens Physiken med sin Hypothese udelukkende holder sig til Phænomener
og betingede Aarsager, gaaer den speculative Immanenslære den omvendte
Vei og tager afgjørende Sigte paa det Ubetingede,
% Spærringen af „Ideen“ er bekræftet paa KB-Exemplaret.
\emph{Ideen}. Den Hegelske Deduction af Solsystemet er i denne
Henseende mærkværdig.

Naturen er mere end en Vrimmel af Kjendsgjerninger; den er Ideen selv i
det udvortes Værende eller i Form af Andetværen. Det ligger i den
guddommelige Idee at oplade sig, for at være Subjectivitet og Aand. Fra
dette Synspunkt maa Solsystemet opfattes; thi dets „absolute Mechanik“
er den Form, hvori Naturideen har realiseret sig selv. De abstracte
% S. 180 ender med det hele Ord „Idee“; Sætningen løber videre paa
% S. 181 („fyldt med Materie ...“). Ingen Orddeling over Bladskiftet.
Former, Rum og Tid, har den sig begribende Idee
% --- p. 181 ---
% S. 181 begynder FLUGTIG (uden Indrykning) og fortsætter altsaa Sætningen
% fra S. 180: „... fyldt med Materie, hvis Væsen er Tyngde“. Kontrolleret paa
% Siden selv (KB-Exemplaret): første Linie staaer helt ude ved Marginen.
% De græske Bogstaver α) β) γ) i denne Opregning ere læste paa KB-Exemplaret
% (Cursiv-Græsk i Texten); OCR'en giver a, B/ß og y/7, som er Feillæsninger.
fyldt med Materie, hvis Væsen er Tyngde; thi Tyngden bestaaer just deri, at
Materien, søgende sit Hvilepunkt udenfor sig, bevæger sig fra Sted til Sted og
ved Sted og Bevægelse ophæver Modsigelsen af Tid og Rum. Bevægelsen er
tredobbelt: α) den mechaniske, udenfra meddeelte, er eensformig; β) den halvt
betingede, halvt frie Bevægelse, nemlig Faldbevægelsen, viser, at et Legemes
Adskillelse fra dets Tyngde endnu er tilfældig, medens Bevægelsen dog allerede
tilhører Tyngden selv; γ) den ubetinget frie Bevægelse, „Himlens store
Mechanik“, sees i Gravitationen.

% „Gravitationen“ er spærret --- bekræftet paa KB-Exemplaret. „Himlens store
% Mechanik“ ovenfor er derimod almindelig Antiqua.
„\emph{Gravitationen} er den materielle Legemligheds sandfærdige og bestemte
Begreb, der er realiseret til Idee. Den almene Legemlighed grunddeler sig
væsenlig i særskilte Legemer og slutter sig som Bevægelsens phænomenale
Tilværen sammen til Momentet af Enkelthed eller Subjectivitet, som herved
umiddelbart er et System af flere Legemer.“\footnote{% trykt Mærke: *)
Vorlesungen über die Naturphilosophie. Berlin 1842. S. 94 flgd.} --- Der kan
ikke med Grund spørges om, hvorfra disse Legemer komme; de ere reale Momenter
af Ideens Væren: de ere til, thi i dem har Ideen Tilværelse. De selvstændige
Legemer, af hvilke Solsystemet bestaaer, referere sig væsenlig til hverandre,
ere tunge, men holde sig selv i denne Relation og sætte deres Eenhed i et Andet
udenfor sig. Saaledes er Fleerheden ikke længer ubestemt, som ved Stjernerne,
men Forskjellen er sat; og Bestemtheden beroer alene: paa
% Hele Vendingen er spærret paa Siden, ogsaa „særskilt Centralitet“ ---
% bekræftet paa KB-Exemplaret; „og paa“ er derimod ikke spærret.
\emph{absolut almindelig Centralitet}, og paa \emph{særskilt Centralitet}. Af
disse to Bestemmelser følge de Bevægelsens Former, i hvilke Materiens Begreb er
udfyldt. Bevægelsen falder paa de relative Centrallegemer, som i sig ere
Stedets almene Bestemthed: og dog er deres Sted tillige ubestemt, forsaavidt
det har sit Centrum i et Andet; denne Ubestemthed maa da ligeledes have
Tilværen, medens det i
% --- p. 182 ---
og for sig bestemte Sted kun er Eet. Det er derfor ogsaa de særskilte
Centrallegemer ligegyldigt, paa hvilket Sted de ere; og det kommer tilsyne
saaledes, at de søge deres Centrum, d. e. forlade deres Sted og sætte sig paa
et andet Sted. Det Tredie er dette: umiddelbart taget, kunde disse Legemer jo
være fjernede lige langt fra deres Centrum; vare de det, saa vare de ikke
fjernede fra hverandre. Bevægede de sig derhos tillige alle i samme Bane, saa
vare de slet ikke forskjellige fra hverandre: de vare da Eet og det Samme, det
ene kun en Gjentagelse af det andet, og deres Forskjellighed selvfølgelig kun
et tomt Ord. Det Fjerde er dette, at, idet de i forskjellig Fjernhed fra
hverandre forandre deres Sted, saa vende de i sig tilbage gjennem en Curve; thi
kun derved fremstille de deres Selvstændighed imod Centrallegemet --- ligesom
ogsaa deres Eenhed med Midtpunktet derved, at de bevæge sig omkring det i samme
Curve. Som selvstændige imod Centrallegemet holde de sig paa deres Sted og
falde ikke længere ind i Centret. Det Dobbeltsidige i Forholdet er, at de
særskilte Legemer samtidig sætte sig et Centrallegeme, og at de blive satte af
Centrallegemet: Centret er meningsløst uden Peripherie, og Peripherien uden
Centret.

Paa dette Grundlag blive de Keplerske Love at deducere af den rene Idee. Hvad
for det Første Banens Form angaaer, saa er Cirkelen kun en Bane for den slet og
ret eensformige Bevægelse; men i den frie Bevægelse maae Rum og Tid gjælde
efter deres Forskjellighed, og herved ophæves den abstracte, simpelthen
eensformige Hastighed. Tænkeligt, som man siger, er det vel, at en eensformig
til- og aftagende Bevægelse kan skee i en Cirkel. Men denne Tænkelighed eller
Mulighed er kun en ubestemt overfladisk Forestilling. Cirkelen er den i sig
tilbagevendende Linie, i hvilken alle Radierne ere lige; d. v. s. den er
fuldkommen bestemt ved Radius, eller: Radius er den hele Bestemthed.
% --- p. 183 ---
% Afsnittet løber ubrudt over Bladskiftet; S. 183 begynder uden Indrykning.
% Spærringerne paa denne Side --- „Forskjellighed“, „Differens“, „to“,
% „Ellipse“, „uafhængig“, „Een“, „Eet Hele“, „Sectoren“ --- ere alle
% bekræftede paa KB-Exemplaret. Det senere „Forskjellighed i Radierne“ og
% „empirisk Størrelse“ ere derimod almindelig Antiqua og faae ingen \emph.
Men i den frie Bevægelse, hvor Rums- og Tidsbestemmelser i deres
\emph{Forskjellighed} træde i et qvalitativt Forhold til hinanden, maa dette
Forhold i det Rumlige selv nødvendig fremtræde som en \emph{Differens} indenfor
samme, hvad der udtrykkelig fordrer \emph{to} Bestemmelser. Fordringen er, at
forskjellige Buer gjennemløbes i samme Tid. Skulle Buerne være forskjellige,
saa maae Radierne ogsaa være det: af den accelererende Bevægelse følger
Forskjellighed i Radierne. Derved bliver Formen for den i sig tilbageløbende
Bane væsenlig en \emph{Ellipse}; --- første Keplerske Lov.

Den abstracte Bestemthed, som er Cirkelens Særkjende, viser sig ogsaa saaledes,
at den Bue eller Vinkel, som indesluttes af tvende Radier, er en af dem
\emph{uafhængig}, empirisk Størrelse. Men i den ved Begrebet bestemte Bevægelse
maa Afstanden fra Centrum og Buen, der gjennemløbes i en vis Tid, sammenfattes
i \emph{Een} Bestemmelse, udgjøre \emph{Eet Hele} (Begrebets Momenter ere ikke
tilfældige for hinanden); saa fremkommer der da en Rumsbestemmelse af to
Dimensioner, \emph{Sectoren}. Buen er paa denne Maade væsenlig en Function af
Radius vector, og fører, som den i lige Tider ulige, Radiernes Ulighed med sig.
At Rumsbestemmelsen formedelst Tiden bliver en Fladebestemmelse, hænger sammen
med, hvad der gjælder om Faldet, at nemlig samme Bestemthed maa forekomme: den
ene Gang som Tid i Roden, den anden Gang som Rum i
% TRYKFEILBLADET retter S. 183, Linie 9 f. n.: „Qvadratet*) --- læs
% Qvadratet*).“ --- altsaa et manglende Punktum efter Fodnotemærket.
% KB-EXEMPLARET har INTET Punktum efter „*)“: der staaer „Qvadratet *) Kun
% at“. ARBEJDSEXEMPLARET (Google-Scanningen) viser derimod et lille Mærke i
% Grundlinien efter Parenthesen --- Blyanten er altsaa kommet dette Sted i
% Forkjøbet, ganske som ved S. 160. Gjengivet SOM TRYKT, uden Punktum.
% Noten strækker sig over Bladskiftet: den begynder ved Foden af S. 183 og
% slutter ved Foden af S. 184 (uden gjentaget Mærke); den er efter Husreglen
% sat i sin Heelhed her ved Mærket.
Qvadratet\footnote{% trykt Mærke: *)
Den løbende Tid er negativ Eenhed, medens Rummet breder sig som et Værende.
Rummet maales med Tiden saaledes, at Eenhederne af Rum give Tælleren, den
valgte Tidseenhed Nævneren (Tid i Roden). Idet nu Rumseenhederne under
Faldbevægelsens accelererede Hastighed reflectere sig i Tidseenheden,
fremkommer et Tal, der er Tidseenhedens eget Product, altsaa Qvadratet af Tid,
et Maal for Rummet. Dette er den fatteligste Omskrivning, vi formaae at give,
af de gaadefulde Ord, om hvilke Hegel (Anfr. Skr. S. 88) siger
% TRYKFEILBLADET retter S. 184, Linie 2 f. n.: „ans --- læs \emph{aus}“.
% KB-EXEMPLARET trykker utvetydigt „ans“ (n, ikke u) i den tyske Sætning;
% ARBEJDSEXEMPLARET har over samme Bogstav et lettere Blyantsmærke, der gjør
% n til u --- atter Blyanten før os. Gjengivet SOM TRYKT: „ans“.
% Ingen Colon efter „siger“, og Tysken er almindelig Antiqua, ikke Cursiv;
% begge Dele læste paa KB-Exemplaret.
Diess ist der Beweis des Gesetzes des Falls ans dem Begriffe der Sache.} Kun at
% --- p. 184 ---
her det Qvadratiske i Rummet er, ved Bevægelsesliniens Tilbagevenden i sig
selv, begrændset til Sector. Disse ere, som man seer, de almindelige Principer,
paa hvilke den anden Keplerske Lov, at der i lige Tider afskæres lige Sectorer,
speculativt beroer.

% „Faldet“, „for de frie Maal“, „Qvadratet“ og „Cubus“ ere spærrede ---
% alle fire bekræftede paa KB-Exemplaret.
Denne Lov angaaer dog kun Buens Forhold til Radius vector; og Tiden er derhos
den abstracte Eenhed, i hvilken de forskjellige Sectorer sammenlignes, fordi
den er det Determinerende som Eenhed. Men det videre gaaende Forhold er et
Forhold af Tid, ikke som Eenhed, men som Qvantum overhovedet, som Omløbstid,
til Banens Størrelse eller, hvad der er det Samme, til Afstanden fra Centrum.
Som Rod og Qvadrat saae vi Tid og Rum forholde sig til hinanden i
\emph{Faldet}, den halvfrie Bevægelse, der vel fra den ene Side er bestemt ved
Begrebet, men dog paa den anden Side er udvortes bestemt. Men i den absolute
Bevægelse, i Sphæren \emph{for de frie Maal}, faaer enhver Bestemthed sin
Totalitet. Som Rod er Tiden en blot empirisk Størrelse, og i qvalitativ
Henseende kun abstract Eenhed. Som Moment i den udviklede Totalitet derimod er
den tillige en derved bestemt Eenhed, en Totalitet for sig, producerer og
refererer sig deri til sig selv. Som det i sig Dimensionsløse kommer Tiden i
sin Production dog kun til formel Identitet med sig selv, til \emph{Qvadratet},
Rummet derimod som det positive Udenforhinanden, kommer til Begrebets
Dimension, til \emph{Cubus}. Dette er den tredie Keplerske Lov, Forholdet af
Afstandenes Cuber til Tidernes Qvadrater; --- „en Lov, som just derfor er saa
stor, fordi den er saa enfoldig og fremstiller Sagens egen Fornuft“.

% Orddelingen løber over Bladskiftet: „Be-“ paa S. 184, „stemmelse“ paa S. 185.
Hvad Legemerne selv angaaer, saa have de, til Be\-%
% --- p. 185 ---
stemmelse af deres forskjellige Natur, de Hovedsider, der udgjøre Momenterne af
deres Begreb. Et af dem er det almene Centrum for den abstracte Relation til
sig selv. Imod dette Extrem staae de \emph{umiddelbare}, udenfor sig værende,
% „umiddelbare“, „Enkeltheder“ og senere „uselvstændige“ ere spærrede ---
% bekræftede paa KB-Exemplaret. De øvrige Kald fra spacing.py paa denne Side
% („selv.“, „Enkeltheder“ i anden Forekomst) ere Justerings-Artefakter og
% forkastede.
centrumløse \emph{Enkeltheder} ligeledes i Form af selvstændig Legemlighed. De
særskilte Legemer staae saavel under Bestemmelsen af Væren-udenfor-sig som af
I-sig-væren; de ere Centra for sig og referere sig dog til hiint første som til
deres væsenlige Eenhed. Centrallegemet fremstiller den abstract rotatoriske
Bevægelse. Centret skal være et Punkt; men det er, idet det er Legeme, tillige
udstrakt og bestaaer af Søgende. Denne uselvstændige Materie, som
Centrallegemet har i sig selv, fordrer, at det maa rotere om sig selv. Thi de
uselvstændige Punkter, tilmed fjernede fra Centret, have intet fast bestemt
Sted --- de ere kun faldende Materie og saaledes kun bestemte efter een
Retning. Den øvrige Bestemthed mangler; ethvert Punkt maa altsaa indtage alle
de Steder, det kan indtage. Det i og for sig-Bestemte er kun Centrum, det
øvrige Udenfor-hverandre er ligegyldigt, thi herved er kun Stedets Afstand
bestemt, ikke Stedet selv. Denne Bestemmelsens Tilfældighed kommer til Existens
paa den Maade, at Materien forandrer sit Sted, og udtrykkes ved Solens
Roteren-i-sig om sit Midtpunkt. Denne axedreiende Bevægelse er ingen
Stedforandring\footnote{% trykt Mærke: *)
At der, ifølge „Himlens Mechanik“, ikke vel kan gives nogen Axedreining uden
translatorisk Bevægelse, kan let blive overseet, naar man kun seer paa den
„absolute Idee“.}; thi alle Punkter beholde samme Sted imod hverandre. Det Hele
er en hvilende Bevægelse. --- De \emph{uselvstændige} Legemer, som tillige have
en fri Existens, og ikke udgjøre sammenhængende Dele af et med Centrum begavet
Legemes Udstrækning, have ogsaa Rotation, men ikke om sig selv, thi de have
intet Centrum. De rotere altsaa om et Midtpunkt, der
% --- p. 186 ---
tilhører et andet Legemindivid, fra hvilket de ere udstødte. Deres Sted er
overhovedet dette eller hiint; og denne det bestemte Steds Tilfældighed udtrykke
de saa ved deres Rotation. Men deres Bevægelse er en træg og stiv Bevægelse om
Centrallegemet, idet de bestandig blive i den samme Stilling imod det, som
f. Ex. Maanen i sit Forhold til Jorden. De uselvstændige himmelske Legemer
danne den Side af Begrebet, der kaldes Særskiltheden; deri ligger, at de med
deres Forskjellighed falde ud fra hinanden, thi i Naturen existerer
% „Eet“, „cometare“, „den lunare“ og „planetare“ ere spærrede --- alle læste
% paa KB-Exemplaret. Det senere „Den cometare Sphære“ og „den lunare derimod“
% ere derimod almindelig Antiqua og faae ingen \emph.
Særskiltheden som Tohed, ikke som i Aanden i Form af \emph{Eet}. For det Første
er da det Moment sat, at den hvilende Bevægelse bliver en urolig Bevægelse, en
Sphære af Udsvæven eller Stræben ud over den umiddelbare Tilværen imod et
Hiinsides. Dette Moment af Væren-udenfor-sig er selv Substansens Moment som en
Masse, en Sphære; thi ethvert Moment erholder her sin egen Tilværen. Denne
Sphære er den \emph{cometare}, og udtrykker en Hvirvel, en bestandig Staaen paa
Spring til at løse sig op og sprede sig i det Uendelige og Tomme.
Udsvævelsesextremet bestaaer nødvendigt i at nærme sig Centrallegemets
Subjectivitet engang og saa at vige for Repulsionen. Men denne Andetværen har
umiddelbart sin Modsætning i sig selv, en Modsætning, som er: at søge sit
Midtpunkt, sin Hvile: den ophævede Fremtid, Forbigangenheden som Moment; dette
er \emph{den lunare} Sphære, der ikke peger hen paa en Udsvæven fra den
umiddelbare Tilværen, men paa det Tilblevne, eller paa For-sig-væren, paa
Selvet. Den cometare Sphære er desaarsag kun henviist til den umiddelbare
axedreiende, den lunare derimod til det nye, i sig reflecterede Midtpunkt, til
Planeten. --- Den \emph{planetare} Sphære er endelig den, der udtrykker
Relation baade til sig selv og til Andet; den er axedreiende Bevægelse og har
dog sit Midtpunkt udenfor sig. Planeten har sit Centrum i sig selv; men det er
kun et relativt; den har ikke sit absolute Centrum i sig, den er forsaavidt
ogsaa
% --- p. 187 ---
uselvstændig. Planeten har begge Bestemmelser i sig og fremstiller dem begge
som Stedforandring. Som selvstændig viser den sig da saaledes, at dens Dele
selv forandre Sted med Hensyn paa den Stilling, de have til den rette Linie,
som forbinder det absolute og det relative Centrum; dette begrunder Planetens
rotatoriske Bevægelse. Planeten er, som det Tredie, den Slutning, hvori vi have
det Hele. Den solare, planetare, lunare, cometare Natur! denne Fiirhed af
Himmellegemer danner det fuldendte System af den fornuftige Legemlighed. Den
planetare Natur er Totaliteten, Modsætningernes Eenhed, medens de andre kun
fremstille dens isolerede Momenter; den planetare Natur er den fuldkomneste,
ogsaa med Hensyn paa Bevægelsen, som her alene kommer i Betragtning.

% spacing.py finder ingen Spærring paa S. 187, og KB-Exemplaret bekræfter det:
% hele Siden er almindelig Antiqua. Latinen er derimod Cursiv (læst paa KB).
% „Aristoteliske“: tesseract giver „Åristoteliske“; Å findes ikke i denne
% Retskrivning.
Grundsvagheden i denne Deduction er iøinefaldende. Ved eensidig at fæste
Opmærksomheden paa Ideen og den ideelle Causalitet kommer Tænkeren i ethvert
afgjørende Vendepunkt til at oversee de virkelige Betingelser og mechaniske
Aarsager. Her argumenteres af Grunde; det er den gamle Aristoteliske Maneer.
Men saa lidt som Naturaarsager lade sig erkjende uden Grunde, lige saa lidt
lade Naturgrunde sig erkjende uden Aarsager: \textit{vere scire est per causas
scire}, det er et frugtbart Ord af Baco.

% Orddelingen løber over Bladskiftet: „Sidebestem-“ paa S. 187, „melser“ paa
% S. 188.
Ikke desto mindre vil Enhver, der forstaaer at læse, faae et Indtryk af, at en
storartet Bestræbelse med en stor Tanke rører sig og aander i denne Deduction,
en Bestræbelse, der gaaer ud paa at ophæve Stykværket, føre de phænomenale
Bestemmelser tilbage til Ideens Eenhed og anskue Verdensaltet i Fornuftens Lys.
At Naturen er Idee, vil jo med andre Ord sige, at den ikke er Sammenstykning af
utallige Fragmenter, men eet uendeligt i sig sammenhængende Væsen med een
Grund, at alle enkelte Kjendsgjerninger, alle særlige Phænomener og Love ere
Sidebestem\-%
% --- p. 188 ---
melser af det ene Grundvæsen, reale Momenter i det ene indholdsrige Begreb.

% Indløbende Overskrift, læst paa Siden selv (KB-Exemplaret): „c)“ i
% almindelig Antiqua, „Natur og Aand.“ HALVFED --- samme Sats som de øvrige
% a)/b)/c)-Overskrifter. Denne Udgave gjengiver ikke Forskjellen mellem
% halvfed og spærret; begge blive \emph.
\runin{c) Natur og Aand.} Physiken og Immanensphilosophien ere enige om, at
Naturen er et i sig sammenhængende Hele, og Verdenssystemet en virksom
Aabenbarelse af den evige Fornuft; men fra ingen af Siderne bydes der os nogen
klar eller indtrængende Fremstilling af, hvad der egenlig skal forstaaes ved
% „Natur“ her er spærret --- bekræftet paa KB-Exemplaret.
Ordet \emph{Natur}.

Physiken viser simpelthen Spørgsmaalet fra sig; den handler om Alt, hvad der er
i Naturen; men hvad Naturen selv er, omhandler den ikke. Naar den gaaer tilbage
til Urtaagen, forstaaer den det saaledes, at Urtaagen oprindelig var i Naturen,
at der var Natur i Urtaagen; men at Naturen i sit Væsen er en Urtaage, eller at
Urtaagen er Naturen selv, paastaaer den ingenlunde. Immanensphilosophien tager
Spørgsmaalet op og besvarer det resolut ved den Erklæring, at Naturen er det
Absolute, Ideen selv i sandselig udvortes Tilværen; men Besvarelsen er
illusorisk.

% „Ingen Idee uden Medium!“ er spærret over to Linier --- bekræftet paa
% KB-Exemplaret.
Immanensphilosophiens Mysticisme viser sig deri, at den uden videre gjør det
Absolute til Idee og Ideen til det Absolute. Ideen for sig som Idee er
ingenlunde det Absolute. Den rene Idee er det huleste, afmægtigste,
tvetydigste Væsen, naar den skal fastholdes i egen Absoluthed, adskilt fra
ethvert tilsvarende Medium. \emph{Ingen Idee uden Medium!}

Ifølge ovenstaaende Deduction seer det jo næsten ud, som var det Planeterne
selv i concret Legemlighed, der af logisk Interesse for et System af tre
Slutninger tilsagde hverandre at møde, ikke paa samme Sted, men, for
Forskjellighedens Skyld, hver paa sit Sted, og kom overeens om at sætte sig et
fælles Centrum, Solen; men da simple Planeter dog i Grunden maae antages at
være for dumme til at indsee Begrebets logiske Pointe, saa kan det ikke være
% --- p. 189 ---
Planeterne simpelthen, men Planetideen, den som Planet i Planeter existerende
Idee, der begriber det Hele. Vel er Materien ogsaa begribende, men i Uklarhed,
i Dunkelhed; den klare Begriben, den rene Tænken er i Rum og Tid. Det er da
ikke den existerende Materie med Egenskaber og Kræfter, men en dialektisk
Proces mellem det speculative Rum og den reflecterende Tid, der har gjort de
% „faae“ uden Accent, læst paa KB-Exemplaret; tesseract giver „faaé“.
Keplerske Love nødvendige. For at to saadanne Begribende som Tiden og Rummet
kunde i Begrebets Navn faae hver Sit, maatte Differensen imellem dem absolut
komme frem, og Planetbanerne blive Ellipser; Tiden maatte blive Rod, at det
Qvadratiske i Rummet kunde begrændses til Sector. Men Tiden vilde ikke nøies
med at være Rod; den stræbte videre frem, den vilde qvadreres, vilde ad Ideens
Vei naae op til at blive Qvadrat. Først da Tiden opnaaede „den formelle
Identitet“ at være Qvadrat, medens Rummet avancerede til Begrebets Dimension,
til Cubus, kom Solsystemet i Orden: den solare, cometare, lunare, planetare
Natur gik op i Ideen.

Som Væsensindhold er Ideen fælles for Natur og Aand; men, det maa udtrykkelig
mærkes, Ideen staaer i et andet Forhold til Aanden end til Naturen; for Aanden
er den et Afledet, thi den bestemmes og formes af almægtig begribende Villie;
for Naturen er den et Oprindeligt, thi den bestemmer og former det naturlige
Hele. Med dette Synspunkt vende vi nu tilbage til den kosmogoniske Hypothese,
til Urtaagen.

% Græsk Underoverskrift, indløbende og SPÆRRET; læst paa KB-Exemplaret ved
% stærk Forstørrelse: α (Cursiv-Alfa), ikke „a“ eller „&“, som OCR'en giver.
% Overskriften er kontrolleret paa Siden selv, ikke i Indholdet.
\subrunin{α) Oprindelse og Begyndelse.} Urtaagen, hvormed der begyndes, maa
selv have en Begyndelse; thi hvis det antages, at den i Rummet udbredte Materie
er til fra Evighed af, da have vi, som ovenfor er viist, i dette
Begyndelsesløse et Tilværende uden Tilblivelse, et Taagephænomen uden Væsen, et
Existerende uden Grund, en drømt Forudsætning, et naturløst, unaturligt Første.
Men med et naturløst, unaturligt Første kan ingen Naturdannelse begynde. Ifølge
% --- p. 190 ---
% Denne Side bærer TO Fodnoter, mærkede *) og **) ved Foden.
% Anførselstegnene: Citatet brydes af „siger han“, lukkes ved „maa“ og
% aabnes paany ved „paa een Gang“ --- to Aabninger og to Lukninger, altsaa
% i Balance. Det er ikke Bogens bekjendte Gjenaabnings-Vane.
Kant skal det hele uendelige Verdensrum fra Begyndelsen af være fyldt med en
Materie, hvis Egenskaber og Kræfter ligge til Grund for alle Forandringer.
„Denne Grundmaterie, en umiddelbar Følge af den guddommelige Tilværelse, maa“
siger han, „paa een Gang være saa rig, saa fuldstændig, at Udviklingen af dens
Sammensætninger i Evighedens Afløb kan udbrede sig over en Plan, som i sig
indeslutter Alt, hvad være kan, som intet Maal antager, kort, som er
uendelig.“\footnote{% trykt Mærke: *)
Allgem. Naturgeschicht. u. Theor. des Himmels. 7tes Hauptstück. S. 319.} Den
theologisk-supranaturalistiske Forestilling om et enkelt Skabelsesmirakel er
den videnskabelige Nødhjælp, Kant endnu%
\renewcommand{\thefootnote}{**)}%
\footnote{% trykt Mærke: **)
1755. Kritik der reinen Vernunft udkom først 1781.}%
\renewcommand{\thefootnote}{*)}
har i Baghaanden.

Nu skulde det vel synes, som om et Mirakel, der ligger en Evighed tilbage i
Tiden, et eneste almindeligt Mirakel, der oven i Kjøbet er skeet med
Nødvendighed og har gjort det Hele af paa een Gang, endda nok kunde tolereres i
Videnskaben --- blot til en Begyndelse; og dog er Modsigelsen let at opdage.
Den existerende Materie kan ikke ved et umiddelbart Magtsprog udspringe af
Selvmagten; Materien tilhører Naturen som Natur og maa have en naturlig
Begyndelse. Materiens Tilblivelse er Gjenstandsmagtens første reale Mulighed,
og Magtens Udslag alene bestemt ved Atomernes Realbegreb. Begyndelsen skeer,
idet Begrebet realiseres. Af hvilken Beskaffenhed den saaledes tilblivende
Materie iøvrigt vil være, hvilke Egenskaber dens Atomer og Masser fra
Begyndelsen af ville antage, afhænger alene af de indenfra bestemmende
Magtbegrebers Eiendommelighed; men disses Eiendommelighed afhænger igjen af den
eiendommelig bestemte Plan, det indholdsrige Anlæg, Ideen.

Ideen er altsaa Naturen iboende; Udviklingen er i sin Begyndelse som i sin
Fortsættelse en betinget-nødvendig
% --- p. 191 ---
Realisering af bestemte Muligheder, af indenfra bestemmende Anlæg. Naturen
giver ikke sig selv sine Muligheder, men den realiserer de givne Muligheder;
den meddeler ikke sig selv Anlæg, men den udvikler de meddeelte Anlæg:
% Hele Vendingen fra „al“ til „Anlæg“ er spærret --- bekræftet paa
% KB-Exemplaret; ogsaa „al“ selv er spærret. Ligeledes „Begyndelse“ og
% „Oprindelse“ nedenfor, og „endskjøndt“ og „fordi“ i næste Afsnit.
\emph{al Naturudvikling er, væsenlig seet, en Udvikling af meddeelte Anlæg}.
Herved føres vi med videnskabelig Conseqvens fra Begrebet \emph{Begyndelse} til
Begrebet \emph{Oprindelse}.

Begyndelsen er naturlig, men Oprindelsen overnaturlig. Naturens meddeelte Anlæg
forudsætte en Anlægsmeddelelse; denne Anlægsmeddelelse stammer ikke fra
Nogetsomhelst i Naturen; den er alene at søge i det, som --- ikke udvortes, men
indvortes --- er over Naturen, altsaa i det Overnaturlige selv, i Aanden.
Naturudviklingens Idee kunde umulig være bestemmende for et sig udviklende
Naturhele, dersom en saadan bestemmende Idee ikke oprindelig havde en indre
Bestemthed i sig selv, dersom den i sit Væsen ikke indeholdt et System af
strenge, ufravigelige Grundbestemmelser. Hvorfra har da Ideen disse
Bestemmelser? En Naturidee kan lige saa lidt bestemme sig selv, som en
Naturtanke kan tænke, eller et Naturbegreb begribe sig selv. Det Alt tænkende,
Alt begribende, alle Ideer bestemmende Magtvæsen er Aanden som Aand, det
almægtige Selvvæsen. Er det da sandt, at Gjenstandsmagten ikke staaer endelig
udenfor Selvmagten; er det sandt, at alle Gjenstandsbestemmelser i
Magtfordoblingen gjennem Realbegreber og Anlæg vise hen til Enemagtens absolut
overgribende Selvbestemmelse: saa er det ogsaa sandt, at denne Verden er skabt,
at Skabelsesacten er Aandens almægtige Villiesact; saa er det sandt, at hele
den virkelige Natur har en naturlig Begyndelse, ikke \emph{endskjøndt}, nei
just \emph{fordi} den har en overnaturlig Oprindelse.

% Græsk Underoverskrift, indløbende og spærret; Bogstavet er læst paa
% KB-Exemplaret: β, ikke „B“, „ß“ eller „P“, som OCR'en giver. Overskriften
% er kontrolleret paa Siden selv, ikke i Indholdet.
\subrunin{β) Grundvillie og Dannelsesdrift.} Endskjøndt Physiken i det Hele er
tilbøielig til at forudsætte en evig
% --- p. 192 ---
% TRYKFEIL: der staaer „i Universitet samme Ælde“, hvor Meningen kræver
% „i Universet“. Begge Exemplarer trykke „Universitet“ (kontrolleret paa
% KB-Exemplaret), og Bogens eget Trykfeilblad berører ikke S. 192.
% Gjengivet som trykt; læs „Universet“.
% „høist“: begge OCR-Vidner læse „heist“ (den kjendte ø→e-Feillæsning);
% ø'et staaer tydeligt paa KB-Exemplaret.
værende Materie, er den dog langt fra at ville tillægge de mange Verdener i
Universitet samme Ælde; den antager tvertimod, at de forskjellige
Verdenssystemer begynde deres Dannelse til høist forskjellige Tidspunkter.
% „Livsudvikling, saaledes“ med KOMMA: Arbejdsexemplaret viser et Punktum-
% lignende Mærke (Tærskningens Fortykkelse), men KB-Exemplaret har utvetydigt
% Komma.
„Ligesom vi i vore Skove“, siger Humboldt, „see den samme Dannelsesmaade
samtidig paa alle Vegetationens Trin, og fra denne Betragtning, fra denne
Coexistens modtage et synligt Indtryk af den fremadskridende Livsudvikling,
saaledes erkjende vi ogsaa de forskjellige Stadier af Stjernedannelsen i den
store Verdenshave. Fortætningsprocessen, som Anaximenes og den hele ioniske
Skole lærte, synes her at foregaae ligesom for vore Øine.“

% „eiendommelig“/„eiendommeligt“ med i, ikke j; læst paa KB-Exemplaret.
% tesseract giver „ejendommelig“, hvilket er en Feillæsning.
For at forklare sig den naturnødvendige Sammenhæng mellem de mange til
forskjellige Tider indtrædende Verdensdannelser, antager Kant, at det uendelige
Hele fra Begyndelsen af maa have et ved Skaberens Viisdom bestemt
Fortætningernes og Bevægelsernes Midtpunkt, et første materielt Centralpunkt,
og at saa Dannelse følger paa Dannelse efterhaanden som den ordnende Bevægelse
udbreder sig og faaer Bugt med Forvirringen i det oprindelige Chaos. Her er i
denne Synsmaade en dobbelt Skævhed. For at bevare den naturlige Aarsagsrække
overseer Kant, at der til hver eiendommelig Verdensdannelse maa svare et
eiendommeligt Anlæg, altsaa en eiendommelig Skabelsesact. Endvidere overseer
han, at en Skabelsesact, videnskabelig forstaaet, ingenlunde afbryder, men
tvertimod muliggjør en lovbestemt Sammenhæng mellem Verdensaarsagerne
indbyrdes: de mange Verdener i Himmelrummet udgjøre et System af Systemer,
fordi alle Verdensideerne udgjøre Momenter af een og samme Universets Idee,
fordi alle skabende Villiesacter udgaae fra een guddommelig Villie, en evig
Grundvillie.

At Solsystemets Indretning ikke umiddelbart kan hidrøre fra Guds Haand, men maa
være et Product af mechanisk virkende Aarsager, vil Kant bevise paa følgende
Maade. Alle Planeterne bevæge sig i samme Retning omkring Solen;
% --- p. 193 ---
% Arkesignaturen „13“ staaer ved Foden af S. 193. Den er ikke Text.
% spacing.py finder ingen Spærring paa S. 193; KB-Exemplaret bekræfter det.
denne Omstændighed kan kun forklares af den Mechanisme, hvoraf Bevægelserne ere
opstaaede. Hvis Gud med egen Haand havde sat Planeterne i Bevægelse, vilde hans
Valg, da her ikke er den ringeste Bevæggrund, hvorfor han skulde binde sig til
een eneste Bestemmelse, have udviist mere Frihed i allehaande Afvigelser og
Forskjelligheder. Hvorfor falde Planetbanerne næsten i eet og samme Plan, men
dog ikke ganske? Hvorfor nærme de sig til at være Cirkler og ere dog lidt
excentriske? I alt Sligt erkjende vi Naturens sædvanlige Fremgangsmaade; den
stræber efter noget Bestemt, men naaer det ikke ganske; den maa anvende
forskjellige Mellemaarsager, men Mellemaarsagerne bringe den til at bøie lidt
af til Siden, saa at den aldrig rammer Sagen paa en Prik. Dersom det var sandt,
hvad en Philosoph har sagt: at Gud bestandig udøver Geometrien, og dersom dette
klarlig kunde sees paa alle Naturlovenes Veie, da vilde der spores en ganske
anden Fuldkommenhed og geometrisk Nøiagtighed i det almægtige Ords umiddelbare
Værker.

Sandhedskjernen i det Kantiske Raisonnement er: at virkelige Naturdannelser
skyldes virkende Naturaarsager, og at det særlig er de mechaniske Aarsager, som
i deres Vexelvirkninger optage de udvortes Tilfælde med alle Tilfældigheder.
Paa det Tilfældige vil Kant netop kjende det Naturlige; herimod er Intet at
indvende. Men den Maade, hvorpaa han i sin „Himlens Theorie“ stiller
Opfyldelsen af de almindelige Love imod Planen selv i den lovgivende Villie,
røber en eensidig Reflexion. Med en menneskelig Plan og dens Udførelse gaaer
det vel undertiden saaledes, at tilstødende Omstændigheder, uberegnelige
Tilfælde, uforudseelige Hindringer kunne bevirke saa gjennemgribende
Forandringer i Udførelsen, at den realiserede Plan aldeles ikke eller dog kun
tildeels svarer til Værkmesterens Øiemed. Men en guddommelig Plan kan umulig
være udsat for slige uforudseelige Hindringer.
% --- p. 194 ---
Der er i Solsystemets Anlæg, i Urtaagen, en indre Mangfoldighed af Anlæg til
Anlæg, af Muligheder, af Betingelser for Betingelser, af Aarsager til Aarsager.
Sammentrækningen er Aarsag til Fortætning og Rotation, Rotationen igjen en
Aarsag til Delenes Sammenhobning ved Æqvator o. s. fremd. Aarsagen til Jordens
Dannelse og dens Bevægelse omkring Solen, Aarsagen til Maanens Dannelse og dens
Bevægelse omkring Jorden, alle slige Aarsager med deres Betingelser, der Tid
efter Tid udvikle sig, ere fra Begyndelsen af indesluttede i Anlægets
Grundaarsag, saa at den Villie, der har villet Anlæget og tilfulde vidst, hvad
den selv vilde, har i Et og Alt villet den hele Udvikling. Idet de mange
udvortes Aarsager med de mange baade indvortes og udvortes Betingelser
efterhaanden træde i Virksomhed, medvirke jo vistnok de indtrædende
Omstændigheder, men der bliver Intet, aldeles Intet overladt til den blotte
Tilfældighed: Lovene lægge Beslag paa de mødende Tilfælde, at de alle uden
Undtagelse punktlig maae tjene til Planens Udførelse.

% „Grundvillien“ og „Dannelsesdriften“ ere spærrede --- bekræftede paa
% KB-Exemplaret; det tidligere „skabende Grundvillie“ i samme Afsnit er
% derimod almindelig Antiqua.
Den skabende Villie bruger Tilfældene, vil det Tilfældige, fordi den i Alvor
vil det Naturlige; den vil en naturnødvendig Dannelse af Delene, fordi den med
Aandens Frihed vil en naturnødvendig Dannelse af det Hele. Det Heles Udvikling
gjennem Millioner Gange Millioner af Secler er i alle Enkeltheder en naturlig
Fuldbyrdelse af den skabende Grundvillie. Mellem \emph{Grundvillien} i det
Indre og de mechaniske Aarsagers Mangfoldighed i det Ydre ligger
\emph{Dannelsesdriften}, aandig og dog naturlig, et væsenligt Bindeled mellem
Natur og Aand.

% Græsk Underoverskrift, indløbende og spærret; Bogstavet er læst paa
% KB-Exemplaret: γ, ikke „y“ eller „7“, som OCR'en giver. Overskriften er
% kontrolleret paa Siden selv, ikke i Indholdet.
\subrunin{γ) Teleologie og Mechanisme.} Den gamle Strid om Forholdet mellem de
mechaniske Aarsager og de saakaldte Hensigtsaarsager maa nu i Principet kunne
hæves. Hvad er vel Hensigten med Solsystemet? En naturlig Udvikling af Alt,
hvad der indesluttes i dets oprindelige Anlæg:
% --- p. 195 ---
% Arkesignaturen „13*“ staaer ved Foden af S. 195. Den er ikke Text.
% „høiere“ og „Hældning“ ere læste paa KB-Exemplaret. Trykfeilbladet retter
% „Hældning“ til „Heldning“ paa S. 136, men berører ikke S. 195.
Anlægets Udvikling er Øiemedets Opnaaelse. Udviklingen skeer i Kraft af
mechaniske Aarsager; men naar Udviklingen er Maalet, og de mechaniske Aarsager
Midlerne, saa kan der jo i Sagen selv ingen Strid være mellem det Mechaniske og
det Teleologiske. Kun naar Betragtningen er hildet i endelige Forestillinger,
kan Striden fremkomme. Enten drager man da „Hensigterne“ ned i det Endelige og
forseer sig paa Vilkaarligheden eller man holder de reelle Aarsager fast i
Udstykningen og forseer sig paa Nødvendigheden.

At drage Hensigterne ned i det Endelige er at fremhæve visse for Livet, især
for Menneskene, gavnlige, heldbringende Indretninger og opfatte dem saa
enkeltviis og udskilte af det Heles Sammenhæng, at det seer ud, som havde
Skaberen vilkaarlig havt Opnaaelsen af slige Resultater for Øie. Det er en
Teleologie saa egoistisk, at den vexler og bliver forskjellig i Forhold til det
Standpunkt, hvorpaa Betragteren er stillet. Hvorfor staaer vel Fuldmaanen
høiere paa Himlen om Vinteren end om Sommeren? Efter Nordboernes Teleologie maa
det være Skaberens Hensigt, at den skal opklare de lange Vinternætter. Skade,
at der ogsaa er lange Nætter uden Fuldmaane! Hvorfor har Jordaxen faaet en
Hældning imod Banen? Det maa fra Jordteleologiens Standpunkt naturligviis være
skeet i den alvise Hensigt, at Adskillelsen mellem hede, tempererede og kolde
Zoner med de Afvexlinger og Forskjelligheder, som derved fremkomme, skulde
afgive en uudtømmelig Rigdom af Betingelser for det organiske Livs Udvikling
paa Jorden. Desværre maae de, der leve under Linien, og de, der boe omkring
Polerne, trods den alvise faderlige Hensigt, føle sig som Stedbørn.

At holde Aarsagerne fast i Udstykningen er at gjøre den betingede Nødvendighed
til Et og Alt. Hertil svarer den aldeles vrange, skjøndt hos adskillige
Physikere ikke
% --- p. 196 ---
% Latinen er Cursiv, med „inqvisitio“ og „tanqvam“ (qv, ikke qu) --- læst
% paa KB-Exemplaret. „Baco's“ med Apostroph, som trykt.
usædvanlige, Forestilling, at et Naturphænomen umulig kan udtrykke et vist
relativt Øiemed, naar det beviislig fremgaaer som lovbestemt Virkning af
mechaniske Aarsager. Er den teleologiske Betragtning paa Endelighedens Trin
eensidig subjectiv, da er den blot mechaniske tilvisse ligesaa eensidig
objectiv. Som forskende Videnskab har Physiken, hvor endelig end dens
Mechanisme kan være, imidlertid den umiskjendelige Fordeel, at dens
Undersøgelser ere frugtbare, og dens Begrundelser videnskabelige. Den endelige
Teleologie derimod, hvad enten den forøvrigt beraaber sig paa den indre eller
den ydre Hensigtsmæssighed, er og bliver uvidenskabelig; om den gjælder Baco's
Ord: \textit{Causarum finalium inqvisitio sterilis est, et tanqvam virgo Deo
consecrata nihil parit.}

Fra Uendelighedens, Principets, d. e. fra Aandens Synspunkt forholder Sagen sig
anderledes. Naturudviklingen er først fuldendt, naar den i Et og Alt
tilfredsstiller Aanden. Vistnok tilfredsstiller den Aanden paa ethvert Trin,
men dog vel kun forsaavidt som ethvert Trin betegner et Skridt mod
Fuldendelsen. Det Spørgsmaal, om Solsystemet har sit Maal alene i sin
mechaniske Form eller er bestemt til at afgive Grundlag for Opnaaelse af andre
og høiere Naturøiemed, har videnskabelig Betydning og falder sammen med det
Principspørgsmaal: hvilke Former Naturudviklingen kræver, for at det hele til
Anlæget svarende Begrebsindhold kan blive grundig, alsidig og fuldstændig
realiseret. Hermed er Overgangen fra den abstract mechaniske til den
elementære, concretere Natur motiveret.

% SLUTNINGEN PAA FØRSTE DEEL. Texten paa S. 196 løber IKKE til Foden: den
% ender omtrent tre Femtedele nede med „concretere Natur motiveret.“, hvorpaa
% følger blank Sats og dernæst en kort, centreret Streg omtrent tre Fjerdedele
% nede paa Bladet; Resten af Siden er blank. Stregen er her sat DOBBELT (to
% tynde, parallele Linier tæt sammen) og altsaa kraftigere end de sædvanlige
% Afdelings-Streger; den Forskjel gjengives efter Husreglen ikke --- \streg
% tjener overalt. Kontrolleret paa begge Exemplarer.
% S. 196 bærer ingen Arkesignatur, og Anden Deel begynder IKKE paa dette
% Blad: „Anden Deel. / Den elementære Natur.“ og „Første Kapitel. /
% Elementerne.“ staae først paa S. 197. S. 196 er altsaa den rene Slutning
% paa Første Deel.
\streg


% ======================================================================
%  Anden Deel — Første Kapitel: Elementerne   (printed pp. 197–249)
% ======================================================================

\deel{Anden Deel.}{Den elementære Natur.}
% Mellem Deel-Overskriften og Kapitel-Overskriften staaer i Satsen en kort,
% centreret Streg. Efterprøvet paa KB-Exemplaret for S. 1 og S. 197.
\streg

\kapitel{Første Kapitel.}{Elementerne.}

%  --- A. Elementernes Oprindelse  (pp. 197–214) ---

% --- p. 197 ---
% \deel og \kapitel staae allerede i Skelettet ovenfor og gjentages ikke her.
% KONTROL AF KAPITELOVERSKRIFTEN: læst paa KB-Exemplaret (Farvescanningen).
% Hovedet paa Siden lyder „Første Kapitel. / Elementerne.“ --- IKKE „Tredie
% Kapitel.“, som Indholdet har. Skelettets Valg staaer altsaa ved Magt.
% Siden aabner reent: Første Deel slutter „... concretere Natur motiveret.“
% nederst paa S. 196; intet af den løber over paa S. 197.
% Mellem „Den elementære Natur.“ og „Første Kapitel.“ staaer en kort centreret
% Streg. Den hører hjemme mellem \deel og \kapitel, altsaa OVENFOR dette
% Fragment; den er derfor ikke sat her --- ganske som ved Første Deel, S. 1.
\afdeling{A.}{Elementernes Oprindelse.}

% De tre spærrede Vendinger nedenfor ere alle kontrollerede paa
% KB-Exemplaret. spacing.py's RUN „forgaae, derfor“ er derimod FORKASTET:
% KB viser almindelig Antiqva.
Ifølge den sædvanlige Forestillingsmaade maa det vel næsten synes modsigende
at tale om Elementernes Oprindelse: Element er jo netop det, som er uden
Oprindelse. Alle Ting kunne opløses i deres Elementer, men Elementerne selv
kunne ikke opløses; de kunne ikke forgaae, og derfor heller ikke opstaae: de
ere evige. Ved nøiere at undersøge Betydningen af \emph{Stof og Element}, af
\emph{Stof og Form} og derefter prøve Stoffets væsenlige Andeel i
\emph{Tingenes Substantialitet} ville vi imidlertid erholde en ny Bekræftelse
paa, hvad Læren om den mechaniske Natur allerede har godtgjort, at
Elementerne, for at have en naturlig Tilværelse, maae have en overnaturlig
Oprindelse --- i Kraft af Aanden.

% Indløbende Overskrift, læst paa Siden selv (KB-Exemplaret): „Stof og
% Element.“ er sat med HALVFED Antiqva, ikke spærret; „a)“ foran den staaer
% med almindelig Skrift. Husmakroen \runin gjengiver hele Overskriften eens.
\runin{a) Stof og Element.} Den græske Philosophie begynder som
Naturphilosophie i al Umiddelbarhed med en Lære om Stoffet. Det
Charakteristiske ved denne Begyndelse er,
% --- p. 198 ---
at Stof og Princip gaae i Eet; Modsætningen af det Reelle og det Ideelle, af
Natur og Aand, er paa dette Bevidsthedstrin endnu ikke til: det er den rene
Immanens i speculativ Uskyldighed. De i Immanensen skjulte Modsigelser dukke
saa efterhaanden op snart under denne snart under hiin Form; de forestillede
Phænomener og de tænkte Begreber geraade i Strid med hinanden; Opgaverne blive
bestandig vanskeligere, og Forhandlingerne mere indviklede.

% „Forvandlingens“ er spærret paa KB. „Ionernes“ (som spacing.py ikke fandt,
% men som Blyanten har understreget i Arbeids-Exemplaret) er derimod
% almindelig Antiqva paa KB --- ligesom „Heraklit“ og „Ild“ længere nede,
% der ogsaa kun bære Blyantsstreger. Ingen af dem er sat som Fremhævelse.
Den almindelige Tanke, der bærer Ionernes Forestilling om et Urstof, er
\emph{Forvandlingens} Grundtanke. Hvorledes Stoffet for sig nærmere blev
forestillet, som Vand, som Luft, som Ild o. s. v., er uvæsenligt;
Stofforvandlingen selv er Hovedsagen. Den Energie, hvormed Forvandlingstanken
er udviklet hos Heraklit, gjør det allerede klart, at intetsomhelst Materielt
har evig Bestaaen, at Stoffet i dets Umiddelbarhed er en Modsigelse. Ilden er
ham ikke et tilværende Stof, ved hvis Fortætning eller Fortyndning der kunde
fremkomme andre tilværende Stoffer; dette er i al Fald kun en sandselig
Forestillingsmaade. Det Afgjørende er, at enhver Ting lige saa meget er et
Ikkeværende som et Værende. Den synlige Ild, der kan slukkes, har i sin Væren
en Ikkeværen; thi hvorledes skulde den ellers kunne slukkes. At en Ild
slukkes, vil sige, at den gaaer over fra at være til ikke at være; og saaledes
med alt Tilværende. Men Ilden selv, der i Sandsetingenes uophørlige Vorden,
Strømmen og Vexlen, er det sande almene Grundvæsen, denne Ild tændes ikke og
slukkes ikke; det er en Tankens Ild, hvis Eenhed er Forskjel, og hvis Forskjel
er Eenhed; det er Verdensfrembringeren, den evig levende Ild, der i sig
% SÆTTERFEIL i Græsken, gjengivet som trykt: der staaer „ἀποσπεννύμενον“
% med π, hvor Heraklits Ord har β („ἀποσβεννύμενον“). Bogstavet er læst paa
% KB-Exemplaret ved høi Forstørrelse og har hverken Løkke eller Underlængde;
% det er samme Type som π i „ἁπτόμενον“ og „ἀπο-“, og adskiller sig tydelig
% fra det β, der staaer i „βασιλεὺς“ paa næste Side. Begge Exemplarer vise
% det, altsaa Sætterens Feil. Læs „ἀποσβεννύμενον“.
% Skriftsnittets ϰ- og ϱ-Former gjengives efter Husreglen med κ og ρ.
antænder Maal og udslukker Maal (ἁπτόμενον μέτρα καὶ ἀποσπεννύμενον μέτρα).
At der paa denne Maade, idet hver Deel af det Hele ikke blot er modsat andre
Dele, men et i sig selv Modsat, maa vise sig Strid, Splid og Spænding overalt
i Naturen, er en Selvfølge. Det Lige bliver et Ulige; det
% --- p. 199 ---
Ulige føier sig sammen. Strid er alle Tings Fader og Herre
% Græsken læst paa KB-Exemplaret. „βασιλεὺς“ har gravis, ikke akut ---
% Trykken sætter gravis foran sluttende Parenthes, se ogsaa „τὸ κενὸν“
% nedenfor, „τὸ στερεὸν“ S. 200 og „μορφὴ“ S. 207.
(πόλεμος πάντων μὲν πατήρ ἐστι πάντων δὲ βασιλεὺς); Livet er en blandet Drik;
Verdens Strenge ere spændte, ligesom Buens og Lyrens.

% Det tostavelses „er“ er SPÆRRET paa KB (spacing.py finder som ventet ikke
% Toboghavsord). Blyantsstregerne under „Eleaternes“, „Benegtelse“,
% „Vorden“, „Opstaaen og Forgaaen“ og „af Intet kommer Intet“ ere ikke Text.
Imod Heraklits gjennemgaaende Opløsning af alt Værende stod Eleaternes skarpe
Benegtelse af al Vorden, Opstaaen og Forgaaen. Skulde Vorden lade sig tænke,
da maatte det enten være det Lige, der kom af det Lige, eller det Ulige, der
kom af sit Ulige. Det Første er umuligt, thi det Lige \emph{er} --- og vorder
ikke --- det Lige. Det Sidste er ogsaa umuligt. Skulde Noget fremkomme af sit
Ulige, da maatte det fremkomme af et Andet, hvori det ikke var, altsaa af
Ikkeværen, af Intet; men af Intet kommer Intet.

% „Elementer“ og „Jord, Vand, Luft og Ild“ ere spærrede paa KB. „Bevægelser
% i Rummet“ og „qvalitativt uforanderlige“, som Blyanten har understreget,
% ere det IKKE.
Imellem de skarpt modsatte Synsmaader forsøgte Empedokles en Mægling. Paa den
ene Side gjør han Eleaterne Indrømmelser og paastaaer med Parmenides, at
Opstaaen og Forgaaen, d. v. s. en qvalitativ Forandring i det oprindelige Stof
er utænkelig; paa den anden Side vil han dog ikke lukke Øinene for den
Kjendsgjerning, at Tingene i Verden, ligesom den hele Verden, ere en bestandig
Vexel undergivne. Der bliver ham da intet Andet tilbage, end at henføre
Omskiftelsesphænomenerne til Bevægelser i Rummet, til Forbindelser og
Adskillelser af uoprundne, uforgængelige og qvalitativt uforanderlige
Substanser eller \emph{Elementer}. Af saadanne Elementer maa der nødvendig
gives flere uligeartede, naar de mangfoldige Ting med deres forskjellige
Egenskaber deraf skulle kunne afledes. Saaledes opstod Læren om de fire
Elementer: \emph{Jord, Vand, Luft og Ild}, en Grundvold for Oldtidens og
Middelalderens Naturbetragtning, som den nyere Physik har opløst.

% spacing.py's „Element“ paa denne Side sigter til denne Linie; KB viser
% almindelig Antiqva. Kaldet er FORKASTET.
Paa en anden Maade blev Forestillingen Element opfattet i Atomlæren. For at
forklare den virkelige Vorden gik Leukipp og Demokrit ud fra Modsætningen
mellem det Fyldte (τὸ πλῆρες) og det Tomme (τὸ κενὸν), mellem
% --- p. 200 ---
Atomet og Rummet. Jord, Vand, Luft og Ild bestaae af Atomer; Atomerne selv ere
altsaa de egenlige Elementer. I ethvert Atom er Stoffet selvstændigt og fast
(τὸ στερεὸν). Atomerne ere uden Vorden; der er i dem intet Ikkeværende: de ere
udeelbare. Der kan ikke tillægges dem særlige Egenskaber; de ere det
umiddelbart, positivt Værende. Det Eneste, hvorved de adskille sig indbyrdes,
er deres Figur, Størrelse, Vægt, Ordning og Leiring, lutter qvantitative
Forskjelligheder. Hvor en Forbindelse af Atomer finder Sted, er det Tomme
altid med som Mellemrum. Af Atomerne sammensættes Tingene. En Ting opstaaer,
naar der danner sig et Complex af Atomer; den forgaaer, naar Atomforbindelsen
løser sig op; den forandrer sig, naar Atomernes Leiring og Stilling vexler;
den voxer, naar nye Atomer indtræde i Forbindelsen; den aftager, naar Atomer
udskilles uden at erstattes. Hvad Bestanddelene angaaer, ere alle Legemer
lige; hvad Egenskaberne angaaer, kunne de derimod være høist ulige.

Uden at opholde os ved Enkeltheder, som mere angaae de individuelle
Forestillingsmaader end Grundtanken i det Hele, og særlig høre ind under
Philosophiens Historie\footnote{% trykt Mærke: *) --- eneste Note paa Bladet,
% helt indenfor S. 200. Titlen staaer med almindelig Antiqva, ikke kursiv.
Jvnfr. Zellers klare og fortræffelige Fremstilling. Die Philosophie der
Griechen. 1ster Theil. Tübingen. 1856.}, kunne vi af det Anførte allerede see,
hvorledes det ved de første naturphilosophiske Forsøg forholder sig med
Immanensen. Hos Heraklit er Tanken uforvansket gjennemgaaende. Den evig
levende Ild, der er Tingenes Fornuft, er ogsaa Fornuften i Sjælen. Væsenet i
Forvandlingerne er kun et Sandsevæsen, forsaavidt det er et Tankevæsen: det
Tænkende i Naturen og det Tænkende i Mennesket er det samme Almene. Medens
Parmenides, for at holde Tankevæsenet reent som det forskjelsløse Ene, lader
Sandseverdenen for sig afgive et tomt, bedragerisk Skuespil af det
% --- p. 201 ---
Ikkeværende, forliger Heraklit derimod det Ikkeværende med det Værende og
luttrer dem begge i Tænkningens Ild, thi det Sandselige afgiver kun Stof til
uendelig Forbrænding. I Forvandlingen kommer det Stofagtige overalt tilsyne
med en Værens Umiddelbarhed som Jord, som Luft, som endelige Ting; men at de
endelige Ting kun ere Tilsyneladelser af det uendelige Tankevæsen, sees deraf,
at Ikkeværen, det Negative, er med overalt. Det Værende er --- i
\emph{Vorden} --- et Ikkeværende: det er umiddelbar Immanens.

Anderledes, naar Stoffet faaer absolut Selvstændighed og bliver Element.
Elementet er intet Tankevæsen; det er netop Tankens haarde Modstand, en for
Tanken uigjennemtrængelig Skranke. Og dog giver Tænkningen ikke uden videre
tabt; den griber Elementet i Forestillingen; Forestillingen bestemmes ved
Sandsning; Sandsning ved Sandseindtryk; Sandseindtryk ved det umiddelbart
Qvalitative; under denne Forudsætning kan det Oprindelige lægges ind i
elementære Qvaliteter, og Empedokles paavise Nødvendigheden af fire Elementer.
I hvilken Grad Forestillingen herved faaer Overhaand over Tænkningen,
fremlyser umiddelbart deraf, at det Værende og det Ikkeværende fastholdes ved
Siden af hinanden saaledes, at de begge, hver paa sin Maade, faae en
umiddelbar Tilværen. Det er et Fremskridt hos Atomikerne, at de ikke vilde
begynde med sandselige Qvaliteter. De vilde begribe Qvaliteterne ved at
deducere dem af Atomernes Forbindelser; de anstrengte sig til det Yderste for
at gjøre de atomistiske Stofelementer til rene Tankemomenter. Atomet er ikke
et ligefrem Sandseligt; det er et kategorisk Udtryk for, hvad Tanken
udtrykkeligt fordrer, et absolut Værende, et Værende, hvis eneste Qvalitet er
at udelukke al Ikkeværen ved at have det Tomme, det Ikkeværende, til Grændse.
Men med alt dette er Tankevæsenet dog oprindelig fornegtet, og Tanken blevet
sig selv et Fremmed. I Forvandlingens Sted kommer nu
% Orddelingen løber over Bladskiftet: „For-“ paa S. 201, „andringens“ paa
% S. 202. Spærringen løber med over; begge Halvdele ere kontrollerede paa KB.
\emph{For\-%
% --- p. 202 ---
andringens} Kategorie til at spille Hovedrollen: Leiring, Ordning,
Adskillelse, Forbindelse betegne lutter Forandringer i det Udvortes. Med en
saadan Udvorteshed er Immanensen i Fare. Alt, ogsaa det Usandselige, bliver
sandseligt; og dog maa der forudsættes et Tænkeligt. Demokrit er ligesaa lidt
som Empedokles i Stand til at aflede Tanken af Elementet eller paavise nogen
væsenlig Forskjelseenhed af Sandsning og Tænkning. Skal Tankevæsenet hævdes
og Immanensen bevares, maa det skee ad en anden Vei.

% Indløbende Overskrift, læst paa Siden selv (KB-Exemplaret): halvfed Antiqva,
% „b)“ med almindelig Skrift.
\runin{b) Stof og Form.} Den simpleste Udvei til Forlig mellem Elementer og
Tankevæsen maa fra et physisk populært Standpunkt vel synes at være den, som
er anviist af Anaxagoras. Anaxagoras antager, hvad det Materielle angaaer, et
Chaos af uligeartede Elementer, som, hvert i sig, ere eensartet-qvalitativt
% Trykken har intet Komma efter Parenthesen; gjengivet som trykt (begge
% Exemplarer). „(Homøomerier)“ staaer i Apposition til „Gulddele“ o. s. v.
bestemte (Homøomerier) Gulddele, Kjøddele, Knokkeldele o. s. v., uendelig
deelbare, i uendelig Blanding; hvad det Ideelle angaaer, forudsætter han en
ordnende Forstand, Aanden (νοῦς), alvidende, almægtig, det fineste, reneste,
absolut enkelte, udeelbare Væsen, der i alle Væsener er og bliver sig selv
ligt. Paa Grundlaget af denne Dualisme kan Immanensen forsaavidt bevares, som
det jo kan forudsættes, at Aanden er med i Alt, bevæger og gjennemtrænger Alt;
at de mechaniske og teleologiske Aarsager udfylde hinanden saaledes, at de
virkelige Dannelser foregaae naturlig ved Delenes Adskillelse og Forbindelse,
medens den altordnende, altbevægende Aand organiserer det Hele og bygger en
Verden deraf, der svarer til Fornuftens, Viisdommens, Øiemed. Men ved nærmere
Betragtning vil man dog snart indsee, at en Dualisme af denne Natur maa løse
sig op i Modsigelser. Allerede den Antagelse, at der udenfor Aanden,
uafhængigt af den, fra Evighed af bestaaer en Mangfoldighed af absolut
selvstændige Grunddele, beviser jo tydelig nok, at Aanden ikke er almægtig,
men afmægtig.
% --- p. 203 ---

Det er den samme Grundmodsigelse, som spøger endnu i halv philosophiske
Hjerner, naar det i eet og samme Aandedræt paastaaes, baade at Materien i dens
Elementer er evig, uden Oprindelse, og at der ikke desto mindre er en
herskende Fornuft og fornuftig Orden i Naturen. Fornuften er et Tankevæsen.
Skal Fornuften herske, maa den ubetinget have Elementerne i sin Magt; men det
er kun muligt under den Forudsætning, at Elementerne selv ere absolut
afhængige af Fornuften. Antager man, at Elementernes evige Tilværelse er
Fornuftens egen evige Realitet, at Fornuften ikke blot existerer i
Elementerne, men at de selvstændige Elementer netop ere den existerende
Fornuft: da har man enten reentud benegtet, at det, vi kalde Fornuft, er et
Tankevæsen, eller bragt en saa urimelig Dualisme ind i Fornuften selv, at det
Fornuftige hører op.

% SÆTTERFEIL, gjengivet som trykt: der staaer „hvcri“ for „hvori“. Bogstavet
% er en utvetydig c --- Skaalen er aaben til høire --- baade paa
% KB-Exemplaret og paa Arbeids-Scanningen, og adskiller sig klart fra det
% lukkede o i „om det aandige“ i Linien ovenfor og i „hvori“ S. 207. Begge
% Exemplarer vise det, altsaa Sætterens Feil (eller en brudt Type). Læs
% „hvori“.
Den Anaxagoriske Dualisme maatte vel fra sin Side bidrage til at styrke
Bevidstheden om det aandige Væsen, hvcri Naturerkjendelsen har sin Sandhed;
men Forklaringen var, hvad Sokrates beklager, i Virkeligheden mechanisk, og
Dualismen utænkelig. Man maatte da vælge mellem Et af To: enten holde paa de
selvstændige Elementer og opgive Tankevæsenet, eller fastholde Tankevæsenet og
opfatte Materien som et Slags Ikkeværende, et formløst Stof, og deducere
Elementerne ved Hjælp af Formen.

Med Platos Ideelære kom Tankevæsenet, i græsk Forstand, til fuld
Anerkjendelse. Uden væsensfyldige Begreber gives der hverken nogen sand Væren
eller sand Viden. Begreberne, som paa een Gang udtrykke Værens og Videns
Sandhed, ere selvgyldige Former af selvgyldigt Indhold, indeslutte al
Virkelighed og bestaae ved sig selv som Ideer. Det Sandselige forandres,
Phænomenerne vexle, Naturtingene opstaae og forgaae; men det i og for sig
Bestaaende, det Almene, det i Sammenligningernes Mangfoldighed Selvlige
% --- p. 204 ---
% „Ideen“ er spærret; kontrolleret paa KB.
(ἓν ἐπὶ πολλῶν) er evigt, uomskifteligt, Videns Gjenstand, det Værende selv,
\emph{Ideen}.

Vel er Ideens Eenhed ogsaa en Fleerhed; det Værende har i sig selv
mangfoldige Bestemmelser af Ikke-Væren; men de Forskjelligheder, som herved
fremkomme, ligge alle i Værens Væsen og stemme overeens med hverandre. Det
Værende kan umulig være det forskjelsløse Ene. Sige vi: det Ene er, saa gjøre
vi Forskjel mellem Eenhed og Væren; vi have altsaa Tvende: det Ene og det
Værende. Sige vi: det Værende er et Hele, saa ligger det i Begrebet Hele, at
det maa have Dele: det Hele er en Flerhed af Dele. Vil Nogen paastaae, at det
Værende ingen Dele har og selvfølgelig heller ikke er et Hele, saa gjør han
det Værende umuligt. Hvorledes skulde det vel kunne være? I sig selv! Men saa
maa det jo være forskjelligt fra sig selv; thi der er jo dog Forskjel mellem
det, hvori Noget er, og det Noget, som er deri. Af disse dialektiske
Antydninger vil det være indlysende, hvorledes det hos Plato forholder sig med
Ideernes Mangfoldighed. Hver Idee er en Værensbestemthed, fordi den er en
Tankebestemthed, en bestemt Tanke, og omvendt. Ideerne maatte altsaa
oprindelig have været betragtede fra to Sider, nemlig: 1) som Tankevæsener,
hvis Væren er en nødvendig Tænken, der forudsætter en evig Tænkende; dernæst:
2) som objectivt Værende, almene Substanser i det Grundværende. Men en ligelig
Udvikling af de to Hovedsider har Plato ikke kunnet give. Ideernes Forhold til
Guddommen, Tankevæsenets subjective Side, berøres kun i Forestillingens Form
under mythiske Indklædninger, medens den objective Side, Ideernes reale
% Bindestregerne i „For-sig-væren“ ere Ordets egne; Linieskiftet falder
% tilfældigvis sammen med den anden af dem.
For-sig-væren, bliver udviklet igjennem den allerfineste Dialektik.

Er nu den sande Verden en Verden af Ideer, hvorfra kommer da Sandseverdenen?
Der maa nødvendig forudsættes et Noget, der kan gjøre en Medværen af det
Sandselige, det Stofagtige, det Vordende forklarlig. Men hvad er det for et
% --- p. 205 ---
Noget? Det skulde jo være Materien, Stoffet. Men et positivt Udtryk for
Materie og Stof findes ikke hos Plato. Materien kan ikke tænkes; thi det, som
i Sandhed tænkes, er Ideen. Materien, det i sig selv formløse Stof, kan heller
ikke sandses; thi det, som sandses, er altid kun et bestemt, formet Stof, ikke
det almene, formløse Grundlag for alt Stofagtigt (kun et τοιοῦτον, ikke τόδε).
% SÆTTERFEIL i Græsken, gjengivet som trykt: der staaer „ἀναιθησίας“ uden
% sigma, hvor Timæus har „ἀναισθησίας“. Læst paa KB-Exemplaret ved høi
% Forstørrelse: mellem ι og ϑ staaer intet σ. Læs „ἀναισθησίας“.
% Skriftsnittets ϑ-Form gjengives efter Husreglen med θ. Tegnet efter „μετ“
% er Elisionsapostrofen som trykt.
Det formløse Stof, Massen (ἐκμαγεῖον) bliver altsaa at opfatte ved Hjælp af en
vis uægte Tænken (μετ’ ἀναιθησίας ἁπτὸν λογισμῷ τινι νόθῳ). Stoffet i og for
sig er et Ikkeværende, men som dog i sandselig Mening maa blive et
Medværende. Skal Stoffet nærmere bestemmes, maa det skee ved dets Modsætning
til Ideen. Medens Ideen er det evigt Værende, er Stoffet det altid Vordende,
det Ubegrændsede, der faaer sin Begrændsning, sin Form, sit Præg af Ideen. Ved
en Forbindelse af Ideerne med det Stofagtige fremkomme de sandselige Ting, den
endelige, dog evige, Natur, den hele udvortes i Rummet tilværende Verden. Der
er da i Naturen ikke en Blanding af to Slags Væsener, Ideernes og Stoffets,
thi Stoffet i Tingene er væsenløst; det er Ideerne, der give Tingene Realitet:
den hele Virkelighed er i Ideen.

At Stoffet, det Ikkeværende, kan, formedelst Ideen, faae Betydning i den
sandselige, Vorden og Forandring underlagte Tilværelse, beroer da paa, at det
formes, at det gaaer op i de for udvortes Væren nødvendige Ideeformer, det vil
sige i geometriske Grundformer. Fra dette Synspunkt vil Platos eiendommelige
Deduction af de fire Elementer være forstaaelig.

Som en legemlig Verden, hedder det i Timæus, maatte Verden være synlig og
haandgribelig. Synlig kunde den ikke være uden Ild; haandgribelig ikke uden
Jord, Grunden til alt det Faste. Mellem Ild og Jord maatte der som Midte være
et Tredie, der kunde knytte dem sammen. Da nu den skjønneste Sammenknytning er
% Orddelingen løber over Bladskiftet: „Propor-“ paa S. 205, „tionen“ paa
% S. 206.
et bestemt Forhold, Propor\-%
% --- p. 206 ---
% TRYKFEILBLADET retter S. 206, Linie 3 f. o.: „tilstrækkelig; --- læs
% tilstrækkeligt;“. KB-Exemplaret trykker „være tilstrækkelig; men nu,“ ---
% med Semikolon og uden -t --- og saaledes staaer det her.
% ARBEIDS-EXEMPLARET er derimod BLYANTSRETTET paa netop dette Sted: et
% haandtegnet t er sat ind over Semikolonet, saa tesseract læser
% „tilstrækkeligt men nu“ og taber Semikolonet. Samme Haand som ved S. 160.
tionen: saa maa dette Tredie staae i Forhold til begge. Havde vi blot med
Flader at gjøre, saa vilde eet Mellemled være tilstrækkelig; men nu, da der
handles om Legemer, maa der nødvendig være to. Vi erholde da fire Elementer,
der indbyrdes danne en sammenhængende Proportion, saaledes at Ilden forholder
sig til Luften som Luften til Vandet, og Luften til Vandet som Vandet til
Jorden.

Hvorledes det i arithmetisk Forstand nærmere forholder sig med denne
Proportion, er i nærværende Sammenhæng uden Interesse; charakteristisk er det
derimod, at Elementerne dannes uden tilværende Stof af geometriske Figurer.
Alle Flader, hedder det, bestaae af Trekanter, alle Trekanter afledes af de
retvinklede, som deels ere ligebenede, deels uligesidede. Blandt de
retvinklede og uligesidede Trekanter er den den skjønneste, hvis mindste
Kathede er halvt saa stor som Hypothenusen. Af sex saadanne Trekanter opstaaer
en ligesidet; men af fire ligebenede Trekanter et Qvadrat. Af Qvadratet dannes
Tærningen, Hexaedret; af ligesidede Trekanter dannes Tetraedret, Oktaedret og
Ikosaedret. Ildens Grunddele ere Tetraedre, Luftens Oktaedre, Vandets
Ikosaedre og Jordens Hexaedre.

Idet Stoffet svinder ind til et Ikkeværende, og Elementerne dannes af rene
Begrændsningsformer i Rummet, er Atomistiken opløst: kun den stofløse
Discretion, den rene mathematiske Punktualitet, bliver tilbage. Man kan nu
vistnok sige, at Immanensen er hævdet, at de sandselige Ting ere immanente i
Ideerne, da jo Ideerne selv ere Tingenes egen Realitet. Men hvorfra det
Sandselige kommer; hvorledes det gaaer til, at et Ikkeværende kan tage sig ud
som Stof, og bevirke, at den evige Ideernes Verden bliver til en udvortes,
sandselig, i Vorden og Forandringer haandgribelig Verden er fra et platonisk
Synspunkt ubegribeligt. Eftertrykket falder paa Stof og Form.
% --- p. 207 ---

Forholdet mellem Stof og Form (ὕλη og μορφὴ) spiller en Hovedrolle hos
Aristoteles, saavel i Naturphilosophien som i Metaphysiken. Stoffet er
Grundlag for al Vorden; thi Vorden bestaaer just deri, at et Stof antager en
bestemt Form, gaaer over fra at være i en Form til at være i en anden. Stof og
Form forholde sig som Mulighed til Virkelighed, som det Passive til det
Active, i særlig Forstand som Tilfældighed og mechanisk Nødvendighed til
Hensigtsmæssighed. Kun Formen som Form er begribelig; det formløse Stof er
uden Begreb og forsaavidt et Ikkeværende. Ikke desto mindre maa der dog
tillægges dette Ikkeværende en vis positiv Væren, thi Stoffet staaer Formen
imod: det i Naturtingene Tilfældige og Uvæsenlige hidrører fra Stoffets
Modstand mod Formen. Den platoniske Deduction af Elementerne er, ifølge
Aristoteles, endnu mere forfeilet end Atomlæren. I Atomerne er der dog et
formet Stof, altsaa en Realitet; men Platos Elementer ere hule Figurer, lige
saa stridende mod Mathematiken som mod Physiken. At sammensætte Legemer af
Flader fører conseqvent til at antage udeelbare Linier, ja til at opløse
Størrelser i Punkter. Som rene mathematiske Figurer kunne Elementerne umulig
opfylde Verdensrummet, og dog vil Plato ikke indrømme noget tomt Rum. Fra et
physisk Synspunkt er Forsøget aldeles urimeligt. Thi hvorledes kan det
Legemlige, som jo dog har Tyngde, bestaae af Flader som ingen Tyngde have?
Hvorfra skulde de enkelte Elementers specifiske Tyngde eller Lethed, under
denne Forudsætning, vel hidrøre? --- Antagelsen af faste, udeelbare Atomer er
% spacing.py's RUN „Figur ikke passede det“ er FORKASTET: KB viser
% almindelig Antiqva.
dog ogsaa falsk. Ifølge den absolute Udeelbarhed maatte det jo kunne skee, at
Atomernes Figur ikke passede i det Rum, hvori de skulde befinde sig. Heller
ikke indseer man nogen Grund, hvorfor der kun skulde gives usynlige smaa, og
ikke ogsaa store Atomer. Bevæges Atomerne af noget Andet, saa undergives de en
Paavirkning, og den dem tillagte Apathie er ophævet;
% --- p. 208 ---
bevæge de sig selv: saa er enten det Bevægende i dem forskjelligt fra det
Bevægede, og de ere da ikke udeelbare; eller, dersom det Bevægende er Eet med
det Bevægede, saa er der i eet og samme Atom modstridende Egenskaber.

Langt fra at sætte Tingenes Opstaaen og Forgaaen i en Sammenhoben og Spreden
af udeelbare Atomer, hvis qvantitative Forskjelligheder ere utilstrækkelige
til at forklare de legemlige Qvaliteter, antager Aristoteles en
gjennemgaaende, af Formen bevirket Stofforandring, en real Vorden i væsenlige
Overgange. Naar Vandet fryser, eller naar Isen smelter, saa skeer det ikke ved
en blot Forandring i Delenes Leiring, men ved en Omformning af Delene selv,
hvorunder Egenskaberne vexle; naar der af Næringsstofferne dannes Knokler og
Kjød i Legemet, da skeer det ikke derved, at enkelte Dele uddrages ligesom
Teglsteen af en Muur, men der danner sig et nyt Stof. Det er den Heraklitiske
Forvandlingslære, udviklet til articuleret Begreb.

Betingelserne for den qvalitative Forandring ligge, som enhver Bevægelse, i
Forholdet mellem det Mulige og det Virkelige, nærmere bestemt som Forhold
mellem det Lidende og det Virkende. Det Virkende og det Lidende maae være
deels ligeartede, deels modsatte. Forandringen bestaaer i, at det Virkende
gjør det Lidende ligt med sig, en Omdannelse, der igjen kan skee paa flere
Maader; thi de mechaniske Stofaarsager ere kun Mellemaarsager; Eftertrykket
ligger paa Bevægelsen. Naturen selv er Bevægelsens indre Grund; enhver
Bevægelse har i sig et Maal, hvorved dens Retning og Grændse bestemmes:
Naturen virker teleologisk. Det er det Eiendommelige ved Bevægelsen nedad og
opad, der danner Grundlaget for Adskillelsen af de fire Elementer. Dog maa der
ikke blot tages Hensyn til Modsætningen mellem det Tunge og det Lette; ogsaa
Legemernes følelige Egenskaber komme i Betragtning. De elementære Egenskaber
ere Varme og Kulde, Tørhed og Fugtighed, de to første
% --- p. 209 ---
% De fire Elementnavne ere spærrede, ikke hele Vendingerne; kontrolleret paa KB.
active, de to sidste passive. Ved at stille disse Egenskaber parviis sammen
erholde vi, med Fradrag af to umulige Forbindelser, fire mulige: Varme og
Tørhed i \emph{Ilden}, Varme og Fugtighed i \emph{Luften}, Kulde og Fugtighed
i \emph{Vandet}, Kulde og Tørhed i \emph{Jorden}; altsaa igjen de fire
Elementer.

Elementerne gaae over i hverandre; thi Alt vorder fra et Modsat til et
Modsat. Jo fuldstændigere Modsætningen er, desto vanskeligere er Overgangen;
jo ufuldstændigere, desto lettere. Naar Vandets Kulde hæves, saa opstaaer
Luft; thi Vand og Luft have Fugtigheden tilfælles og ere ufuldstændig
modsatte. Vand og Ild ere derimod fuldstændig modsatte; Overgangen maa
desaarsag skee middelbart, idet Vandet først gaaer over til Luft og saa, naar
Fugtigheden hæves, over til Ild.

Uden at lægge Vægt paa, at Phænomenernes Udtydning i væsenlige Punkter strider
mod Nutidens Physik, skulle vi fra Naturphilosophiens Standpunkt nu særlig
undersøge, hvorvidt Aristoteles ved sin Forklaring af Elementernes Oprindelse
virkelig har været i Stand til at forlige Formen med Stoffet, hævde
Tankevæsenet og gjennemføre Immanensen.

% Indløbende Overskrift, læst paa Siden selv (KB-Exemplaret): halvfed Antiqva,
% „c)“ med almindelig Skrift.
\runin{c) Stof og Substans.} Atomikerne fastholde Stofvæsenet i ubetinget
Selvstændighed og udelukke Tankevæsenet; Plato hævder Tankevæsenet ved at
nedsætte Stoffet til Skin af et Ikkeværende. Aristoteles gjør Tankevæsenet til
Form og Sandselighedscomplementet til et positivt Substrat, et Stof, der lader
sig forme. Han anerkjender med en vis Ligelighed baade det Ideelle og det
Reelle; hos ham skulde vi da vente at finde, hvad den nyere Speculation saa
ivrig søger, det ideelt-reelle Princip, der kan begrunde den materielle
Verdens immanente Eenhed med den objective Fornuft.
% Nederst paa Bladet staaer Arksignaturen „14“; den er ikke Text.
% --- p. 210 ---
Men netop hos Aristoteles lære vi, hvor urimelig og selvmodsigende den hele
speculative Bestræbelse er, naar Immanensen mangler sin transcendente
Forudsætning, og Eenheden af Natur og Idee umiddelbart og paa første Haand
skal søges --- ikke i Aanden, men i Naturen. Der er i de legemlige Ting en
egen Forening af et Sandseligt med et Tænkeligt; paa denne Forening beroer
ethvert Legemes eiendommelige Bestaaen eller Substantialitet. Abstrahere vi
fra det Tænkelige d. v. s. fra Tankevæsenet, saa have vi kun en flygtig,
drømmeagtig Tilsyneladelse, men ingen Substans, intet virkeligt Legeme.
Abstrahere vi fra det Sandselige, det Stofagtige, saa have vi et almindeligt
Legemsbegreb, men den tilværende Substans, det virkelig bestemte Legeme er
borte. Saa er det da uimodsigelig vist, at Stoffet som Stof maa have en
ligesaa afgjørende Betydning for et Legemes Substantialitet, som Formen. Denne
Kjendsgjerning sætter Aristoteles i Forlegenhed, en Forlegenhed, som endnu
ingenlunde er overvunden enten af Schelling eller Hegel eller Feuerbach.

% „reale“ er spærret; kontrolleret paa KB.
Almeenvæsenet --- Slægten, Arten --- har, ifølge Aristoteles, ingen Bestaaen
for sig; det bestaaer kun i og ved det Enkelte. Substansen i fuldeste
Betydning, det sande, virkelige Væsen, der er Subject, men aldrig Prædicat, er
et Enkeltvæsen. Da nu Formen som Form er det Almene, et Prædicat, der kun i
abstract, uegenlig Forstand kan blive Subject: saa følger jo heraf, at Formen
kun kan være en Side af Substansen, ikke Substansen selv. Ligger Grunden, den
\emph{reale} Grund, til Enkeltvæsenet ikke i Formen, men i Stoffet; er det
Stoffet, der giver Tingene særlig Bestaaen ved at give dem Individualitet, da
maa det egenlig Substantielle ligge i Stoffet. Men det vil Aristoteles ikke
indrømme. Hvis Stoffet var Substans, maatte det Substantielle absolut falde
udenfor Viden, thi al Viden gaaer paa Formen. Substansen maa da begribes
saaledes, at Stoffet kun efter
% --- p. 211 ---
sin Mulighed er et Enkeltvæsen, men at Enkeltvæsenets Virkeliggjørelse skyldes
Formen. Og dog maa ogsaa denne Antagelse forkastes. At Stoffet har Mulighed
til at blive legemligt Enkeltvæsen formedelst Formen, maa nødvendig tænkes
saaledes, at det samme Stof, der er Enkeltvæsenets Mulighed, ogsaa er Formens
Mulighed. Stoffet faaer altsaa ikke Formen som et Andet, en fuld færdig
Virkelighed, der træder til; det udvikler paa givne Betingelser sin Form, sin
Virkelighed af sin egen Mulighed: Stoffet lader sig forme, fordi det er dets
Natur at have Formen i sig, at forme sig. Det er en reen Indbildning, at der
skulde være et almindeligt Underlag (ὑποκείμενον), der faldt udenfor Tanken,
fordi det faldt udenfor Formen! det formløse Stof er, hvad der da ogsaa kan
kjendes paa, at det i sin Formløshed opfattes som noget Almindeligt, selv en
Tanke, en abstract, ufuldbaaren Tanke. Det formløse Stof er en Betegnelse for
noget Usandt, for et ikkeværende og dog værende Formeligt, der venter paa
Formen. Det formløse Stof er ingen Mulighed; det er tvertimod en Umulighed,
thi det er en Modsigelse. Modsigelsen vedkommer slet ikke Stoffet, men Tanken.
Det i sig formløse Stof er en lige saa modsigende Tanke, som den i sig
formløse Form, det i sig stofløse Stof.

% „til det Tænkelige eller det Ikketænkelige.“ staaer med Punktum, ikke
% Spørgsmaalstegn, skjøndt Spørgsmaalet fortsættes; gjengivet som trykt.
Den Maade, hvorpaa Aristoteles vil deducere Elementerne, og i det Hele Alt,
hvad han lærer om Legemernes Egenskaber, viser tydelig, hvor modsigende det
er, at gjøre Formen uden videre til det Tænkelige, og det Stofagtige til det
Ikketænkelige. Hvorhen høre vel saadanne elementære Egenskaber som Varme og
Kulde, Tørhed og Fugtighed? til det Tænkelige eller det Ikketænkelige.
Forsaavidt de opfattes som Former, maae de jo udelukkende falde ind under
Viden, altsaa under det Tænkelige. Men ifølge Deductionen svare de til
Følelsen, altsaa til det Stofagtige, det Ikketænkelige. Hvorledes skulde vel
Tørheden lade sig forvandle til Tanke? Det maatte da være en tør Tanke. Men
% Nederst paa Bladet staaer Arksignaturen „14*“; den er ikke Text.
% --- p. 212 ---
Tanker ere i Virkeligheden lige saa lidt enten tørre eller fugtige, som de i
naturlig Forstand ere enten varme eller kolde.

% spacing.py's RUN „blot viist“ er FORKASTET: KB viser almindelig Antiqva.
Stoffet i al Ubestemthed er et Usandseligt; det er Formbestemtheden, der gjør,
at Stoffet kan sandses. I Conseqvens af den Aristoteliske Opfattelse kan man
da med lige saa megen Ret paastaae, at Formen er sandselig, og Stoffet
tænkeligt, som omvendt. Man kan lade Stoffet være Moment i Formens Begreb og
gjøre Formen til den hele Virkelighed; men man kan ogsaa lade Formen være
Moment i Stoffets Begreb, og sætte den hele Virkelighed i Stoffet. Denne
Svingen mellem en eensidig Idealisme, der i Formen holder paa Begrebet, og en
eensidig Realisme, der søger Realiteten i Stoffet, har ikke blot viist sig i
Oldtiden, men --- en nødvendig Følge af Immanenslærens Tvetydighed --- den har
paa en ret charakteristisk Maade gjentaget sig i den nyere Tid. Hegel, der
gjør Materie og Stof til lige saa logiske Kategorier som Form og Substans,
culminerer i Begrebsidealisme; Feuerbach, som af Forargelse over det naturløse
Hule i Begrebsidealismen gaaer til den modsatte Yderlighed, sætter Formen som
Moment i Stoffets Selvindividualisering og ender med et sandseligt
Detteværende, et Uudsigeligt.\footnote{% trykt Mærke: *) --- eneste Note paa
% Bladet, helt indenfor S. 212. Mærket staaer i Texten efter Punktummet,
% adskilt ved et Mellemrum: „et Uudsigeligt. *)“.
Eftersnakkerne forsikkre, at Hegel er eensidig, fordi han overseer det Reelle,
og Feuerbach eensidig, fordi han overseer det Almeen-Ideelle; saa kan man da
blive alsidig ved at sætte det Ideelle immanent i det Reelle og det Reelle
immanent i det Ideelle; saa har man det Ideelle og det Reelle i den rene
Immanens, forstaaer sig --- i den immanente Snak.}

Uklarhed over, hvad Stof, Element og Stofform er, maa nødvendig falde sammen
med en tilsvarende Uklarhed over, hvad Natur er. At Aristoteles, trods sit
% Orddelingen løber over Bladskiftet: „Er-“ paa S. 212, „faringsblik“ paa
% S. 213.
indtrængende Er\-%
% --- p. 213 ---
faringsblik, sin overlegne Skarpsindighed og beundringsværdige Udholdenhed i
omfattende naturphilosophisk Tænkning, ikke veed, hvad Natur er, ikke er enig
med sig selv om, hvorledes han egenlig skal bestemme Naturens Væsen, er
beviisligt af hans eget System. Overgangen fra Metaphysiken til
Naturphilosophien betegner han som en Overgang fra Læren om det Ubevægede til
Læren om det Bevægede. I Kredsen af Bevægelserne og alt det Bevægede indstille
sig umiddelbart de legemlige Enkeltvæsener, bestaaende af Stof og Form. Hvad
er nu Stoffet? Det vides ikke og kan ikke vides, thi Stoffet som Stof er ingen
Gjenstand for Viden. Hvad er da Formen? Det er just det, der vides; og dog,
naar vi see nøiere til, vides det heller ikke. Formen, hedder det, er
Virkeligheden. Den fuldkomne Form er den fuldkomne Virkelighed, den høieste
Energie, det er Guddommen. Dette Formvæsen er i absolut Forstand et
Tankevæsen; det er den evige Tænken, der kun tænker sig selv, den salige
Fornuft, der begriber sig selv, en Virkelighed uden alle Muligheder. Men
Naturformerne, hvad er det for Væsenheder? Ere de Guddommens Tanker?
Ingenlunde! At antage dette vilde i lige Grad være stridende mod Guddommens
Væsen og mod Formernes Natur. Mod Guddommens Væsen! Thi Fuldkommenheden i den
guddommelige Tænken beroer netop paa, at den alene tænker det Fuldkomne. For
at tænke det Fuldkomne maa den udelukke Tanken om alt Andet; thi Intet er
fuldkomment uden Guddommen selv. Mod Formernes Natur! Thi dersom Naturformerne
vare Guddommens Tanker, saa havde de en Væren i et Andet, nemlig i den
tænkende Guddom, og vare afledede af et Andet, nemlig af den guddommelige
Tænken. Under denne Forudsætning kunde Formen umulig være Tingenes Substans.
Substans er jo netop det, som hverken er i et Andet eller kan afledes af et
% Orddelingen løber over Bladskiftet: „Be-“ paa S. 213, „greb“ paa S. 214.
Andet, men ubetinget er i sig selv. Hvilken Spaltning af Formens Be\-%
% --- p. 214 ---
greb! Den høieste Form er ikke blot tænkelig af os, men i sig selv tænkelig:
dens Tænken er Væren og dens Væren Tænken. Naturformerne derimod ere
tænkelige af os, men i sig selv utænkelige, thi i Tankens Form ere de ikke
Substanser.

Formbestemtheden er i sig selv da ogsaa utænkelig. Naturtingenes Egenskaber og
Bestemmelser ligge i de bevægende og bestemmende Former. Stoffet bestemmes af
den i sig bestemte Form; men Formen selv bestemmes ikke; Formen er bestemt.
Har Formen da forud formet og bestemt sig selv? Ingenlunde! Skulde
Naturformerne forud have formet sig selv, før de kunde begynde at forme
Stoffet: da maatte de jo under denne Selvformen have gjennemgaaet en Vorden;
de maatte allerede have havt det Mulige, det Passive, det Stofagtige i sig,
før de kom i Berøring med Stoffet; der maatte da igjen have været en forud
færdig Form, som kunde give Stofformerne Form. Formbestemtheden forudsætter
umiddelbart sig selv: det hele i Naturriget virkende System af Former bliver
at tage som en given uforklarlig Kjendsgjerning. Denne Kjendsgjerning har saa
igjen det Betænkelige ved sig, at den lider af uopløselige Modsigelser.
Forsaavidt Formerne opfattes i Viden, opfattes de for sig i abstraheret
Almindelighed og kunne ikke være Substanser; forsaavidt de gaae sammen med
Stoffet, besværes de af det Stofagtige og undergaae en væsenlig Forandring; de
bestemmes idet de bestemme; de blive passive ved at være active, ligesom
Stoffet formedelst sin Modstand bliver activt ved at være passivt. Man kan ved
at fortsætte Analysen finde Modsigelse paa Modsigelse. Hemmeligheden er, at
den umiddelbare Immanens af Stof og Form beroer paa en Grundmodsigelse; men at
en Grundmodsigelse, formælet med dialektisk Reflexion, yngler stærkt, er en
bekjendt Sag.

% Sætningen løber over Bladskiftet og fuldføres paa S. 215 (næste Batch):
% „... som Element; det er i sig selv et Uselvstændigt, ...“
Udbyttet af den gamle Immanenslære bliver da følgende: Stoffet er ikke
umiddelbart at opfatte som Element;
%  --- B. Elementer og Naturgrund  (pp. 215–232) ---
% --- p. 215 ---
% Afdeling A løber seks Linier ind paa denne Side; Sætningen fortsættes
% fra S. 214 („... at opfatte som Element;“) uden Afsnitsskifte.
det er i sig selv et Uselvstændigt, der først faaer Qvalitet og qvalitativ
Bestaaen, naar det bestemmes ved en til Viden svarende Form. Elementerne
have en Oprindelse; Fordringen er, at deres Oprindelse skal paavises, at
Elementerne skulle deduceres. Fordringen i sig selv er sand; men endnu er
den ikke opfyldt.

% Kort, centreret Streg midt paa S. 215, mellem Afdeling A og Afdeling B;
% læst paa begge Exemplarer. Efter Husreglen gjengives Ornamentets Længde
% ikke: \streg tjener overalt.
\streg

% Overskriften er læst paa Siden selv (KB-Exemplaret), ikke i Indholdet:
% „B.“ paa egen Linie, derunder „Elementer og Naturgrund.“ sat med fed,
% spærret Skrift. Overskriften staaer omtrent midt paa Bladet.
\afdeling{B.}{Elementer og Naturgrund.}

Den haarde Modsætning af Element og Tankevæsen har sin oprindelige Grund
i Modsætningen af umiddelbar Gjenstandsmagt og overgribende Selvmagt. Den
absolut overgribende Selvmagt begriber Naturgrundens Indhold som en
Totalitet af Momenter i eet Tankevæsen; derpaa beroer
% Spatieringen efterprøvet paa KB-Exemplaret.
\emph{Elementernes Idealitet}. I Gjenstandsmagten komme de reflecterede
Momenter til Existens som Grundforskjellighedernes umiddelbare Udslag, som
selvstændige Led i udvortes Virkelighed; deri bestaaer
\emph{Elementernes Realitet}. For at grunde i sig selv maa Naturen
begrundes af Aanden: den uendelige Omsætten af Idealitet i Realitet, og
omvendt, foregaaer i \emph{Grundmediet}.
% „Grundmediet“ er spatieret paa Siden, men blev ikke fanget af
% spacing.py; læst paa KB-Exemplaret.

% Indløbende Overskrift, sat med halvfed Skrift, paa samme Linie som
% Afsnittet den aabner. Læst paa Siden selv (KB-Exemplaret).
\runin{a) Elementernes Idealitet.} Tankevæsenet, bestemt alene ved
menneskelig Viden, kan umulig opfattes som det reale, i Naturelementerne
virksomme Væsen. Eet af To: enten maa Bevidstheden standse for en Skranke
af uigjennemtrængelig Realitet, bøie sig for den udvortes mægtige Natur og
tillægge dens Elementer og Kræfter en evig, ab\-%
% --- p. 216 ---
solut Selvstændighed; eller Tanken maa i speculativ Forglemmelse af sin
Afmagt, erklære Stoffet, den utænkelige Grund til al udvortes Væren, for
et Ikkeværende, sætte \emph{sit} Væsen som Tingenes Væsen og gjøre sig selv
% „sit“ staaer spatieret i Satsen, ogsaa paa KB-Exemplaret, skjøndt Ordet
% er kort og staaer sidst paa Linien; Spatieringen er tydelig og
% Linien iøvrigt ikke løst udslaaet.
til Guddom. Skal Tankevæsenet magte Elementerne, da maa det være et
overmenneskeligt Væsen, hvilende i guddommelig Viden, begrundet af en
almægtig Tænken.

Ideerne, hedder det hos Plato, ere ikke Tanketing
% Græsk, sat med kursiv græsk Skrift: νοήματα (akut over eta), læst paa
% KB-Exemplaret.
(\textit{νοήματα}), der blot have Tilværelse i Sjælen. Rigtigt! Men naar
saa heraf udledes, at Ideerne ikke kunne være Guddommens Tanker, er
Slutningen falsk. Sjælens Tanker, opfattede for sig, ere Afmagtstanker,
abstracte Tanketing, der have Realiteterne udenfor sig; men Guddommens
Tanker ere Almagtstanker med indre, væsenlig Realitet. At Ideeformerne
ere nødvendige Tankeformer, beviser den Platoniske Dialektik; men
Tankeformer ere jo Tænkningens egne, af Tænkningen selv bestemte, Former,
Former, som kun et tænkende Selv kan frembringe. Er Ideernes evige Væren
ikke begrundet i evig Tænken, saa gives der ingen sand Viden, saa er den
Platoniske Dialektik et tomt Ordspil.

Hvad vil det egenlig sige, at Tænkningen vel er afhængig af Ideen, men at
Ideen, det i og for sig Værende, er uafhængig af Tænkningen? Det vil
udtrykkelig sige, at Ideen i sig selv er et utænkeligt Væsen. Et Værende,
der gaaer forud for al Tænken, er et Værende, hvis Bestemthed, langt fra
at beroe paa nogen bestemmende Tanke, tvertimod falder udenfor enhver
Tankebestemmelse. For at et saadant Værende kan blive Gjenstand for
Tanken, maa det forestilles; det kan ikke smelte sammen med Tanken, men
maa i værende Selvstændighed fastholdes som Tankens Modstand, dens
uigjennemtrængelige Andet. Tanken kan nu gjerne i Et og Alt rette sig
efter denne i Forestillingen usynlige Gjenstand og forvisse sig om, at den
noiagtig opfatter alle
% --- p. 217 ---
i det Værende indeholdte Bestemmelser; men hvad er skeet? De evige af
Tænkningen uafhængige Værensbestemmelser ere forvanskede; de ere
forvandlede til Noget, som de dog i sig selv ikke ere, nemlig til lutter
Tankebestemmelser. Det i Tanken Værende er altsaa ikke Ideen selv ---
Ideebestemmelserne ere jo ikke Tankebestemmelser --- men en Efterligning,
en intellectuel Copie af en udenfor værende Original; Ideernes virkelige
Væren falder udenfor deres tænkte Væren: det er en haard Dualisme af
Tænken og Væren.

For at danne Verden og bringe den bedst mulige Orden ind i det
Uordenlige, maatte, siger Timæus, den øverste Bygmester bestandig see hen
% Spatieringen af „Urbillederne“ og „Forbilleder“ efterprøvet paa
% KB-Exemplaret; „Ideerne“ derimellem staaer almindeligt.
til \emph{Urbillederne}, Ideerne, Tingenes evige \emph{Forbilleder};
herved, paastaaer man, er Ideernes Substantialitet uimodsigelig hævdet.
Rigtigt! Men naar man saa deraf vil udlede, at de Urbilleder, til hvilke
Guddommen selv seer hen, umulig kunne være Frembringelser af Guddommens
egen Tænkning, er Slutningen falsk. Den lange, uklare Strid mellem Ideerne
og den guddommelige Forstand, maa dog vel omsider kunne bringes til Ende.

Saalænge den menneskelige Bevidsthed, være sig i berusende Speculation
eller fastende Kritik, anmasser sig en absolut Viden af det Absolute, vil
den mystiske Forestilling om „ein unvordenkliches Sein“ bestandig dukke op
% Den tydske Vending staaer i Anførselstegn, men i almindelig Antiqva:
% hverken kursiveret eller spatieret (efterprøvet paa KB-Exemplaret).
og bringe Forstyrrelse ind i den lovpriste Eenhed af Tænken og Væren. I
det psychologisk begrændsede Selv gaaer den reale Væren, hvorledes man saa
end dreier det, nødvendig forud for enhver nok saa speculativ Tænken; kun
i det absolute Selv er den endelige Tidsfølge hævet. Den guddommelige
Væren og den guddommelige Tænken forudsætte hinanden i det Uendelige:
Væren følger af Tænken, og Tænken af Væren i Aandens evige substantielle
Væsen. At gjøre Væren til det absolut Første og lade en Tænken følge bag
efter, er, med Hensyn paa Guddommens Begreb, lige saa tankeløst, som at
gjøre Tænken til det Første og deraf aflede Væren.
% --- p. 218 ---

For at bevise, at der i den evige Væren er Noget, som den guddommelige
Tænken ikke kan forandre, og hvorefter den altsaa maa rette sig, har man
anført saadanne Almeensætninger som: at To er mindre end Tre, at en
Trekant maa have tre Vinkler, at det Hele maa have Dele, at en Virkning
maa have en Aarsag o. s. v. Men hvo seer nu ikke strax, at slige saakaldte
„evige Sandheder“ kun ere nødvendige Værensbestemmelser, fordi de ere
nødvendige Tankebestemmelser. Dersom Gud nogensinde kunde blive træt af at
tænke, da vilde han med det Samme være træt af at være.

Ifølge den Aristoteliske Definition, at Substans er det, som ikke kan være
i et Andet, kunne de guddommelige Tanker vistnok ikke være Tingenes
Substanser; men af Definitionen følger et Enten---Eller. Enten gives der
% Tankestregen i „Enten---Eller“ staaer uden Mellemrum paa begge Sider,
% som trykt.
slet ingen legemlige Substanser, hvad Aristoteles ikke kan indrømme, da
han jo paa det Eftertrykkeligste indskærper, at de legemlige
Enkeltvæsener ere de egenlige Substanser. Eller de legemlige Substanser
maae være absolute; hvad han heller ikke kan indrømme, efterdi det staaer
i udtrykkelig Strid med hans Lære om Elementernes Overgange, om legemlige
Dannelser, Forbindelser og Opløsninger. Var den legemlige Substans et
Enkeltvæsen i den Forstand, at dens Væren i sig gjorde dens Væren i Andet
umulig, da maatte et saadant Enkeltvæsen jo være et evigt, uopløseligt
Atom: de legemlige Substanser maatte da være hævede over alle
Forandringer. At Tingene undergaae Forandringer, gaae til Grunde og
opstaae ved en Forvandling af Stoffet, er et Beviis for, at deres
Substantialitet ikke er absolut, men relativ. Substansernes Relativitet
gjør det uimodsigeligt, at deres Bestaaen i sig er en Bestaaen i Andet.

Fra et almindeligt Begreb om Guddomstanker til en sandselig Forestilling
om saadanne bestemte Stofsubstanser som Kulstof, Svovl, Guld o. desl., er
der, udvortes seet,
% --- p. 219 ---
unegtelig et stort Spring; men naar Sagen gjennemtænkes, skal det dog vise
sig, at Kategorierne: Stofsubstans og substantielt Tankevæsen, i Grunden
selv forudsætte hinanden.

Det substantielle Tankevæsen objectiveres i en formaaende Tænken. I den
uendelige Tænken maae de indbyrdes forskjellige Tankebestemmelser
gjennemskinne hverandre: hvert særligt Begreb er et Tankespeil, der
reflecterer det hele System af Begreber, i hvilket det væsensfyldige
Indhold er udformet til Idee. Hvad giver nu dette Tankevæsen
Substantialitet? Den tænkende, vidende, altformaaende Selvmagt! At Tanken
er Magt, og Tænken uendelig Formaaen, vil sige, at Tænken er Virken og har
i sig selv alle Virkelighedens Momenter som Reflexer i den uendelige
Reflexions overgribende Energie. Men naar alle Virkelighedsmomenter uden
Undtagelse ere indesluttede i det substantielle Tankevæsen, saa følger jo
heraf, at intet Stof, intet Realmoment, kan komme udenfra, men at al
Realitet, al Virkelighed, den være nok saa stofagtig og haandgribelig
udvortes, maa komme indenfra, fra Tankevæsenet, den væsenlige Grund til
den reale Verdens, den ydre Naturs, Elementer og Stofsubstanser.

Vi have seet, hvorledes Tankevæsenet i og ved sig selv forudsætter
Stofsubstanserne; vi skulle nu see, hvorledes Stofsubstanserne i deres
Realitet forudsætte Tankevæsenet.

Hvorvidt Chemiens saakaldte Grundstoffer virkelig ere Enkeltstoffer og
angive den yderste Grændse for Legemernes mulige Opløsning, er et
Spørgsmaal for sig; paa Videnskabens nuværende Trin forestille de
Elementerne, og det er os nok.

De elementære Stofsubstanser have nu for det Første det Idealitetens
Særkjende, at deres uforandrede Bestaaen kommer tilsyne i lutter
Forandringer. Naar Ilt og Brint forene sig og danne Vand, saa forandre de
hinanden gjensidig; de opgive deres Selvstændighed; deres særlige Egen\-%
% --- p. 220 ---
skaber omdannes og forsvinde i nye Egenskaber: de gaae til Grunde, idet
Vandet kommer til Existens. Stofsubstansernes Væren i sig er en Væren i
Andet; deres umiddelbare Selvstændighed er en eensidig og forbigaaende
Tilstand; for at kjende deres Væsen, maa man kjende deres Forbindelser.
Hvad de ere i sig selv, ere de for hverandre indbyrdes; de spille i hver
ny Forbindelse en ny Rolle, de blive uophørlig et Andet ved Andet; der
bliver i Forandringernes Strøm ikke Spor eller Rest af absolut umiddelbar
Selvstændighed tilbage: for deres Ændringer og Forandringer gives der
ingen udvortes endelige Grændser. Og dog er der Grændser, væsenlige,
lovbestemte Grændser; thi de elementære Substanser have i al denne
Andetværen hver sit særegne substantielle Væsen; den ene kan ikke
forvandles til den anden, Ilten ikke til Brint, Brinten ikke til Kulstof.
Men at de elementære Stofvæsener netop ved at undergaae Forandringer holde
sig uforandrede, beviser, at deres Grundværen er et i sig reflecteret
% Spatieringen af „substantielt“ efterprøvet paa KB-Exemplaret.
Væsen, et \emph{substantielt} Tankevæsen.

Det næste Idealitetens Særkjende er, at de elementære Stofsubstanser, hvad
Qvaliteter og Egenskaber angaaer, baade ere uafhængige af hverandre og dog
væsenlig bestemte for hverandre indbyrdes. Kaste vi et Blik paa de noget
over tredsindstyve Grundstoffer, Metaller og Metalloider, som for Tiden
ere bekjendte, da sees det jo vel, at nogle iblandt dem staae hinanden
nærmere, andre fjernere; men det viser sig ogsaa, at ethvert enkelt har
sine Egenskaber, sit eiendommelige Væsen for sig. Der er, naar vi undtage
de faa, som, endskjøndt de ere sammensatte, dog spille en Rolle som
Grundstoffer, ikke Spor til, at den ene Elementærsubstans har laant Noget
fra den anden. Fra dette Synspunkt ere Grundstofferne altsaa uafhængige af
hverandre; deres Forskjelligheder tage sig ud, som vare de fremkomne ved
en reen Tilfældighed. Hvo skulde vel see paa Kulstoffet, at dets
eiendommelige Tilværelse paa nogen Maade beroer paa
% --- p. 221 ---
Brintens, eller omvendt? Men Kulbrinterne lære os dog, at der maa være et
indre, væsenligt og ret eiendommeligt Slægtskab mellem de to Stoffer; de
passe for hinanden, ja, det seer livagtigt ud, som om de vare afpassede
efter hinanden. Men hvorledes er det at forstaae?

Det er en Kjendsgjerning, at Egenskaberne ved det ene Stof, uagtet de
umiddelbart tilhøre det paagjældende Stof, kun faae Betydning ved de
Egenskaber, der lige saa umiddelbart tilhøre et andet Stof. Denne
Kjendsgjerning viser, at Egenskaberne ved det ene Grundstof oprindelig,
uafhængigt af ydre Sammentræf, maae være reflecterede i Egenskaberne ved
de andre; i Dannelsen af det ene er der udtrykkelig et Hensyn til
Dannelsen af de andre: her er imellem dem alle en i guddommelig Viden
% Den latinske Vending staaer med Kursiv i Satsen (Antiqva-Kursiv),
% ikke spatieret; æ-Ligaturen er trykt i „præstabilita“.
forudbestemt Harmonie (\textit{harmonia præstabilita}). Men dette beviser
igjen, at de elementære Qvaliteter ere Momenter af det ene og samme
Grundvæsen, det substantielle Tankevæsen.

Endelig, og dette er Idealitetens fuldt udviklede Særkjende, har det hele
Indbegreb af elementære Substanser Charakteren af et i alle Retninger
% Mellem „af“ og „et“ sidder i Arbeidsscanningen en Trykplet, som
% spacing.py læste som Punktum; KB-Exemplaret viser almindeligt Mellemrum.
gjennemført, til Naturideens Indhold oprindelig svarende Begrebssystem.
Stofsubstansernes Realbegreber frembyde nemlig for Betragtningen to
Hovedsider, en endelig og en uendelig. Fra Endelighedssiden ere de som
Existensbegreber bundne, hver især, til sit eiendommelige Stof; fra
Uendelighedssiden ere de i Ideen hævede over Udstykning til evig klar og
gjennemsigtig Eenhed.

Den indre reelle Sammenhæng mellem de existerende Substanser lære vi i
Chemien at kjende af deres Affiniteter. Grunden til Affiniteternes
Forskjellighed, hvorpaa den hele Mangfoldighed af Forbindelser og
Adskillelser væsenlig beroer, er en reel Grund, Grunden som Stofvæsen. Men
at denne reelle Grund maa være ideelt articuleret, altsaa oprindelig
bestemt af en rationel Grund, kan udtrykkelig paa\-%
% --- p. 222 ---
vises. Affiniteten er i sit Væsen dobbeltsidig: Chloret har Affinitet til
Guld, fordi Guldet har Affinitet til Chlor, og omvendt. Ilten har
Affinitet til Brint, fordi Brinten har Affinitet til Ilt, og omvendt. Og
hvorfor har saa Ilten ingen umiddelbar Affinitet til Fluor? Ganske
ligefrem, fordi Fluor ikke har umiddelbar Affinitet til Ilt. Men i denne
udvortes Henviisning fra det Ene til det Andet er der jo ingen
Begrundelse. Affinitetens formelle, sig selv ophævende Fordi---Fordi
% Tankestregen i „Fordi---Fordi“ staaer uden Mellemrum, som trykt, og
% falder paa Linieskiftet i Satsen.
fordrer et dybere liggende Motiv, en for Væsenet afgjørende, rationel
Grund.

Affiniteten forudsætter, at de paagjældende Stofbegreber nødvendig maae
være Sidebegreber af eet og samme omfattende Totalbegreb. Iltens og
Brintens gjensidige Affinitet, ifølge hvilken de forene sig og danne Vand,
er Existensudtrykket for den Natursandhed, at Realbegreberne af Ilt og
Brint ere som Sidebegreber af Vandets Realbegreb indbegrebne i hinanden.
Hvad der gjælder om to Stofbegreber, gjælder om dem alle: de ere alle
Begreber af samme Begriben. Fra Uendelighedssiden ere Stofbegreberne
Ideebegreber; i hvert enkelt af disse Ideebegreber er Ideens hele Indhold
paa særskilt Maade nærværende. Den reelle Grund grunder i den rationelle,
Stofvæsenet i Tankevæsenet; det er Elementernes Idealitet.
% „Elementernes Idealitet“ staaer her IKKE spatieret (modsat S. 215);
% efterprøvet paa Siden.

% Indløbende Overskrift, halvfed, paa samme Linie som Afsnittet; læst paa
% Siden selv. Ingen Streg foran.
\runin{b) Elementernes Realitet.} Stoffet er lige saa lidt et Tilværende i
og for sig som den tomme, abstracte Form. Stof og Form ere
Sidebestemmelser af den reelle Magt, Gjenstandsmagten, som oprindelig er
Selvmagtens Idealitet modsat. I Magtbegrebets ideelle Form er Elementet et
reflecteret Moment; i Gjenstandsmagtens reelle Form er Momentet et formet
Stof, et tilværende Element. Der føies ikke Noget til Momentet udenfra,
for at det kan blive til Element; heller ikke drages der Noget fra
Elementet, for at det kan blive til Moment: Moment og Element ere, hvad
Indholdet angaaer,
% --- p. 223 ---
væsenlig det Samme. Og dog ere de hinanden modsatte, ja saa væsenlig
modsatte, at det Ene er et Tænkeligt, det Andet et i og for sig
Utænkeligt.

Kun et Tankevæsen lader sig tænke; men et existerende Element --- Kul,
Guld, Svovl --- er intet Tankevæsen, det er som tilværende Gjenstand
Tankens virkelige Modstand. For den afmægtige Tanke er Modstanden en
Hindring, en udvortes Realitet, som ved Sandsning kan opfattes i Billede
og Forestilling, men ikke begribes. For den almægtige Tanke er Gjenstanden
den substantielle Fyldestgjørelse, det Andet, hvormed Tanken mættes,
Objectet, hvori det objectiverende Selv gjenkjender sin egen Magt. At der
i den guddommelige Tænken er ideel Magt, bliver kategorisk udtrykt derved,
at der i Elementerne er en reel Magt, en virkelig Gjenstandsmagt. Aanden
tænker ikke Elementerne; den skaber dem; den magter dem i formgivende
% „begribe“ og „magte“ ere spatierede, „er at“ derimellem ikke;
% efterprøvet paa KB-Exemplaret.
Begreber: at \emph{begribe} er at \emph{magte}.

Fra Elementbegreber til existerende Elementer kan den menneskelige Tænken
ikke paavise nogen endelig Overgang. Det kan i Videnskaben indsees, at
Elementerne udtrykke Begreber, og at Begreberne begribes af Aanden. Tages
Begreberne bort, saa ere Elementerne umulige; tages Elementerne bort, saa
ere Begreberne uvirkelige: Omsætningen af Begreb i Element er et Udslag af
Magten. Men Overgangen selv, den virkelige Skabelsesact kan ikke deduceres
af den almindelige, i Magtfordoblingen liggende Forudsætning.
Magtfordoblingen udtrykker det evige Grundforhold mellem Begrebernes
ideelle og Elementernes reelle Væren i Almindelighed; men Frembringelsen
% Femdobbelt spatieret „disse“; det FØRSTE „disse“ („af disse, just“) er
% ikke spatieret. Alle sex Forekomster efterprøvede paa KB-Exemplaret.
af disse, just \emph{disse} Elementer med \emph{disse} Qvaliteter, af
\emph{disse} Substanser med \emph{disse} Kræfter i \emph{disse}
Egenskaber, er for hvert særeget Verdenssystem en sig selv betingende,
ubetinget Villiesact, en Skaberact, der i Anlægets Eenhed bestemmer de
særegne Anlæg.
% --- p. 224 ---

Grundstoffernes eiendommelige Beskaffenhed maa udfindes ad Erfaringens
Vei; men Erfaringslæren for sig har ikke til Opgave at forene
Forestillingen om Elementernes Mangfoldighed med videnskabelig Erkjendelse
% TRYKFEIL-BLADET: „S. 224, Lin. 5 f. o.: mechaniske, læs: chemiske.“
% Ordet gjengives som trykt. KB-Exemplaret trykker „mechaniske“ reent og
% urettet; ARBEIDSEXEMPLARET har Ordet overstreget med Blyant og
% „chemiske“ skrevet i Margenen af den tidligere Eiers Haand — Rettelsen
% er altsaa allerede naaet hertil i Arbeidsscanningen og er ikke Sats.
af Naturgrundens Eenhed. De mechaniske Analyser have, især i dette
Aarhundrede ved Voltastøtten, ledet til Opdagelse af en stor Mængde
forhen ubekjendte Grundstoffer. Spørgsmaalet er nu dette: vil en fremtidig
Adskillelse af Substanser, som hidtil ere anseete for Grundstoffer, ende
med en Forøgelse eller en Formindskelse af Grundstoffernes Antal? Ved
Udskillelsen af Kalium, Natrium, Norium, Thorium, Erbium, Terbium o. s. v.
er Antallet unegtelig blevet større og større. Men der er da ogsaa
Antydninger af, at det Vendepunkt engang kunde indtræde, da Resultaterne
vendte sig om, og Bevægelsen gik fra Mangfoldigheden hen imod Eenheden.
Det har viist sig, at f. Ex. Salpeter, langt fra at være, hvad man i
Middelalderen antog, et modificeret Element, er et Salt (salpetersuurt
% Berzelius' Iltprikker: N'et bærer fem Prikker, to over tre. Gjengivet
% $\ddot{\dddot{\mathrm{N}}}$. Læst paa KB-Exemplaret ved stærk
% Forstørrelse; Arbeidsscanningen viser kun en ubestemt Klat.
Kali) og Salpetersyren en Forbindelse af Qvælstof og Ilt
($\ddot{\dddot{\mathrm{N}}}$); det har viist sig, at Cyanet, skjøndt det
forholder sig som Grundstof til de andre Grundstoffer, er en Forbindelse
% De chemiske Formler ere trykt opret med sænkede Tal og Mellemrum
% mellem Tegnene: „C2 N“, „N H4“; saaledes gjengivet.
af Kulstof og Qvælstof (C$_2$ N), og Ammoniet, om hvilket det Samme
gjælder, en Forbindelse af Qvælstof og Brint (N H$_4$). Der er nu dem, der
antage, at man ved at eliminere det ene foregivne Grundstof efter det
andet, tilsidst vil kunne opløse dem alle i nogle ganske faa, ja muligviis
i et eneste. Paa en saa storartet Reduction maatte der da følge en
tilsvarende Deduction. Ligesom man allerede nu kan omdanne Ilten til Ozon,
ja lade den omdannede Ilt optræde baade som Ozon og Antozon, saaledes
skulde man da, naar Idealet var naaet, kunne lade de mange qvalitativt
forskjellige Stoffer opstaae og danne sig af eet og samme fælles
Grundstof.

Hvor høit vi end prise den analytiske Konst, hvem Æren for saa mange store
og frugtbare Opdagelser tilkommer,
% --- p. 225 ---
saa tør vi dog aldrig glemme, hvad den forskende Videnskab bestandig
minder om, at blot empiriske Analyser, om de end fortsættes med nok saa
vidunderligt Held, umulig ved sig selv kunne føre til sand Indsigt i
Naturgrundens Væsen. Standser man ved Forestillingen om en Mangfoldighed
af evigt værende Grundstoffer, saa fornegter man Grundens Eenhed, ja
Grunden selv. Ere Grundstofferne evige, saa ere de jo til uden at være
begrundede, saa træde atomistiske Adskillelser og Forbindelser ved alle
Dannelser overalt i Begrundelsens Sted; Tingene blive da sammenstykkede,
men ikke begrundede. Uden oprindelig Grund bliver Naturen kun et overmaade
fiint mechanisk Flikværk, en naturløs Natur, en Modsigelse. Ender man i
Forestillingen om et fælles, forskjelsløst, almindeligt Grundstof og
sætter dette værende Stof i Grundens Sted, er Modsigelsen endnu mere
iøinefaldende. Dersom forskjelligartede Stoffer skulde fremkomme af et
første Stof, hvori der ingen Forskjelligheder var, da maatte jo
Forskjellighederne opstaae af Intet. Dersom der i det forudsatte Grundstof
virkelig var en Mængde forskjellige Bestanddele, saa var det i
Virkeligheden ikke eet, men flere Stoffer; men naar der slet ingen
Forskjelligheder er deri, saa er der heller ikke Noget, hvoraf
Forskjelsbestemmelser enten umiddelbart eller middelbart kunne afledes.
Forskjellighederne høre lige saa oprindelig og nødvendig med i
Forudsætningen som Eenheden. Forskjelsbestemmelsernes indholdsrige Eenhed,
de uligeartede Stoffers reale Grund, er intet Grundstof, ingen Erfaringens
Gjenstand. Skal Naturgrunden ikke erkjendes, saa blive vi i Modsigelsen;
skal Naturgrunden erkjendes og de elementære Forskjelligheder begrundes,
saa maa Grændsen for den empiriske Viden overskrides.

En anden, Erfaringslæren ledsagende Biforestilling vil ligeledes fordunkle
Opfattelsen af Grundforholdet mellem det Oprindelige og det Afledede.
Enkeltstofferne, de elementære Realiteter, der ikke kunne opløses, vise
sig for Iagttageren
% Arksignaturen „15“ staaer nederst paa S. 225. Den er ikke Text.
% --- p. 226 ---
som det Bestandige, det Oprindelige, medens de sammensatte Stoffer
forestille det Ubestandige og Afledede. Et Dobbeltsalt har da fra dette
Synspunkt mindre Selvstændighed end de særskilte Salte, hvori det kan
adskilles; disse kunne bestaae uden hiint, men hiint ikke uden disse. Af
samme Grund have Saltene mindre Selvstændighed end deres Syrer og Baser,
disse igjen mindre Selvstændighed end Elementerne selv. Denne Synsmaade
% Trykt Mærke: *) — det staaer EFTER Punktummet, som gjengivet.
% Eneste Note i dette Batch; Noten staaer helt paa sit eget Blad.
kan indenfor Realiteternes egen Kreds være fuldkommen berettiget.\footnote{%
Jvnfr. Forelæsninger over „Philosophisk Propædeutik“ fra
Universitetsaaret 1860--61. S. 138.} Men for Erkjendelsen af Elementernes
Forhold til Naturgrunden er den vildledende. Grunderkjendelsen nedsætter
den relative Forskjel mellem enkelte og sammensatte Stoffer til et Afledet
og fastholder det elementære Kredsløb. I Kredsløbet gjælder den uophørlige
Vexel af Opstaaen og Forgaaen lige saa vel for Enkeltstofferne som for de
af dem sammensatte Stofsubstanser. At det Sammensatte opstaaer, betyder,
at Elementerne gaae til Grunde; at det Sammensatte gaaer til Grunde,
betyder, at Elementerne opstaae. Ethvert Stof er et Afledet; Intet er i
Grunden oprindeligt uden Grunden selv.

Den overfladiske, Forskningen uvedkommende, Biforestilling om Elementernes
ubetingede Prioritet forvirrer Begreb og Begriben. Det kommer bestandig
til at see ud, som om de første, faste, uforanderlige Naturbegreber
udelukkende laae i Elementerne, medens de særlige, mere indholdsrige, der
udtale sig i Dannelserne af sammensatte Stoffer, kun vare afledede,
tilfældige, af ydre Omstændigheder afhængige Blandingsbegreber. Det er en
Grundvildfarelse. Elementerne ere, ifølge deres Begreber, udtrykkelig
anlagte paa særlige Dannelser. Først af de særlige Dannelser kan det sees,
hvad de egenlig betyde som Elementer. Vandets Begreb,
% --- p. 227 ---
Qvartsens Begreb, Feltspathens Begreb ere lige saa oprindelige
Naturbegreber, som Begreberne: Ilt, Brint, Silicium, Aluminium, Kalium.
Det er netop de alsidigere, concretere Begreber, der i sig indeholde de
eensidigere; ikke omvendt.

Paa Indsigt i Forholdet mellem abstractere og concretere Begreber beroer
al væsenlig Indsigt i det Teleologiske. Erfaringslæren er i sin gode Ret,
naar den afviser løse Raisonnementer om Lavere og Høiere, om Midler og
Øiemed i Naturens Huusholdning; men Begreber som Naturteleologie og
Naturfornuft kan den ikke afhandle. Naturfornuft og Teleologie beroe paa
Naturanlæg, Naturanlæg paa Naturgrund; men det Spørgsmaal, som
Erfaringslæren ifølge sin Charakteer maa lade ubesvaret, er netop
Grundspørgsmaalet, Spørgsmaalet om Grunden.

% Indløbende Overskrift, halvfed, paa samme Linie som Afsnittet; den
% staaer omtrent midt paa Bladet. Læst paa Siden selv.
\runin{c) Grundmediet.} Magtens Grundfordobling medfører nødvendig en
Fordobling af Grunden. I Selvmediet er Grunden Tankevæsen, Fornuftgrund; i
Gjenstandsmediet er den Stofvæsen, Realgrund. Vel har Fornuftgrunden sin
Sandhed, Bekræftelsen paa sin Idealitet, i den reale Grund, og denne igjen
sin Realitets Bekræftelse i hiin; men hiin er reflecteret i Inderlighed,
denne gaaer op i existentiel Umiddelbarhed: de danne Extremer.
Grundmediet, det immanente Mellemvæsen, hvori Modsætningen af Idealitet og
Realitet uophørlig opstaaer og lige saa uophørlig forsvinder, er
Extremernes Midte.

Den ydre, stofagtige Virkelighed er et Extrem af positiv Umiddelbarhed, en
antilogisk Realitet, der udelukker al Negation, al Reflexion og dermed al
Idealitet. Ethvert Element udtrykker sit Begreb paa en saa realistisk
Maade, at Begrebet forsvinder i Udtrykket. Sige vi: dette Stykke Kul har
en Grændse, Grændse er Negation; altsaa har det i sin Tilværen en
Negation, saa sammenblande vi det, der tilhører Tanken, med det, der
tilhører Gjenstanden. Kulstykkets
% Arksignaturen „15*“ staaer nederst paa S. 227. Den er ikke Text.
% --- p. 228 ---
Grændse er det i Overfladen Positive; de Dele, der ligge i Overfladen,
have en lige saa umiddelbar Existens, som de, der findes i det Indre.
Grændsen som Negation tilhører Tanken: Overfladen er Legemets Ophør, dets
Ikkeværen, det er en Dom af den reflecterende Bevidsthed. Legemets
Ikkeværen vedkommer ikke det værende Legeme: Elementets Negation falder
udenfor Elementet.

Det Virkeliges antilogiske Umiddelbarhed findes antydet i ældgamle
Problemer. Eleaterne negtede Vorden, fordi det Værende maa udelukke det
Ikkeværende. Megarikerne gik videre og negtede Muligheden, fordi det
Mulige i Virkeligheden er en Modsigelse, et Værende, som dog ikke er. At
% „kunne“, „virke“ og det første „er“ ere spatierede; „sagde de“ og
% „slet ikke“ ere det ikke. Alt efterprøvet paa KB-Exemplaret.
% Arbeidsexemplaret har en Blyantsstreg i venstre Margen ud for dette
% Afsnit; den er ikke Text.
\emph{kunne}, sagde de, er at \emph{virke}; hvad Naturen kan, det virker
den; hvad den ikke virker, det kan den ikke
% Græsk, sat med kursiv græsk Skrift; polytonisk, læst paa
% KB-Exemplaret. „φασιν“ staaer uden Accent, som trykt.
(\textit{ὅταν ἐνεργῇ μόνον δύνασθαι, ὅταν δὲ μὴ ἐνεργῇ οὐ δύνασθαι,
φασιν}): kun det Virkelige \emph{er}, det Mulige er slet ikke. Aristoteles
forklarer Vorden og Mulighed ved at oversee Pointen og bryde Spidsen af
Problemet. For at forklare Vorden, henviser han til Muligheden; for at
forklare Muligheden, henviser han til Stoffet, --- det Uforklarlige! Hvad
er Stoffet? Det er Muligheden. Hvad er Muligheden? Det er Stoffet. I den
nyere Tid har Hegel givet et dialektisk Stød til den reale Muligheds
Opklaring. Den immanente Negation er alle Muligheders Hemmelighed; det er
Sagen paa en Prik. Skade, at Negationens Immanens, hvad Naturtingene
angaaer, selv er forblevet en mystisk Hemmelighed.

For at Elementerne kunne virke, maa der i ethvert virkeligt Element være
utallige Muligheder. I isoleret Tilstand maa det existerende Kulstof jo
udtrykke sit Begreb; og kan dog i Virkeligheden ikke udtrykke hele sit
Begreb. Hvor Kulstoffet optræder som Diamant, er kun en Side af Begrebet
realiseret, thi Muligheden er, at det ogsaa kan optræde som Graphit; hvor
det i Virkeligheden optræder
% --- p. 229 ---
% „som i Anthracit“ staaer saaledes i begge Exemplarer; „i“ er altsaa
% Sats og ikke en Plet. Gjengivet som trykt.
som Graphit, er Muligheden den, at det ogsaa kan optræde som i Anthracit.
I Kulstoffets Forbindelse med Brint er Kulsyrens Mulighed udelukket. At
det hele Begreb umulig kan realiseres paa een Gang og paa eet Sted i een
og samme Virkelighed, beroer paa, at Existensen er exclusiv, at den ene
Virkelighedsform nødvendig maa udelukke den anden. Efter sit Begreb er det
% Spatieringen af „baade“, „og“, „enten“, „eller“ og „denne“ efterprøvet
% paa KB-Exemplaret.
rene Kulstof \emph{baade} Graphit \emph{og} Diamant, men i Virkeligheden
er det \emph{enten} Graphit \emph{eller} Diamant. Det er altsaa ikke, hvad
Aristoteles antager, Stoffet, men Begrebet, der er Muligheden; naar
Begrebsformen er blevet Stofform, have vi Virkeligheden. Ethvert Element
er altid noget Andet i Virkeligheden, end det er i Muligheden;
Virkeligheden er hver Gang kun \emph{denne}, kun een, men Mulighederne ere
mangfoldige. Begrebets mangfoldige Muligheder have en negativ Væren i den
ene positive Virkelighed, hvormed et Element existerer. Denne
Mulighedernes negative Væren er det Virkeliges immanente Negation, dets
Idealitet, dets Indre.

Men efter, hvad vi ovenfor have seet, har det Virkelige i Naturen jo intet
Indre; Naturtingen er i antilogisk Umiddelbarhed, hvad den er, hverken
mere eller mindre. Saa strides der om Forholdet mellem Naturens Indre og
dens Ydre; Striden gaaer baade paa Vers og Prosa. Imod Hallers strenge
% De to tydske Verscitater staae i almindelig Antiqva, uden
% Anførselstegn og uden Kursiv; efterprøvet paa KB-Exemplaret.
Adskillelse af det Indre og det Ydre: Ins Inn're der Natur dringt kein
erschaffner Geist, sætter Goethe sin intuitive Genius i Pant paa det
Indres umiddelbare Eenhed med det Ydre: Natur hat weder Kern noch Schaale,
Alles ist sie mit einem male. Hvem har nu Ret, Haller eller Goethe? I
Realiteten er der overalt Eenhed af det Indre og det Ydre, thi Realiteten
er, at det Indre gaaer op i sit Ydre. Men Realiteten er endelig; imod
Realitetens Endelighed staaer Idealitetens Uendelighed, og saa er her dog
en afgjørende Modsætning mellem det Indre og det Ydre.
% --- p. 230 ---

Idealiteten, det Indre, svarer til Tankevæsenet, Fornuftgrunden; i
Fornuftgrunden ere Elementbegreberne Ideebegreber: der er i enhver
Naturexistens et uendeligt Overskud af Idealitet. Realiteten i sin
Umiddelbarhed svarer til Stofgrunden, Realgrunden; thi Existensens reale
Grund gaaer med de givne Betingelser umiddelbart op i den begrundede
Existens.

For Anskuelse og Forestilling kommer det da let til at see ud, som om
Fornuftgrund og Realgrund afgave to paa hinanden liggende Grundlag, saa at
Realgrunden gik sammen med Virkeligheden, medens Fornuftgrunden var at
betragte som en Slags usynlig Beholder af Mulighederne. Virkeliggjørelsen
af en Mulighed blev da at opfatte saaledes, at Muligheden kom ved en
Tilstrømning fra det Indre, og fik ved at gribe ind i de ydre Betingelser
Magt til at forandre det Bestaaende, tage Virkeligheden fra det forud
Givne og derved give sig selv Virkelighed. Paa denne Maade skulde altsaa
enhver ny Mulighed fra Idealiteternes Dyb komme som en Erobrer, alliere
sig med en Kreds af Betingelser og rive Virkeligheden til sig. Et
unaturligt Maskineri! Man forestiller sig Elementerne svømmende i et Hav
af Muligheder; og saa er der imellem Mulighed og Virkelighed ingen
substantiel Eenhed. Man kunde da lige saa godt forestille sig, at de
elementære Realiteter ikke vare Andet end Anledningsaarsager, løse
Brikker, som en almægtig Villie stillede op og flyttede efter Behag.

Al Naturvirksomhed beroer paa en substantiel Eenhed af Mulighed og
Virkelighed, men denne Eenhed beroer igjen paa en gjennemgaaende
Modsætning af ideel og reel Magt. Paa Forestillingens Maade veed jo
vistnok Enhver, hvad der saadan i Almindelighed menes med en
Naturvirksomhed. Saa antager man, at Virken og Virksomhed er Noget, der
ingen Forklaring behøver, og forklarer uden videre Overgangen fra Mulighed
til Virkelighed ved Virksomhed, idet
% --- p. 231 ---
man særlig henviser til Aarsager, Kræfter, Betingelser som det Virksomme.
Man glemmer, at netop det, hvorved man forklarer, i høi Grad trænger til
en Forklaring. Hvad nytter det at beraabe sig paa „det Virksomme“ og sige:
% Fire spatierede „virke“ i dette Afsnit; efterprøvet paa KB-Exemplaret.
% Anførselstegnene om „det Virksomme“ ere lav-høie, som trykt.
Aarsagerne \emph{virke}, Kræfterne \emph{virke}, Betingelserne
\emph{virke}, naar man ikke er i Stand til at gjøre sig Rede for, hvad der
da egenlig ligger i det Begreb: at \emph{virke}.

Realiteten for sig, den antilogisk bestemte Virkelighed, kan ikke virke,
thi den mangler Mulighed. To Stykker Træ kunne f. Ex. ikke gnides mod
hinanden uden at det ene maa, hvad Indtrykket angaaer, afgive Mulighed for
det andets Virkelighed; idet de gjensidig ere baade mulige og virkelige
for hinanden, virke de, og der frembringes Varme. Ilt og Brint kunne
blandes uden at opgive deres umiddelbare Virkelighed; der bevirkes da ikke
nogen ny Frembringelse. Men dersom en elektrisk Gnist slaaer igjennem
Blandingen, saa oplades Muligheden: Elementerne virke og frembringe Vand.
Muligheden for sig, den rene Idealitet, kan heller ikke virke, thi den
mangler Virkelighed. Den i Vandbegrebet liggende Mulighed kan ikke
realisere sig selv; den maa have Ilten og Brinten til Virkelighed for at
kunne virke. Virken er under alle Former en uendelig Omsætten af Ideelt i
Reelt, af Mulighed i Virkelighed, og omvendt.

Denne Omsætning forudsætter, at det Ideelle er Magt. Dersom Idealiteten
ikke var Magt, kunde den umulig have Virkelighedsmomentet i sig og udvirke
sin egen Realitet. Al Virken er, naar den sees fra Mulighedens Side, et
kategorisk Udtryk for, at det Ideelle magter det Reelle. Men Omsætningen
forudsætter ogsaa, at det Reelle fra sin Side er Magt og magter det
Ideelle. Hvis den umiddelbare Realitet ikke havde Idealiteten, Muligheden,
som et skjult Moment, en iboende Negation, i sig selv, var den jo
afmægtig overfor Muligheden; den kunde da ikke være med i Virksomheden,
den kunde slet ikke virke. Fra Virkelighedens
% --- p. 232 ---
Side er al Virken et kategorisk Udtryk for, at det Reelle magter sit
Ideelle.

% Græsk og Latin staae begge med Kursiv, adskilte af et almindeligt
% Komma; læst paa KB-Exemplaret.
Mulighed er en Formaaen (\textit{δύναμις}, \textit{potentia}), dens Væsen
er en ideel Magt. Virkelighed, antilogisk bestemt, er en Værens
Selvstændighed, er Modstand, Gjenstand; dens Væsen er en reel Magt.
Virkelighed er Udslag af Magt; Mulighed er Idealitetens Overskud. Vand er
efter Mulighed Iis, og dog er her en qvalitativ Forskjel. Overgangen fra
Vand til Iis er en brat, afgjørende Omsætning; Kuldegraderne med de
forberedende Mellembestemmelser tabe sig i Disjunctionen: enten Vand eller
Iis. Saalænge Vandet er virkeligt, er Isen kun mulig; den nye
Tilstandsform, hvori Isen er virkelig og Vandet muligt, beroer paa et
Udslag af Magt. Magtens Udslag kan iøvrigt være af høist forskjellig Art;
% spacing.py flagede „Art“ paa denne Side; KB-Exemplaret viser
% almindelig Antiqva. Kaldet forkastet.
men den er altid kjendelig paa den qvalitative Omsætning. Den reale
Mulighed kræver Betingelser. Betingelserne kunne indtræde, een efter een;
saalænge en eneste Betingelse mangler, er her kun en Mulighed og et Antal
Betingelser; men med den sidste Betingelses Indtræden forandres Alt;
Muligheden er ikke længer Mulighed, Betingelserne ikke Betingelser; det
Hele er en ny Virkelighed, det er Udslag af reel Magt. Betingelsernes
Forvandling til Momenter i Grunden er Udslagets Væsen. Den reelle Magt
magter Betingelserne, thi den forvandler dem; og dog er dette kun den ene
Side af Sagen. Det Modsatte er, at den reelle Magt er afhængig af
Betingelserne, idet den er afhængig af Muligheden og den ideelle Magt. Paa
den urolige, dog i Grunden rolige, Eenhed af reel og ideel Magt beroer al
Virken og Virksomhed. Elementerne forbindes, og de sammensatte
Stofsubstanser fremkomme ved Udslag af reel Magt --- paa Grund af Magtens
Idealitet. Det Sammensatte adskilles, den mægtige Idealitet oplader sine
Muligheder, og Elementerne vende tilbage --- i Magtens reelle Udslag.
% S. 232 slutter her, med fuldt afsluttet Sætning og Afsnit, en Linie
% kortere end Satsspeilet; ingen Streg ved Bladets Fod. Afdeling C
% begynder paa S. 233 (næste Batch).

%  --- C. Elementer og Legemer  (pp. 233–249) ---

% --- p. 233 ---
% S. 233 aabner med et NYT Afsnit --- første Linie er indrykket paa Siden
% (kontrolleret paa KB-Exemplaret) --- saa Afdeling B slutter her med to hele
% Afsnit i Bladets øverste Trediedeel.
Grundmediet, det almene Væsen, hvori Fornuftgrunden er Eet med
Realgrunden, reflecterer i sit Indre, hvad der udstykkes i det Ydre, og
forener i det Ydre, hvad der er forskjelligt i det Indre. Grundmediet
udtrykker Forskjelseenheden af Natur og Aand; i Grundmediet er det
Naturlige aandigt og det Aandige naturligt. At et existerende Element
kan bære alle sine Muligheder, sine immanente Negationer, i umiddelbar
Existens, er kun muligt derved, at den reelle Magt, hvis Udslag
bestemmer dets indskrænkede Virkelighed, er i Grundmediet Eet med den
indre ideelle Magt, der bærer Idealitetens Overskud.

Saa er her altsaa en Immanens, en oprindelig Identitet af Tankevæsen og
Stofvæsen. Ja! Men Eenheden er grundet paa en Modsætning, Immanensen
paa en Transcendens; Naturgrunden selv er en afledet Grund. Naturen
begrundes af Aanden.

% Kort centreret Streg (enkelt) mellem Afdeling B og den følgende Afdeling,
% omtrent midt paa Bladet. Læst paa KB-Exemplaret.
\streg

% AFDELINGSOVERSKRIFT MED FORKERT BOGSTAV --- og Overskriften aabner IKKE
% Siden, men staaer omtrent midt paa S. 233, efter Stregen.
% Trykfeilbladet retter S. 233, Linie 13 f. n.: „B. --- læs C.“; ogsaa
% Rækkefølgen (A. paa S. 197, B. paa S. 215) kræver C. Trykken har „B.“,
% læst paa KB-EXEMPLARET.
% ARBEIDS-EXEMPLARET ER HER RETTET MED BLYANT: ved Siden af det trykte „B.“
% staaer et haandtegnet „C“, lysere end Trykken, og selve B'et er overstreget
% med en lodret Blyantsstreg; derfor giver OCR'en „B.C“. Dette er det syvende
% Sted, hvor Blyantshaanden har anvendt Nielsens Errata --- jvnfr. det ganske
% tilsvarende Tilfælde ved den indløbende Overskrift paa S. 180 (c) --> b)).
% Gjengivet som trykt, ikke omnummereret.
% Overskriften „Elementer og Legemer.“ er sat med halvfed Antiqua.
\afdeling{B.}{Elementer og Legemer.}

% Spærringen af „Legemlighed“ er bekræftet paa KB-Exemplaret.
Den naturlige Forskjel mellem Element og Legeme beroer paa en
umiddelbar Adskillelse af Selvstændigt og Uselvstændigt, en bevægelig
Grændsebestemmelse, som særlig kommer tilsyne, naar Begrebet
\emph{Legemlighed} undersøges. At de elementære
Selvstændighedsbegreber ere Magtbegreber, giver sig i Kjendsgjerninger
tilkjende derved, at Begrebsbestemmelserne, uagtet deres indre
væsenlige Sammenhæng nødvendig maa forudsættes, ikke paa rationel
Maade lade sig aflede af hverandre; den factiske Sammenhæng erkjendes
% Hele Vendingen „af de legemlige Egenskaber“ --- ogsaa de to smaa Ord ---
% er spærret paa Siden; bekræftet paa KB-Exemplaret.
\emph{af de legemlige Egenskaber}. Legemernes Selvstæn\-%
% --- p. 234 ---
% Spærringen af „Legemernes Individualitet“ er bekræftet paa KB-Exemplaret.
dighed er en stofagtig Substantialitet, en selvløs Selvstændighed, som
i den elementære Natur har sit Maal og sin Maalestok i \emph{Legemernes
Individualitet}.

% Indløbende Overskrift, læst paa Siden selv (KB-Exemplaret) og ikke i
% Indholdet: sat med halvfed Antiqua, paa samme Linie som Afsnittet.
\runin{a) Legemlighed.} Af det Foregaaende sees, at Ordet Element kan
tages i meget forskjellig Betydning; snart er det Grundatomerne, snart
Enkeltstoffernes Masser, snart igjen de almeenlegemlige
Tilstandsformer og Virksomheder, hvortil der sigtes. Ethvert Element
er et Legeme, hver mindste Deel af Materien udfylder et Rum; men ikke
ethvert Legeme er et Element. Ved Bestemmelsen af Elementet er man
oprindelig gaaet ud fra, at det maa være det absolut Enkelte, det i
sig Selvstændige; men Enkelthed og Selvstændighed svare ikke ligefrem
til hinanden. Atomet, det abstract
% TRYKFEIL: „Eukelte“ for „Enkelte“ --- et omvendt n. Bogstavet er læst paa
% KB-Exemplaret, hvor u'et staaer utvetydigt, og kalibreret mod det rigtigt
% satte „Enkelte“ tre Linier ovenfor og „det Enkelte i Grunden“ nedenfor paa
% samme Side. Trykfeilbladet tier. Gjengivet som trykt; læs „Enkelte“.
Eukelte, er et Uselvstændigt; Atomerne bestaae i Massen og kjendes kun
af de Vægt- og Maalforhold, hvori Masserne forbinde sig.%
\footnote{% trykt Mærke: *) --- staaer EFTER Punktummet: „forbinde sig.*)“
Om chemiske Proportioner, Æqvivalenter og Atomer jvnfr.
Forelæsninger over Phil. Propæd. 1860--61. Side 200 flgd.}
Om ogsaa Grundstofferne i sig selv ere chemisk uadskillelige, saa ere
de dog derfor ingenlunde absolut selvstændige; et Blik paa det
elementære Kredsløb viser, at det Enkelte i Grunden ikke har større
Selvstændighed end det Sammensatte. Hegel protesterer mod Chemiens
udelukkende Ret til at bestemme, hvad et Element er; hos ham er
Grændsen mellem Selvstændigt og Uselvstændigt
% Spærringen af „logisk“ er bekræftet paa KB-Exemplaret.
\emph{logisk} flydende. Den eiendommelige Maade, hvorpaa han dialektisk
har opfrisket Læren om de fire gamle Elementer er for charakteristisk
til at den her tør forbigaaes.

Materien, hedder det, har Individualitet; den har i Solsystemet lagt
for Dagen, at den ud af sit eget Væsen kan bestemme og forme sig selv.
Hvad Tyngden er i det Mechaniske, er Lyset for de elementære
Qvaliteter. Lyset er den første qvalificerede Materie, den
reflecterede Isigværen, Naturens umiddelbare, almene Selv. Lysets ydre
Tilværelse
% --- p. 235 ---
% Spærringen af „Stjernen“ og af det FØRSTE „Solen“ er bekræftet paa
% KB-Exemplaret; det andet „Solen“ (Linie 2) er almindelig Antiqua.
er \emph{Stjernen}, som Moment i en Totalitet \emph{Solen}. Solen er
det tilværende, af Planeterne selv forudsatte Lys. Lyset selv er intet
Element, men Principet for Elementernes Udvikling. Lyset har nemlig
ikke blot Mørket, det har som qvalificeret Materie i Idealitetens Form
tillige en Fleerhed af qvalificerede Materier i Realitetens Form til
Modsætning. Det er disse Realiteter, Lysvæsenets immanente Negationer
og derhos Materier for sig i udvortes Tilværen, der optræde som fire
Elementer, svarende til de fire, i Solsystemet udtalte Naturer: Luften
til den solare, Ilden til den lunare, Vandet til den cometare og
Jorden til den planetare Natur.

Luften er Almindelighedens Element, ikke i Selvhedens Form, men
umiddelbart som Andetværen; Luften, det til Væren for Andet
degraderede Lys, er tung; Lyset selv er uden Tyngde. Passiv mod
Lyset, gjennemsigtig, mechanisk elastisk, en snigende, tærende Magt,
er Luften overalt i Begreb med at opløse og forflygtige det
Individuelle.

% Kun „Modsætningens“ er spærret; „Særskilthedens,“ staaer i almindelig
% Antiqua. Kontrolleret ved Forstørrelse paa KB-Exemplaret.
Særskilthedens, \emph{Modsætningens}, Elementer ere Ilden og Vandet.

I Maanen er den umiddelbare Forsigværen stivnet til et svovlagtigt
Forsigværende, et Ildprincip uden Ild --- Læseren vil maaskee bedst
kunne tænke sig det som en frossen Ild --- denne Stivhed (Starrheit)
gaaer nu i Elementernes Region over til for sig værende Uro, til
brændende Ild. Luften er, hvad der viser sig, naar den comprimeres, i
Grunden Ild. I Ilden bliver altsaa Luften udtrykkelig sat som det, den
i Begrebet er, som negativ Almindelighed eller sig til sig refererende
Negation. Ilden er den materialiserede Tid eller Selvhed (Lys identisk
med Varme), det slet og ret Urolige og Fortærende --- en Fortæren af
Andet, der med det Samme fortærer sig selv og saaledes gaaer over i
Neutralitet.

Det cometare Element, Modsætningens andet Led, det Neutrale, er
Vandet. Som den i sig sammensunkne Mod\-%
% --- p. 236 ---
sætning har Vandet ingen for sig værende Enkelthed, ingen Fasthed
eller Bestemthed i sig selv. Her er i dette Element en gjennemgaaende
Ligevægt af mulige Tilstande; enhver mechanisk Bestemthed, som sættes
deri, opløses. Formgrændsen modtages fra Andet og søges i Andet
(Adhæsion). Vandet er flydende, ligesom Luften, men ikke elastisk
flydende, saa at det expanderer sig til alle Sider. Det er mere
jordisk end Luften, søger et Tyngdepunkt, staaer det Individuelle
nærmest og stræber hen derimod, fordi det væsenlig er den concrete
Neutralitet, der endnu ikke er
% Spærringen af det enkelte Ord „sat“ er bekræftet paa KB-Exemplaret.
\emph{sat} som concret.

% Spærringen af „Jordagtige“ er bekræftet paa KB-Exemplaret.
Det planetare, concrete Element, Jorden, er for det Første det endnu
ubestemte \emph{Jordagtige} overhovedet; nærmere bestemt, er det den
Totalitet, der i sig optager de andre Elementer og sammenfatter
Forskjellighederne til individuel Eenhed. Det er Jordens Bestemmelse
at tilegne sig de almene astrale Magter, der som himmelske Legemer
have Skin af Selvstændighed, fordøie dem, bringe dem ind under
Herredømmet af en Individualitet, i hvilken disse Kæmpelemmer
nedsættes til Momenter. Den første store Fordøielsesact er den
elementære Proces.

De physiske Elementer, tilføies der, ere ikke individualiserede i sig,
men skikkelsesløse; de ere Legemlighedsmomenter, men ikke bestaaende
Legemer. Det uorganisk Individuelle er for Begrebet det mest
Haardnakkede, fordi Materien her er specificeret til Individualitet;
og dog er Individualiteten umiddelbar, og som Qvalitet umiddelbart
identisk med det Almene.

Dette er nu Hegels Grundtanke; det er den sidste og fuldstændigste
Udvikling af Immanenslæren, vi her have for os. Det gjælder om den
Hegelske Deduction, hvad der ovenfor er fremhævet i Anledning af den
Aristoteliske. Spørgsmaalet er ikke nærmest, hvorvidt den givne
Forklaring tilfredsstiller Physiken --- her er jo med Undtagelse af
nogle
% --- p. 237 ---
Smaafeil egenlig saa godt som Intet, der vedkommer Physiken ---
Grundspørgsmaalet er, om en saadan naturphilosophisk Methode har
Gyldighed, og hvilket Værd der kan tilkjendes den i Videnskaben.

% Schellings Ord ere satte med CURSIV (Antiqua-Cursiv mod den omgivende
% Antiqua), læst paa KB-Exemplaret. Ingen Spærring.
\textit{Ueber die Natur philosophiren heisst die Natur schaffen} ---
disse Ord af Schelling, som Udgiveren af Hegels Naturphilosophie har
valgt til Motto%
\footnote{% trykt Mærke: *) --- staaer FORAN Kommaet: „til Motto*),“
G. W. F. Hegels Vorlesungen über die Naturphilosophie als der
Encyclopädie der philosophischen Wissenschaften im Grundrisse zweiter
Theil. Herausgegeben von D. C. L. Michelet. Berlin 1842.}, ere ret
egnede til at sætte det immanente Standpunkt i lueklar Belysning. Der
er noget Orakelagtigt, noget Ypperstepræsteligt, noget høit
Dogmatisk-Bispeligt i den speculative Selvvished, hvormed Naturens
Hemmeligheder her blive aabenbarede. Det seer unegtelig ud, som havde
Hegel selv skabt de fire Elementer, for at indføre dem, om ikke i
Naturen, saa dog i Naturphilosophien.

Det er i Sandhed ikke Mangel hverken paa videnskabelig Alvor eller
logisk Dybde, men en i Immanenslæren rodfæstet Misforstaaelse af det
objective Begrebs Selvbegriben, som i denne abstruse Elementskabelse
lader sig tilsyne. Opfattelsen af Legemlighedsbegrebets Forhold til
den existerende Legemlighed er i Bund og Grund forfeilet. Her er intet
Grundmedium for rationel-reel Nødvendighed, thi her er ingen
qvalitativt bestemte Extremer; her er, med andre Ord, intet
tilforladeligt Udgangspunkt, hverken i det Aprioriske eller i det
Empiriske. I det Aprioriske ikke; thi Begreber som Luft og Ild, Vand
og Jord lade sig ikke med logisk Nødvendighed aflede af den rene
Tænken. Elementerne ere ikke Tankevæsener, men Stofvæsener. Saa
indsmugles Bestemmelser, laante fra Sandsningen og den almindelige
Dagligdagserfaring. I det Empiriske ikke; thi
% --- p. 238 ---
det skal jo dog lade, som om de indsmuglede Sandsebestemmelser ikke
skyldtes Erfaringen, men bleve paa det Strengeste afledede af det
immanente Begreb. Det logiske Begreb gjentager sig med subjectiv
Nødvendighed i tre Hovedformer: det Almene, det Særskilte og det
Enkelte; ergo skal det med objectiv Nødvendighed gjentage sig i fire
Elementer. Og hvorfor? Fordi Naturen, der ikke kan fastholde
Begrebet, maa lade Særskilthedens Momenter falde ud fra hinanden i to
Led og desaarsag behøver to Elementer, Ild og Vand, til at udtrykke
det Særskilte. Man faaer et Indtryk af, at den sidste Tilsætning
skriver sig, ikke fra nogen høiere Begriben, men fra en slet og ret
Hverdagserfaring. Med Anvendelsen af de fire Naturer i Solsystemet
seer det endnu betænkeligere ud; her er slet Intet, hverken sandseligt
eller tænkeligt, at holde sig til. At Maanen er Svovl, er vel omtrent
lige saa videnskabeligt, som at den er en grøn Ost. Hvilket Spring
fra Ildprincipet i Maanen til en brændende Ild paa Jorden!
Allerværst er det dog med Cometvandet. Sæt, at Cometerne vare
uvandede, og det kunde endda nok være, saa er der jo slet intet
Vandprincip. Men naar der intet Vandprincip er i Cometerne, maa det
enten gaae ud over „det immanente Begreb“, som da ved den nye
Comettheorie lider et Knæk, det aldrig forvinder; eller, hvis
„Begrebet“ absolut skal have Ret, maa det gaae ud over den stakkels
Natur, thi hvorledes skal Naturen kunne holde sit Vand, naar der intet
Vandprincip er.

Hegel gjør gjældende, at det ved slige Deductioner egenlig ikke
kommer an paa Beskaffenheden af den ydre, materielle Existens;
Hovedsagen er, at man overalt i den legemlige Verden kan see Begrebet.
Paracelsus, bemærker han, har sagt, at alle jordiske Legemer bestaae
af fire Elementer, af Mercurius, Svovl, Salt og den jomfruelige Jord,
ligesom man jo ogsaa havde fire Cardinaldyder. Sligt maa ikke
forstaaes efter Bogstaven; den høiere Mening er, at den
% --- p. 239 ---
reale Legemlighed har fire Momenter. --- Man lærer da heraf, at det
under disse Forudsætninger ikke hjælper at lægge Vind paa Realismen
ved at beraabe sig paa Eenheden af det Reelle med det Ideelle. Naar
det Ideelle ikke forlanger mere, end at den reale Legemlighed skal
have fire Momenter, saa kan man jo lige saa godt søge Momenterne hos
Paracelsus og Jacob Böhme som hos Schelling og Hegel.

% SPÆRRING, DER BRYDES VED LINIESKIFTET. Paa Siden staaer „E x i s t e n s-“
% spærret som sidste Ord i Linien, medens „grændse;“ øverst i næste Linie er
% almindelig Antiqua. Kontrolleret ved Forstørrelse paa KB-Exemplaret.
% Gjengivet som trykt: kun Forleddet er sat som Fremhævelse.
Grændsen for det i Legemverdenen Selvstændige og Uselvstændige er en
bevægelig, men dog lovbestemt \emph{Existens}grændse; Momenterne i
Legemlighedens Begreb ere virkelige Kjendsgjerninger.

% Indløbende Overskrift, læst paa Siden selv (KB-Exemplaret): halvfed Antiqua.
\runin{b) Legemlige Egenskaber.} Forholdet mellem
Legemlighedsbegrebet og det existerende Legeme er, hvad vi alt fra
flere Sider have seet, et Magtforhold. I Grundmediet er Eenheden af
Begreb og Virkelighed reflecteret som Virken og Virksomhed; men i
Gjenstandsmediet er Begrebet saa umiddelbart Eet med sin Existens, at
Gjenstanden staaer i Begrebets Sted. Af hvad Natur et Legeme er, maa
kjendes paa dets Oprindelse; men Legemets Beskaffenhed, den hele
Eiendommelighed, hvormed det existerer, maa kjendes paa dets
Egenskaber. Legemet med sine Egenskaber er et umiddelbart Givet; det
Givne vil, i antilogisk Umiddelbarhed, fra alle Sider i alle Maader
sees og sandses, saaledes som det er givet. Denne
Gjenstandsnødvendighed gjør uvilkaarlig Indtryk af Magt; men lige saa
uvilkaarlig sætter Gjenstandsmagten igjen Reflexionen i Bevægelse. Det
er Iagttageren umuligt at sandse uden at tænke: dette er et Legeme,
dette Legeme er haardt, er guult o. s. v. Saa bliver Legemet logisk
Subject (Substrat), og Egenskaberne Prædicater; thi der erfares kun,
idet der tænkes og dømmes. Det kommer nu an paa, hvad Retning
Reflexionen tager. Den reflecterende Bevidsthed kan, hvad
Videnskabens Historie lærer,
% --- p. 240 ---
springe Grundmediet over og kaste sig paa et af Magtens Extremer,
enten paa Selvmagten i al Vilkaarlighed eller paa Gjenstandsmagten i
al Tilfældighed. Den scholastisk theologiske Metaphysik er et Exempel
paa den første Eensidighed, den empiriske Skepticisme paa den sidste.

% Latinen er sat med Cursiv, læst paa KB-Exemplaret; Bogen skriver her
% „distingvendum“ med qv-Formen (jvnfr. Vexlen mellem qv og qu i Latinen
% paa S. 145). Gjengivet som trykt.
Legemets Existens er et Magtens Ord, en Aabenbarelse af Selvmagtens
Villie; det er Scholastikens theologiske Forudsætning. Hvorledes denne
Forudsætning i sin Tid blev anvendt, sees af den katholske
Forvandlingslære. Brød og Viin forvandles til Legeme og Blod og
beholde dog deres sandselig-legemlige Egenskaber. Hvorledes er nu
dette muligt? Her maa skjelnes --- \textit{distingvendum est} ---
mellem Substans og Accidenser. Den vilkaarlige Selvmagt, der af Intet
har skabt Substansen og meddeelt den visse Accidenser, kan efter Behag
lade to uligeartede Substanser være fælles om eet og samme Indbegreb
af Accidenser eller ogsaa omskabe Substansen, men lade dens
Accidenser forblive uforandrede. Naar en Muus æder det
Høiærværdige, hvad faaer saa Musen? Brød eller Brødhed
(\textit{panis} eller \textit{paneitas})? Gud veed det, svarer Petrus
Lombardus. Et træffende Svar. I en naturløs Natur er Baandet mellem
Substans og Accidenser saa vilkaarligt, at Almagten selv umiddelbart
løser og binder det, efter som Præsten messer.

% Brøken er trykt som ægte Brøk (1 over 20000) midt i Linien.
Den empiriske Skepticisme tager sit Udgangspunkt i det umiddelbart
Givne; det Givnes Nødvendighed er skjult i Tilfældighed:
Gjenstandsmagten er antilogisk i Dette og Dette og Dette. Viden beroer
paa Erfaring, men det, der erfares, er hver Gang et Enkelt, og hvert
Enkelt et Tilfældigt. Vi mene at vide, at der er en Substans, der
kaldes Guld; men, at det er en Indbildning, kan sees af den Maade,
hvorpaa Guldets Egenskaber sammenstilles. Guldet, hedder det, er et
guult Metal, der kan krystallisere i Oktaedre; det er saa smidigt, at
det lader sig udvalse til Blade af $\frac{1}{20000}$ Linies Tykkelse
o. s. v., lutter empiriske Kjendemærker,
% --- p. 241 ---
Skulde det lade sig godtgjøre, at Guldet var en Substans, og
Guldsubstansen disse Egenskabers (Accidensers) nødvendige Eenhed, da
maatte det kunne paavises, hvorledes hver enkelt Egenskab fulgte af de
andre og af den Forbindelse, de danne, ligesom det f. Ex. ved en
Trekant kan paavises, at hver enkelt Side afhænger af de to andre og
af den Vinkel, de indeslutte. Men dette er saa langt fra at være
Tilfældet, at tvertimod hver enkelt Egenskab erfares for sig uden
mindste Antydning af, hvad den egenlig har med de andre at gjøre. At
være guult, at være smidigt, at krystallisere i Oktaedre er i
Virkeligheden tre saa uligeartede Ting, at Ingen let falder paa at
forestille sig dem i uadskillelig Eenhed. At Guldet iblandt alle
Metallerne har den mindste Affinitet til Ilt, anføres i Chemien som en
almeengyldig Sætning. Men hvorledes beviser man denne Sætning?
Begrundet paa Erfaring kan Sætningen ikke være; thi man
% De tre Spærringer i denne Sætning ere alle bekræftede paa KB-Exemplaret;
% „Tilfælde“ er hver Gang almindelig Antiqua.
kan erfare en \emph{Mængde enkelte} Tilfælde, men ikke \emph{alle}
Tilfælde, endnu mindre \emph{et almeengyldigt} Tilfælde. Man beraaber
sig paa Inductionen; det vil sige: man sætter nogle Tilfælde for alle
Tilfælde; man kunde lige saa godt sætte $1 = 2$, $2 = 3$,
$3 = \infty$.

For at undgaae Vilkaarlighedsbestemmelser af Selvmagten og
forstyrrende Tilfældigheder i Gjenstandsmagten har Immanenslæren, det
Absolutes Philosophie med den absolute Viden, rask væk slaaet en Streg
over begge Extremer; den vil i Kraft af sit ideel-reelle Princip
articulere den rationelle Nødvendigheds Indhold; men Articulationen
strander paa de legemlige Egenskabers vilkaarlig-tilfældige
Bestemthed.

% Ved Foden af S. 241 staaer Arksignaturen „16“. Den er ikke Text.
Immanenslæren staaer i principiel Strid med Physiken. Vel er Physiken
hævet over Skepticismens Løsagtighed, thi den hævder Fornuftens Ret og
erkjender Nødvendigheden i „de uforanderlige Love“. Men den rationelle
Nødvendighed er i Naturen reel, og den reelle Nødvendighed
uadskillelig fra
% --- p. 242 ---
Tilfældighed. I hver legemlig Egenskab har det Nødvendige sin
umiddelbare Væren i det Tilfældige. Det Tilfældige er, at Guldet er
guult, at det er smidigt, at det har stærkere Affinitet til Chlor end
til Ilt. Men dette Tilfældige er netop Existensens Udtryk for en
articuleret Nødvendighed, en lovbestemt Sammenhæng: Physiken stoler
paa Inductionen, fordi den stoler paa Loven; den erkjender i Guldet en
Substans, fordi den i Egenskabernes Forbindelse erkjender en factisk
Nødvendighed. At Nødvendigheden er factisk, uigjennemsigtig, udtaler
Physiken ved at erklære, at vi kun kjende Tingenes Egenskaber, ikke
deres Væsen.

% Spærringen af „magnetisk“ og af „hele Theorie“ er bekræftet paa
% KB-Exemplaret; „Magnetismens“ er almindelig Antiqua.
% ANFØRSELSTEGN OVER BLADSKIFTET: Hegel-Citatet aabnes her („Denne
% Tilsyneladelse“ --- Indskud --- „er Magnetismens hele Theorie ...“) og
% lukkes først paa S. 243 ved „Hemmelighedsfuldt.“. Siden er derfor +1 i
% Anførselsbalancen og S. 243 tilsvarende -1; det er en Anførsel over
% Bladskiftet, ikke en Trykfeil.
Imod en saadan Udtalelse mener Immanenslæren at kunne sætte en
Forklaring som den, at Tingens Væsen reflecterer sig i dens
Egenskaber, saa at en Erkjendelse af Egenskaberne nødvendig maa gaae i
Eet med en Erkjendelse af Væsenet. En saadan Forklaring er god nok i
en Logik; men hvordan tager den sig ud, naar den skal gjennemføres i
en Physik? Lad os f. Ex. vælge den Egenskab ved et Legeme, at det er
\emph{magnetisk}, for at see, hvad Immanensphilosophien formaaer at
oplyse om denne Egenskab. Magnetismens Væsen er, ifølge Hegel, aldeles
gjennemsigtigt i sin Tilsyneladelse (Erscheinung). Det Væsenlige er,
at Begrebets differente Momenter ikke paa udvortes Maade blive satte
identiske, men at de sætte sig selv identiske. Det sprøde Punkt
oplader sig til Adskillelse i Begrebet, saaledes have vi Polerne. Paa
den physiske Linie, der har Formens Forskjel i sig, er det de to
levende Ender, af hvilke enhver er sat saaledes, at den kun er i
Relation til sit Andet, og ingen Betydning har, naar det Andet ikke
er. Begge ere een Determination. „Denne Tilsyneladelse“, tilføier
han, „er Magnetismens \emph{hele Theorie}. Physikerne sige, at man
endnu ikke veed, hvad Magnetisme er, om det er en Strøm o. s. v. Alt
dette hører til den Slags Metaphysik, der ikke
% --- p. 243 ---
% Ved Foden af S. 243 staaer Arksignaturen „16*“. Den er ikke Text.
bliver anerkjendt af Begrebet. Magnetismen er intet
Hemmelighedsfuldt.“

Efter en saadan Udtalelse skulde man nu vente, at det i Begrebet
opklarede Væsen maatte udbrede et næsten guddommeligt Lys over den
ellers saa gaadefulde Forskjel mellem magnetiske og ikkemagnetiske
Legemer; men Forventningen skuffes. „Paa hvilke Legemer Magnetismen
kommer til at lægge sig for Dagen, er“, hedder det, „Philosophien
fuldkommen ligegyldigt. Fornemmelig findes den i Jernet, men ogsaa i
Nikkel og Kobolt.“ Vi lære af denne Erklæring, at de virkelige
Legemers virkelige Egenskaber staae i et aldeles udvortes, tilfældigt
og ligegyldigt Forhold til Legemlighedens Væsen; og endda er Legemets
Væsen immanent i dets Egenskaber! Tilfældighedsmærket er uadskilleligt
fra Kjendsgjerningen; det er Facticitetens umiddelbare Bestemthed,
dens eget antilogiske Moment. Immanenslæren opløser ikke
Tilfældigheden, fordi den fornemt overseer det Tilfældige; den knækker
ikke Virkelighedens haarde Nød, fordi den kaster Kjendsgjerningerne
bort som ubrugelige Skaller.

At Immanensphilosophien ikke kan komme til Rette med
Gjenstandsmagten, har i dybere Forstand sin Grund i, at den har
udelukket Selvmagten; den kan ikke fordøie det Tilfældige, fordi den
ikke kan tilegne sig det i Naturbegreberne Vilkaarlige; istedenfor
Naturens indenfra fastsatte Kategorier underskyder den sine egne
selvvalgte Begreber, og bliver falsk vilkaarlig: dens Idealisme er
lige saa huul som dens Realisme.

Vi skulle, siger Physikeren, gaae i Skole hos Naturen; det er ikke os,
der skal lære Naturen Noget, men det er Naturen, der skal lære os
Noget. Physikeren erkjender altsaa, at han i Naturtingene har med
bestemte, udtrykkelig foreskrevne Tanker og Begreber at gjøre. Og deri
har han Ret. De legemlige Egenskaber vise, at Naturbegreberne
% --- p. 244 ---
have deres virkelige Udtryk i et kategorisk Magtsprog. Guld er et
Metal, Jern er et Metal; men af det almene Begreb Metal er det umuligt
at udlede de særlige Metallers særlige Egenskaber. Ligesom ethvert
Metal har i udvortes Tilværen et Præg af Tilfældighed, saaledes har
det i Væsen og Begreb Præg af en almægtig Vilkaarlighed. Hvorfor er
Guldet ikke magnetisk lige saavel som Jernet? Ja hvorfor er der i
Jorden baade Guld og Jern? Paa slige Spørgsmaal vil ingen Physiker
svare. Realbegreber som Guld og Jern ere i deres kategorisk ideelle
Bestemthed Villiesbegreber. Der er i hver Tankebestemmelse et Moment
af Vilkaarlighed, ligesom der i hver Egenskab er et Moment af
Tilfældighed.
% ATTER EN SPÆRRING, DER BRYDES VED LINIESKIFTET, som paa S. 239:
% „r a t i o n e l l e“ staaer spærret sidst i Linien, medens
% „Nødvendighed,“ øverst i næste Linie er almindelig Antiqua. „r e e l l e“
% og „N a t u r n ø d v e n d i g h e d“ ere spærrede i deres Heelhed.
% Alt kontrolleret ved Forstørrelse paa KB-Exemplaret; gjengivet som trykt.
Vilkaarlighedsmomentet forholder sig til den \emph{rationelle}
Nødvendighed, som Tilfældighedsmomentet til den \emph{reelle};
Eenheden af begge er \emph{Naturnødvendighed}, en Eenhed, der
articuleres i Grundmediet.

% Indløbende Overskrift, læst paa Siden selv (KB-Exemplaret): halvfed Antiqua.
\runin{c) Legemlig Individualitet.} At ethvert Legemes
Eiendommelighed maa lægge sig for Dagen i dets Egenskaber, forstaaes
af sig selv; den, der tilfulde kjender alle et Legemes saavel udvortes
(physiske), som indvortes (chemiske) Egenskaber, maa med Sandhed siges
at kjende Legemets Væsen. Men fra et menneskeligt Standpunkt vil
Erkjendelsen, hvad enten man forøvrigt holder sig til den reelle eller
den ideelle Hovedside, altid beroe paa en Tilnærmelse.

Erfaringslæren er realistisk; dens Adskillelse mellem haarde og bløde,
seige og skjøre, spændige og uspændige Legemer grunder sig paa
Forestillingen om de usynlige Deles indbyrdes Forhold:
Tilstandsformerne ere af væsenlig Betydning for Legemernes
Individualitet. Og dog kan Tilstandsformen forandre sig uden at
Legemet derfor taber sin substantielle Eiendommelighed. Iis, Vand og
Damp kunne vel siges at være forskjellige Ting, fordi der med hver
Aggregationsmaade følger en Forandring af Egenskaber; men Sub\-%
% --- p. 245 ---
stansen, den chemiske Charakteer, er for Vanddampen den samme som for
Isen og Vandet. De physiske Egenskaber kunne forandre sig, og de
chemiske forblive uforandrede. Individualitetsforandringer kunne
foregaae uden Forandring af Substans; men en Substansforandring kan
ikke foregaae uden at forandre Individualiteten.

Hvorpaa beroer nu egenlig en Substansforandring? Dette Spørgsmaal
faaer en anden Besvarelse i Physiken end i Philosophien.

Det ligger i Sagens Natur, at Erfaringslæren med sine exacte Analyser
maa lægge særlig Vægt paa den materielle Discretion. „I vore Tider“,
siger H. C. Ørsted%
\footnote{% trykt Mærke: *) --- staaer FORAN Kommaet: „Ørsted*),“
Naturl. mech. Deel. 1859. S. 14.}, „have Undersøgelserne over de
Love, efter hvilke ueensartede Stoffer indgaae de sandselig fuldkomne
Foreninger, som vi kalde chemiske, ledet til den nødvendige
Forudsætning, at Stofferne under disse Foreninger ikke fuldstændig
gjennemtrænge hinanden, som Kant og hans Meningsfølgere antage, men at
Foreningshandlingen bliver staaende ved visse Legemers Grunddele, som
med stor Inderlighed hænge sammen uden at gjennemtrænge hinanden.
Herved føres vi da til en Lære om hemmelige Grunddele i Legemerne, som
har saa megen Lighed med den gamle Atomlære, at Anledning herved blev
tagen til at give de her omhandlede chemiske Grunddele Navn af
Atomer.“

% Spærringen af „isomere“, „empiriske“ og „rationelle“ er bekræftet paa
% KB-Exemplaret; det mellemstaaende „og“ er almindelig Antiqua.
Ved Dannelsen af nye Substanser kommer det naturligviis an paa
Bestanddelenes Beskaffenhed, Mængde og Leiring. Et nyt Legeme opstaaer
ikke alene derved, at tilsvarende Grunddele forene sig i et vist
Mængdeforhold; men Legemer med ny Individualitet kunne ogsaa
fremkomme ved samme Bestanddele i samme Mængde. Ved disse saakaldte
\emph{isomere} Legemer maa der da, hvad Sammensætningsmaaden angaaer,
gjøres Forskjel mellem \emph{empiriske} og \emph{rationelle} Formler.
% --- p. 246 ---
% CHEMISKE FORMLER. Grundstofbogstaverne ere trykte som opreist Antiqua med
% sænkede Tal; de sættes her som almindelig Text med Tallene i Mathematik-Mode,
% saa Bogstaverne beholde den opreiste Form, Trykken har. Alle Talrækker ere
% læste ved Forstørrelse paa KB-Exemplaret --- OCR'en er intet Vidne her.
% NB. Trykken skriver „12 Vægtdele Brint“ mellem „6 Atomer Kulstof“ og
% „4 Atomer Ilt“; begge Exemplarer have det saaledes, og Trykfeilbladet tier.
% Spærringen af „metamere“ og „polymere“ er bekræftet paa KB-Exemplaret.
Et givet Product være f. Ex. C$_6$ H$_{12}$ O$_4$, altsaa ifølge den
empiriske Formel bestaaende af 6 Atomer Kulstof, 12 Vægtdele Brint, 4
Atomer Ilt; hvad er det nu for en Substans? Her er i dette Stof en
Mulighed til to Substanser, en Dobbelthed, som finder sin Forklaring
deri, at den empiriske Formel lader sig adskille i to rationelle
Formler: C$_4$ H$_6$ O$_3$ + C$_2$ H$_6$ O, Substansen er da
eddikesuurt Metylilte, og C$_2$ H$_2$ O$_3$ + C$_4$ H$_{10}$ O, men
saa er Substansen myresuurt Æthylilte. Den substantielle Forskjel
beroer paa en Delenes Omordning, hvorfor slige isomere Dannelser da
ogsaa særlig have faaet Navn af \emph{metamere}. Paa tilsvarende Maade
udledes Substansforandringen i visse Tilfælde deraf, at
Forbindelserne vel indeholde de samme Stoffer i samme relative, men
ikke i samme absolute Mængde, saa at der altsaa fremkommer
\emph{polymere} Dannelser. Saaledes giver C$_8$ H$_{16}$ O$_4$ samme
procentiske Sammensætning, som C$_4$ H$_8$ O$_2$, men da Atomernes
Mængde i hvert af Elementerne er forskjellig, kan Forskjellen mellem
de to Producter være substantiel.

Atomernes forskjellige Leiringsmaade kan naturligviis lige saa lidt
være Gjenstand for umiddelbar Iagttagelse, som deres Størrelse, Figur
eller Vægt; men hvilken substantiel Betydning der maa tillægges den
Maade, hvorpaa Bestanddelene tænkes grupperede, kan alene den
Indflydelse, som Krystalformen udøver paa Bestemmelsen af de physiske
Egenskaber, lære os. Krystalformerne bestemmes fornemmelig ved
Axerne, den centrale Linie, hvorom den hele Dannelse regelmæssigt er
anlagt. At nu Axen virkelig afgiver det væsenlige Støttepunkt for
Leiringen, der altsaa maa være forskjellig, eftersom den skeer
parallelt med Axen eller lodret paa den, bekræfter sig fra mange
Sider. De bestemte Retninger, hvorefter Krystallen spaltes, staae i
Forhold til Axen; det Samme gjælder om Lysets Gjennemgang, som, da den
beroer paa Ætherdelenes Svingninger, netop er afhængig af
% --- p. 247 ---
Atomernes Leiring. Parallelt med Axen er f. Ex. Turmalinen næsten
uigjennemsigtig, hvorimod den, naar den sees i en Retning lodret paa
Axen, er gjennemsigtig; Dichroiten, der er blaa i Retning af
Hovedaxen, er bruunguul lodret paa samme. Noget Lignende gjælder om
Varmen. Krystaller af Terningsystemet (det sphæroedriske System), der
har tre Slags Hovedaxer, lede Varmen lige godt i alle Retninger;
hvorimod Kalkspathen og Qvartsen, der høre til Rhomboedersystemet, som
kun har een Axe, vise en Ledningsevne, der er størst parallelt med
Hovedaxen, mindst lodret paa denne. --- Af alt dette sees, at
Erfaringslæren ifølge sin Charakteer i det Hele maa gaae i atomistisk
Retning.

I Modsætning til Physiken, der afleder de sammensatte Legemers
substantielle Individualitet af uforandrede Grunddeles eiendommelige
Gruppering, gaaer den speculative Immanensphilosophie til den
idealistiske Yderlighed, at den gjør Stofsubstanserne til umiddelbare
Selvvæsener og tillægger de elementære Legemer individuel
Subjectivitet. Saaledes Schelling i intuitiv Mystik; saaledes Hegel i
overspændt Logik.

% ANFØRSELSTEGN OVER TO BLADSKIFTER: det store Hegel-Citat aabnes her paa
% S. 247 og lukkes først paa S. 249 („i sig selv.“). S. 247 er derfor +1 og
% S. 249 -1 i Anførselsbalancen; S. 248 har hverken Aabning eller Lukning.
% Det er en Anførsel over Bladskiftet, ikke en Trykfeil.
% Spærringen af „Vægtfylde“, „Cohæsion“ og „Klangen“ er bekræftet paa
% KB-Exemplaret. „Cohæsion“ er spærret i BEGGE Halvdele over Linieskiftet
% („C o h æ-“ / „s i o n.“), i Modsætning til S. 239 og S. 244.
„De legemlige Egenskaber, Prædicater af den sig individualiserende
Materie, ere reale Momenter af en iboende, mod Subjectiviteten
henstræbende Formvirksomhed. Den af Subjectivitetsprincipet
udspringende Formbevægelse gjennemløber to Stadier: fra Tyngden til
Klangen, og fra Varmen til det elektriske Lys. Tyngden qvalificeres
til \emph{Vægtfylde}; dens articulerende Momenter: Punctualitet,
Linearitet og Fladeagtighed, blive særlig satte i den qvalificerede
\emph{Cohæsion}. Cohæsionens Idealitet bestaaer deri, at de
materielle, udenfor hinanden værende Dele sættes i hinanden; den
uophørlige Vexlen af Delenes Væren i Eet og Væren udenfor hinanden er
den Svingen og Zittren, der lader sig høre i \emph{Klangen}. Klangen
for sig er Individualitetens Selv. Tonen berører vort Inderste, fordi
den selv er et Indre, et
% --- p. 248 ---
% Spærringen af „Varmen“, „Magnetismen“, „Krystal“, „Farver“, „Lugt“ og
% „Smag“ er bekræftet paa KB-Exemplaret.
% „Homogeneitet“ og „Formløshed“ ere læste paa KB-Exemplaret; tesseract
% giver her „Homogeueitet“ og „Formløeshed“, som begge ere Feillæsninger.
Subjectivt. Til Individualitet hører der Materie og Form; Klangen er
den totale Form, der kundgjør sig i Tiden --- den hele Individualitet,
der ikke er Andet, end at denne Sjæl nu er sat i Eet med det
Materielle. Det næste er, at den foreløbige, endnu ufuldkomne
Individualitet opløses i \emph{Varmen}. Ved Varmen restitueres
Materien til dens oprindelige Formløshed, dens Flydenhed; det er dens
abstracte Homogeneitet, der feirer en Triumph over de første
specifiske Bestemtheder. Fremskridtet bestaaer i, at Materiens kun i
Væsen værende Continuitet udtrykkelig bliver sat som Negationens
Negation, som Activitet, som tilværende Opløsen. Ud fra denne Opløsen
begynder i dybere Forstand den individualiserende Formvirksomhed
paany. Den umiddelbare, i sig formløse Skikkelse er i sit ene Extrem
Punctualiteten, i sit andet Extrem den kugleagtige Flydenhed, ---
Skikkelsen som indre Skikkelsesløshed. Punktet maa nu gaae over i
Linien, og Formen sætte sig paa Linien i polære Extremer med en
indifferent Midte. Denne Slutning udgjør Formgivningens
skikkelsedannende Princip --- \emph{Magnetismen}. Idet den magnetiske
Virksomhed afsætter sit Product, er Skikkelsen bestemt som
\emph{Krystal}. De individuelle Legemer ere nu udviklede til physiske
Totaliteter, selviske Subjecter, hvis særlige Prædicater vise sig som
physiske Egenskaber i Forhold til de almene Elementer, med hvilke de
indgaae i reale Processer. I Forhold til Lyset, den med sig identiske
Selvhed, udtale de legemlige Individualiteter deres eiendommelige
Selvværen i Uigjennemsigtighed, Gjennemsigtighed og de mange
Fordunklinger, der give dem \emph{Farver}; i Forhold til Luften er
Individualiteten \emph{Lugt}, i Forhold til Vandet \emph{Smag}.
Formudviklingen af de endnu umiddelbart selvstændige Individualiteter
ender med en virkelig Conflict, en physisk Spænding af deres
Særskilthed, deres reelle Selvviskhed, der udlignes i det elektriske
Lys. Det begyndte med, at Legemerne havde Lyset udenfor sig; det ender
med,
% --- p. 249 ---
% To Noter paa dette Blad, trykte *) og **); den anden sættes med et lokalt
% \renewcommand, som strax efter stilles tilbage.
% Mærket for den første Note staaer EFTER det lukkende Anførselstegn:
% „i sig selv.“*)“.
at hver legemlig Individualitet har sit Lys, sit eget subjective Væsen
i sig selv.“%
\footnote{% trykt Mærke: *) --- første af to Noter paa S. 249.
Hegels Vorlesungen über die Naturphil. S. 129--360.}
--- Det Vilkaarlige i en saadan Deduction er iøinefaldende.

Først ved Læren om Grundmediet kan Naturphilosophien bringes til
væsenlig sand Overeensstemmelse med Erfaringslæren uden at opgive sin
Selvstændighed som særegen Videnskab. I Grundmediet have vi Principet
for den fulde Anerkjendelse af Erfaringens Kjendsgjerninger; thi
Gjenstandsmagten er forudsat, og Begrebernes Udslag i Magt skeer i
Masser og Grunddele, saa at disse naturlig maae spille Hovedrollen ved
Stofsubstansernes, de legemlige Individualiteters Sammensætning og
Opløsning. I Grundmediet have vi Principet for en virkelig
Articulation af Legemsbegreberne; thi Selvmagten er forudsat: den
articulerende Tænkning gaaer i Eet med en articulerende Villie.
Magtbegrebernes System i sin uendelige, eiendommelige Sammenhæng er et
Villiessystem, hvori Begrebernes Overgange punktlig afgive de
Forvandlinger, der i Virkeligheden begrunde de legemlige
Forandringer.%
\renewcommand{\thefootnote}{**)}%
\footnote{% trykt Mærke: **) --- staaer EFTER Punktummet: „Forandringer.**)“
Om „Forandringsproblemet i Philosophien“ jvnf. Forelæsn. over
„Phil. Propædeut.“ 1860--61. S. 143--177.}%
\renewcommand{\thefootnote}{*)}

Hvad Elementer og Legemer saavel i ideel som i reel Forstand ere for
hverandre indbyrdes, maa erkjendes af deres indbyrdes Vexelvirkninger,
deres lovbestemte Virksomheder; de lovbestemte Virksomheder ere
Functioner.

% KAPITLETS SLUTNING. Texten paa S. 249 hører op omtrent tre Femtedele nede
% paa Bladet; derunder staaer blankt Rum og saa en kort centreret Streg,
% TRYKT DOBBELT (to parallele Linier tæt over hinanden), omtrent to
% Trediedele nede --- som ved S. 61 og S. 125. Efter Husreglen sættes den
% som en enkelt \streg; Forskjellen i Ornamentet gjengives ikke.
% Under Stregen følger kun Notefeltet med de to Noter. Intet af Andet
% Kapitel staaer paa dette Blad: Kapiteloverskriften begynder S. 250.
\streg


% ======================================================================
%  Anden Deel — Andet Kapitel: Elementære Functioner   (printed pp. 250–317)
% ======================================================================

\kapitel{Andet Kapitel.}{Elementære Functioner.}

%  --- A. Dynamiske Functioner: Lys og Varme  (pp. 250–273) ---

%  MATHEMATICAL — work from the image; the OCR is not a witness here.
% --- p. 250 ---
% \kapitel staaer allerede i Skelettet ovenfor og gjentages ikke her.
% KONTROL AF KAPITELOVERSKRIFTEN, læst paa KB-Exemplaret (Farvescanningen):
% Hovedet „Andet Kapitel. / Elementære Functioner.“ AABNER S. 250 --- Siden
% begynder med Folioen 250, derunder Overskriften, og derunder strax
% Kapitlets første Afsnit. Der staaer INGEN Streg ovenfor Hovedet (til
% Forskjel fra S. 1 og S. 197, hvor en Streg skiller \deel fra \kapitel:
% her gaaer ingen Deel-Overskrift forud). Intet af S. 249 løber over paa
% S. 250. Anden Linie er sat halvfed.
% Afdelingsoverskriften A. staaer IKKE øverst paa Bladet, men omtrent tre
% Fjerdedele nede: først løber Kapitlets Indledningsafsnit til Ende, derpaa
% følger en kort centreret Streg, derpaa „A.“ paa egen Linie og derunder
% „Dynamiske Functioner: Lys og Varme.“ (halvfed, spærret).

% Spærringen paa dette Blad er efterprøvet paa KB-Exemplaret Ord for Ord:
% „Almene“, „Særskilte“, „det Individuelle“ (her er ogsaa „det“ spærret),
% „dynamiske“, „Lys og Varme“, „leddelende“, „Magnetisme og Elektricitet“,
% „legemdannende“, „Chemisme og chemiske Processer“. „Functioner“ er hver
% Gang almindelig Antiqva og staaer derfor udenfor \emph.
I den elementære Legemverden ere de existerende Legemer at opfatte:
deels i reflecteret Tilværen som Momenter for det legemlige Hele, altsaa
undergivne det \emph{Almene}; deels som Dele, bestaaende Led, der under
mangesidige Conflicter holde sig mod hverandre i det \emph{Særskilte};
deels endelig som substantielle Producter, eiendommelige Frem\-%
bringelser af den i \emph{det Individuelle} virksomme Natur\-%
grund. Forsaavidt Legemerne efter deres særegne Beskaffenhed staae i
Forhold til en qvalitativt bestemt, almeenlegemlig Væsenhed, vise de
heraf flydende Virksomheder sig som \emph{dynamiske} Functioner,
svarende til \emph{Lys og Varme}. Forsaavidt det i Legemerne Selvegne
hævder sig under ind\-%
byrdes spændende Modsætninger, vil den igjennem det dynamiske
Almeenvæsen gaaende Vexelvirkning komme tilsyne i \emph{leddelende}
(articulerende) Functioner, svarende til \emph{Magnetisme og
Elektricitet}. Forsaavidt endelig den hele Formvirksomhed er saa
gjennemgaaende, at Leddene blive Momenter og Momenterne Led af
individuelle Dannelser, have vi de \emph{legemdannende} Functioner,
svarende til \emph{Chemisme og chemiske Processer}.

% Kort, centreret Streg mellem Kapitlets Indledningsafsnit og Afdeling A,
% midt paa S. 250. Efter Husreglen gjengives Ornamentets Længde ikke.
\streg

% Overskriften er læst paa Siden selv (KB-Exemplaret), ikke i Indholdet:
% „A.“ paa egen Linie, derunder „Dynamiske Functioner: Lys og Varme.“
% med Kolon, sat med fed, spærret Skrift.
\afdeling{A.}{Dynamiske Functioner: Lys og Varme.}

Den sædvanlige Forestilling om det Dynamiske, at Stoffer omdannes til
Kræfter, gjennemtrænge hinanden o. desl., er uklar og uden Betydning i
Videnskaben. Det Dynamiske har det Mechaniske til Grundlag og bestaaer i,
% --- p. 251 ---
at de virksomme Deles endelige Discretion i den Grad nærmer sig det
Uendelige, at den gjør Indtryk af Continuitet.%
\footnote{% trykt Mærke: *) --- staaer efter Punktummet
Jvnfr. Grundideernes Logik I. S. 354--69.}
% tesseract sætter en Tankestreg foran „Naar“; det er Randstøi paa
% Arbeidsscanningen (se Hovedet), ikke Sats. Google-Laget har den ikke.
Naar Lys og Varme henføres til Svingninger i Ætheren, da er
Virksomheden mechanisk; men herved er dog det Særegne, at Ætherdelene
og deres Bevægelser smelte sammen i Forestillingen saaledes, at de
enkelte Forholdsled ikke fast\-%
holdes for sig, men gaae op som Momenter i Virksomheden. Den nærmere
% „Qvalitative“: Google-Laget læser „Qualitative“ (den kjendte
% qv/qu-Forvexling); Typen har v.
Forskjel bestemmes ved det Qvalitative. For Gravitationens abstract
mechaniske Tyngdeforhold er al Materie eensartet; her er kun de tunge
Masser og det tomme Rum. Det Dynamiske i Lys og Varme betinges derimod
% „væsenlig“: tesseract læser „vægenlig“; Typen har s (efterseet paa
% Arbeidsscanningen og KB-Exemplaret).
af en qvalitativ Forskjel i Materien selv, en væsenlig Modsæt\-%
ning mellem ponderabel og imponderabel Materie. At kalde Ætheren en
ideel Materie er imidlertid overfladisk, thi Ætheren fylder sit Rum
lige saa vel som de veielige Masser. Naar det hedder, at Ætheren
„gjennemtrænger de gjennem\-%
sigtige Legemer“, da er det Dynamiske i denne Gjennemtrængen en Reflex
af det Mechaniske. Om ogsaa ethvert Legemsatom er „omgivet af en
Æthersphære“, saa er det dog utænkeligt, at Ætheratomet skulde
gjennemtrænge Legemets Atom og punktlig indtage det selvsamme Rum som
dette. At Ætheren er uden Vægt, betyder kun, at den paa Grund af sin
yderst ringe Tæthed er uveielig; men det beviser ingenlunde, at den i
sig selv er uden Tyngde. An\-%
tages dens Tæthed at være 845000 Millioner Gange mindre end Luftens,
% Brøken staaer i Satsen som en skraa Brøk: 1 og derefter 1/3.
% Fodnotemærket **) staaer FORAN Kommaet, som trykt.
saa at en Cubikmiil Æther ikke vilde indeholde mere Masse end
$1\tfrac{1}{3}$ Pund%
\renewcommand{\thefootnote}{**)}%
\footnote{% trykt Mærke: **)
C. Holten. Lysets Naturlære. Kbhvn. 1861. S. 167.}%
\renewcommand{\thefootnote}{*)}%
, saa er det heraf forstaaeligt, at Ætheren meget vel kan være
uveielig, uden at være Tyngdeloven unddragen. Det ideelle Moment i
Ætheren er, at den umiddelbart forestiller den almene Materie, et
hypo\-%
thetisk Grundlag for gjennemgaaende dynamiske Virksom\-%
% --- p. 252 ---
heder. Men Modsætningen mellem det Ideelle og det Reelle ligger ikke i
Materien selv, den beroer paa Functionernes Eiendommelighed. Den
% Spærringen efterprøvet paa KB-Exemplaret: „Lyset“, „Varmen“ og
% „Forhold mellem Lys og Varme“ ere spærrede; Kommaet og Semikolonet
% efter de to første ere det ogsaa, men medtages ikke i \emph.
dynamiske Idealitet er i \emph{Lyset}, hvad den tilsvarende Realitet er
i \emph{Varmen}; Vexelbestem\-%
melsen af det Ideelle og det Reelle articuleres i \emph{Forhold mellem
Lys og Varme}.

% Indløbende Overskrift, læst paa Siden selv (KB-Exemplaret): sat med
% halvfed Skrift, ikke spærret. Denne Udgave gjengiver ikke Forskjellen
% mellem halvfed og spærret; se Hovedet.
\runin{a) Lyset.} For at erkjende, hvad Lyset er, maae vi begynde med
at undersøge det Virksomme, hvorpaa dets Til\-%
væren beroer. Her maa da forudsættes to ueensartede Virksomheder, en i
Andethedens og en i Selvhedens Form, en synliggjørende og en seende, en
materiel og en ideel Virksomhed: Lyset selv er de to ueensartede
Virksomheders reflecterede Eenhed. Dette kan fra Naturlærens Standpunkt
udtales i følgende to Sætninger: 1) Lyset beroer paa en Zittren i
Ætheren; naar Ætheren er i Ro, er der intet Lys. 2) Lyset beroer paa en
Synsvirksomhed; hvor der ingen Seen er, er der intet Lys. --- Lysets
Væsen er altsaa paa een Gang baade objectivt og subjectivt. Men uagtet
de to Virk\-%
somheder umulig kunne skilles fra hinanden, uden at Lyset forsvinder,
er der dog Intet til Hinder for, at jo enhver af dem kan betragtes for
sig; kun at man ved Betragtningen af den ene bestandig maa være sig
bevidst, at man under\-%
forstaaer den anden. Vi afhandle først det objective Væsen og dernæst
det subjective, for at naae til en udviklet Reflexion af Lysets
reel-ideelle Natur.

% Det græske Bogstav er læst paa KB-Exemplaret: α, ikke „a“ eller „&“,
% som OCR'en giver. Overskriften er indløbende og spærret.
\subrunin{α) Lysets objective Væsen.} Seet fra Natursiden maa Lyset
nødvendig have et til sin Virksomhed svarende materielt Grundlag. Ved
at opfatte Aarsagen, forsaavidt den reflecterer sig i sine Virkninger,
søger Naturlæren at tydeliggjøre den af Grundlaget betingede
Virksomhed fra alle Sider. Det, som ikke sees, reflecterer sig i det,
som sees: Lysets objective Væsen kan alene bestemmes ved dets
Objectivering. Det Eiendommelige ved Svingningshypothesen
% --- p. 253 ---
er, at det underforstaaede Sandselige forsvinder for Iagt\-%
tagelsen, thi det, som iagttages, er ikke de svingende Ætherdele, det
er Phænomener, der som Virkninger henføres til Svingninger som til
usynlige Aarsager; det Flygtige i Forestillingen forvandler sig til
Reflex af Bevægelsen: Begreber og Love indgaae i Reflexioner af
% S. 253 bærer to Noter, *) og **). Note **) løber videre paa S. 254,
% hvor dens Slutning staaer ØVERST i Fodnoteblokken, foran S. 254's egen
% Note *). Bogen gjentager ikke Mærket paa Fortsættelsen. Efter Husreglen
% staaer hele Noten her ved sit Mærke.
Aarsagen.%
\footnote{% trykt Mærke: *) --- staaer efter Punktummet
Det ligger nær at sammenligne Vandbølger, Luftbølger og Ætherbølger med
hverandre; men de ere dog meget forskjellige. Cauchy har --- i sit
Skrift: Memoire sur la dispersion de la lumière --- med Atomistiken til
Grundlag behandlet Ætherbølgerne og udledet rigtige Formler for
Farveadspredelsen. Heraf har Fechner sluttet: ergo er Atomistiken
% TRYKFEIL, som Trykfeilbladet IKKE retter: „Fcrmler“ for „Formler“ ---
% en gal Type (c for o). Begge Exemplarer og begge OCR-Vidner vise det,
% altsaa Sætterens Feil. Gjengivet som trykt; læs „Formler“. Til
% Sammenligning staaer det rigtige „Formler“ tre Linier ovenfor.
beviist. Men de samme Fcrmler kunne udledes med Continuitetslæren til
Grundlag; man kunde da med samme (U-)Ret slutte: ergo er Continuiteten
beviist. Det, som under begge Forudsætninger kan godtgjøres, er, at det
Mechaniske her gaaer over i det Dynamiske.
% „Memoire sur la dispersion de la lumière“ staaer i almindelig
% Antiqva, IKKE i Kursiv (efterprøvet paa KB-Exemplaret); derfor intet
% \textit. Accenten paa „lumière“ er gravis.
}
Svingningens Aarsag skiller sig oprindelig i to Aarsager, en bevægende
og en levende Kraft, der uophørlig omsættes i hinanden, uophørlig
fortære og frembringe hinanden. Denne Virken er en gjennemgaaende
Reflecteren; at Aarsagen virker dynamisk vil sige: den reflecteres i
sine mechaniske Momenter.%
\renewcommand{\thefootnote}{**)}%
\footnote{% trykt Mærke: **) --- staaer efter Punktummet
% Noten begynder paa S. 253 og slutter øverst i Fodnoteblokken paa S. 254
% („Naar y = o ... forsvundet.“); den er samlet her ved Mærket.
% ± er et Plus over en Streg, uden Prikker, læst paa KB-Exemplaret;
% Minus er derimod en almindelig lang Tankestreg-Minus. Nielsens Tegn
% harmoniseres ikke. φ er det krøllede phi. „v = o“ og „y = o“ ere sat
% med Bogstavet o, ikke med Nullet; saaledes gjengivet.
% De trykte Variable staae opret antiqua; efter Husreglen gjengives det
% ikke, men den almindelige mathematiske Kursiv beholdes.
At en Ætherdeel er i Svingning, vil sige, at den bevæger sig frem og
tilbage lodret over en given ret Linie; i Linien selv har den sin
Hvilestilling. Betegnes Afstanden fra Hvilestillingen ved $\pm y$
svarende til Tiden t, saa kan Maximum for y betegnes ved $+a$ til den
ene og ved $-a$ til den anden Side af Linien. Kaldes Ætherdelens Masse
$m$, $m\varphi$ den bevægende, $mv^{2}$ den levende Kraft, saa er
Forholdet mellem $\varphi$ og $v^{2}$ følgende:
$\varphi = -ky = \dfrac{d^{2}y}{dt^{2}}$;
$\dfrac{dy^{2}}{dt^{2}} = v^{2} = k(a^{2}-y^{2})$.
Heraf sees, at naar $y = a$, saa er $\varphi$ i sit Maximum, men
$v = o$, d. v. s. naar den bevægende Kraft er størst, saa er den
levende Kraft forsvundet. Naar $y = o$, altsaa Ætherdelen i
Hvilestillingen, saa er $v^{2} = ka^{2}$, men $\varphi = o$, d. v. s.
naar den levende Kraft er størst, er den bevægende Kraft forsvundet.}%
\renewcommand{\thefootnote}{*)}

% --- p. 254 ---
Dersom Modsætningen af Lys og Mørke blot var umiddelbart qvalitativ:
enten Lys eller Mørke, saa vilde Svingningsbegrebet i Almindelighedens
Form være tilfreds\-%
stillende. Men allerede den Kjendsgjerning, at Belysningen svækkes med
Udbredelsen, forudsætter en Lov og fordrer en Begrundelse.
Begrundelsen følger naturlig af Begrebets Ud\-%
vikling, thi Svingningsbegrebet forudsætter, at Bølgerne ud\-%
brede sig som Kugleflader omkring det lysende Punkt, og at Ætheren er
fuldkomment spændig. I to concentriske Kugleflader ville Summerne af
de levende Kræfter altsaa være lige store. Da nu Kuglefladernes Masse
voxer med Afstanden, og Kuglefladerne forholde sig ligefrem som de
qvadrerede Radier, maae de levende Kræfter, altsaa Lys\-%
% S. 254 bærer to Noter, *) og **), foruden Slutningen af S. 253's Note
% **), der staaer øverst i Fodnoteblokken. Note **) løber videre paa
% S. 255 og er samlet her ved sit Mærke.
styrken, forholde sig omvendt som Radiernes d. e. Afstandenes
Qvadrater.%
\footnote{% trykt Mærke: *) --- staaer efter Punktummet
% Hele Opregningen m, n, r, v og deres mærkede Led sættes i Mathematik,
% saa at Rækken ikke skifter Skriftsnit undervejs (Husreglen fra Batch 5).
% I Formlen staaer $r_{1}$ med sænket Ettal, medens den løbende Text
% skriver $r'$ med Apostrof; begge Exemplarer vise det, og det samme
% Forhold gjentager sig S. 256 ($x'$ i Texten, $x_{1}$ i Formlen).
% Gjengivet som trykt.
For en Kugleflade med Massen $nm$ være Radien $r$, og
Svingningshastigheden $v$, for en Kugleflade med Massen $n'm$ Radien
$r'$ og Svingningshastigheden $v'$; man vil da have
$nmv^{2} = n'mv'^{2}$. Da nu $\dfrac{n'}{n} = \dfrac{r_{1}^{2}}{r^{2}}$,
saa følger heraf $\dfrac{mv^{2}}{mv'^{2}} = \dfrac{r_{1}^{2}}{r^{2}}$.}
Men nu Modsætningen mellem Lysstyrke og Farve? Ogsaa for Begrundelsen
af denne Modsætning indeholder Svingningsbegrebet de nødvendige
Betingelser, nemlig saadanne særlige Momenter som Udsving, Fjernelse
fra Hvilestillingen i den løbende Tid og Svingningstiden.%
\renewcommand{\thefootnote}{**)}%
\footnote{% trykt Mærke: **) --- staaer efter Punktummet
% Noten begynder paa S. 254 og løber videre øverst i Fodnoteblokken paa
% S. 255 („svarer y = o ... hvori a er Udsvinget.“); samlet her.
% „arc (sin = y/a)“ er Nielsens egen Skrivemaade for Arcus sinus og
% beholdes. Alle „o“ i Ligningerne ere Bogstavet o, ikke Nullet.
Af Ligningen $\dfrac{dy^{2}}{dt^{2}} = k(a^{2}-y^{2})$ erholdes
$\sqrt{k}\,dt = \dfrac{dy}{\sqrt{a^{2}-y^{2}}}$ og heraf igjen ved
Integration
$\sqrt{k}\,t + C = \operatorname{arc}\left(\sin = \dfrac{y}{a}\right)$,
altsaa $y = a\sin(\sqrt{k}\,t + C)$. Til $t = o$ svarer $y = o$,
% Satsen har her ikke noget Lighedstegn foran det sidste o: „og C --- i
% første Qvadrant --- o.“ staaer saaledes paa begge Exemplarer.
% Gjengivet som trykt.
følgelig er $o = a\sin C$, og C --- i første Qvadrant --- o. Man har
altsaa
$\sqrt{k}\,t = \operatorname{arc}\left(\sin = \dfrac{y}{a}\right)$.
Betegnes Svingningstiden ved T, saa svarer $t = \tfrac{1}{4}T$ til
$y = +a$, altsaa er
$\sqrt{k}\,\tfrac{1}{4}T = \operatorname{arc}(\sin = 1) = \dfrac{\pi}{2}$,
$\sqrt{k} = \dfrac{2\pi}{T}$, og Svingningsformlen
$y = a\sin\dfrac{2\pi t}{T}$, hvori a er Udsvinget.}%
\renewcommand{\thefootnote}{*)}
% --- p. 255 ---
Lysstyrken beroer paa Udsvinget, Farven paa Svingningstiden. Heraf
indsees, at Udsvinget kan være forskjelligt og Svingningstiden den
samme, d. v. s. Lysstyrken kan være for\-%
skjellig, og Farven den samme; omvendt kan Udsvinget være det samme og
Svingningstiden forskjellig, d. v. s. Lysstyrken kan være den samme og
Farven forskjellig.%
\footnote{% trykt Mærke: *) --- staaer efter Punktummet
% Det græske Bogstav er læst paa KB-Exemplaret: β, ikke „B“, „ß“ eller
% „3“, som OCR'en giver. ± er et Plus over en Streg, uden Prikker.
% Punktummet mellem 2πa/T og cos er Nielsens Gangetegn og beholdes.
% Citatet aabner „Da det ...“, lukkes efter „lige stærke“, aabnes paany
% efter Indskuddet „(eensfarvede)“ og lukkes efter „af dem.“ ---
% Bogens sædvanlige Vane; Anførselstegnene gaae op.
Betegnes Lysbølgernes Hastighed ved V, Svingningstiden ved T,
Bølgebreden ved $\beta$ og Svingningstallet ved N, saa faaer man
$T = \dfrac{\beta}{V}$ og $N = \dfrac{V}{\beta}$, hvoraf sees, at
Svingningstiden er i ligefremt og Svingningstallet i omvendt Forhold
til Bølgebreden. De røde Straaler have den største, de violette den
mindste Svingningstid; naar 500 Billioner Svingninger i Secundet give
Rødt, saa give 733 Billioner Violet. Lysstyrken derimod har til Maal
Qvadratet af den Hastighed, hvormed Ætherdelen gaaer igjennem
Hvilestillingen. Af Formlen $y = a\sin\dfrac{2\pi t}{T}$ følger nemlig
$\dfrac{dy}{dt} = \dfrac{2\pi a}{T}\,.\,\cos\dfrac{2\pi t}{T}
 = \dfrac{2\pi}{T}\sqrt{a^{2}-y^{2}} = v$. I Hvilestillingen er
$y = o$, Hastigheden $\dfrac{dy}{dt} = v = \pm\dfrac{2\pi a}{T}$,
følgelig $v^{2} = \dfrac{4\pi^{2}a^{2}}{T^{2}}$. „Da det imidlertid
ikke kommer an paa en absolut Bestemmelse, kan man lige saa godt bruge
$a^{2}$ som Maal for Lysstyrken, saa at to lige stærke“ (eensfarvede)
„Lysstraaler, som træffe sammen under saa gunstige Forhold som muligt,
ville give en Belysning, 4 Gange saa stærk, som een af dem.“
C. Holten. Lysets Naturlære. S. 173.}
Endelig har
% --- p. 256 ---
Begrebet i sig Mulighed til at optage alle de særlige Betingelser af
hvilke Lovene for Ætherbølgernes fremad\-%
% Fodnotemærket staaer her uden foregaaende Tegnsætning, midt i Sætningen.
skridende Bevægelser%
\footnote{% trykt Mærke: *)
% TRYKFEIL, rettet af Bogens eget Trykfeilblad:
%   S. 256, Linie 12 f. n.: „vil --- læs ville“.
% Efter Husreglen gjengives Texten som TRYKT: „saa vil --- under“.
% VIGTIGT: Arbeidsscanningen (Google-Exemplaret) er blyantsrettet af en
% tidligere Eier, og Blyanten HAR naaet dette Sted --- der staaer et
% paaskrevet „le“ efter „vil“, som Google-Textlaget læser med og gjengiver
% som „ville“. KB-Exemplaret, som ikke er rettet, trykker klart „vil“.
% Læsningen er altsaa taget paa KB-Exemplaret.
% θ er lille Theta, læst paa KB-Exemplaret; begge OCR-Vidner læse „0“.
% β er lille Beta. I Texten staaer m', x' med Apostrof, i Formlerne
% $a_{1}$, $x_{1}$, $t_{1}$ med sænket Ettal --- som trykt, se S. 254.
% „y = 0“ er derimod her et NUL (ikke Bogstavet o), efterprøvet paa
% KB-Exemplaret mod „o“ i „og“ paa samme Linie.
Forplanter en Bølge sig fra Punktet $m$, hvori Svingningen betegnes ved
$y = a\sin\dfrac{2\pi t}{T}$ til $m'$, er Veien $x$ og Tiden, som
derpaa anvendes, $\theta$, saa faaer man for Svingningerne af
Ætherdelen $m'$ Formlen: $y = a\sin\dfrac{2\pi(t-\theta)}{T}$.
Skulle to Bølger samtidig gjennemløbe den samme Linie, men fra
forskjellige Udgangspunkter, saa at den ene tilbagelægger Veien $x$,
den anden $x'$, inden de ankomme til $m'$, saa vil --- under
Forudsætning af, at Bølgebreden $\beta$ er for begge den samme, men
Udsvinget for den ene $a$, for den anden $a_{1}$ --- Svingningerne af
$m'$ udgjøre Summen af begges Svingninger og udtrykkes ved Formlen:
\[
y = a\sin 2\pi\left(\frac{t}{T}-\frac{x}{\beta}\right)
  + a_{1}\sin 2\pi\left(\frac{t}{T}-\frac{x_{1}}{\beta}\right).
\]
Betegnes $\dfrac{t}{T}-\dfrac{x}{\beta}$ ved $\dfrac{t_{1}}{T}$,
$x'-x$ ved $e$ og antages $a = a_{1}$, saa viser Beregningen, at
$y = 2a\cos\dfrac{\pi e}{\beta}
 \sin\left(\dfrac{2\pi t_{1}}{T}-\dfrac{\pi e}{\beta}\right)$;
heraf sees, at man for et ulige Antal halve Bølgebreder faaer
$y = 0$, saa at Lys med Lys giver Mørke; for et lige Antal derimod
$y = 2a\sin\dfrac{2\pi t_{1}}{T}$, saa at Lys med Lys giver forstærket
Lys.}
og gjensidige Indvirkninger paa hinanden lade sig udlede. Det kan
nøiagtig beregnes, hvorledes to Lysbølger, som samtidig forplante sig
langs hen ad samme Linie, men fra forskjellige Udgangspunkter, ville
træffe sammen, i hvilke Tilfælde de ville forstærke og i hvilke
Tilfælde de ville ophæve hinanden; det kan heraf forklares, at Lys med
Lys kan give, ikke blot forstærket Lys, men ogsaa Mørke.

Af disse Antydninger maa det vel kunne skjønnes, i hvilken Forstand
Svingningsbegrebet objectivt svarer til
% --- p. 257 ---
Lysbegrebet og bestemmer Lysets Naturaarsag paa en saadan Maade, at det
hele System af særskilte Aarsager til særskilte Phænomener af
% Fodnotemærket staaer uden foregaaende Tegnsætning, foran „o. s. v.“
Interferens, Reflexion, Refraction, Plansætning%
\footnote{% trykt Mærke: *)
% TRYKFEIL, rettet af Bogens eget Trykfeilblad:
%   S. 257, Linie 13 f. n.: „Reflection --- læs Reflexion“.
% (Linietællingen paa dette Blad regner Fodnotelinierne med og lader
% Arkmærket „17“ ude; Linie 13 f. n. er netop Notens første Linie.)
% Efter Husreglen gjengives Texten som TRYKT: „Reflection“.
% LÆSNINGEN ER TAGET PAA KB-EXEMPLARET og er POSITIVT afgjort. Dette er
% ellers netop den x/ct-Sondring, som TRANSCRIPTION-PLAYBOOK §2 advarer
% imod; men her er Spørgsmaalet ikke en Hales Tilstedeværelse paa een
% Pixel, men eet Bogstav mod to: efter „Refle“ staae i Noten TO Typer ---
% et rundt, aabent c og derpaa et t med tydelig Overlængde og Tværstreg
% --- medens Textens „Reflexion“ paa samme Side (Linie 3 f. o.:
% „Interferens, Reflexion, Refraction“, og atter Linie 4 f. n. i Brødtexten:
% „den rene Reflexion i sig“) har EEN smal Type uden Overlængde.
% Ordene ere ogsaa maalbart forskjellige i Bredde. Calibreringen er
% altsaa foretaget paa Ord, som Udgaven allerede læser den anden Vei.
% Arbeidsscanningen duer IKKE her: Blyanten har ogsaa naaet dette Sted og
% skrevet et x hen over „ct“, hvad der paa Scanningen staaer som en Klat.
% Graderne ere sat med det trykte Gradtegn; „0“ i „0°“ er et Nul.
Reflection (Tilbagekastning) og Refraction (Brydning) skee altid efter
samme Love, hvad Retningen, men ikke hvad Intensiteten angaaer. To
parallele Speile A og B med en Heldning af lidt over $33^{\circ}$, saa
at Indfalds- og Udfaldsvinkel ved begge blive omtrent $57^{\circ}$,
vise, at de fra A tilbagekastede Straaler enten ganske eller kun
tildeels eller aldeles ikke tilbagekastes af B, eftersom Vinklen mellem
Speilenes Planer er $0^{\circ}$, $90^{\circ}$, $180^{\circ}$ eller
mellemliggende. Herved bestemmes Lysets Plansætning eller Polarisation.
% Spærringen i Noten er efterprøvet paa KB-Exemplaret: „Almindeligviis“
% er spærret, „hedder det,“ er det IKKE, dernæst er alt spærret til og
% med „plansat“, medens „(polariseret)“ igjen staaer i almindelig
% Antiqva. spacing.py's fire Løb her ere altsaa ægte.
\emph{Almindeligviis}, hedder det, \emph{foregaae de paa Straalen
lodrette Svingninger i alle Retninger, men under visse Betingelser
bindes Svingningerne til et vist Plan, og Lyset siges da at være
plansat} (polariseret).}
o. s. v. med alle derunder hørende Ændringer og Afskygninger lader sig
udvikle af een Grundaarsag. I Lysphænomenernes reflecterede Eenhed med
det naturlige Aarsagsbegreb erkjende vi Lysets objective Væsen.

% Det græske Bogstav er læst paa KB-Exemplaret: β, ikke „6“ eller „B“.
% Overskriften er indløbende og spærret.
\subrunin{β) Lysets subjective Væsen.} Uden Syn intet Lys. Opfattet
fra denne Side bliver Lyset altsaa ikke en Function af den svingende
Æther, ikke en Andetheds-, men en Selvhedsfunction, en Function af det
seende Selv. Dette, betragtet for sig, er Lysets subjective Væsen, dets
Idealitet. Denne Idealitet, skjøndt den vitterlig er subjectiv, har den
speculative Naturphilosophie ikke desto mindre forvandlet til noget i
Naturen Objectivt og derved givet et nyt Beviis paa, hvilke
Abstractionstvetydigheder Immanenslæren maa gribe til, naar den i Alvor
forsøger at indlade sig med Virkelig\-%
heden. Lyset, siger Hegel, er Materiens rene Selvhed, mod hvilken den
selv efter sin umiddelbare Tilværen er det Selv\-%
løse, Mørket. Lyset er den uendelige Oscilleren i sig selv, den rene
Reflexion i sig, det Samme, hvad Aandens Jeg er
% Arkmærket „17“ staaer nederst paa S. 257; det er ikke Text og gjengives
% ikke.
% --- p. 258 ---
i høiere Form. Rummet, Lyset og Aanden betegne forsaa\-%
vidt en trinviis Formudvikling af det samme Væsen. „Jeg er det
uendelige Rum, Selvbevidsthedens uendelige Lighed med sig,
Abstractionen af den tomme Selvbevidsthed, den rene Identitet af mig
med mig. Jeg er kun Identiteten i Forholdet af mig selv som Subject til
mig som Object. Med denne Identitet er Lyset parallelt, heraf er det et
tro Af\-%
billede. Det er kun derfor ikke Jeg, fordi det ikke fordunkles og
brydes i og ved sig selv, men er kun den abstracte Skinnen.“

Ved at ivre for Lysets objective Idealitet har Immanenslæren meent, at
den virkelig stred for Aandelighedens Ret, og har i denne Anledning
reist en hæftig, man kan gjerne sige: fanatisk, Polemik mod Physiken.
Især er det Newton, der maa holde for; Hegel vil i hans Forsøg
paa\-%
vise Fadhed, Urigtighed, ja Uredelighed. Der tillægges Physikerne de
% „horribile dictu“ staaer i Kursiv (efterprøvet paa KB-Exemplaret);
% Tankestregerne omkring det ere almindelig Antiqva.
plumpeste Forestillinger, at de tale om Straale\-%
knipper, fiirkantede Straaler, at de --- \textit{horribile dictu} ---
betragte Lyset som noget Materielt og Sammensat og lade hver enkelt
Straale have sin enkelte Farve i det hvide Lys, ligesom hvert enkelt
Horn har sin enkelte Tone i en russisk Hornmusik o. s. v. Med alt dette
bliver dog det, som egenlig er Hovedsagen i Physikens exacte
Undersøgelser, fuldstændig misforstaaet og det til stor Skade --- ikke
for Physiken, men for Philosophien.

Physiken lærer os ikke, at den svingende Æther i og for sig er Lys; den
gjør Øiet til mæglende Organ mellem det Virksomme i Naturen og
Synsvirksomheden, og opfatter Lysaarsagen som en lovbestemt
Vexelvirkning af to ueens\-%
artede Aarsager, en objectiv i Ætheren og en subjectiv i Sjælen. At
Physiken ikke vil indlade sig paa særegen Forklaring af Lysets ideelle
Væsen, er ikke noget Beviis for, at den fornegter Lysets Idealitet, det
er kun et Beviis paa, at den hverken er eller vil være Philosophie.
% S. 258 slutter midt i Ordet „fastholder“: Orddelingen løber over
% Batchgrændsen, og S. 259 begynder med „holder“.
Physiken fast\-%

% [text to be added: pp. 259--273]

%  --- B. Leddelende Functioner: Magnetisme og Elektricitet  (pp. 274–293) ---

%  MATHEMATICAL — work from the image; the OCR is not a witness here.
% [text to be added: pp. 274--287]

% [text to be added: pp. 288--293]

%  --- C. Legemdannende Functioner: Chemisme  (pp. 294–317) ---

% [text to be added: pp. 294--307]

% [text to be added: pp. 308--317]


% ======================================================================
%  Anden Deel — Tredie Kapitel: Det elementære Hele: Jordlegemet   (printed pp. 318–374)
% ======================================================================

\kapitel{Tredie Kapitel.}{Det elementære Hele: Jordlegemet.}

%  --- A. Jordlegemets Sammensætning  (pp. 318–345) ---

% [text to be added: pp. 318--331]

% [text to be added: pp. 332--345]

%  --- B. Jordklodens historiske Udvikling  (pp. 346–361) ---

% [text to be added: pp. 346--361]

%  --- C. Jordlegemet som elementær Organisme  (pp. 362–374) ---

%  MATHEMATICAL — work from the image; the OCR is not a witness here.
% [text to be added: pp. 362--374]


% ======================================================================
%  Tredie Deel — Første Kapitel: Det organiske Livs Oprindelse   (printed pp. 375–444)
% ======================================================================

\deel{Tredie Deel.}{Den organiske Natur.}
% Mellem Deel-Overskriften og Kapitel-Overskriften staaer i Satsen en kort,
% centreret Streg. Efterprøvet paa KB-Exemplaret for S. 1 og S. 197.
\streg

\kapitel{Første Kapitel.}{Det organiske Livs Oprindelse.}

%  --- A. Liv og Organisme  (pp. 375–400) ---

% [text to be added: pp. 375--388]

% [text to be added: pp. 389--400]

%  --- B. Organiske Grundformer  (pp. 401–423) ---

% [text to be added: pp. 401--414]

% [text to be added: pp. 415--423]

%  --- C. Grundformernes Systematik  (pp. 424–444) ---

% [text to be added: pp. 424--437]

% [text to be added: pp. 438--444]


% ======================================================================
%  Tredie Deel — Andet Kapitel: Organiske Naturer   (printed pp. 445–508)
% ======================================================================

\kapitel{Andet Kapitel.}{Organiske Naturer.}

%  --- A. Plantenaturen  (pp. 445–465) ---

% [text to be added: pp. 445--458]

% [text to be added: pp. 459--465]

%  --- B. Dyrenaturen  (pp. 466–490) ---

%  MATHEMATICAL — work from the image; the OCR is not a witness here.
% [text to be added: pp. 466--474]

%  MATHEMATICAL — work from the image; the OCR is not a witness here.
% [text to be added: pp. 475--483]

% [text to be added: pp. 484--490]

%  --- C. Menneskenaturen  (pp. 491–508) ---

%  MATHEMATICAL — work from the image; the OCR is not a witness here.
% [text to be added: pp. 491--508]


% ======================================================================
%  Tredie Deel — Tredie Kapitel: Slutning: Natur og Aand   (printed pp. 509–549)
% ======================================================================

%  NOTE: this chapter head does NOT open a page. It stands
%  part-way down printed p. 509, below the close of the
%  preceding chapter. The head is therefore NOT pre-placed here:
%  the agent transcribing p. 509 writes it, in its printed
%  position, as  \kapitel{Tredie Kapitel.}{Slutning: Natur og Aand.}

%  --- A. Naturen skabes af Aanden  (pp. 509–522) ---

%  MATHEMATICAL — work from the image; the OCR is not a witness here.
% [text to be added: pp. 509--517]

% [text to be added: pp. 518--522]

%  --- B. Naturen bæres af Aanden  (pp. 523–534) ---

% [text to be added: pp. 523--534]

%  --- C. Naturen forklares af Aanden  (pp. 535–549) ---

% [text to be added: pp. 535--549]


% ============================================================
%  TRYKFEIL — the book's own errata leaf, unfoliated, working PDF 571.
%  Transcribed as printed. Nielsen's corrections are NOT applied to the
%  text above; each is noted in a % comment at its site instead, so that
%  a reader sees both what was printed and what the author wanted.
% ============================================================

% ============================================================
%  TRYKFEIL --- the book's own errata leaf.
%
%  WHERE IT IS. An unfoliated leaf following printed p. 549.
%      working scan (Google)   PDF 571
%      KB colour scan          PDF 569
%  Its verso is blank in both copies (KB PDF 570). The leaf carries the
%  errata table and NOTHING else: no colophon, no printer's name, no
%  advertisement, no folio.
%
%  THE TWO COPIES AGREE THROUGHOUT, entry for entry and glyph for glyph.
%  Every reading below was read off the image at high zoom on BOTH scans.
%  The OCR of this leaf is not a witness and was not used: tesseract loses
%  the column alignment, mangles the ditto marks, and misreads the four
%  mathematical entries.
%
%  TYPOGRAPHY OF THE LEAF.
%    - "Trykfeil." is centred and BOLD; it is the only bold line on the
%      leaf. None of the book's heading macros fits an unfoliated errata
%      leaf, so it is set here directly rather than with \kapitel etc.
%    - A short centred rule stands under the heading, and a DOUBLE rule
%      (two parallel lines, set right of centre) closes the table. Per the
%      house rule \streg serves for both and the difference in the printed
%      ornament is not reproduced. Two rules, therefore two \streg.
%    - The leaf is set in a size smaller than the text: \small.
%    - LETTERSPACING. The corrected reading is letterspaced wherever the
%      corrected WORD differs from the printed one, and is NOT letterspaced
%      where only punctuation or a single letter/label changes (pp. 180,
%      183, 233, 443) or where the correction is mathematical (pp. 60, 79,
%      160, 169, and the "= 0, o. s. v." of p. 163). Letterspacing is
%      \emph{} here as everywhere in this edition. The pattern is regular
%      and was checked entry by entry on both copies.
%
%  DITTO MARKS. The Side and Linie columns repeat by ditto, printed „ --
%  the low double comma, the same sort as the book's opening quotation
%  mark. They are transcribed AS PRINTED, and the EDITORIAL EXPANSION of
%  each is given in the % comment heading its row, so that this fragment
%  can be grepped as a lookup table without emending the page. Three rows
%  print NO ditto at all in the Side column -- it is simply left blank --
%  and are transcribed so: p. 163 lines 10 and 8 f. n., p. 178 line 9 f. o.
%
%  ABBREVIATIONS. f. o. = fra oven (from the top); f. n. = fra neden (from
%  the bottom); læs = read. Line counting on the printed page INCLUDES the
%  footnote lines and EXCLUDES the signature mark -- verified on p. 257,
%  where "13 f. n." lands on the first line of the footnote, and on p. 357,
%  where "2 f. n." lands on the last text line but one.
%
%  p. 257 AND p. 357 ARE TWO DIFFERENT ENTRIES, not a folio confusion. The
%  book misprints p. 357's folio as "257" (see the file header), so the
%  pair invites suspicion; both were checked against their pages and both
%  are right. p. 257 line 13 f. n. is the footnote opening "Reflection
%  (Tilbagekastning) og Refraction (Brydning)"; p. 357 line 2 f. n. reads
%  "i den tertiære Perioder (Molassegruppen)".
%
%  PRINTER'S DEFECTS ON THE LEAF ITSELF, both copies, transcribed as printed:
%    - p. 79 entry: the em-dash between the wrong and the corrected reading
%      is MISSING -- the line reads "Q_k læs Q'_k". Every other entry on the
%      leaf carries the dash.
%    - p. 60 entry: the minus sign in the subscript of the corrected reading
%      prints as two short strokes, "n--1", in both copies. The sense-reading
%      n-1 is set here; the doubled stroke is a battered or doubled rule, not
%      a second character.
% ============================================================

\clearpage

% The heading is centred and bold; it is not one of the book's numbered heads.
{\centering\large\textbf{Trykfeil.}\par}

\streg  % short centred rule under the heading

\begingroup
\small
\setlength{\tabcolsep}{4pt}
\begin{center}
\begin{tabular}{@{}l r l r l l@{}}
% p. 25, line 11 f. o.: sin -> læs sine
Side & 25 & Linie & 11 & f.\ o. & sin --- læs \emph{sine}\\
% p. 60, line 1 f. n.: sidste Led i Tælleren (x-x_n) -> læs (x-x_{n-1})
% „ = Side, „ = Linie (ditto in both columns)
„ & 60 & „ & 1 & f.\ n. & sidste Led i Tælleren $(x-x_{n})$ --- læs $(x-x_{n-1})$\\
% p. 79, line 1 f. o.: Q_k -> læs Q'_k   [NO em-dash printed; see header]
% „ = Side, „ = Linie
„ & 79 & „ & 1 & f.\ o. & $Q_{k}$ læs $Q'_{k}$\\
% p. 113, line 8 f. n.: blive til -> læs blive til.
% „ = Side, „ = Linie
„ & 113 & „ & 8 & f.\ n. & blive til --- læs \emph{blive til.}\\
% p. 136, line 14 f. o.: Hældning -> læs Heldning
% „ = Side, „ = Linie
„ & 136 & „ & 14 & f.\ o. & Hældning --- læs \emph{Heldning}\\
% p. 140, line 2 f. n.: quiqvam -> læs quidqvam
% „ = Side, „ = Linie
„ & 140 & „ & 2 & f.\ n. & quiqvam --- læs \emph{quidqvam}\\
% p. 160, line 13 f. n.: c/mu i Tælleren -> læs c^2/mu
% „ = Side, „ = Linie
„ & 160 & „ & 13 & f.\ n. & $\dfrac{c}{\mu}$ i Tælleren --- læs $\dfrac{c^{2}}{\mu}$\\
% p. 163, line 13 f. o.: det Antilogiske -> læs det Antilogiske,
% „ = Side, „ = Linie
„ & 163 & „ & 13 & f.\ o. & det Antilogiske --- læs \emph{det Antilogiske,}\\
% p. 163, line 10 f. n.: være med -> læs være, med
% Side column BLANK on the leaf (no ditto printed); „ = Linie
 & & „ & 10 & f.\ n. & være med --- læs \emph{være, med}\\
% p. 163, line 8 f. n.: = 0; -> læs = 0, o. s. v.
% Side column BLANK on the leaf (no ditto printed); „ = Linie
 & & „ & 8 & f.\ n. & $=0$; --- læs $=0$, o.\ s.\ v.\\
% p. 164, line 15 f. o.: pertuberede -> læs perturberede
% „ = Side, „ = Linie
„ & 164 & „ & 15 & f.\ o. & pertuberede --- læs \emph{perturberede}\\
% p. 169, line 19 f. o.: E = 2e, -> læs ε = 2e,
% „ = Side, „ = Linie. The corrected symbol is a Greek small epsilon,
% lunate on both copies -> \varepsilon.
„ & 169 & „ & 19 & f.\ o. & $E=2e$, --- læs $\varepsilon = 2e$,\\
% p. 176, line 4 f. n.: Ingriben -> læs Indgriben
% „ = Side, „ = Linie
„ & 176 & „ & 4 & f.\ n. & Ingriben --- læs \emph{Indgriben}\\
% p. 178, line 5 f. o.: dem. -> læs den.
% „ = Side, „ = Linie
„ & 178 & „ & 5 & f.\ o. & dem. --- læs \emph{den.}\\
% p. 178, line 9 f. o.: hældende -> læs heldende
% Side column BLANK on the leaf (no ditto printed); „ = Linie
 & & „ & 9 & f.\ o. & hældende --- læs \emph{heldende}\\
% p. 180, line 13 f. n.: c) Systemets -> læs b) Systemets
% „ = Side, „ = Linie
„ & 180 & „ & 13 & f.\ n. & c) Systemets --- læs b) Systemets\\
% p. 183, line 9 f. n.: Qvadratet*) -> læs Qvadratet*).
% „ = Side, „ = Linie. The correction adds the full stop after the
% footnote mark; the mark itself is unchanged.
„ & 183 & „ & 9 & f.\ n. & Qvadratet*) --- læs Qvadratet*).\\
% p. 184, line 2 f. n.: ans -> læs aus
% „ = Side, „ = Linie
„ & 184 & „ & 2 & f.\ n. & ans --- læs \emph{aus}\\
% p. 206, line 3 f. o.: tilstrækkelig; -> læs tilstrækkeligt;
% „ = Side, „ = Linie
„ & 206 & „ & 3 & f.\ o. & tilstrækkelig; --- læs \emph{tilstrækkeligt;}\\
% p. 224, line 5 f. o.: mechaniske -> læs chemiske
% „ = Side, „ = Linie
„ & 224 & „ & 5 & f.\ o. & mechaniske --- læs \emph{chemiske}\\
% p. 233, line 13 f. n.: B. -> læs C.
% „ = Side, „ = Linie. A division label, not a word.
„ & 233 & „ & 13 & f.\ n. & B. --- læs C.\\
% p. 256, line 12 f. n.: vil -> læs ville
% „ = Side, „ = Linie
„ & 256 & „ & 12 & f.\ n. & vil --- læs \emph{ville}\\
% p. 257, line 13 f. n.: Reflection -> læs Reflexion
% „ = Side, „ = Linie. This is the real p. 257 (the footnote), NOT the
% leaf whose folio misprints as "257"; see the header.
„ & 257 & „ & 13 & f.\ n. & Reflection --- læs \emph{Reflexion}\\
% p. 287, line 15 f. n.: Polentialet -> læs Potentialet
% „ = Side, „ = Linie
„ & 287 & „ & 15 & f.\ n. & Polentialet --- læs \emph{Potentialet}\\
% p. 355, line 11 f. o.: Organismens -> læs Organismernes
% „ = Side, „ = Linie
„ & 355 & „ & 11 & f.\ o. & Organismens --- læs \emph{Organismernes}\\
% p. 357, line 2 f. n.: Perioder -> læs Periode
% „ = Side, „ = Linie. This is the leaf whose folio misprints as "257".
„ & 357 & „ & 2 & f.\ n. & Perioder --- læs \emph{Periode}\\
% p. 407, line 6 f. o.: anticulerede -> læs articulerede
% „ = Side, „ = Linie
„ & 407 & „ & 6 & f.\ o. & anticulerede --- læs \emph{articulerede}\\
% p. 414, line 4 f. n.: tl -> læs til
% „ = Side, „ = Linie
„ & 414 & „ & 4 & f.\ n. & tl --- læs \emph{til}\\
% p. 418, line 4 f. o.: betegne -> læs betegner
% „ = Side, „ = Linie
„ & 418 & „ & 4 & f.\ o. & betegne --- læs \emph{betegner}\\
% p. 430, lines 2--3 f. n.: Phænomer -> læs Phænomener
% „ = Side, „ = Linie. The only entry spanning two lines; the printed
% line reference is "2-3" with a short dash.
„ & 430 & „ & 2--3 & f.\ n. & Phænomer --- læs \emph{Phænomener}\\
% p. 443, line 13 f. o.: Synsvidde -> læs Synsvidde.
% „ = Side, „ = Linie. The correction adds the full stop only.
„ & 443 & „ & 13 & f.\ o. & Synsvidde --- læs Synsvidde.\\
% p. 479, line 2 f. o.: Cascuta -> læs Cuscuta.
% „ = Side, „ = Linie. Last entry on the leaf; the full stop after
% "Cuscuta" is printed and may simply close the table.
„ & 479 & „ & 2 & f.\ o. & Cascuta --- læs \emph{Cuscuta.}\\
\end{tabular}
\end{center}
\endgroup

\streg  % the leaf closes with a DOUBLE rule; \streg per the house rule

\end{document}
